% !TeX program = xelatex
\documentclass[11pt,a4paper,oneside]{ctexart}
\usepackage[margin=25mm,headheight=15pt]{geometry}
\usepackage{amsmath,amssymb,amsthm,mathtools,bm,mathrsfs}
\usepackage{booktabs,longtable,array,graphicx,calc}
\usepackage{textcomp}
\usepackage{etoolbox}
\AtBeginEnvironment{longtable}{\small}
\allowdisplaybreaks[2]
\providecommand{\real}[1]{#1}
\providecommand{\pandocbounded}[1]{#1}
\usepackage{xcolor,enumitem,fancyhdr}
\usepackage[unicode,colorlinks=true,linkcolor=blue!45!black,citecolor=blue!45!black,urlcolor=blue!45!black]{hyperref}
\usepackage{bookmark}
\setcounter{tocdepth}{1}
\setcounter{secnumdepth}{2}
\newtheoremstyle{cnplain}{6pt}{6pt}{\normalfont}{}{\bfseries}{。}{0.5em}{}
\theoremstyle{cnplain}
\newtheorem{theorem}{定理}[section]
\newtheorem{proposition}[theorem]{命题}
\newtheorem{lemma}[theorem]{引理}
\newtheorem{corollary}[theorem]{推论}
\theoremstyle{definition}
\newtheorem{definition}[theorem]{定义}
\newtheorem{example}[theorem]{例}
\newtheorem{problem}[theorem]{问题}
\theoremstyle{remark}
\newtheorem{remark}[theorem]{注}
\renewcommand{\proofname}{证明}

\newcommand{\Q}{\mathbb Q}
\newcommand{\C}{\mathbb C}
\newcommand{\Z}{\mathbb Z}
\newcommand{\Gm}{\mathbb G_m}
\newcommand{\Gr}{\operatorname{Gr}}
\newcommand{\Ext}{\operatorname{Ext}}
\newcommand{\End}{\operatorname{End}}
\newcommand{\Hom}{\operatorname{Hom}}
\newcommand{\rank}{\operatorname{rank}}
\newcommand{\im}{\operatorname{im}}
\newcommand{\id}{\operatorname{id}}
\newcommand{\Res}{\operatorname{Res}}
\newcommand{\dR}{\mathrm{dR}}
\newcommand{\tightlist}{\setlength{\itemsep}{0pt}\setlength{\parskip}{0pt}}
\setlength{\emergencystretch}{3em}
\pagestyle{fancy}
\fancyhf{}
\fancyhead[L]{\small 费曼周期的可见结构与重整化}
\fancyhead[R]{\small 长论文稿}
\fancyfoot[C]{\thepage}
\hypersetup{pdftitle={费曼周期的可见结构与重整化},pdfauthor={Yu Zhehao}}

\usepackage{tikz}
\usetikzlibrary{arrows.meta,positioning}
\numberwithin{equation}{section}
\newtheorem*{maintheorem}{结构定理}
\ctexset{section={format=\Large\bfseries},subsection={format=\large\bfseries},subsubsection={format=\normalsize\bfseries}}
\setlist{itemsep=2pt,topsep=4pt}
\newcommand{\otherref}[2]{\href{#1.pdf}{#2}}
\title{费曼周期的可见结构与重整化\\[5pt]\large Visible Structures and Renormalization of Feynman Periods}
\author{Yu Zhehao\\[4pt]\small 北京大学物理学院}
\date{2026年10月13日}
\begin{document}\maketitle
\begin{abstract}
费曼积分所需的几何不仅由图决定，也由边界、积分链和物理读出决定。本文在实际相对模型中建立一个结构判据：留数像给出背景之外的组成上界，链相容的割给出下界，而由背景对象承载的加性变化保持其余简单组成及重数。锚定的有限重整化将这一稳定性用于方案比较。若微分算子消去完整相对留数类，其非齐次源属于固定环境的周期空间。
四传播子banana给出纯K3核、质量Kummer商、regulator与最小标架的具体实现；指定有理纤维的秩八结论使用明确列出的几何和有限认证输入。切割、森林端点下降及二次减除sunset将这一框架连接到有限质量两圈MOM的RG核，其紫外系数为$3$与$-17/3$。独立标量散射基准检验归一化、交叉和相应阶数的幺正性。一般图的完整分类和基准小数的窄区间认证仍为开放问题。
\end{abstract}
\noindent\textbf{关键词：}费曼周期；相对动机；物理可见性；K3；重整化；MOM。

\tableofcontents


\part{结构框架}
\section{引言}\label{ch:intro}
\subsection{从一个积分到它所需要的几何}
一个费曼积分可以有许多不同的表示。对物理而言，这些表示描述同一个量；对几何而言，它们可能显露出不同的部分。最大割可以给出一条椭圆曲线或一个K3曲面，而完整积分还包含边界和扩张。减除又会加入参考点的周期。因此，我们要问的不是某张图是否“属于椭圆类”，而是：给定质量、运动学、分子、物理链和减除条件，哪些几何组成必然进入这个积分，它们在物理操作之后是否仍然存在？

动机周期为这个问题提供了自然的语言。Brown的形式周期及最小对象将积分关系与实际态射联系起来\cite{brownperiods,browncosmic}；Bloch--Kreimer和Brown--Kreimer将重整化置于相对几何与森林公式中\cite{bk2008,bk2011,bk2012}。sunset的椭圆描述及三圈banana的高阶正规函数说明，纯几何与混合扩张必须一起研究\cite{bv,bkv}。另一方面，Kompaniets--Panzer的六圈MS计算提供了局部RG系数与临界指数的精确参照\cite{kp}。从几何对象到这些物理读出，中间仍须指定链、边界和方案。

本文建立一个共同的结构判据。在实际相对模型存在的范围内，留数正合列给出完整周期的几何上界，链相容的割给出物理下界；二者相等时，可见组成及其重数便被确定。加入由指定背景承载的周期，不改变背景之外的组成。锚定重整化的有限方案关系使这一代数事实成为物理保护定理。四传播子banana给出包括纯核、质量Kummer商、regulator和最小标架的完整实例；森林端点下降、二次减除sunset与有限质量两圈MOM给出相同判据的操作和物理实现。

全文分为四部分。第一部分陈述并证明结构判据，以一圈结构和边界椭圆检验其含义。第二部分从组合入口进入banana的实际对象、比较与参数导数。第三部分构造切割、收缩和减除所需的几何态射。第四部分给出两圈MOM的RG出口及独立散射基准。扩展图族和较长的操作计算收于同一论文的补充材料。

\subsection{结构定理及其含义}
写一个形式周期为$p=[E,\sigma,\lambda]$：$E$是实际对象，$\sigma$是Betti链，$\lambda$是de Rham读出。令$R_\sigma$为包含链的最小子对象，$I_\lambda$为被读出消去的最大子对象，则
\[
 M(p)=R_\sigma/(R_\sigma\cap I_\lambda).
\]
这是矩阵系数的最小对象。它同时保留链与读出，因而区别于承载整个积分族的环境对象。

\begin{maintheorem}[完整模型中的可见组成]
在特征零、有理系数及第\ref{ch:unified}节规定的实际范畴中，设光滑除子族$Y_s\subset X$横截固定正常交叉边界$B$，并设物理周期属于
\[
 M_s=H^n(X\setminus Y_s,B\setminus Y_s),\qquad E_s=M_s^\vee.
\]
记$J=H^n(X,B)$、$A_s=\operatorname{im}(M_s\xrightarrow{\rm res}H^{n-1}(Y_s,B\cap Y_s)(-1))$。取包含$J^\vee$全部组成的Serre背景$\mathcal S$。

\begin{enumerate}[label=(\roman*),leftmargin=2em]
\item 背景外的物理组成受$A_s^\vee$控制。若实际管态射$f:C\to E_s$及割链$\gamma$满足$R_{f\gamma}\subset R_\sigma$，则对每个背景外简单对象$V$，
\[
 [M(C;\gamma,f^*\lambda):V]\leq[M(p):V]\leq[A_s^\vee:V].
\]
上下界相等时，该重数被精确确定。
\item 若$\ell$由$\mathcal S$中的对象承载，Serre商中有$\pi M(p+\ell)\simeq\pi M(p)$。特别地，在本文构造的相容MOM图类中，共同锚点后的方案变化保护严格低圈背景之外的全部简单组成及重数。
\item 若$L\in k(S)[\partial_s]$消去实际相对留数类$\operatorname{res}(\omega_s)$，则
\[
 L I(s)=\sum_j a_j(s)P_j,\qquad a_j\in k(S),
\]
其中$P_j$是固定环境$J$的周期。
\end{enumerate}
关于固定Nori纤维的陈述要求相应实际态射；仅由单值得到的保护首先是变分层面的结论。
\end{maintheorem}

证明分为四个短步骤：留数像与环境分离、最小标架的加性稳定性、链相容割的单射，以及连接下的线性提升。第\ref{ch:unified}节逐步给出完整证明。相容MOM的低圈方案公式在第\ref{ch:finitemom}节独立构造，故方案应用并不反过来充当代数定理的前提。

这个定理首先回答组成问题。Serre商会忘掉背景内部的扩张；要理解完整周期，还须恢复这些扩张怎样粘合。banana实例使这一点具体化：纯K3核、质量对数和顶层链共同构成一个带过滤的对象。指定$m=(1,2,4,8)$、$s=-64$、基域$\mathbb Q$时，第\ref{ch:core}--\ref{ch:fixed}节在所列有限证书输入下给出最小秩八。四参数很一般纤维的秩十和等质量很一般纤维的秩四另行证明，不能由这一固定点外推。

同样，微分方程只记录连接的一部分信息。消去紧致纯割的算子可能留下相对边界源；只有消去完整相对留数类时，定理才将源限制到固定环境。这是本文组织Picard--Fuchs方程和重整化导数的共同出发点。

\subsection{已有工作、证明依赖与约定}
\paragraph{标准工具与本文的使用。}
Serre商的存在性和正合性是标准同调代数\cite{stacks}；形式周期的最小对象沿用Brown\cite{brownperiods}。本文所需的加性结论进一步追踪两个核在扰动载体中的位置，从而保留背景外的\emph{重数}。我们将这一结论与实际森林方案差连接，不把Serre商本身作为新构造。

一圈积分的混合Hodge结构和双对数背景来自Bloch--Kreimer\cite{bkone}。第\ref{ch:unified}节的一圈三角图与盒图沿用这一几何起点，列出未分裂基域上的算术角色，并把实际物理链、固定环境及留数源纳入同一模型。sunrise的椭圆曲线与任意质量微分方程已有系统研究\cite{abwsun,bkvsunset}；补充材料的串联细分部分使用这些端点对象，所要证明的是实际相对下降、阈值分离后的可见重数以及分子对这些组成的选择。

Brown--Kreimer已建立参数重整化中的尺度、角度与森林结构\cite{bk2011,bk2012}。本文的森林部分给出规定图类和参考条件下的共同辅助参数、相对端点态射及其与物理BPHZ规则的比较。二次减除sunset单独处理幂次发散。两圈MOM部分则在同一有限质量方案中计算读出；Kompaniets--Panzer的MS结果用于核对紫外系数与比较方案，不能直接替代有限质量核\cite{kp}。上述比较说明各项结果所承担的任务；一般图的优先权与完备分类并不由这些实例推出。

\begin{figure}[ht]
\centering
\begin{tikzpicture}[node distance=7mm and 6mm,box/.style={draw=black!45,rounded corners=1pt,align=center,font=\small,text width=43mm,minimum height=11mm},arr/.style={-Latex,draw=black!65}]
\node[box] (geo) {相对模型与留数正合列\\第\ref{ch:unified}节};
\node[box,right=of geo] (alg) {最小标架与Serre保护\\第\ref{ch:unified}节};
\node[box,right=of alg] (ren) {共同森林与相容MOM\\第\ref{ch:zero}、\ref{ch:finitemom}节};
\node[box,below=of geo] (ban) {banana对象、比较、群\\第\ref{ch:core}--\ref{ch:fixed}节};
\node[box,below=of alg] (cut) {物理割与导数读出\\第\ref{ch:connection}、\ref{ch:cut}节};
\node[box,below=of ren] (sun) {端点下降与sunset\\第\ref{ch:descent}、\ref{ch:sunsetjet}节};
\node[box,below=of cut] (protect) {新组成的检测与保护\\推论\ref{unified:cor:prediction}};
\node[box,below=of sun] (rg) {有限质量两圈MOM RG\\第\ref{ch:momrg}节};
\draw[arr] (geo)--(ban);\draw[arr] (geo)--(cut);
\draw[arr] (alg)--(cut);\draw[arr] (cut)--(protect);
\draw[arr] (ren.west) to[out=225,in=40] (protect.north east);
\draw[arr] (ren)--(sun);\draw[arr] (sun)--(rg);
\draw[arr] (ban)--(cut);
\end{tikzpicture}
\caption{主要证明依赖。箭头表示被调用的输入。banana的固定纤维认证与四维两圈MOM是两条独立实例链；两者共享结构方法，后者不以秩八K3对象为前提。}
\label{fig:dependencies}
\end{figure}

\paragraph{范畴与归一化。}\label{sec:notation}
基域$k\subset\mathbb C$及其指定嵌入始终保留，动机取有理系数。$\mathbb Q(j)$的权重为$-2j$，标准Tate周期为$(2\pi i)^j$。物理链采用同调方向，几何留数常先以上同调书写，再取对偶。族中使用几何产生的可容许混合Hodge变分及其子商；固定代数纤维使用实际Nori对象。六运算与实现的相容性采用Tubach的已发表结果\cite{tubach}，一动机商采用Ayoub--Barbieri-Viale\cite{abv}。这两项基础输入及实际使用方式在附录\ref{app:inputs}重述。

本文不假定数值周期映射单射。一个数在一点为零、一个族函数恒为零、一个形式周期为零和一个实际子对象被读出消去，有不同的证明条件。指定banana纤维的有限证书作为明确输入列出；散射基准的普通小数与严格区间分别标注。读者可以先读第\ref{ch:unified}节及banana实例，再沿图\ref{fig:dependencies}进入所需的操作。

\section{统一结构框架}\label{ch:unified}
本节从最小标架出发。先证明加入背景周期时的稳定性，再把几何留数与物理链放入同一组上下界。最后，连接将这组结构关系变成非齐次微分方程。

\begingroup
\providecommand{\Q}{\mathbb Q}\providecommand{\cC}{\mathcal C}
\providecommand{\cS}{\mathcal S}\providecommand{\dR}{\mathrm{dR}}
\providecommand{\rk}{\operatorname{rk}}\providecommand{\im}{\operatorname{im}}
\subsection{带标架的实际对象}
\subsubsection{输入与分类输出}
输入为
\[
 \mathfrak P=(G,D,\nu,\mu,\kappa,N,\sigma,k,R,\mathcal O),
\]
分别包含图、维数、传播子幂次、质量、运动学、分子、链、基域、减除方案及可观测量。
输出首先是带标架的实际对象，不是仅由边表决定的“MPL/椭圆/K3”单选标签。

在已构造的有理Nori范畴或明确指定的实现范畴$\cC$中，采用同调方向：
$E$承载链$\sigma\in E_B$，形式为读出$\lambda\in E_{\dR}^{\vee}$。
de Rham系数可在数域$k$，实际子对象仍在原范畴中求，保留下降条件。
周期记作$p=[E,\sigma,\lambda]$，此记号省略了与上同调约定之间的对偶。定义
\[
 R_\sigma=\bigcap_{\sigma\in J_B}J,\quad
 I_\lambda=\sum_{\lambda|J_{\dR}=0}J,\quad
 M(p)=R_\sigma/(R_\sigma\cap I_\lambda).
\]
有限长度与实现正合性保证存在。最小矩阵系数对象的表示无关性是经典结果
\cite[\S2.4]{brownperiods}，不是本文新创的定义定理。

完整分类包含三组数据：简单组成、权重、Hodge数、极化与算术下降；
各层间的实际粘合与多层相容关系；标架及每个操作的源、靶、实际态射。
只列简单组成会漏掉Kummer和regulator；只列相邻$\operatorname{Ext}^1$
也未必确定多层对象。精确输出仍是带过滤和标架的实际对象或等价的已证明表示。

\subsubsection{结构表的分层语义}
\begin{center}\small
\begin{tabular}{p{24mm}p{47mm}p{72mm}}
\toprule 名称&合法结论的层次&不能由名称单独推出\\\midrule
MPL&指定字母、路径和端点的迭代积分表示&任意代数字母的值都是MZV\\
MZV&指定归一化的常数及其形式周期&全部mixed Tate周期属于该子类\\
mixed Tate&实际对象的组成均为相应Tate对象&扩张分裂或自动有指定MPL表达式\\
椭圆/高亏格&曲线$H^1$及可见因子&同亏格等于同对象或全部因子物理可见\\
K3&指定超越$H^2$及边界粘合&秩六方程必为K3或整个周期只有纯层\\
较高维CY&指定几何、次数与上同调证书&任意高阶微分方程都是CY\\
Kummer&单位及其实际扩张类&不需基域、分支和比较数据\\\bottomrule
\end{tabular}
\end{center}
例如质量单位$2$的Kummer扩张与特定MZV端点对象有不同的算术输入。
这一说明不需要断言$\log2$与某组普通数值之间的未知代数独立性。
本表规定输出语义；充分判据仍须由实际几何给出。


\subsection{背景周期的加性稳定性}
\subsubsection{背景子范畴与正确的商}
设$\cS\subset\cC$为Serre子范畴，即对子对象、商与扩张封闭。
记精确商函子$\pi:\cC\to\cC/\cS$，采用标准存在性定理
\cite[Tag 02MS]{stacks}。该商只比较哪些组成超出背景结构，
一般不是带纤维函子的动机范畴。

若背景含单位对象，就不能同时要求它是对任意环境对象张量封闭的
“张量理想”：$1\otimes E=E$会使其等于全部$\cC$。
下文用给定背景张量子范畴的\emph{Serre闭包}，没有这一要求。
在有限长度范畴中，它恰由全部简单组成属于背景清单的对象构成。

\begin{theorem}[背景周期的加性稳定性]\label{unified:thm:addition}
设$p$为实际形式周期，$\ell=[L,\tau,\eta]$，$L\in\cS$，
$q=p+\ell$。则
\begin{equation}\label{unified:eq:quotient}
 \pi M(q)\simeq\pi M(p).
\end{equation}
对任何不属于$\cS$的简单对象$S$，有
\[
 [M(q):S]=[M(p):S].
\]
另有与给定扰动载体有关的秩界
\[
 |\rk M(q)-\rk M(p)|\leq\rk L .
\]
无需计算$p,\ell$的复数值，也不需要周期映射单射性。
\end{theorem}
\begin{proof}
取$P=M(p)$的最小标架$(\sigma,\lambda)$。
于是$\sigma$生成$P$，且$P$中没有被$\lambda$消去的非零子对象。
$q$在$P\oplus L$上由链$(\sigma,\tau)$和读出$(\lambda,\eta)$表示。
令$R$为该链生成的子对象，$I\subset R$为该读出消去的最大子对象，
则$M(q)=R/I$。

投影$R\to P$的像承载$\sigma$，故满射；其核$K$是$L$的子对象，
所以$K\in\cS$。另一方面，
\[
 I\cap(P\oplus0)=0,
\]
因为该交是$P$的实际子对象，在其上$\lambda$为零，违背最小性。
因此投影$I\to L$为单射，故$I\in\cS$。两条短正合列
\[
 0\to K\to R\to P\to0,\qquad
 0\to I\to R\to M(q)\to0
\]
在$\pi$下成为同构，证明\eqref{unified:eq:quotient}。
组成重数的可加性给
$[M(q):S]-[P:S]=[K:S]-[I:S]$；
对背景外的$S$右侧为零。
秩差为$\rk K-\rk I$，二者都介于$0$和$\rk L$，得秩界。
这些均是实际子商，不是数值相加后猜测对象。
\end{proof}

该结论强于“加入背景后生成同一张量范畴”：
后者不决定新简单因子在对象中出现几次，本定理保留精确重数。
背景可代表低圈图、已知边界或其它允许结构，
但须先证明扰动确由该背景内的实际对象承载。

\subsubsection{为何仍必须单独研究混合扩张}
若$\cS$为所有组成均为Tate的对象，则非分裂Kummer对象和其分裂替代都映为零。
例如$p=\log^{\mathfrak m}2$、$\ell=-p$允许$q=0$：
它们在商中都是零，原最小对象不同。
所以该商不替代regulator、单位和完整coaction，
只精确回答背景之外的新结构能否被加性变化消掉。
完整分类要把这一稳定纯组成与实际粘合并列保留。


\subsection{锚定方案下的新组成保护}
\subsubsection{实际图类与低圈背景}
采用本文已证明的正质量欧氏$D=4$、单位传播子、至多对数发散图类，
使用已构造的零动量或相容代数非零MOM参考方案。
令$a,b$为两方案，$h_G=h>0$，固定共同物理锚点$\kappa_*$：
\[
 p_a=\widehat\Phi_a^{\mathfrak m}(G;\kappa,\kappa_*)
 =\Phi_a^{\mathfrak m}(G;\kappa)-\Phi_a^{\mathfrak m}(G;\kappa_*).
\]
当前点、锚点及参考数据在已证正质量代数欧氏范围。
改变基域或代数规范时，所需算术对象须显式纳入背景。

第\ref{ch:finitemom}节的实际有限方案定理给
\begin{equation}\label{unified:eq:forest}
 p_b-p_a=\sum_{\varnothing\ne\gamma\subsetneq G}
 g_{ba}^{\mathfrak m}(\gamma)\,
 \widehat\Phi_a^{\mathfrak m}
 (G/\gamma;\pi_\gamma\kappa,\pi_\gamma\kappa_*).
\end{equation}
子图可为森林余积允许的不连通并，每个因子均来自严格低圈实际载体。
经典森林背景见\cite{bk2012}；实际周期提升与截断物理比较见第\ref{ch:zero}、\ref{ch:finitemom}节。

令$\mathcal B_{<G}$为本节转换所需有限子图--收缩闭包、
参考点和当前/锚点评价的全部严格低圈载体生成的满张量子范畴，
包含所需Tate因子；$\cS_{<G}$为其在共同环境中的Serre闭包。
“新几何”总是相对于这个明确背景，不能只按圈数口头指定。

\begin{theorem}[新结构的方案不变性]\label{unified:thm:ren}
在上述范围，
\[
 \pi_{<G}M(p_a)\simeq\pi_{<G}M(p_b).
\]
每个背景外简单对象$S$的重数相同。
若全部低圈背景mixed Tate，而某方案的锚定最小对象含非Tate简单组成，
另一方案仍含同一组成及同样重数。
\end{theorem}
\begin{proof}
式\eqref{unified:eq:forest}每项乘积由实际载体的张量积承载，
有限和由这些张量积的直和$L$承载。
故$L\in\mathcal B_{<G}\subset\cS_{<G}$。
应用定理\ref{unified:thm:addition}即可。
若低圈均mixed Tate，其张量、子商及扩张仍只有Tate组成，
因此非Tate简单因子在背景外。
\end{proof}

这直接接合结构表与物理操作：一旦几何证明某个新纯因子进入指定标架，
允许的锚定方案变化不能将其改成背景类型。
没有求完整积分，也没有把局部性误解为所有常数消失。

\subsubsection{锚点条件及物理求和的区别}
未锚定差含整图常数$g_{ba}(G)$，
可含同圈数参考周期，不能归到低圈。
本原图已经有
\[
 \Phi_b(G;\kappa)-\Phi_a(G;\kappa)
 =\Phi_0(G;\rho_a)-\Phi_0(G;\rho_b).
\]
bubble在自己的减除点为零，改参考后该点可严格非零，
所以孤立最小对象可由零变为非零。
物理分子或完整图和的改变也不是方案变化；
它们可能在同一新几何内部真实相消，本定理没有禁止这种现象。


\subsection{几何上界与物理下界}
\subsubsection{有限实际操作的上界}
\begin{proposition}[结构上界]\label{unified:prop:upper}
给定有限组实际种子，允许有限直和、张量、对偶、Tate扭转、
实际态射的核像商及已给出的短正合列。
若其组成属于对子商、张量、对偶及扩张封闭的背景类，
则全部输出及物理最小对象的组成仍属于该类。
\end{proposition}
\begin{proof}
直和与正合列的重数可加，子商只删组成。
张量积对有限长度过滤正合，其相继商是组成的张量积；
再取组成，连同对偶和扭转即可归纳。
物理最小化最后也是子商。
\end{proof}
扩张仍可增加regulator与幺幂深度。
切割、收缩及Gauss--Manin导数也不都是固定对象的自同态：
新几何必须给其实际对象与比较，而非从形式相似推断。

\subsubsection{链相容的割给强制下界}
设实际管态射$f:C\to E$已包含正确Tate扭转与取向。
割读出为$f^*\lambda$，割链为$\gamma$，
其像在$E$中生成$R_{f\gamma}$。
\begin{proposition}[割的物理保护]\label{unified:prop:lower}
若$R_{f\gamma}\subset R_\sigma$，则有实际单射
\[
 M([C,\gamma,f^*\lambda])\hookrightarrow M([E,\sigma,\lambda]).
\]
故割最小对象的每个简单组成在物理最小对象中至少有同样重数。
\end{proposition}
\begin{proof}
实际标架关系识别左侧与
$R_{f\gamma}/(R_{f\gamma}\cap I_\lambda)$，
再由包含关系注入$R_\sigma/(R_\sigma\cap I_\lambda)$。
\end{proof}
原链循环生成$E$时条件自动满足。
否则，非零广义割可能来自原链没有激发的方向，只能证明环境几何。
完整物理求和必须先形成总读出，再检验割，不能逐图忽略干涉。

\begin{corollary}[不求完整积分的不可消去证书]\label{unified:cor:prediction}
若实际几何上界与标架已构造，某个割满足上述链条件，
割最小对象含简单因子$S$，且$S\notin\cS_{<G}$，
则物理输出必须含$S$，并在定理\ref{unified:thm:ren}的全部锚定方案中保留。
若上界中$S$的重数等于割下界，则该重数已精确确定。
\end{corollary}
\begin{proof}
依次用命题\ref{unified:prop:lower}及定理\ref{unified:thm:ren}，再夹逼整数重数。
\end{proof}
这是可以贯通本文的主判据。每个前提都要有证据；
未知的链生成性或regulator非零不能藏在“自动编译”中。
banana指定纤维的相应输入已有独立证明；
一般图可以停在上界或条件下界，但须标明输出强度。

\subsubsection{由物理变分验证链条件：族与固定纤维分开}
令$S$为连通光滑参数曲线，$\mathcal E$为已经构造的相对上同调
同调对偶变分，具有有理Betti局部系统和代数de Rham连接。
下面在一个有限长度的实际变分类别中取子对象，且其Betti实现为局部系统；
可使用几何产生的可容许混合Hodge变分及其子商。
不据实现结果宣称存在尚未构造的Nori族态射。
固定基点$s_*$、链$\sigma$及单值代数读出$\lambda$。
$R_\sigma^{\rm var}$表示包含$\sigma$的最小子变分，
$I_\lambda^{\rm var}$表示被$\lambda$恒等消去的最大子变分。

\begin{theorem}[物理变分给出的割保护]\label{unified:thm:variation}
设实际管态射$f:\mathcal C\to\mathcal E$和路径$\ell$满足
\[
 (T_\ell-1)\sigma=c\,f_B\gamma,\qquad c\in\Q^\times.
\]
则在变分类别中
\[
 M_{\rm var}(\mathcal C;\gamma,f^*\lambda)
 \hookrightarrow M_{\rm var}(\mathcal E;\sigma,\lambda).
\]
若实际管像包含一个Betti局部系统不可约的子变分$\mathcal A$，
且$f_B\gamma\in\mathcal A_{s_*}$非零、$\lambda|_{\mathcal A}\not\equiv0$，
则$\mathcal A$单射物理最小变分。
\end{theorem}
\begin{proof}
$R_\sigma^{\rm var}$的Betti纤维被所有单值保持，所以同时含$\sigma$
和$T_\ell\sigma$，从而包含$f_B\gamma$。最小性给
$R_{f\gamma}^{\rm var}\subset R_\sigma^{\rm var}$，应用命题\ref{unified:prop:lower}。
第二项中，交$\mathcal A\cap R_\sigma^{\rm var}$的局部系统非零，
不可约性使其等于$\mathcal A$；而
$\mathcal A\cap I_\lambda^{\rm var}$要么为零，要么全部，
后者被非零读出排除。
\end{proof}
这消除了族中链包含的一个实质待证条件，但没有消除实际管比较和交数的责任。
若变分不是单一管周期而是多个割之和，应先对总割求最小对象，
不能把各项分别当作已被保护。

\begin{proposition}[无需完整积分的微分阶数下界]\label{unified:prop:order}
令$F$为物理周期函数，并设其某个变分等于非零割函数$c\,g$。
若$g$的全部解析延拓所张成的解空间$W$不可约且维数为$r$，
则任意以单值亚纯函数为系数、满足$LF=0$的标量齐次算子都有
$\operatorname{ord}L\ge r$。
若$W$由$K[\partial_s]$中的首一$r$阶算子$L_{\rm cut}$给出，
则每个$L\in K[\partial_s]$满足
\[
 L=A\,L_{\rm cut}.
\]
\end{proposition}
\begin{proof}
解析延拓与单值系数算子交换，故$L$杀死$F$的变分和它的所有延拓。
于是它杀死整个$r$维$W$；普通点处$m$阶方程的局部解空间至多$m$维。
对$L$用$L_{\rm cut}$作右除，余项阶数小于$r$且仍杀死$W$，只能为零。
\end{proof}
不可约割局部系统与非零代数读出通常提供这里的$r$：
其平坦向量到周期函数的核被单值保持，不可约性迫使核为零。
如扩大到代数系数，须在相应有限覆盖上仍有不可约性；
连通几何单值群已不可约时满足这一点。

\textbf{固定纤维边界。}
$T_\ell$一般不是$(\mathcal E_{s_*},F^\bullet,W_\bullet)$的Hodge自同态。
所以本定理的$M_{\rm var}$不能替换成任意指定点的$M$。
例如既有banana固定纤维$\operatorname{End}_{\rm Hdg}T=\Q$，
非标量节点反射就不是固定Hodge态射。
指定$s=-64$的秩八最小性仍用原regulator与群证明，
不从本条微分下界重证。
更一般地，某一纤维数值周期偶然为零，也不抹去非零族函数的解空间。

\subsubsection{一个任意圈数的实际物理入口：二维banana正常阈值}
设$h\ge2$、$n=h-1$、$m_i>0$，$\mu_i=m_i^2$。
采用明确的正参数归一化
\[
 I_h(s)=\int_{\alpha_i>0,\ \alpha_h=1}
 \frac{d\alpha_1\cdots d\alpha_n}
 {(\sum_i\mu_i\alpha_i)U_h-s\prod_i\alpha_i}.
\]
这里不含动量测度的外部$\pi,i$因子；
$h=4$时正取向与本文欧氏参数周期相同。
在图表中置$y_i=\log(m_i\alpha_i/m_h)$、$y_h=0$，则
\[
 I_h(s)=\int_{\mathbb R^n}\frac{dy_1\cdots dy_n}{f(y)-s},
 \qquad
 f(y)=\Big(\sum_{i=1}^h m_i e^{y_i}\Big)
       \Big(\sum_{i=1}^h m_i e^{-y_i}\Big).
\]
令$M=\sum_i m_i$、$s_0=M^2$。

\begin{theorem}[正物理链的阈值非消失]\label{unified:thm:threshold}
$I_h$在$\mathbb C\setminus[s_0,\infty)$上全纯。
$f$严格凸、恰当，唯一极小点为$y=0$，极小值为$s_0$。
对$t>s_0$，正向边界差为
\[
 I_h(t+i0)-I_h(t-i0)=2\pi i\,\rho_h(t),\qquad
 \rho_h(t)=\int_{\mathbb R^n}\delta(f(y)-t)\,dy>0.
\]
在$\varepsilon\downarrow0$时，
\[
 \rho_h(s_0+\varepsilon)=
 \frac{\pi^{n/2}}
 {\Gamma(n/2)\,M^{(h-2)/2}\sqrt{\prod_i m_i}}\,
 \varepsilon^{(h-3)/2}\bigl(1+O(\varepsilon)\bigr).
 \tag{T}
\]
局部链变分为正常Morse消失球的单次管，取向由上式固定，
交数绝对值为一；尤其物理分支确实激发该割。
\end{theorem}
\begin{proof}
展开为
\[
 f=\sum_i m_i^2+2\sum_{i<j}m_i m_j\cosh(y_i-y_j).
\]
沿任意非零方向$v$（$v_h=0$），Hessian二次型为
$2\sum_{i<j}m_i m_j\cosh(y_i-y_j)(v_i-v_j)^2>0$。
全体差同时为零只在$y=0$，由Cauchy--Schwarz得到极小值。
若$R=\max_i y_i-\min_i y_i$，则
$f\ge m_{\min}^2e^R$且$R\ge\|y\|_\infty$；
故$f$恰当，积分在无穷远有指数可积控制。
对避开割线的任意紧$s$集合，分母有统一相对下界；
可在积分下任意求$s$导数，得全纯性。

严格凸性使每个$t>s_0$的紧水平面均无临界点。
用余面积公式推前正测度，再对$1/(u-t\mp i0)$作一维Cauchy边界值，
得到跳跃式。水平面非空且其密度处处正，故$\rho_h(t)>0$。
这是本参数积分的直接证明，不调用任意维多圈Cutkosky作为前提。

在极小点，令$m=(m_1,\ldots,m_n)^t$，有
\[
 f-s_0=y^tAy+O(\|y\|^4),\quad
 A=M\,\operatorname{diag}(m_1,\ldots,m_n)-mm^t,
 \quad \det A=M^{h-2}\prod_{i=1}^h m_i.
\]
矩阵行列式引理给最后等式。
实解析Morse坐标把$f-s_0$化为$\|u\|^2$，
体积密度在零点为$(\det A)^{-1/2}$。
球面平均消去所有奇次项，故首项为
\[
 (\det A)^{-1/2}
 \int_{\mathbb R^n}\delta(\|u\|^2-\varepsilon)\,du
 =\frac{\pi^{n/2}\varepsilon^{n/2-1}}
 {\Gamma(n/2)\sqrt{\det A}},
\]
下一相对阶为$O(\varepsilon)$，包括$n=1$情形。

最后，取极小点的小球；球外$f-s_0$有正下界，
外部积分在$s_0$邻域全纯，不能贡献这次局部变分。
球内实链与Morse消失胞的对偶局部胞仅横截相交一次；
亦可在径向变量中绕过唯一简单极点，直接得到一次Cauchy管。
其留数正密度恰为上式，因此倍数不是未定整数。
经典局部管与Picard--Lefschetz解释可比较\cite[\S2.3]{muhlbauer}。
\end{proof}
由正$y$图表对应未来向在壳支，以及既有光锥/参数割簇同构，
此处给出了正常阈值处的物理激发。
它没有计算全体异常阈值或所有全局轮廓，也未补完全部动量空间测度比较。

\begin{corollary}[从阈值到不可约几何核]\label{unified:cor:banana-threshold}
若此正常消失球的紧割周期属于一个已证明不可约的$r$维纯局部系统，
则$I_h$的标量齐次微分方程阶数至少为$r$，并含对应割算子作为右因子。
已构造实际管变分时，同时得到定理\ref{unified:thm:variation}的对象保护。
\end{corollary}
\begin{proof}
式(T)保证割函数非零，不可约性使其延拓生成全纯核，
再用命题\ref{unified:prop:order}与定理\ref{unified:thm:variation}。
\end{proof}
$h=3$在不可约椭圆变分上给$r=2$；
$h=4$、$m=(1,2,4,8)$射线利用已证纯单值群给$r=6$。
两者结论是物理函数的下界；完整混合对象和齐次阶数不因此分别等于$2,6$。
$h>4$的同一阈值公式成立，但不能凭阈值幂次就宣布某个CY核不可约。


\subsection{结构判据的适用范围}
下表分开全积分、最大割和有限系数展开，不将它们的秩混用。
\begin{center}\small
\begin{longtable}{>{\raggedright\arraybackslash}p{29mm}>{\raggedright\arraybackslash}p{51mm}>{\raggedright\arraybackslash}p{61mm}}
\toprule 输入类&已有结论&当前边界\\\midrule\endfirsthead
\toprule 输入类&已有结论&当前边界\\\midrule\endhead
泛独立能量根的单位标量图&
dual-critical充要支撑、边束奇偶压缩&
字符支撑不是动机直和\\
非退化圈星图&
实际二次覆盖、字符宿主、标量全字符与分子规则&
紧割已补实际三层边界；完整图到曲线的比较、链和等源特化另判\\
正质量一圈三角/盒图&
完整混合Artin--Tate载体；一般位置秩7/8及二次角色表&
完整参数链和代数源已接入；载体秩不是未经最小化的物理秩\\
接受的Euler/MZV图规则&
原v2的有限Euler展开与MZV形式余模&
解析图恒等式不自动给高维原图Nori最小对象同构\\
二维四传播子banana&
实际相对对象、K3、质量Kummer、regulator；指定点最小秩八、coaction及群&
指定点、很一般四参数与等质量层分别陈述\\
任意$h$的二维正质量banana&
正常阈值唯一夹挤、交数绝对值一、非零割首项；不可约核给物理微分下界&
较高维纯核的不可约性另证；不覆盖所有异常阈值\\
留数零的实际相对曲线周期&
端点给代数非齐次源；割保护下最小齐次阶数为$r$或$r+1$&
原函数/移动端点及实际比较不可省略；固定点动机秩另判\\
正质量对数图及相容MOM&
实际森林与端点；本节新结构的方案保持&
一般幂次、IR和任意方案不包含\\
四维sunset及两圈$\phi^4$&
实际二点载体及同一方案有限阶RG出口&
作为物理实例，不继续以低阶系数扩展代替总分类\\
单耦合MS型RG&
条件内局部系数与临界展开可逆及结构保持&
保留方案和热方向条件，MOM质量不机械代入\\
较高维CY和一般混合图&
几何上界、割下界及扩张检查的统一规则&
尚无覆盖任意边表的完整判定\\
串联细分sunrise，二维正质量、零内部动量注入 &
完整周期实际下降到端点sunrise直和；各端点至多一个椭圆组成 &
质量互异且正常阈值分离时，每个未被分子消去的椭圆组成恰一次；同质量合流时至多一个。\otherref{feynman_periods_supplement}{补充材料S1.1}。\\
四边两圈三点图，允许内部非零注入 &
有理极点曲面与两个椭圆端面；完整相对载体的非AT半单秩四 &
单标量读出恰含两椭圆核；收缩分子各选一项。显式切片齐次阶数至少四；绝对消去可有椭圆源。第\ref{unified:sec:boundaryelliptic}--\ref{unified:sec:scalarboundary}节。\\
物理图求和与阈值背景&
整个变分与半单分级的秩相等；实际森林载体跨过选定退化分量延拓。&
续接跳跃秩给函数核重数；背景外重数受方案比较保护。三图交叉组合保留三椭圆通道，非AT秩6--12。\otherref{feynman_periods_supplement}{补充材料S1.2}。\\
零内部注入的三路径theta细分族&
完整相对Stokes下降；每个质量三元组给一个椭圆载体上界。&
阈值分离下，分子网格非零项各给一次椭圆核；秩$2r_1r_2r_3$可无界。有限骨架界与适用条件见\otherref{feynman_periods_supplement}{补充材料S1.3}。\\
\bottomrule
\end{longtable}
\end{center}
不同入口须经过实际比较才能拼接；不能以环境对象取代最小物理对象。
表内依据是已保存内部证明及原始材料，不将其称为外部同行评议结论。


\subsection{留数像、固定环境与非齐次源}
\label{unified:sec:pencil}
\subsubsection{固定环境是必须提交的几何输入}
第\ref{ch:graph}节的曲线规则要求完整周期已被识别为相对曲线周期。
一般图尚未完成这一步时，可以直接在其高维参数模型上工作。
设$k\subset\mathbb C$为特征零域，$X/k$光滑，
$B\subset X$为固定的简单正常交叉边界。
固定极点集合已从$X$中删去。
令$Y_s\subset X$为代数参数曲线$S$上的光滑除子族，
与$B$每个交叉层横截。
必要时缩小$S$，使所用相对上同调成为带Gauss--Manin连接的局部自由对象。
相对形式和物理链都必须实际属于下列模型：
\[
 J=H^n(X,B),\qquad
 M_s=H^n(X\setminus Y_s,B\setminus Y_s),\qquad
 C_s=H^{n-1}(Y_s,B\cap Y_s)(-1).
\]
这里$M_s$用上同调方向，前文最小物理对象承载于$E_s=M_s^\vee$。
$J$是固定几何的对象；这不表示$J$自动为Tate。
$\omega_s\in(M_s)_{\dR}$是代数相对类，
$\sigma_s\in(M_s)_B^\vee$在一个单连通参数开集上平坦，
$I(s)=\langle\sigma_s,\omega_s\rangle$。
物理边界的移动项已包含于相对Gauss--Manin连接。

\begin{theorem}[留数像与环境背景的精确分离]\label{unified:thm:pencil-separation}
在上述实际模型中有带连接的正合段
\begin{equation}\label{unified:eq:pencil-gysin}
 J\longrightarrow M_s\xrightarrow{\mathrm{res}}C_s
       \longrightarrow H^{n+1}(X,B).
\end{equation}
令$K_s=\ker(\mathrm{res})$、$A_s=\im(\mathrm{res})$。
则$K_s$是$J$的商，$A_s$是$C_s$的子对象，且
\[
 0\longrightarrow K_s\longrightarrow M_s\longrightarrow A_s
 \longrightarrow0.
\]
在每个指定代数纤维，取含$J$全部简单组成的Serre背景$\cS$，有
\[
 \pi M_s\simeq\pi A_s.
\]
因此完整模型在该背景以外的简单组成及其重数，恰等于实际留数像的组成及重数。
物理最小对象在对偶Serre背景$\cS^\vee$以外的重数至多为$A_s^\vee$中的重数；
若实际割和物理标架另给达到此上界的下界，则这些重数被精确确定。
\end{theorem}
\begin{proof}
对$Y_s\subset X$写局部化三角，余维一纯性给
$R\Gamma(Y_s)(-1)[-2]\to R\Gamma(X)\to
R\Gamma(X\setminus Y_s)\to$。
在$B$的每个交叉层上写相同三角。
横截基变换使它们与限制相容，取有限交叉复形及相对映射纤维得到
\[
 R\Gamma(Y_s,B\cap Y_s)(-1)[-2]\longrightarrow
 R\Gamma(X,B)\longrightarrow
 R\Gamma(X\setminus Y_s,B\setminus Y_s)\longrightarrow.
\]
取第$n$上同调即得\eqref{unified:eq:pencil-gysin}；
实际六运算及其Nori实现给逐纤维态射，Gauss--Manin自然性给族连接相容。
这里不额外假设已构造一般“Nori变分类别”。
正合性给$K_s$为$J$的商；Serre商消去$K_s$，得到同构。
对偶后$E_s$是$A_s^\vee$与$K_s^\vee$的扩张，
而最小物理对象是$E_s$的子商，故在$\cS^\vee$外逐简单因子重数不超过相应上界。
达到上界时再用有限长度性。
\end{proof}

这将两个容易混淆的判断分开：$C_s$可能比留数像$A_s$大；
即使$A_s$完全已知，物理标架也可能只看到其一部分。
若$J$是mixed Tate，所有非Tate简单组成都来自\emph{相对}留数对象$C_s$。
但$C_s$自身仍含边界，它不自动等于紧致纯割核。
族中的物理变分可给连接/变分层面的下界，
不能不经比较改称指定纤维的Nori下界。

\subsubsection{非齐次源不需通过完整积分猜出}
取$K=k(S)$，配一个非常数局部参数的导子$\partial$，
假设常数域为$k$。
在$J_{\dR}$中固定$k$基$e_1,\ldots,e_d$，$d=\dim J_{\dR}$。

\begin{theorem}[实际留数湮灭后的固定环境源]\label{unified:thm:ambient-source}
设$L\in K[\partial]$满足
$L\,\mathrm{res}(\omega)=0$，这里是实际de Rham类的零。
则存在$a_j\in K$，使
\begin{equation}\label{unified:eq:ambient-source}
 L I(s)=\sum_{j=1}^d a_j(s)P_j,\qquad
 P_j=\langle\tau,e_j\rangle,
\end{equation}
其中$\tau$为$\sigma_s$在固定对$(X,B)$中的像，与$s$无关。
每个$P_j$是固定环境$J$的周期。
等式在相应实际形式周期中同样成立，
其中求导由Gauss--Manin连接定义，纤维在$k$或明确的扩域上取值。
特别地，这个结论不使用数值周期映射的单射性。
\end{theorem}
\begin{proof}
留数与连接交换，故$L\omega\in(K_s)_{\dR}$。
由\eqref{unified:eq:pencil-gysin}，泛点处的线性满射
$J_{\dR}\otimes_k K\to(K_s)_{\dR}$可为此向量选一个$K$提升
$\theta=\sum a_je_j$。
$\sigma_s$映入固定对的相对同调，映射水平，故其像$\tau$平坦且为常数。
因此
\[
 LI=\langle\sigma_s,L\omega_s\rangle
   =\langle\tau,\theta(s)\rangle=\sum_j a_j(s)P_j.
\]
提升不唯一，但两个提升之差在映射核中，与$\tau$配对为零。
以上每次等号分别来自连接配对、实际态射函子性及线性关系，
所以也给形式周期等式，而非仅给解析数值相等。
\end{proof}

源的求取现在是有限de Rham约化与线性提升问题，
不要求先积分$I$。
不过，本定理本身没有提供适用于任意图的终止约化算法，
也不把$P_j$自动认作代数数、MZV、低圈周期或局部反项。
哪种固定结构出现，仍由$(X,B)$决定。

\begin{corollary}[不求固定周期值的齐次阶数上界]\label{unified:cor:ambient-order}
设$r=\operatorname{ord}L$，
$q=\dim_k\operatorname{span}_k\{a_1,\ldots,a_d\}$。
存在$L_0\in K[\partial]$、$\operatorname{ord}L_0=q$，
使$L_0LI=0$。故最小齐次阶数满足
\[
 \operatorname{ord}_{\min}(I)\le r+q\le r+d.
\]
若源为零，可直接使用上界$r$。
\end{corollary}
\begin{proof}
$q>0$时选$a_j$张成空间的一组$k$基$b_1,\ldots,b_q$。
微分域的常数恰为$k$，故其Wronskian非零。
用Wronskian的最后一列放入待作用函数，除以最高项系数，
得系数在$K$中的首一$q$阶算子$L_0$，杀死所有$b_i$及$a_j$。
$P_j$对$\partial$为常数，因而$L_0$杀死\eqref{unified:eq:ambient-source}的右端。
$q=0$时取$L_0=1$。
\end{proof}

这是上界，不是最小秩公式，也没有假设各$P_j$独立。
不同提升可改善$q$。
若源是非零$a(s)\in K$，则一阶算子$\partial-a'/a$足够；
若源为$P_1+sP_2$，却未证明$P_1/P_2\in k$，
就不能声称其对数导数在$k(s)$中。
$\partial^2$仍给无需这些常数关系的合法湮灭算子。
扩大常数域到$\mathbb C$与在原算术域$k$上求方程，是不同问题。

\begin{corollary}[有物理保护时的精确两支]\label{unified:cor:ambient-exact}
若$L$同时是被非零物理变分激发的$r$维不可约割的最小算子，
且\eqref{unified:eq:ambient-source}的源属于$K$，
则完整函数的最小齐次阶数为$r$（源为零）
或$r+1$（源非零）。
\end{corollary}
\begin{proof}
前文物理变分定理使任何湮灭$I$的算子都以$L$为右因子。
零源给$r$阶上界。
非零源时，$r$阶算子只能为非零函数乘$L$，不可能杀死$I$；
另一方面$(\partial-a'/a)L$给$r+1$阶上界。
\end{proof}

\subsubsection{完整相对留数与纯核方程之间的差别}
一个纯核算子$L_{\rm pure}$可能只在$C_s$的某个纯商中杀死留数像。
此时$L_{\rm pure}\mathrm{res}(\omega)$仍可在$C_s$的边界中非零，
定理\ref{unified:thm:ambient-source}不能直接用于$L_{\rm pure}$。
需要先为这个实际边界类找算子$L_{\rm bdry}$，
然后对$L=L_{\rm bdry}L_{\rm pure}$应用定理。
这也解释为何一些完整积分在消去最大割后，仍出现随参数变化的下sector源。

同理，多重动量最大割与参数超曲面的一次留数不是定义上相同。
只有给出实际变量变换、管与相对链比较，才能沿用同一个$C_s$。
本文banana有专门比较；一般圈星仍须提交其完整sector的实际模型，
不能把紧曲线宿主及其亏格直接当作完整积分的结构表。

\subsubsection{图多项式铅笔的接受条件}
在一参数切片上，若第二图多项式写为$F_s=A-sB_0$，
可先研究有理映射$A/B_0$及固定边界的消解。
特征零下可在固定空间上消解其不定性，并在稠密开集删去坏纤维，
为上述固定环境模型提供几何候选。
但这一步不自动证明原物理链避开新极点，也不消去例外面的周期。

接受一个图或sector需实际提交四项：
固定$(X,B)$与所有被删除的固定极点；移动除子的横截性及坏参数集合；
原形式的相对de Rham代表；原物理链、边界及其随参数的比较。
完成后，定理\ref{unified:thm:pencil-separation}给结构上界，
定理\ref{unified:thm:ambient-source}给源，
物理管变分给下界，锚定森林定理追踪方案变化。
这是可直接拼接的推导接口。
仅知道$X$有理或具有环面坐标，不能证明删去固定图超曲面后的$J$为Tate。

\subsubsection{banana：从正割密度恢复完整函数}
沿用前文正质量二维$h$传播子归一化，
$s_0=(\sum_i m_i)^2$、$I_h(s)=\int_{\mathbb R^{h-1}}(f(y)-s)^{-1}dy$。

\begin{proposition}[全物理谱的无减除色散表示]\label{unified:prop:dispersion}
对$s\in\mathbb C\setminus[s_0,\infty)$，有
\[
 I_h(s)=\int_{s_0}^{\infty}\frac{\rho_h(t)}{t-s}\,dt,
 \qquad
 \rho_h(t)=\int_{\mathbb R^{h-1}}\delta(f(y)-t)\,dy .
\]
因此这个完整参数函数由其正物理谱密度确定，无需额外减除常数。
\end{proposition}
\begin{proof}
先取实$s<s_0$。正性允许Tonelli，恰当函数$f$的余面积公式把
$\mathbb R^{h-1}$按水平集分解，直接给等式。
前文指数尾界保证左端收敛，故右端也收敛。
对复平面的任意紧集$K_0\subset\mathbb C\setminus[s_0,\infty)$，
固定实$s_*<s_0$，
$(t-s_*)/|t-s|$在$t\ge s_0,s\in K_0$上一致有界；
右端由已收敛的$\rho_h(t)/(t-s_*)$控制，
从而全纯。左端也全纯，由恒等定理得到结论。
\end{proof}
这里积分遍历整个参数区间上的不同割纤维。
函数重建不表示完整相对动机与单个纯割纤维同构；
沿参数积分正会保留混合扩张。
证明依赖此二维正质量收敛模型，不覆盖四维发散图的减除多项式。
也不以这一解析恒等式代替一般动量与参数形式周期的实际比较；
该问题的范围在Brown的讨论中已被明确区分\cite[\S10.2]{browncosmic}。

\subsubsection{一个已知完整例子在统一规则中的位置}
等质量$m_i=1$的四传播子banana，BKV给出\cite[Theorem2.2.1]{bkv}
\begin{align*}
 L_3={}&s^2(s-4)(s-16)\partial^3
       +6s(s^2-15s+32)\partial^2\\
      &+(7s^2-68s+64)\partial+(s-4),
 &L_3I_4=-24 .
\end{align*}
这里算子和常数源均为已有结果，不能计作本文首次求得。
该文的模参数描述将纯解空间识别为椭圆秩二表示的对称平方；
在有限模覆盖上，单值包含有限指标模群。
两个相反的非平凡幺幂一参数闭包生成$SL_2$，
其对称平方三维表示不可约，故纯连接绝对不可约。
前文正常阈值证明使$s=16$的割非零；
因此右因子保护和非零常数源给
\[
  \operatorname{ord}_{\mathbb C(s),\min}(I_4)=4,
  \qquad \partial L_3 I_4=0 .
\]
三阶以内的齐次方程不可能湮灭完整物理函数。
这一例子把纯K3核、物理激发、非齐次源和完整函数阶数接在一起。
它不意味着不等质量族只有四阶，也不据函数阶数认定每个代数纤维的动机秩。
本文指定不等质量秩八的结论仍由其独立相对对象及扩张证明承担。


\subsection{完整一圈图的接受域：二次曲面、算术角色与物理链}
\label{unified:sec:oneloop}
本节使前述规则在banana之外获得完整积分的实际输入。
一圈图的Tate性质与四维对数、双对数表达是经典结果，
尤其三角图的Hodge维数已见Bloch--Kreimer\cite[\S7--10]{bkone}；
偶数维多边形的二次曲面递归见Spradlin--Volovich\cite{svone}。
下面重新从实际参数对推导，保留原基域上的二次Artin角色，
并接入本节的物理最小化、操作规则与失败边界。
这些整合不作未经检索的原创性声明。

\subsubsection{从物理积分到完整相对对象}
取特征零数域$k\subset\mathbb R$，正质量平方$\mu_i$及欧氏不变量
$\kappa_{ij}=(r_i-r_j)^2$属于$k$。定义
\[
 U=\sum_{i=1}^N\alpha_i,\quad
 F=U\sum_i\mu_i\alpha_i+\sum_{i<j}\kappa_{ij}\alpha_i\alpha_j
   =\alpha^t C\alpha,
\]
其中$C_{ii}=\mu_i$、$2C_{ij}=\mu_i+\mu_j+\kappa_{ij}$。
Minkowski延拓使用$\kappa_{ij}=-s_{ij}$，不另外改变$F$的符号。
设$B$为坐标超平面之并，$Q=V(F)$，$H=V(U)$。

四维单位传播子的正欧氏归一化为
\begin{equation}\label{unified:eq:onephysical}
 I_N=\int_{\mathbb R^4}\frac{d^4\ell}{\pi^2}
           \prod_i\frac1{(\ell+r_i)^2+\mu_i}
     =\Gamma(N-2)\int_{\sigma}
        \frac{U^{N-4}\Omega_{N-1}}{F^{N-2}},\qquad N\ge3.
\end{equation}
$\sigma$为正射影单纯形；取$U=1$时$\Omega_{N-1}$的取向使其成为正体积形式。
特别地，三角图与盒图分别由
\begin{align}\label{unified:eq:onepairs}
 M_3&=H^2\bigl(\mathbb P^2\setminus(Q\cup H),
             B\setminus(Q\cup H)\bigr),\notag\\
 M_4&=H^3\bigl(\mathbb P^3\setminus Q,B\setminus Q\bigr)
\end{align}
承载。它们是完整参数对象，不是只对最大割取的紧宿主。

\begin{proposition}[实际链比较与一般一圈上界]\label{unified:prop:onecarrier}
在正质量欧氏点，式\eqref{unified:eq:onephysical}是收敛积分，
形式与正单纯形给式\eqref{unified:eq:onepairs}中的实际Nori标架。
更一般地，偶数维$D=2d$、整数幂$a_i\ge1$、$A=\sum a_i>d$
的收敛一圈标量积分，由
\[
 H^{N-1}\bigl(\mathbb P^{N-1}\setminus(Q\cup H),
                    B\setminus(Q\cup H)\bigr)
\]
承载；若形式在$H$上无极点可不删除$H$。
这个对象经某个有限代数扩域后为混合Tate，原域上为混合Artin--Tate。
同一极点集合内的分子读出及参数导数不产生新的非Artin--Tate简单组成。
\end{proposition}
\begin{proof}
欧氏Schwinger公式的被积函数非负，故Tonelli允许先完成平方并积分$\ell$。
高斯积分给$U^{-d}\exp(-F/U)$，再写$\alpha_i=t x_i$、$\sum x_i=1$，
尺度积分给$\Gamma(A-d)$和射影形式
\[
 \frac{\Gamma(A-d)}{\prod_i\Gamma(a_i)}
 \frac{\prod_i\alpha_i^{a_i-1}\,U^{A-2d}\Omega_{N-1}}{F^{A-d}}.
\]
在闭单纯形上$F\ge\min_i\mu_i>0$，$U=1$；
单纯形与其全部面避开极点。尺度零端由$A>d$控制，
无穷端由正质量控制，因而没有尚待认证的边界正则化。
同一次高斯/尺度推导给出完整物理积分等于该参数周期。
这里构造参数Nori载体，不宣称动量对象与它已经有全局动机同构。

每个坐标交叉层是射影空间，$Q$在其上仍为二次式，$H$仍为超平面。
选择共同有限扩域使这些有限多个二次型全部分裂。
光滑分裂射影二次曲面有仿射胞腔分解；
退化二次曲面去掉其线性顶点后，是非退化低维二次曲面上的仿射空间丛。
因此二次曲面及其超平面截面均由Tate对象构成。
对$Q\cup H$的闭并交叉列、开补空间局部化以及$B$各交叉层的有限复形，
依次应用正合性，便得相对对象在扩域后混合Tate。
原域上取有限Galois扩域的限制与迹；除以扩域次数给原对象的直和投影。
故其组成是有限Artin表示的Tate扭转，而非仅在Hodge实现上对角。

有理形式的整数幂提高只改变读出；在常秩参数开集上，
Gauss--Manin求导仍在同一相对连接内。
实际子商不会引入环境中没有的简单因子，得到最后结论。
\end{proof}

有限扩域的选择是上界，不是最小根式域。
分裂后的混合Tate性也不等于MZV性：质量比、运动学支点、分支与扩张类仍在。
本命题没有覆盖维数正则化的全部$\varepsilon$系数、任意UV分子或质量零点的IR极限。

\subsubsection{三角图与盒图的完整算术分级}
为得到精确载体秩，以下取一般位置：$Q$与各坐标层横截，
三角图的$Q$还与$H$及其坐标交叉横截；顶点质量非零。
记$\chi_d$为平方类$d\in k^\times/k^{\times2}$对应的一维二次Artin角色，
$d$为平方时它就是单位角色。对每条坐标边$ij$置
\begin{equation}\label{unified:eq:edgedisc}
 d_{ij}=(\mu_i+\mu_j+\kappa_{ij})^2-4\mu_i\mu_j .
\end{equation}
它是$Q$与该边的二点交的判别式。
三角图另置
\[
 d_\infty=(\kappa_{12}-\kappa_{13}-\kappa_{23})^2
                          -4\kappa_{13}\kappa_{23}.
\]
因为$F|_H=\sum_{i<j}\kappa_{ij}\alpha_i\alpha_j$，
这是$Q\cap H$的二点判别式；等价于$-\det(F|_H)$的平方类。
盒图令$d_Q=\det C$，即四元二次型的有符号判别式。

\begin{theorem}[可直接读出的结构表]\label{unified:thm:onegraded}
上述一般位置中，完整相对对象的权重分级为
\begin{align}\label{unified:eq:trianglegraded}
 \operatorname{Gr}^W_0M_3&=\Q(0),\notag\\
 \operatorname{Gr}^W_2M_3&=\Q(-1)^2\oplus
                 \bigoplus_{1\le i<j\le3}\chi_{d_{ij}}(-1),\notag\\
 \operatorname{Gr}^W_4M_3&=\chi_{d_\infty}(-2),
 \qquad \rk M_3=7;
\end{align}
以及
\begin{align}\label{unified:eq:boxgraded}
 \operatorname{Gr}^W_0M_4&=\Q(0),\notag\\
 \operatorname{Gr}^W_2M_4&=\bigoplus_{1\le i<j\le4}\chi_{d_{ij}}(-1),\notag\\
 \operatorname{Gr}^W_4M_4&=\chi_{d_Q}(-2),
 \qquad \rk M_4=8.
\end{align}
其它权重为零。式中7和8是所列\emph{完整载体}的秩；
尚未证明某组特化标架生成全对象时，它们只是物理最小秩的上界。
\end{theorem}
\begin{proof}
先处理三角图。令$V=\mathbb P^2\setminus(Q\cup H)=\mathbb A^2\setminus Q^\circ$。
光滑射影圆锥曲线$Q$的有理上同调为$\Q(0)\oplus\Q(-1)$，
即使没有$k$有理点亦如此：有限点的零循环可除以其次数。
从$Q$删去$Q\cap H$的二点，留数正合列给
\[
 H^0(Q^\circ)=\Q(0),\qquad H^1(Q^\circ)=\chi_{d_\infty}(-1).
\]
仿射平面的Gysin列因而给
$H^1(V)=\Q(-1)$、$H^2(V)=\chi_{d_\infty}(-2)$。

每个边界分量$B_i^\circ$是仿射直线删去二点，故
\[
 H^1(B_i^\circ)=H^0(Q\cap B_i)(-1)
               =\Q(-1)\oplus\chi_{d_{jk}}(-1).
\]
三分量在三个未删除顶点组成一个三角环。
其规范化交叉复形给$H^1(B^\circ)$的权重零部分$\Q(0)$，
以及权重二部分$\Q(-1)^3\oplus\bigoplus\chi_{d_{ij}}(-1)$；
$H^2(B^\circ)=0$。
限制映射$H^1(V)\to H^1(B^\circ)$由$d\log F$表示，
在每一边的两个根上留数均为$+1$，故在三个平凡角色中是非零对角向量。
它是单射，恰消去一个$\Q(-1)$。相对正合列
\[
 H^1(V)\longrightarrow H^1(B^\circ)\longrightarrow M_3
            \longrightarrow H^2(V)\longrightarrow0
\]
和权重严格性给\eqref{unified:eq:trianglegraded}。

再处理盒图。$V=\mathbb P^3\setminus Q$，$Q$为光滑二次曲面。
扩域后两族直线给$H^2(Q)$的两个生成元；
两族之和为超平面类，差由$\chi_{d_Q}$作用。
Gysin的推前在和上为非零、在差上为零。
所以$H^1(V)=H^2(V)=0$，$H^3(V)=\chi_{d_Q}(-2)$。

边界的四个面各为$\mathbb P^2$删去光滑圆锥曲线，
有理上同调只有$H^0=\Q$；六条交边为$\mathbb P^1$删去二点，
其$H^1=\chi_{d_{ij}}(-1)$；四个三重交为未删除顶点。
交叉复形的$H^0$行是四面体边界的复形，给$H^2$中的一个$\Q(0)$。
六个交边$H^1$位于总次数二，给六个权重二角色。
没有其它项能进出这些位置，故$H^3(B\setminus Q)=0$，且
\[
 \operatorname{Gr}^W H^2(B\setminus Q)
   =\Q(0)\oplus\bigoplus_{i<j}\chi_{d_{ij}}(-1).
\]
相对列$0\to H^2(B\setminus Q)\to M_4\to H^3(V)\to0$
给\eqref{unified:eq:boxgraded}。所有映射均由实际几何诱导，
不是仅由Betti维数猜出的分解；只断言权重分级，不将扩张强行拆开。
\end{proof}

这个结果无需积分，只需二次型与有限个限制判别式。
一般位置失效时，不能保留7、8照抄；须重新计算实际交叉层。
判别式为零首先表示此模型发生退化，并不自动证明所选物理分支发散。
三角图的7维分裂Hodge结果归于\cite[Proposition10.1]{bkone}；
这里显式写出算术下降和与最小物理对象的区别。

\subsubsection{物理出口、操作闭合与源的范围}
\begin{corollary}[已接受图类中的物理上界]\label{unified:cor:oneoperations}
在定理\ref{unified:thm:onegraded}的常秩参数域内，完整三角图和盒图的物理最小对象
均为混合Artin--Tate；没有椭圆$H^1$、K3超越核或其它非Artin--Tate简单因子。
沿代数参数曲线，其标量齐次微分阶数分别至多7、8。
分裂上述有限Artin角色后，连接具有长度至多三的平凡分级过滤，
几何幺幂作用的连续三次严格降过滤乘积为零。

几何边界、留数/管及同一环境中的分子读出、参数导数保持此组成上界。
若有限方案改变量另有混合Artin--Tate载体，则改变方案后的最小对象仍属此类。
这些结论既不保证完整载体都物理可见，也不保证全部扩张类不随方案变化。
\end{corollary}
\begin{proof}
物理最小对象是$M_N^\vee$的子商，给第一结论。
一个秩$r$连接中的$r+1$个连续导数线性相关，给阶数上界；
平凡分级过滤来自式\eqref{unified:eq:trianglegraded}--\eqref{unified:eq:boxgraded}及有限根式覆盖。
三个严格降阶算子的乘积越过全部权重。
几何操作使用实际正合列和Tate扭转，
读出变化用子商；有限方案用直和载体及本节加性比较定理。
\end{proof}

\textbf{必须区分的收缩接口。}
$\alpha_e=0$确实给收缩图的二次多项式$F|_{\alpha_e=0}$，
所以参数面上的几何仍为二次型。
但盒图的面载体是$H^2(\mathbb P^2\setminus Q,\ldots)$，
真正四维三角图还需要删除$H$，形式从$\Omega/F^2$改为$\Omega/(UF)$。
后者由前者的限制以及$H$上的局部化补项连接；
该补项也是Artin--Tate，但二者不是未经比较就相同的物理对象。
同样，三角图缩成四维bubble时出现整体UV减除，
须使用实际减除载体，不能把一个边界周期直接当作未经减除振幅。
因此这里证明的是操作的\emph{类别上界闭合}，而非割、收缩和积分的无条件互换。

\textbf{与完整参数源的连接。}
三角图可在$X=\mathbb P^2\setminus H\simeq\mathbb A^2$上应用前节，
其固定边界为三条仿射直线，构成三角环。
它们自身及交点无高上同调，故
\[
 J_3=H^2(X,B\cap X)=\Q(0).
\]
盒图取$X=\mathbb P^3$及四个坐标面，其交叉复形给
$H^2(B)=\Q(-1)\oplus\Q(0)$，
$H^2(X)=\Q(-1)$同构到前一项，因而
\[
 J_4=H^3(X,B)=\Q(0).
\]
两个固定环境的正单纯形比较值均按取向为$1$。
故在这两个\emph{完整}模型中，任何消去实际\emph{相对留数类}的算子$L$
都给代数源$LI_N=a(s)\in k(s)$（在所取参数曲线上）。
这兑现了前节的代数源条件，而不是把它额外假设为“所有图都成立”。
只消去顶权重的秩一角色通常不够：还剩边界中的对数扩张。

\textbf{三角图的函数形式可以直接解释。}
在$U=1$图表写$\alpha_1=x,\alpha_2=y,\alpha_3=1-x-y$，则
\[
 I_3=\int_0^1dx\int_0^{1-x}\frac{dy}{a y^2+b(x)y+c(x)},
 \quad r^2=b(x)^2-4ac(x),\quad\deg_x(r^2)\le2.
\]
一般$a\ne0$时对$y$的原函数为
\[
 \frac1r\log\frac{2ay+b-r}{2ay+b+r}.
\]
代入两个端点，只剩$\int(dx/r)\log R(x,r)$。
对二次根曲线作有限代数扩域后的有理参数化，
它成为有理一形式乘有理函数对数。
部分分式中高阶极点由分部积分降阶；简单极点的原函数由
$\operatorname{Li}_2$、对数乘积及对数组成。
零次分母或$a=0$在一般开集上直接线性积分，或由无夹点极限处理。
由欧氏连续分支起步再延拓，保留所有$2\pi i$分支常数。
这给三角图的实际双对数充分判据；
盒图的经典双对数结论可用\cite[\S8--9]{bkone}，
不能仅从“权重0、2、4”三个数字推断任意变分都有同样的参数化公式。

\subsubsection{为什么不能从所有子图为Tate推出完整图为Tate}
\begin{proposition}[纤维结构表不对积分推前封闭]\label{unified:prop:fiberwarning}
即使一个代数映射的每个光滑纤维都在有限扩域后为Tate，
它的总空间仍可含非Tate纯核。
因此仅列出子积分在固定外部数据处的结构，再将它们张量相乘，
不能替代完整图的变分及推前。
\end{proposition}
\begin{proof}
取$\lambda\ne0,1$及光滑射影曲线
\[
 E_\lambda:\ y^2=x(x-1)(x-\lambda),\qquad
 \pi:E_\lambda\longrightarrow\mathbb P^1.
\]
四个分支点是$0,1,\lambda,\infty$。
每个非分支纤维只是二点，其动机为$\Q\oplus\chi_x$；
每个分支纤维的约化动机为$\Q$。
然而Riemann--Hurwitz给$2g(E_\lambda)-2=2(-2)+4=0$，故$g=1$。
更精确地，令$j:V=\mathbb P^1\setminus\{0,1,\lambda,\infty\}\hookrightarrow\mathbb P^1$，
在$V$上$\pi_*\Q=\Q\oplus\mathcal L$，
$\mathcal L$在每个分支点的单值为$-1$。
延拓后$\pi_*\Q=\Q\oplus j_*\mathcal L$，所以
\[
 H^1(E_\lambda)=H^1(\mathbb P^1,j_*\mathcal L).
\]
各穿孔无不变向量，故$j_*=j_!$；全局$H^0,H^2$均为零，
Euler特征为$\chi_c(V)=-2$，右边维数为二。
它是权重一的椭圆核，不是固定二点纤维的Tate张量积。
\end{proof}

这里的根式随着下一积分变量变化；使\emph{每个固定纤维}分裂的域
不能当作一个共同\emph{常数}域。
在整个底上进行分裂覆盖，正好产生了$E_\lambda$本身。
这解释了“先积出一个双对数子图”为什么不会自动保住MPL类别：
下一次推前还看见判别式的分支局部系统及其上同调。
对应的结构预测规则应输出
\[
 \text{纤维核}\quad+\quad\text{判别式局部系统}\quad+
       \quad\text{底与边界的实际推前},
\]
而不能只有纤维类型标签。
同样，二维三传播子sunrise的任意真圈子图至多一圈、收缩该圈后的余图亦为一圈，
但完整光滑参数三次曲线有属一核。
这并不违反森林的低圈张量规则：森林规则控制减除项与方案差，
从未断言未减除完整图就是子图和余图的张量积。

本节由此闭合一类实际入口：完整一圈积分的物理链、算术角色、边界源及操作上界
都无需完整积分即可给出；同时指出推广到多圈必须新增的几何信息。
一般图的推前与最小标架非消失仍是下一项理论责任，
而不是通过增加一圈数值精度就能消除的缺口。



\subsection{边界携带的椭圆核：非零动量注入与实际收缩}
\label{unified:sec:boundaryelliptic}
\subsubsection{从绝对几何到完整相对结构的补充规则}
\otherref{feynman_periods_supplement}{补充材料S1}的差商下降要求内部顶点零外动量。现在允许该顶点注入非零动量，
不再使用串联部分分式。新的入口是边界的权重，而非进一步求解积分。
下面先给适用于一类三维参数空间的规则，再在四边两圈图上落实其几何及物理读出。

在变分层，本文称一个分级组成为Artin--Tate（简称AT），
是指它在某个有限\emph{代数}基变换后成为常值Tate变分。
整个对象允许这些组成间存在扩张。有限单值的代数分支不应误记为椭圆核。
固定数域纤维则采用相应的Artin--Tate对象；这两层不能互相替代。

\begin{theorem}[有理极点曲面的边界结构规则]\label{unified:thm:boundaryweight}
设$X$是光滑射影三维簇，其上同调为AT；$S\subset X$为光滑有理曲面，
也允许$S$为若干光滑有理曲面组成的SNC除子，其双交有理、三交有限。
设$B=\bigcup B_i$为SNC除子，各$B_i$为光滑有理曲面。
允许先作由有限点与有理曲线组成的相容消解。
假设$S+B$为SNC；$B_i\cap S$的正亏格分支正规化记为$C_\alpha$，
其余分支有理；$B_i\cap B_j$的分支均有理，三重交有限。
令
\[
 U=X\setminus S,\quad B^\circ=B\setminus S,\quad
 M=H^3(U,B^\circ;\Q),\quad V_\alpha=H^1(C_\alpha;\Q)(-1).
\]
则$H^*(U)$为AT，且在有理混合Hodge结构中有
\begin{equation}\label{unified:eq:boundaryweight}
 \operatorname{Gr}^W_3 M\simeq\bigoplus_\alpha V_\alpha;
 \qquad\text{其余权重分级均为AT。}
\end{equation}
同一几何正合列在实际有理Nori对象的AT Serre商中给
$M\simeq\bigoplus_\alpha V_\alpha$。
这不是完整混合对象的直和分裂，也不是任意标架都看见全部右端的断言。
在光滑相容参数开集上，相同结论成立于变分层。
\end{theorem}
\begin{proof}
$X$与$S$的Gysin列（多分支时用逐交层的权重谱序列）
表明$H^*(U)$的组成均为AT，故$H^2(U)$与$H^3(U)$均无权重三。
单一光滑有理$S$时，还可具体看出$H^2(U)$纯权重二，
$H^3(U)\hookrightarrow H^2(S)(-1)$纯权重四；
多分支时不需要沿用这两个额外的纯性断言。
对每个有理曲面$B_i$取去除曲线的权重列。
其$H^2(B_i\setminus S)$的权重三部分恰为
$\bigoplus_{C_\alpha\subset B_i}H^1(C_\alpha)(-1)$；
曲线交点以及有理分支只贡献偶权重AT对象。
这一事实可由逐分支Gysin列与交点的Mayer--Vietoris列得到，
因此不要求被去除的曲线并集本身光滑。

在$B^\circ$的覆盖谱序列中，双交是去掉有限点的有理曲线，
三交是有限点。它们所有上同调均为AT，特别没有权重三。
于是从单个面的$H^2$发出的任何可能微分都不能消去权重三部分，
得到$\operatorname{Gr}^W_3H^2(B^\circ)=\bigoplus V_\alpha$。
再使用实际相对列
\begin{equation}\label{unified:eq:boundaryrelative}
 H^2(U)\longrightarrow H^2(B^\circ)
 \overset{\delta}{\longrightarrow}M
 \longrightarrow H^3(U)\longrightarrow H^3(B^\circ).
\end{equation}
两侧环境项没有权重三，权重严格性给出式\eqref{unified:eq:boundaryweight}。
全部其它组成由上述AT项的核、商及扩张构成。

所用Gysin、覆盖与相对列均来自实际几何对；在Nori范畴中先取这些列，
再作AT Serre商，环境项和所有有理交层消失，得到所述同构。
这里没有用Hodge实现的同构反推一个未经构造的动机态射。
族中须保留有理曲面除子类的有限代数单值：
光滑有理曲面的$H^2$由有限个代数除子生成，在泛点有限扩域后可定义，
缩小参数开集使其延伸，便得到所需有限覆盖。故这里的AT断言
具有实际代数除子依据，不是从任意逐纤维Tate性推出全推前。
\end{proof}

\subsubsection{四边两圈三点图及接受开集}
顶点为$A,B,C$；边$x,y$平行连接$A,B$，质量均为$m>0$；
边$u$连接$A,C$，质量$a>0$；边$v$连接$C,B$，质量$b>0$。
三个入射动量满足$p_A+p_B+p_C=0$，令$s_i=p_i^2$，
并允许$p_C\ne0$。在二维复运动学中，三个$s_i$没有额外多项式Gram关系。
先在真实欧氏区域定义正链积分，再作解析延拓。

生成树及二森林直接给
\begin{align}
 U&=xy+(x+y)(u+v),\\
 M_0&=m^2(x+y)+a^2u+b^2v,\\
 F&=M_0U-s_Axyu-s_Bxyv-s_Cuv(x+y),\\
 I(s_A,s_B,s_C)&=\int_{\sigma_3}\frac{U\Omega_3}{F^2}.
 \label{unified:eq:injectedgraph}
\end{align}
每个二维环动量采用$d^2k/\pi$归一化，故本式的$\Gamma(2)=1$。
正质量欧氏积分UV和IR均收敛。图在文献中称为$(2,1,1)$ ice cream cone；
Doran--Harder--Vanhove已研究其绝对几何及局部系统
\cite[\S7.3, 式(120)--(121), Proposition7.6]{dhv2023}。
变量对应为其$(y_1,y_2,x_1,z)=(u,v,x,y)$，
$(p_1^2t,p_3^2t,p_2^2t)=(s_A,s_B,s_C)$。
下文重新证明本节所需的曲面与\emph{相对边界}结论，
不把图超曲面的Tate结论称为新的发现。

取以下非空Zariski开集：$a,b,m\ne0$；
\begin{equation}\label{unified:eq:injecteddiscriminant}
 \Delta_C=(a^2+b^2-s_C)^2-4a^2b^2\ne0;
\end{equation}
两个端面三次曲线光滑；下述两个节点普通，曲面在它们之外光滑；
消解后的交层具有上述正常交叉性质。
这是接受条件，特殊参数必须重新分析，不能按连续性自动继承结论。

\begin{proposition}[极点曲面有理，两个端面椭圆]\label{unified:prop:injectedgeometry}
在上述开集上，$S_0=V(F)\subset\mathbf P^3$有两个普通双点，
其光滑消解为有理曲面。存在由有理曲线与有限点的吹起得到的共同模型$X$，
使正链与极点曲面分离，并满足定理\ref{unified:thm:boundaryweight}的条件。
其中仅有两个正亏格边界分支：
\begin{align}
 E_a(s_A):\ &(m^2x+m^2y+a^2u)\{xy+(x+y)u\}-s_Axyu=0,\\
 E_b(s_B):\ &(m^2x+m^2y+b^2v)\{xy+(x+y)v\}-s_Bxyv=0.
 \label{unified:eq:twofaceelliptics}
\end{align}
因此该实际完整相对载体$M$满足
\begin{equation}\label{unified:eq:injectedcarrier}
 \operatorname{Gr}^W_3 M=H^1(E_a)(-1)\oplus H^1(E_b)(-1),
\end{equation}
其它分级为AT，而忘记边界后的$H^3(U)$全部为AT。
\end{proposition}
\begin{proof}
直线$\ell=\{x=y=0\}$包含于$S_0$。其上梯度为
\[
 (F_x,F_y,F_u,F_v)=(d,d,0,0),\quad
 d=a^2u^2+(a^2+b^2-s_C)uv+b^2v^2.
\]
式\eqref{unified:eq:injecteddiscriminant}保证$d$有两个不同根$P$。
在根处取直线局部坐标$w$，二次首项为
\[
 d'(P)w(x+y)+A_Pxy+m^2(u+v)(x+y)^2,
 \qquad A_P=a^2u+b^2v-s_Au-s_Bv.
\]
$d'(P)\ne0$且$A_P\ne0$时此二次型非退化，故为普通双点。
参数族在基点之外由Bertini光滑；其基点集为$\ell$及有限点。
有限点除四个坐标顶点外，位于某坐标边的非零质量线性因子的根，
或$[1:-1:0:0]$。沿该边或质量方向的导数非零，故均光滑；
四顶点也由$F$的非零线性项光滑。
$A_P\ne0$、端面光滑等条件仅去掉真代数闭集，证明接受开集非空。

任选一个节点，并在有限扩域后将其写为$[1:0:0:0]$。
三次方程成为$X_0Q_2(Y)+Q_3(Y)=0$。
从该点投影到$\mathbf P^2_Y$有有理逆
$Y\mapsto[-Q_3(Y):YQ_2(Y)]$，在$Q_2\ne0$上成立。
故曲面及其光滑消解有理。这个论证保留可能的有限算术扩域。

实际正链首先要求处理整条$\ell$，不能只消解两个复节点。
吹起$\ell$，在$x=\rho\xi,y=\rho\eta$图中有
\[
 U=\rho(u+v)(\xi+\eta)+O(\rho^2),\qquad
 F=\rho d(u,v)(\xi+\eta)+O(\rho^2).
\]
$\Omega_3$也含一个$\rho$因子，因此$U\Omega_3/F^2$沿新例外面
没有径向极点。正链方向$\xi,\eta\ge0$满足$\xi+\eta>0$，
且欧氏区域正$u,v$上$d>0$，故链与严格极点曲面分离。
例外面是有理直纹面；其与严格曲面的交为一条截面$\xi+\eta=0$
和$d=0$的两条纤维。在相交点处，严格方程的$\rho$导数
为$A_P\xi\eta\ne0$，所以此吹起也消解了两个节点。

剩余$x,y$坐标顶点吹起时$U,F$的径向阶均为一，
$\Omega_3$阶为二，故拉回形式阶为一，仍无例外极点。
它们的例外曲面与极点曲面交于直线，正方向由正质量首项避开。
为使复杂边界正常交叉所需的进一步吹起，只在极点上的有限点及由此产生的
有理交线进行，不沿整条光滑椭圆曲线吹起。
例如例外面的截面与纤维交点处，局部形式为$S:\rho+dh=0$与$E:\rho=0$；
点吹起后再吹起有理交线即可作嵌入消解。
这些后续中心在极点上，取极点的\emph{约化全逆像}，
则新的补空间与先前$U$同构，不会错误地把新增极点面纳入可积边界。
新增极点分支和其交层均有理，适用定理的多分支版本。
所有新曲面、双交曲线均有理，$X$上同调为AT。

原端面$v=0,u=0$分别给式\eqref{unified:eq:twofaceelliptics}的光滑平面三次曲线。
其它两个面，例如$x=0$，限制为
\[
 F|_{x=0}=y\bigl[(u+v)(m^2y+a^2u+b^2v)-s_Cuv\bigr].
\]
它由直线与二次曲线组成；$y=0$面相同。
上述吹起仅在端面曲线的有限点作修改，保持它们的$H^1$。
应用定理\ref{unified:thm:boundaryweight}即得式\eqref{unified:eq:injectedcarrier}。
\end{proof}

有理曲面并不等于所有定义域上的纯Tate对象；这里保留有限Artin下降。
也不能从绝对载体为AT推出整个标量积分属于MPL，
因为式\eqref{unified:eq:injectedcarrier}中的椭圆层恰由相对边界进入。
这个结论针对几何载体；单个形式的可见性由第\ref{unified:sec:scalarboundary}节另作标架比较。

\subsubsection{收缩端面的实际态射与联合物理读出}
设$B_v$为$v=0$面的严格变换，$B_{\ne v}$为其余边界并集，
用上标$\circ$表示去掉极点。闭覆盖的相对切除和连接态射给
\begin{equation}\label{unified:eq:facemap}
\begin{split}
 H^2(B_v^\circ,B_v^\circ\cap B_{\ne v}^\circ)
 &\simeq H^2(B^\circ,B_{\ne v}^\circ)\\
 &\longrightarrow H^2(B^\circ)
 \overset\delta\longrightarrow M.
\end{split}
\end{equation}
记复合为$\delta_v$，类似定义$\delta_u$。
这是完整带角点面对象的态射，不是把裸椭圆曲线直接嵌入完整积分。

\begin{theorem}[实际边界读出的精确组成]\label{unified:thm:jointsector}
在式\eqref{unified:eq:injectedgraph}的接受开集上，取联合物理读出
\[
 \mathcal I=(I,\ I[D_v],\ I[D_u]),
\]
其中$D_u,D_v$为相应欧氏传播子分母。则
\begin{equation}\label{unified:eq:cancelledphysical}
 I[D_v]=S_{a^2}(s_A),\qquad I[D_u]=S_{b^2}(s_B).
\end{equation}
两者通过$\delta_v,\delta_u$和同一正相对链实现。
联合读出的最小变分对象，以上同调约定记录其非AT半单组成时，恰为
\[
 H^1(E_a(s_A))(-1)\oplus H^1(E_b(s_B))(-1),
\]
各出现一次，非AT半单秩为四。仅取单个$I$时，本节只证明它的非AT组成
是上述两项的子集；第\ref{unified:sec:scalarboundary}节再以独立的边界链变分证明两项均可见。
\end{theorem}
\begin{proof}
在欧氏动量中置$q=p_A-k-\ell$，可取
\[
 D_u=q^2+a^2,\qquad D_v=(q+p_C)^2+b^2.
\]
分母$D_v$消去后，剩余三条传播子为质量$(m,m,a)$、外动量$p_A$
的sunrise；消去$D_u$后，因$p_A+p_C=-p_B$，得到质量$(m,m,b)$、
外动量$p_B$的sunrise。两积分收敛，同一归一化给式\eqref{unified:eq:cancelledphysical}。
这也由各自的三传播子Schwinger参数化直接给出对应端面形式，
不依赖尚未构造的一般动量/参数动机同构。

连接态射与链边界的配对公式为
\[
 \langle\sigma_3,\delta_v\omega_a\rangle
   =\langle(\partial\sigma_3)_v,\omega_a\rangle,
\]
取$\delta_v$的取向号使右端为正sunrise参数链；$u$面同理。
相对面对象包含其它边和例外边界，故没有遗漏角点或填充项。
定理\ref{unified:thm:boundaryweight}的证明还表明$\delta_v,\delta_u$在权重三上
分别识别$E_a,E_b$的两项。

将$\omega,\delta_v\omega_a,\delta_u\omega_b$作为同一相对载体的三个读出，
联合最小对象是$M$的子商，故由式\eqref{unified:eq:injectedcarrier}得到重数上界一。
两条sunrise物理函数的正常阈值割均非零，
\otherref{feynman_periods_supplement}{补充材料S1}的非恒定椭圆变分与不可约单值证明给各自完整$H^1$下界。
在独立参数$s_A,s_B$上，它们互不同构：
沿$s_A=(a+2m)^2$绕行而固定一般$s_B,s_C$，第一个有非平凡幺幂单值，
第二个无单值；交换$A,B$同理。故两项都必须出现，匹配上界。
这里使用的是变分与实际物理面读出，未宣称所有指定代数纤维都互不同源。
\end{proof}

三个读出可进一步作可观测量选择。
只保留$I[D_v]$时，$E_b$方向在该读出中确实消失，只剩$E_a$；
只保留$I[D_u]$则相反。这给出非零内部动量注入下的实际分子筛选与收缩规则，
不需要先求三个函数的特殊函数表达式。
若加入已证由AT对象承载的局部项，加性Serre定理继续保持联合非AT组成。
这不赋予任意局部反项自动为AT的性质。

\subsubsection{为什么绝对微分方程可能有椭圆边界源}
\begin{corollary}[绝对消去不等于完整消去]\label{unified:cor:boundarysource}
将上述族限制到一条光滑代数参数曲线，令$L$为消去$[\omega]\in H^3_{\dR}(U)$
的有理系数微分算子。则$LI$的实际周期载体可取
\[
 K=\operatorname{coker}\{H^2(U)\longrightarrow H^2(B^\circ)\}.
\]
$K$的非AT组成恰由两个端面椭圆核给出；$LI$的标架可能只看见其中一部分。
特别，绝对载体的有限代数单值，不能单独推出$LI$是代数函数。
\end{corollary}
\begin{proof}
相对列\eqref{unified:eq:boundaryrelative}把$K$识别为
$M\to H^3(U)$的核，且该列与Gauss--Manin连接相容。
因此$L[\omega]$在相对de Rham实现中属于$K$，
正链读出限制到$K$即给$LI$。其组成由定理\ref{unified:thm:boundaryweight}确定。
在形式层若$L\omega=d\eta$，相对锥中的恒等式为
$[(d\eta,0)]=[(0,-\eta|_{B^\circ})]$，明确展示边界源。
该负号采用相对微分$d(\alpha,\beta)=(d\alpha,\alpha|_B-d\beta)$；
改变锥约定只同时改变连接态射和取向号，不改变实际周期。
\end{proof}

上述推论是前文完整\emph{相对}留数定理的补充：
只消去绝对部分时，边界还可能带非Tate源；
消去包含边界的完整相对留数像后，才可使用已证的环境背景结论。
这也说明文献中的绝对Picard--Fuchs算子与完整物理函数的齐次最小算子
不可直接等同。本节不把忘边界的上同调映射等同于尚未比较的全部物理最大割。

若退回$p_C=0,s_A=s_B$，本族落在特殊几何层，\otherref{feynman_periods_supplement}{补充材料S1}的实际Stokes公式重新适用；
不能在未检查奇异层及附近循环的情况下，用该特殊层的最小性证明一般参数标架最小性。
本节已闭合的输出为：一般非零注入下的完整载体上界、两个端面核的来源、
实际收缩态射、联合读出的精确组成，以及绝对消去后的源载体。
单个标量$I$的两个椭圆方向由下一节的专属链变分证明均受保护；
不能仅用本节的载体组成替代该证明。



\subsection{从边界核到单个物理读出：一维变分保护}
\label{unified:sec:scalarboundary}
\subsubsection{把非零跳跃升级为指定简单因子的判据}
上一节的两个椭圆端面是完整载体的组成。
本节证明：在声明的光滑参数开集上，不必添加收缩分子，
原标量积分$I$本身就看见两者。证明只用局部几何与渐近首项，
不求完整积分、连接矩阵或数值单值。

首先在同调方向表述一个可重复使用的判据。
设$E$是连通参数开集上的有限秩有理局部系统，
带单值不变过滤$W$；相关分级半单。
设$V$是其中恰出现一次的简单局部系统。
令$T$为一个具体闭路的单值，$N=T-1$。
形式$\lambda$为单值代数de Rham读出，$\sigma$为一个平坦Betti芽，
$f=\langle\lambda,\sigma\rangle$。
子商$R/P$承载此标架，是指$\sigma\in R$且$\lambda|_P=0$；
这里的等式均在相应实现及参数开集上理解。

\begin{theorem}[一维变分保护]\label{unified:thm:rankonevisibility}
假设$\rk N=1$是在整个$E$上成立，且$N$在某个
$\operatorname{Gr}^W_kE$中的$V$上诱导非零算子。
若实际函数的跳跃
\[
 \Delta_T f=\langle\lambda,N\sigma\rangle
\]
不是恒零函数，则每个承载此标架的局部系统子商$R/P$都含有$V$。
因而任何相容的实际最小变分对象也保留$V$。
\end{theorem}
\begin{proof}
记$\Gamma=\im N$。由于$N$在$V$上非零，
可以在$W_kE$中选一个提升，使其$N$像在$\operatorname{Gr}_k^W$中非零。
整个像只有一维，故
$\Gamma\subset W_kE$、$\Gamma\cap W_{k-1}E=0$，
而其非零分级像位于$V$。
跳跃非零说明$N\sigma\ne0$，于是
$\Gamma\subset R$且$\Gamma\not\subset P$。
分级半单性与$V$的简单性表明$\operatorname{Gr}_k^WR$包含整个$V$。

反设$R/P$不含$V$。对诱导过滤取分级仍有正合列，
且$V$在$\operatorname{Gr}E$中重数一，
所以$\operatorname{Gr}_k^WP$必须包含该$V$。
选$p\in P\cap W_kE$，使其分级像$v\in V$满足$Nv\ne0$。
单值保持$P$，故$Np\in P$；它在分级上的像非零，
因而$Np$生成$\Gamma$。这与$\Gamma\not\subset P$矛盾。
\end{proof}

判据中的总秩一不可换成“$V$上的变分秩一”。
例如将一个具有非零退化单值的椭圆变分与一个对数扩张取直和，
只读出后者，仍有非零跳跃却完全看不见椭圆层。
同时，本定理不把一个参数闭路当成固定纤维Nori群的元素：
它首先是变分及实际函数的结论。

\subsubsection{边界普通二重点的总变分秩}
考虑一个光滑三维复流形中的相对对$(U,B^\circ)$，
其中$U=X\setminus S$。设围绕参数$\epsilon=0$只有一处失去横截性，
发生在一个光滑极点分量$S$与一个光滑边界分量$B_v$之间；
其它边界均不经过此点，其余已紧化、解消的对在参数盘上局部平凡。
假设该点的横截局部模型为
\begin{equation}\label{unified:eq:boundarypinchmodel}
 S=\{w=0\},\qquad
 B_v=\{w=z_1^2+z_2^2+\epsilon\}.
\end{equation}
非零单位与参数的可逆解析换元不改变下述结论。

\begin{proposition}[边界夹挤的秩一性]\label{unified:prop:boundaryrankone}
局部闭路对$H_3(U,B^\circ;\Q)$的变分满足$\rk(T-1)\leq1$。
若它在边界曲线的$H^1$组成上是非零Picard--Lefschetz变分，
则总秩恰为一。
\end{proposition}
\begin{proof}
选适配Milnor邻域$D$及固定外圈。
对的单值可取为$D$外恒等，沿外圈也恒等；
把任意相对链细分后，其差分支撑于$D$，
故全局变分像因子化经过
$H_3(U_D,B_D^\circ;\Q)$的全局像。
这一局部化是实际链的切除，不是只对权重分级作估计。

由\eqref{unified:eq:boundarypinchmodel}，$U_D$收缩至一个圆，
所以$H^2(U_D)=H^3(U_D)=0$。
$B_D^\circ$则为复二维Milnor邻域去掉
$C_\epsilon=\{z_1^2+z_2^2+\epsilon=0\}$。
后者的Milnor纤维收缩至一个圆。局部Gysin列与相对正合列依次给出
\[
 H^2(B_D^\circ)\simeq H^1(C_\epsilon)(-1),\qquad
 H^3(U_D,B_D^\circ)\simeq H^2(B_D^\circ).
\]
这里只需两个空间均为一维，不把局部极限的混合Hodge结构
误认为光滑全局纤维的权重分解。
由有理系数的对偶及切除，全局变分秩至多一。
若边界曲线的全局组成已有非零变分，则总变分不为零，秩即为一。
在本节使用的普通椭圆节点情形，分级上的变分平方为零；
总像的一维线在该分级中投影非零，故它与更低过滤的交为零。
于是$N$再次作用于像线必为零，即总算子也满足$N^2=0$。
\end{proof}

局部管、变分与交数的标准背景见\cite[\S2.3]{muhlbauer}。
这里另给了带边界的相对链论证；没有用单圈物理定理替代本两圈问题。
在同调方向，上一节的$V_a=H^1(E_a)(-1)$对应其对偶，
位于权重$-3$；下文陈述几何组成时仍用上同调记号$V_a,V_b$。

\subsubsection{实际正链的非零阈值首项}
仍采用上一节的四边图、正质量$m,m,a,b$和每圈$d^2k/\pi$归一化。
取正射影链的$x+y=1$截面，写$y=1-x$、$z=x(1-x)$。
在欧氏区域，齐次换元给
\begin{equation}\label{unified:eq:scalaricechart}
 I=\int_0^1dx\int_0^\infty du\int_0^\infty dv\,
       \frac{z+u+v}{F^2}.
\end{equation}
沿已固定的正链取向，$dx\,du\,dv$为正测度。
此式直接来自齐次参数形式，Gamma因子为$\Gamma(2)=1$。
多项式准确分解为
\begin{align}
 F&=F_a+vB(x,u)+b^2v^2,\label{unified:eq:scalariceexpand}\\
 F_a&=(m^2+a^2u)(z+u)-s_Azu,\nonumber\\
 B(x,u)&=m^2+(b^2-s_B)z+(a^2+b^2-s_C)u.\nonumber
\end{align}
固定$s_B,s_C<0$，令$h_a=a+2m$、$\tau_a=h_a^2$。
$F_a$在$s_A=\tau_a$的正域中恰有一个零点：
\[
 x_0=\tfrac12,\qquad u_0=\frac{m}{2a},\qquad v_0=0.
\]
例如在$z>0,u>0$时，
$F_a/(zu)=m^2/u+m^2/z+a^2+a^2u/z-s_A$；
先对$u$、再对$z\leq1/4$取最小值，即得$\tau_a$及唯一等号点。
又由$B>0$与$b^2>0$，全正链上没有其它同时达到此阈值的极点。

设$\epsilon=\tau_a-s_A$、$\xi=x-\tfrac12$、$\eta=u-u_0$。
在$v=0$处有精确展开
\begin{equation}\label{unified:eq:scalaricehessian}
 F_a=a^2\eta^2+\frac{m^2h_a}{a}\xi^2
       +4m(m+a)\eta\xi^2+\epsilon zu.
\end{equation}
记
\begin{align}
 U_a&=\frac{h_a}{4a},&
 c_a&=\frac{m}{8a},&
 A_a&=\frac{m^2h_a}{a},\nonumber\\
 B_a&=m^2+\frac{b^2-s_B}{4}
             +\frac{m}{2a}(a^2+b^2-s_C)>0.\label{unified:eq:scalariceconstants}
\end{align}
由于$\partial_vF=B_a\ne0$，极点曲面在夹挤点本身仍光滑。
非退化的正切Hessian和参数横截项$c_a\epsilon$，
由带参数的解析Morse坐标化给出\eqref{unified:eq:boundarypinchmodel}。

\begin{proposition}[不求完整积分的非零跳跃]\label{unified:prop:scalaricelog}
在上述归一化及从正欧氏链续接的分支上，
\begin{equation}\label{unified:eq:scalaricelog}
 I=-C_a\log\epsilon+O(1),\qquad
 C_a=\frac{\pi U_a}{B_a\sqrt{A_aa^2}}
     =\frac{\pi h_a}{4a B_a m\sqrt{a h_a}}>0
 \quad(\epsilon\longrightarrow0^+).
\end{equation}
局部奇异部分可写成$H_a(\epsilon)\log\epsilon$，
其中$H_a$解析且$H_a(0)=-C_a$。
因此$\epsilon$逆时针一周的实际跳跃非零，极限为$-2\pi iC_a$。
交换$(a,s_A,u)$与$(b,s_B,v)$得到另一正常阈值处的非零跳跃。
\end{proposition}
\begin{proof}
先限制到夹挤点的一个小正半邻域。
利用$\partial_vF\ne0$，以$w=F$作法向坐标，
其下端为$w=F_a$。写解析密度
$J(x,u,w,\epsilon)=(U/\partial_vF)(x,u,v(w),\epsilon)$。
对$J/w^2$作法向积分，常数项产生
$J(x,u,0,\epsilon)/F_a$；
其余项至多产生$\log F_a$或解析项。
在二维切向中，$\log F_a$的积分只有
$\epsilon\log\epsilon$等次领先项，因此不改变首项系数。
同理，$J(x_0,u_0,0,0)=U_a/B_a$以外的密度项不改变该系数。

现在使用\eqref{unified:eq:scalaricehessian}的解析Morse坐标，
切向Jacobian的常数为$1/\sqrt{A_aa^2}$，
临界值为$c_a\epsilon+O(\epsilon^2)$。
径向积分
\[
 \int_{X^2+Y^2<r^2}
       \frac{dX\,dY}{c_a\epsilon+X^2+Y^2}
       =-\pi\log\epsilon+O(1)
\]
即给\eqref{unified:eq:scalaricelog}。
解析密度逐次展开后，角向奇项积分为零，
偶项径向积分为解析项加$\epsilon^j\log\epsilon$；
法向的次领先项具有相同性质。
缩小邻域即可使展开一致收敛，故对数系数是解析函数。
这一步也证明了非零单值，而非仅从实轴上的$O(1)$记号猜测分支。

其余链在先前构造的紧化、吹起模型中不接触极点，
在足够小的参数盘上给解析函数；汇合到无穷远的坐标端
由该模型的收敛性与横截性一并控制。
所以局部首项没有被其它链片抵消。
\end{proof}

式\eqref{unified:eq:scalaricelog}的系数只需要一个局部Hessian和一个旁观传播子的
非零法向系数。它是实际函数的结构证据，不是数值拟合。
本节固定了参数闭路的方向；没有将此符号直接称为所有物理通道的
Cutkosky约定，后者仍须相应Minkowski链比较。

\subsubsection{单标量积分的精确椭圆组成}
取上一节已接受的连通光滑参数开集。沿两个正常阈值的一般点，
要求解消后的对除了所述端面夹挤外无其它同时退化。
该条件是几何适用条件，不能省去。
在本族中可满足：固定$s_A=\tau_a$后，
$s_B,s_C$仍提供线性系，Bertini排除基集外的新奇点；
在$v=0$的固定端面节点处，
$\partial_vF=B_a$对负的$s_B,s_C$严格非零。
旧基线及坐标顶点的解消照上一节进行，
再排除有限个代数坏条件即可。另一个端面同理。

\begin{theorem}[非零内部注入图的单标量最小组成]\label{unified:thm:scalaricevisibility}
对于此参数族及其正物理链芽，完整标量积分$I$的最小变分对象，
以上同调方向记录时，其非Artin--Tate半单组成恰为
\begin{equation}\label{unified:eq:scalariceexact}
 \bigl(M_{\min}(I)_{\mathrm{nonAT}}\bigr)^{\mathrm{ss}}
       =H^1(E_a(s_A))(-1)\oplus H^1(E_b(s_B))(-1),
\end{equation}
各出现一次，非AT半单秩为四。
这是单个$I$的结论，无需同时加入$I[D_v]$和$I[D_u]$。
\end{theorem}
\begin{proof}
上一节的相对权重列给出载体上界：非AT组成只有$V_a,V_b$，
重数各一。两条sunrise椭圆族均非等常，
其$H^1$几何局部系统绝对不可约，由\otherref{feynman_periods_supplement}{补充材料S1}的椭圆单值论证说明。
沿$s_A=\tau_a$的一般点，$V_a$有普通单节点退化的非零变分，
而$V_b$横截延拓。故两者作为全参数局部系统不相同。

先选$s_B,s_C<0$的一般点，在$s_A$方向绕$\tau_a$。
命题\ref{unified:prop:boundaryrankone}给完整同调载体上的总变分秩一，
它在$V_a^\vee$上非零；命题\ref{unified:prop:scalaricelog}给实际$I$的非零跳跃。
定理\ref{unified:thm:rankonevisibility}因而强制任何最小标架子商保留
$V_a^\vee$。对另一个正常阈值重复论证，得到$V_b^\vee$。
这些基点可在同一光滑参数开集内平行运输到同一芽，
所以两个下界属于同一物理变分对象。
与重数一的上界比较并对偶，即得\eqref{unified:eq:scalariceexact}。
\end{proof}

此处的“秩四”是非AT半单层之秩，不能改写为完整混合对象恰为四维。
Tate层与扩张仍可增加完整最小秩。
同样，本定理不声称任意两个固定代数点的曲线互不同源，
也不以参数变分排除某个孤立数值偶然等于MZV。
与上一节的分子公式合用，得到同一个图中的精确选择表：
标量$I$保留两项；收缩分子$D_v$只保留$V_a$；
$D_u$只保留$V_b$。这给出了实际物理操作对结构的筛选，
而不是仅对载体作标签分类。

\subsubsection{显式一参数切片的齐次阶数下界}
取$\Lambda>0$、$c<0$，并令
\begin{equation}\label{unified:eq:scalariceray}
 s_A=t,\qquad s_B=-t-\Lambda,\qquad s_C=c.
\end{equation}
在$-\Lambda<t<0$，三个不变量均负，正参数链无极点。
如要求一个实欧氏动量起点，可取$t=-\Lambda/2$、
$-2\Lambda<c<0$；Gram条件允许这样的非共线动量。
其余$t$按该芽解析续接。
选择一般$\Lambda,c$使上述正常阈值处满足单一夹挤条件。
下述有理函数系数域$k(t)$包含固定质量及$\Lambda,c$；
若讨论数域上的算术对象，另取这些常数为代数数。

\begin{corollary}[无需完整积分的微分阶数障碍]\label{unified:cor:scalariceorder}
沿\eqref{unified:eq:scalariceray}，$I(t)$的每个非零齐次有理系数
微分消去算子都至少四阶。有限代数系数扩张后仍有此下界。
\end{corollary}
\begin{proof}
$E_a(t)$的有限退化值属于
\[
 D_a=\{0,a^2,(a-2m)^2,(a+2m)^2\}\subset[0,\infty),
\]
而$E_b(-t-\Lambda)$的有限退化值属于
$-\Lambda-D_b\subset(-\infty,-\Lambda]$，两组不交。
两个仿射拉回椭圆局部系统仍非等常且绝对不可约，
且一个正常阈值处仅其中一个有非零单值。
在$t=\tau_a$和$t=-\Lambda-\tau_b$，
另一通道参数为负，故本节的两个正链跳跃证明均适用。
秩一保护使该一参数标量芽的变分子商也含两个秩二因子。

更明确地，将平坦链$r$映为函数芽$\langle\lambda(t),r(t)\rangle$。
在由$\sigma$的单值轨道生成的有限秩局部系统中，
其恒零函数核是单值不变子系统。
为讨论复函数关系，可将此有限秩局部系统复化；
所用两个椭圆因子绝对不可约，保护论证仍成立。
对这个商可直接用定理\ref{unified:thm:rankonevisibility}，
不需假定所有函数关系都来自动机关系。
若$L$为$r$阶齐次算子，所有续接芽均在其$r$维局部解空间中，
因此$r\geq2+2=4$。
有限代数扩张只将局部幺幂单值换成非零整数次幂，
总秩一与非零跳跃仍保留；椭圆连通几何单值亦不失不可约性。
\end{proof}

本节没有给出四阶上界，因而不能把“至少四阶”写成“恰四阶”。
其作用是用边界类型、普通夹挤与局部正链，
在不计算全函数的前提下把结构表的上界提升为精确简单组成。
加入实际AT背景项后，这些非AT组成还受前文加性比较定理保护。



\endgroup
\subsection{从判据到实例}
上述证明把任务分成两步。首先确定一个实际载体允许哪些组成；其次证明物理链与读出确实激发了其中哪些方向。第二步有时可由正常阈值的非零变分完成，未必需要完整积分。

本节的一圈与边界椭圆实例给出两个层次的检验。正质量欧氏一圈积分在所列整数维数和收敛条件下有Artin--Tate载体；一般位置四维三角图与盒图的载体秩分别为七和八。这是环境秩，仍须由标架决定是否取满。\otherref{feynman_periods_supplement}{补充材料S1}的串联细分图则通过相对Stokes公式下降到有限个sunrise端点；质量阈值分离时，不同椭圆变分在物理读出中各出现一次。改变传播子分子会使部分端点权重消失。

边界例说明了相对对象的重要性。一个非零内部动量注入的四边两圈图可以有有理极点曲面，却在两个边界面上携带椭圆核。接受开集内，整个相对载体的局部秩一变分与正链阈值首项共同证明，单个标量积分看见两份不同的椭圆变分。绝对曲面的有理性因而不足以判定完整积分。

物理求和还要求新的重数检查。若若干图共享一个阈值，必须对总读出的完整跳跃求秩；只检验某个首项系数，不能排除更高阶的对数跳跃。补充材料同时给出这种判据以及一个三通道交叉对称组合。它们说明结构上界、阈值保护和实际求和可以相互配合，也说明何时只能得到重数区间。

\part{核心实例}

\section{能量字符与覆盖几何}\label{ch:graph}

\begingroup

本节研究泛独立能量下的组合支撑。其结论先处于有理函数层；从支撑到实际周期对象的比较另行提出。Dual-critical 图的图论背景见\cite{dualcritical}。

\subsection{对象、输入与输出}\label{sec:graph:1}

图 \(G=(V,E)\) 有 \(n\geq1\) 个顶点，边以输入顺序编号 \(0,\ldots,m-1\)，允许平行边和自环。根变量 \(\omega_e\) 彼此独立；顶点能量 \(p_v\) 除 \(\sum_vp_v=0\) 外独立。先在正实根和实欧氏能量处定义收敛能量积分，再使用其有理延拓。传播子幂均为一，无分子和反项。原稿另论高次幂，本节不把它加入算法定理。

令 \(R_e\) 翻转 \(\omega_e\)，

\begin{equation}
e_S=2^{-m}\sum_{\epsilon\in\{\pm1\}^{E}}
       \left(\prod_{e\in S}\epsilon_e\right)R_\epsilon,\qquad
 \mathcal V(G)=\{S\subset E:e_S\mathscr E_G\not\equiv0\}.
\end{equation}

这里``不恒为零''是有理函数意义。固定外腿能量、质量关系或其他物理约束后，某些字符可能消失。一般支撑给出特化支撑的上界，不能逐点保证非零。

记 \(\delta(A)\) 为割边，\(W(A)=\sum_{\delta(A)}\omega_e\)，\(p(A)=\sum_{v\in A}p_v\)。以下固定能量核的规范：

\begin{equation}
\mathscr E_G=
 \frac1{\prod_e2\omega_e}
 \sum_{\pi\in S_n}\prod_{j=1}^{n-1}
 \frac1{W(A_{\pi,j})+ip(A_{\pi,j})}.                 \label{eq:graph:1}
\end{equation}

单顶点图的空乘积为一。自环只贡献 \(1/(2\omega_e)\)，故每个非零字符必须含全部自环。

程序接受 JSON 的 n 与 edges，拒绝不连通图、无效端点及超出给定资源界的请求。默认顶点上界12，总奇偶掩码/子集状态预算 \(2^{22}\)；这不意味着接受所有十二顶点图。预算拒绝不是数学上的零支撑，也不是不可判定结论。输出包含边束、自环强制标签、每种保留奇偶图、全部可达子集位表、存在时的好排序、支撑总数和按删边数计数的因式化多项式。


\subsection{支撑定理与极点证明}\label{sec:graph:2}

无自环多重图 \(H\) 称为 dual-critical，如果存在排序 \(v_1,\ldots,v_n\)，使每个 \(j>1\) 与此前顶点之间有奇数条边，按重数计。单顶点空图满足定义。该术语、好排序及奇割递归属于既有图论；见 Király--Kisfaludi-Bak \cite{dualcritical} 第1--3页。不能把某个奇度顶点的存在当作充分条件。

\begin{theorem}[泛能量字符支撑]

\label{thm:graph:1} 在上述一般参数范围内，

\begin{equation}
e_S\mathscr E_G\not\equiv0
 \quad\Longleftrightarrow\quad G\setminus S\text{ 是 dual-critical}.       \label{eq:graph:2}
\end{equation}

右侧若留有自环即判否。

\end{theorem}

\subsubsection{极点公式}\label{sec:graph:2.1}

设 \(V=A\sqcup B\)，两侧诱导图均连通，\(C=\delta(A)\)，\(u=W(C)+ip(A)\)。在 \(u=0\) 上令

\begin{equation}
\widetilde p_v=
 \begin{cases}
 p_v-i\sum_{e\in C:e\ni v}\omega_e,&v\in A,\\
 p_v+i\sum_{e\in C:e\ni v}\omega_e,&v\in B.
 \end{cases}
\end{equation}

两侧分别守恒。在这一规范下，留数表示 \(u^{-1}\) 的系数，不含 \(2\pi i\)：

\begin{equation}
\operatorname{Res}_{u=0}\mathscr E_G
 =\left(\prod_{e\in C}(2\omega_e)^{-1}\right)
   \mathscr E_{G[A]}(\widetilde p_A)
   \mathscr E_{G[B]}(\widetilde p_B).                         \label{eq:graph:3}
\end{equation}

在排序式中，产生此极点的排序把 A 放在 B 之前；去掉这一个分母后，两侧的前缀分母恰是移位能量下的式\eqref{eq:graph:1}。内部根翻号与上述移位交换。跨边根只有全不翻、全翻能够产生同一个一般超平面；全翻项对应 B 在 A 前，另有分母的负号。因此

\begin{equation}
\operatorname{Res}_{u=0}(e_S\mathscr E_G)=
 \frac{1+(-1)^{|C|+|S\cap C|+1}}{2^{|C|}}
 \left(\prod_{e\in C}(2\omega_e)^{-1}\right)
 (e_{S\cap E(A)}\mathscr E_{G[A]})
 (e_{S\cap E(B)}\mathscr E_{G[B]}).                         \label{eq:graph:4}
\end{equation}

所有下图能量取 \(\widetilde p\)。系数非零等价于 \(|C\setminus S|\) 为奇数。

\subsubsection{不连通侧的伪极点相消：一个有理桥极限证明}\label{sec:graph:2.2}

补充原稿弱指数正规化论证，避免用分布极限代替有理恒等式。对有 \(c>1\) 个分量的图 H，仍以式\eqref{eq:graph:1}的有限排序和定义其形式核，顶点能量只作总守恒。用 \(c-1\) 条桥把分量连成树，每条桥根取 \(\eta\)，得到连通 \(H_\eta\)。

在一般分量能量处，排序式每个分母都有非零的 \(\eta=0\) 极限，所以

\begin{equation}
\mathscr E_H^{\rm formal}
 =\lim_{\eta\to0}(2\eta)^{c-1}\mathscr E_{H_\eta}.          \label{eq:graph:5}
\end{equation}

另一方面，新边均为桥，其能量由桥一侧的外能量总和 \(P_b\) 固定，不形成新圈。沿桥逐个分解原能量积分得

\begin{equation}
\mathscr E_{H_\eta}
 =\prod_b(P_b^2+\eta^2)^{-1}
   \prod_\alpha\mathscr E_{H_\alpha}(p^\alpha).
\end{equation}

\(p^\alpha\) 由桥能量向附件顶点的注入得到，分别守恒，且不依赖 \(\eta\)。在一般 \(P_b\ne0\) 处乘 \((2\eta)^{c-1}\) 后极限为零。故形式核恒为零。这是有限有理函数的极限，无须交换不绝对收敛的时间积分。将此结论用于式\eqref{eq:graph:3}的一侧，证明不连通分割的候选留数为零。

\subsubsection{归纳}\label{sec:graph:2.3}

充分性：令 \(H=G\setminus S\) dual-critical，取好排序的最后顶点 v。\(H-v\) 仍 dual-critical 且跨边保留数为奇数。特别地 \(G-v\) 连通。式\eqref{eq:graph:4}的系数及按归纳所得下图投影均非零。极点超平面上的下图守恒能量仍是独立参数，所以留数非零。

必要性：当 \(n>1\)，式\eqref{eq:graph:1}关于独立外能量只有线性候选极点，且沿各一般超平面至多一阶。若所有留数为零，则投影是这些外能量的多项式；但沿一般方向 \(p\mapsto\lambda p\)，每个排序项为 \(O(\lambda^{-(n-1)})\)，只能恒零。非零投影因而有一个非零留数。其极点可能含带符号根；利用投影在任意根翻号下为特征向量，可把它移到式\eqref{eq:graph:4}的标准极点。第\ref{sec:graph:2.2}节排除不连通侧，式\eqref{eq:graph:4}及归纳遂给出两个 dual-critical 诱导子图和奇数保留跨边。文献\cite{dualcritical} Theorem3 的奇割递归等价性给出 H dual-critical。单顶点为归纳起点。证毕。

证明决定恒零项，有限域证书只额外验证实现和某个特化上的非零性。


\subsection{平行边压缩与精确计数}\label{sec:graph:3}

令 q 为自环数，非空无序端点边束为 \(B_1,\ldots,B_b\)，重数 \(r_j\ge1\)，于是 \(\sum_jr_j=m-q\)。对删边字符 S 定义

\begin{equation}
p_j=|B_j\setminus S|\bmod2.
\end{equation}

令 \(H_p\) 在第 j 对端点上有一条边当且仅当 \(p_j=1\)。

\begin{theorem}[压缩输出]

\label{thm:graph:2} S 属于支撑，当且仅当它包含全部自环且 \(H_p\) dual-critical。因此

\begin{equation}
|\mathcal V(G)|=2^{m-q-b}
    \#\{p\in\{0,1\}^b:H_p\text{ dual-critical}\},             \label{eq:graph:6}
\end{equation}

并有删边度数生成函数

\begin{equation}
\sum_{S\in\mathcal V(G)}z^{|S|}
 =z^q\sum_{p:H_p\ {\rm dual\mbox{-}critical}}
   \prod_{j=1}^b
   \frac{(1+z)^{r_j}+(-1)^{r_j-p_j}(1-z)^{r_j}}2.           \label{eq:graph:7}
\end{equation}

\end{theorem}

\begin{proof}

对任意固定顶点顺序，每个向前边数的奇偶性只取决于对应边束保留数的奇偶性；这也就是文献\cite{dualcritical} Proposition4 的删除平行边对不变性。应用定理\ref{thm:graph:1}即可。固定 \(p_j\) 时，束内删边子集大小须满足 \(|S\cap B_j|\equiv r_j-p_j\pmod2\)，恰有 \(2^{r_j-1}\) 种。各束独立，自环强制，给式\eqref{eq:graph:6}。偶/奇二项式筛选给式\eqref{eq:graph:7}。

\end{proof}

式\eqref{eq:graph:6}按字符计数一次，不按可能不唯一的好排序计重。证明不是声称平行边不影响能量核的数值或几何。

例如，m 条平行边的 banana 只需两个奇偶图，接受奇数保留边的一个分支，支撑数 \(2^{m-1}\)。\(m=100\) 时为 \(633825300114114700748351602688\)。显式列出全部字符不可能具有关于 m 的多项式输出长度。


\subsection{判定递推与复杂度}\label{sec:graph:4}

对每个 \(H_p\) 置 \(D[\varnothing]=0\)、\(D[\{v\}]=1\)，其余

\begin{equation}
D[A]=\bigvee_{v\in A}
   \bigl(D[A\setminus\{v\}]
     \land(\deg_{H_p}(v,A\setminus\{v\})\equiv1\bmod2)\bigr). \label{eq:graph:8}
\end{equation}

从最后顶点删除与逆向添加立即证明正确性。每一步子集大小下降，因此终止；记录最后顶点即可重建好排序。拒绝证书用完整递推位表，独立重放也可穷举排序。

在一个顶点掩码可装入机器字的计算模型中，每个奇偶图耗 \(O(n2^n+b)\) 个图运算，完整压缩输出耗

\begin{equation}
O(m+b\log b+2^b(n2^n+b)).
\end{equation}

只保留当前图时工作空间 \(O(m+2^n+n)\)；现有程序还保留全部证书，空间/输出大小包括 \(O(2^b(2^n+bn))\) 比特及标签和计数所需的多项式附加项。可达位表按字节打包，避免逐项增长的大整数求和。任意精度整数计数和根数值计算不能免费计作机器字操作；位复杂度仍有界于

\begin{equation}
2^{O(n^2)}\operatorname{poly}(m,n),
\end{equation}

因 \(b\le n(n-1)/2\)。所以压缩问题关于顶点数 n 是 FPT，\textbf{不声称关于圈数、树宽或任意其他参数的 FPT}，也不声称全支撑显式列表有同一复杂度。

默认预算预先检查 \(2^{n+b}\)，再开始枚举；超限抛出明确错误，不返回伪``空支撑''。压缩查询还检查字符标签唯一且在范围内。图输入不包含质量和动量，因此不能承诺专化支撑。



\subsection{圈星图的多重二次覆盖}
原圈星文稿提供一个具有明确传播子支撑的无限图族。取 $k+1$ 个圈变量，中心圈标为 $0$，叶圈标为 $1,\ldots,k$；每圈恰有三条只依赖自身圈动量的传播子，另有 $k$ 条共享传播子分别依赖 $(0,i)$。圈数为 $k+1$，内边数为 $4k+3$。从生成树及基本圈矩阵检查这一证书，可识别该类，而非先知道积分结果再命名图族。

在三重割有理化所需的有限代数运动学覆盖上，每个圈先化为有理二次曲线，共享传播子给双次数 $(2,2)$ 的关系。因此最大割函数域具有形式
\begin{equation}
 k(C)=k(z)(\sqrt{\Delta_1(z)},\ldots,\sqrt{\Delta_k(z)}),
 \qquad\deg\Delta_i\le4.
\end{equation}
以下亏格计算假设四次判别式平方自由、两两没有公共根且平方类独立。覆盖群为 $(\Z/2)^k$，每个非零字符 $S$ 对应
\begin{equation}
 C_S:\quad y_S^2=\prod_{i\in S}\Delta_i(z).
\end{equation}
字符投影由实际覆盖群作用定义，所以在这一特定几何中确有上同调分解。它与能量核翻号角色的抽象投影属于不同构造，只有给出两种作用之间的实际比较后才能相互识别。

\begin{proposition}[一般覆盖的亏格]
在上述假设下，$g(C_S)=2|S|-1$，且
\begin{equation}
 g(C)=\sum_{S\ne\varnothing}g(C_S)=1+2^k(k-1).
\end{equation}
\end{proposition}
\begin{proof}
$C_S\to\mathbf P^1$ 的分支点有 $4|S|$ 个，偶次数使无穷远不再分歧。Riemann--Hurwitz 给 $2g(C_S)-2=-4+4|S|$。全部非平凡字符组成 $H^1(C)$；利用 $\sum_S|S|=k2^{k-1}$ 求和即得第二式。也可直接对次数 $2^k$ 的覆盖应用 Riemann--Hurwitz：每个分支点的惯性群阶为二。
\end{proof}

标量最大割的规范形式只选出全乘积字符，其宿主为
\begin{equation}
 Y^2=\prod_{i=1}^k\Delta_i(z),\qquad
 \omega=c\,\frac{z^{k-1}\,dz}{Y},\qquad g=2k-1.
\end{equation}
因此整个覆盖的亏格并不是标量读出的最小几何秩。分子可改变所选字符；共享根、平方类依赖及无穷远分歧则改变正规化后的分支计数。程序或理论规则必须先处理这些退化，不能把泛参数公式直接套到特殊点。

原圈星稿还给出充分独立离壳外部源下的参数占优及一般单值论证，主张所选超椭圆部分具有 $\mathrm{Sp}_{4k-2}$ 几何单值和一般绝对简单 Jacobian。本文将其保留为原稿中的族结果；把它接入统一实际动机范畴、固定四点物理切片及完整未割振幅仍有独立证明责任。它不意味着任意运动学专化都保持相同群。


\subsection{极化障碍与不能跨越的逻辑缺口}
极化给候选几何子商施加双线性形式、权重和签名约束。若某局部模型的全部可用子商都不满足目标 K3 的极化要求，则该模型不能承担所声称的纯 K3 核。这是排除一个比较方案的证据，不是对完整积分所有可能表示的否定。

原质量分室文稿讨论了局部八维立方体模型与六维 K3 极化子商之间的障碍。由于其结论依赖具体局部模型、配对及允许子商范畴，本综述不把``八维中没有六维子商''写成一般线性代数断言；统一稿的后续任务是逐一列出这些假设，给出与真实几何对的比较或反例。正文已严格使用的 K3 纯核来自边界 Chow 投影与相对长正合列，独立于该局部候选模型。


\endgroup

\begingroup
\providecommand{\Q}{\mathbb Q}\providecommand{\cC}{\mathcal C}
\providecommand{\cS}{\mathcal S}\providecommand{\dR}{\mathrm{dR}}
\providecommand{\rk}{\operatorname{rk}}\providecommand{\im}{\operatorname{im}}
\subsection{从边表进入几何：圈星类的结构规则}
\subsubsection{可终止的图级入口}
原\nolinkurl{conic_star_classification}定义任意圈数的可识别类：
在一个整数基本圈基中，有中心圈与$k$个叶圈，每圈恰有三条私有边，
另有中心分别与各叶共享的$k$条边；$L=k+1$、$E=4k+3$。
枚举生成树并检查基本圈支撑，是对该定义完备的有限识别。
至多检验$\binom{E}{V-1}$棵候选树，每项为有限图运算；
不声称任意图的多项式分类。

每圈三个私有割之差给两条线性方程。
秩二、剩余二次曲线光滑且有两个不同无穷远点时，
经明确有限代数扩域可参数化为$\ell_j=v_j+a_jt_j+b_j/t_j$。
置中心$t_0=z$，共享割给
\[
 F_i=A_i(z)t_i^2+B_i(z)t_i+C_i(z),\quad
 \deg A_i,\deg B_i,\deg C_i\leq2,
\]
\[
 y_i=2A_it_i+B_i,\quad y_i^2=\Delta_i(z)=B_i^2-4A_iC_i,\quad
 \deg\Delta_i\leq4.
\]
这由传播子支撑推导，未按图名预设椭圆答案。
秩、判别式及平方类条件都是明确的精确代数前提。

\subsubsection{紧割分解和分子选择}
先取平方类独立、四次平方自由且两两不相交的一般参数层。
紧化覆盖$C$的群为$H=(\mathbb Z/2)^k$，
$C_S:y_S^2=\prod_{i\in S}\Delta_i$。

\begin{proposition}[曲线核与标架选择]\label{unified:prop:star}
代数对应给出
\[
 H^1(C)\simeq\bigoplus_{\varnothing\ne S\subset[k]}H^1(C_S),
 \quad g(C_S)=2|S|-1,\quad g(C)=1+2^k(k-1).
\]
单位标量最大割仅选全字符：
\[
 \omega_{\rm sc}=c\,z^{k-1}dz/(y_1\cdots y_k),\quad c\ne0.
\]
其紧宿主亏格为$2k-1$。若$N=\sum_U n_U(z)y_U$，
则$n_Uy_U\omega_{\rm sc}$进入补集字符$[k]\setminus U$。
\end{proposition}
\begin{proof}
群幂等元$e_S=2^{-k}\sum_h\chi_S(h)h$为代数自同构图的有理线性组合。
商映射到$C_S$的拉回与迹复合为覆盖次数；除以该次数，
像即字符$S$部分，平凡字符$H^1(\mathbb P^1)$为零。
Riemann--Hurwitz对$4|S|$个分歧点给亏格，再求和。
共享留数将$dt_i/t_i$除以$F_i/(zt_i)$，得到$\pm z/y_i$；
乘中心$dz/z$即得标量形式。
分子规则来自$y_U/(y_1\cdots y_k)=1/y_{[k]\setminus U}$。
在一般层该标量形式非零且全纯。
\end{proof}
曲线群幂等元分解为标准机制\cite{KR}，
最大割亏格的相关先行工作见\cite{HHMZ}。
本节吸收原稿的传播子与留数证明，不主张经典公式优先权。

\subsubsection{宿主不自动等于最小对象}
若选中各$H^1(C_S)$在当前范畴中简单且两两不同构，
最小纯对象恰为链投影与读出同时非零的因子直和。
因为不等型简单直和的子对象按因子分解，链闭包先选非零链分量，
读出再删去零分量。这时才是精确分类，而不仅是亏格上界。

这些条件可在充分泛参数族中通过占优分支变化和单值群证明；
原稿论证使用完整独立离壳外源，不能套用每个指定物理切片。
本节不凭亏格重证其泛单值定理。
在等源特殊点，不同字符可成为同构对象，发生标架相消：
同一简单对象$A$的两份中，$(v,v)$与$(\lambda,-\lambda)$给零周期。
两份各自非零不足以保证总和可见；同型重数须单独处理。

分支碰撞时，先在代数闭常数域删除平方类依赖，得秩$r$。
每个非零几何类$v$有
\[
 b(v)=\#\{p:\sum_i v_i\operatorname{ord}_p\Delta_i=1\bmod2\},
 \qquad g_v=b(v)/2-1 ,
\]
其中包括无穷远。非零几何类在$\mathbb P^1$至少有两个分支点；
常数非平方仅为算术扭曲，不可算出几何亏格$-1$。
原域可按不可约因子次数加权，同时保留常数域与连通分量。
约去平方因子仍在微分中留下极点，因此相对边界扩张另记。
实际系数为数域时，因子分解与二元线性代数部分有限终止；
这不是指定Jacobian全部等源或一般动机的判定算法。

\subsubsection{补齐纯核之外的相对边界}
下面给接受曲线核后的实际混合对象，避免把紧亏格直接当作完整周期。
设$C/k$为光滑几何连通射影曲线，$D,B$为互不相交的有限约化点集，
$U=C\setminus D$，$d=\#D(\bar k)$、$b=\#B(\bar k)$。
所有点均带实际代数定义；必要时只作有限扩域。

\begin{proposition}[曲线相对对象的三层结构]\label{unified:prop:curve-boundary}
实际同调对象$E=H_1(U,B)$仅有如下权重分级：
\[
 \begin{split}
 \operatorname{Gr}_{-2}^W E
   &=\operatorname{coker}\{\Q(1)\to H_0(D)(1)\},\\
 \operatorname{Gr}_{-1}^W E&=H_1(C),\\
 \operatorname{Gr}_{0}^W E
   &=\ker\{H_0(B)\to\Q(0)\}.
 \end{split}
\]
空$D$或空$B$相应项记零，其维数为
\[
 \max(d-1,0),\quad 2g(C),\quad\max(b-1,0).
\]
原域上的端点项为Artin及其Tate扭转；点全部分裂后才写为Tate直和。
\end{proposition}
\begin{proof}
光滑曲线除点的局部化及其同调对偶给
\[
 H_2(C)\longrightarrow H_0(D)(1)\longrightarrow
 H_1(U)\longrightarrow H_1(C)\longrightarrow0,
\]
第一映射在几何点基中为对角映射。
相对同调正合段给
\[
 0\longrightarrow H_1(U)\longrightarrow H_1(U,B)
 \longrightarrow H_0(B)\longrightarrow H_0(U).
\]
$H_0(U)=\Q(0)$；各项纯权重及权重严格性给上述分级。
这些先是实际几何对的映射，Hodge实现只用于读取权重；
不从一个抽象分级矩阵推测扩张存在。
\end{proof}
因此在\emph{已构造的这一个曲线模型}内，
开放链和极点不会创造新的紧$H^1$因子，却可能增加
Abel--Jacobi、第三类微分及端点Kummer扩张。
其除子类、支点与路径决定粘合；仅列$b,d,g$不能决定这些扩张。
给定图若还需更高维积分、其他sector或重整化边界，
本命题不能将它们全部压缩成该曲线对象。

\subsubsection{非齐次源由端点给出：一个精确的微分分类规则}
在参数曲线$S$上考虑上述光滑射影曲线族与互不相交有限平坦$B,D$。
删去有限坏点，再经有限覆盖把点分裂为代数截面；
记函数域为$K$、导数为$\partial$。
令$\lambda\in H^1_{\dR}(U,B)$为实际相对形式，
其在$D$的全部留数为零。
相对类可以含必要的端点分量，移动端点不能在求导时冻结。

令$V\subset H^1_{\dR}(U,B)$为留数核。
由上两条正合列的对偶得到带连接的正合列
\[
 0\longrightarrow H^0_{\dR}(B)/K
 \xrightarrow{j}V\longrightarrow H^1_{\dR}(C)\longrightarrow0.
 \tag{C}
\]
空$B$时左项记零。
把$j(h)$的配对约定为$\sum_i n_i h(B_i)$，
其中$\partial\gamma=\sum_i n_i[B_i]$；在常用映射锥中$j(h)$由$(0,-h)$表示。
令$\bar\lambda$为$\lambda$在$H^1_{\dR}(C)$中的像，
$r$为$\bar\lambda$生成的循环连接模的维数，
$L_{\rm pure}\in K[\partial]$为对应首一$r$阶算子。

\begin{theorem}[端点源与一阶附加界]\label{unified:thm:endpoint}
对任意平坦相对链$\gamma$，$F=\langle\gamma,\lambda\rangle$满足
\[
 L_{\rm pure}F=a(s),\qquad
 a(s)=\sum_i n_i h_i(s)\in K,\qquad
 L_{\rm pure}\lambda=j(h).
\]
若$a=0$，齐次阶数至多$r$；若$a\ne0$，则
\[
 \bigl(\partial-a'/a\bigr)L_{\rm pure}F=0,
 \qquad \operatorname{ord}\le r+1.
 \tag{E}
\]
这里$a$由相对de Rham约化和代数端点求值决定，
不用先求出$F$的完整积分。
\end{theorem}
\begin{proof}
$\bar\lambda$满足$L_{\rm pure}\bar\lambda=0$。
式(C)为带连接的正合列，故$L_{\rm pure}\lambda$落在$j$的像。
点分裂后$H^0(B)$为带平凡连接的$K^b$，
其对平坦零链的配对是$K$中的代数函数。
配对与Gauss--Manin连接相容，所以得到非齐次式。
再对$a$施以$\partial-a'/a$即得(E)。
常数同时加到全部$h_i$不影响结果，因为$\sum_i n_i=0$。
\end{proof}
若用纤维上的恰当形式$df$实现一次约化，其端点贡献为
$f(B_{\rm end})-f(B_{\rm start})$；
一般移动端点还包含相对Cartan修正，均由实际连接确定。
只看紧曲线上的恰当性、把这些端点项置零会错误地降低阶数。
有非零留数时还存在第三类微分的Tate商，
必须先处理其连接；本定理不自动覆盖该情形。

\begin{corollary}[由割和端点精确决定阶数]\label{unified:cor:exact-order}
若$\bar\lambda$所在的$r$维纯连接不可约，
且物理变分已满足命题\ref{unified:prop:order}的$r$维割保护，
则$F$的最小齐次阶数恰为
\[
 \begin{cases}
 r,&a=0,\\
 r+1,&a\ne0.
 \end{cases}
\]
\end{corollary}
\begin{proof}
割保护给$r$阶下界。若$a=0$，由上界得等号。
若$a\ne0$而有$r$阶湮灭算子，则命题\ref{unified:prop:order}
的右整除使它为$L_{\rm pure}$的非零标量倍数，
与$L_{\rm pure}F=a\ne0$矛盾。式(E)给剩余上界。
\end{proof}
对于一般圈星标量割，形式为紧全字符曲线$C_{[k]}$上的非零全纯形式。
若其纯连接已在所用参数曲线上证明不可约，则$r=4k-2$。
实际闭链周期的源为零；实际开链表示且满足割保护时，
齐次阶数恰为$4k-2$或$4k-1$。
这是一条\emph{带实际曲线表示和物理链条件的精确规则}；
尚未构造完整费曼积分到该相对曲线的比较时，
只能将它用于割或已识别的曲线子sector。
同样的表述适用于椭圆$r=2$的闭/开Abel积分。
这些微分阶数不等于任意指定代数点的Nori最小秩。


\endgroup

\section{Banana的相对对象、扩张与最小性}\label{ch:core}

\begingroup

四传播子banana说明纯核怎样与边界和物理链共同决定完整积分。我们先构造实际相对对象，再计算分级与扩张，最后用标架确定最小子商。指定纤维所需有限输入集中列于第\ref{ch:comparison}节和附录\ref{app:inputs}。

\subsection{对象、参数与主结果}\label{sec:core:1}

令 k 为带指定实嵌入的特征零域，采用有理系数的 Nori 动机范畴。固定

\begin{equation}
\mu_i=m_i^2>0\quad(1\leq i\leq4),\qquad s<0,
\end{equation}

并置

\begin{equation}
U_G=abc+abd+acd+bcd,\qquad
 M=\mu_1a+\mu_2b+\mu_3c+\mu_4d,\qquad
 F_s=M U_G-sabcd.                                      \label{eq:core:1.1}
\end{equation}

本文的物理量是去除动量积分外部归一化后的参数积分

\begin{equation}
I(s)=\int_{\sigma_+}\frac{\Omega_{\mathbb P^3}}{F_s},
 \qquad \sigma_+=\{[a:b:c:d]:a,b,c,d\geq0\}.             \label{eq:core:1.2}
\end{equation}

积分形式和链采用使欧氏积分为正的相容取向。

记 \(\mathbb T=\mathbb P^3\setminus\{abcd=0\}\)、 \(U_s=\{F_s=0\}\cap\mathbb T\)。令 \(X_s\) 是相应节点四次曲面的最小解消，D 为其二十条有理边界曲线。以 \(H_{\rm res}\) 表示从 \(H^2(X_s)\) 剔除边界除子所张成十六维空间后的实际纯因子。暂不把它写成完整超越因子；后者还要求 \(\rho(X_s)=16\)。置

\begin{equation}
A=H_{\rm res}(3),\qquad B=\mathbb Q(1),\qquad
 \kappa=\dim_{\mathbb Q}\langle[\mu_1/\mu_4],
 [\mu_2/\mu_4],[\mu_3/\mu_4]\rangle\subset k^\times\otimes\mathbb Q.
                                                               \label{eq:core:1.3}
\end{equation}

\begin{theorem}[实际对象及比较]

\label{thm:core:A} 上述模型给出秩十五的有效对象 \(E_{\rm env}=H^3_{\rm N}(\mathbb T,U_s)(3)\)。其最大一动机商是

\begin{equation}
K_{\rm bd}=\left[\mathbb Z\longrightarrow\mathbb G_m^8,\quad
 1\longmapsto(\mu_1/\mu_4,\mu_2/\mu_4,\mu_3/\mu_4,1,1,1,1,1)\right]_{\mathbb Q},
                                                               \label{eq:core:1.4}
\end{equation}

核实际同构于 A。拉回该点生成的最小子环面给出

\begin{equation}
0\subset A\subset N\subset E=E_{\rm act},\qquad
 \operatorname{Gr}^W E=(A,B^\kappa,\mathbf1),\qquad
 \operatorname{rank}E=7+\kappa.                           \label{eq:core:1.5}
\end{equation}

正物理链属于 \(E_B\)，物理读出属于 \(F^3E_{\rm dR}^{\vee}\)。下层扩张的各个边界方向等于相应高阶 Chow 循环的投影 Nori regulator。这里一动机的同调余字符基取负的留数基；第\ref{sec:core:4}节固定这个符号。

\end{theorem}

\begin{theorem}[最小化与幺幂群]

\label{thm:core:B} 对第\ref{sec:core:6}节所列三步扩张假设，令 \(n=\dim A>1\)、\(D_A=\operatorname{End}(A)\)，下层扩张的 \(D_A\) 秩为 \(r>0\)，有理关系空间维数为 d。则

\begin{equation}
\dim R_uG(E)=nr+\kappa+n,\qquad
 \operatorname{End}(E)=\mathbb Q.                         \label{eq:core:1.6}
\end{equation}

对顶层非零的链和在 A 上非零的物理 Hodge 读出，最小对象是 \(\bar E=E/\widetilde B_0\)，其中 \(\widetilde B_0\simeq B^d\) 是下层关系方向的唯一 Tate 提升。特别地

\begin{equation}
\operatorname{rank}\bar E=n+\kappa+1-d,
 \qquad 1\longrightarrow\mathbb V(B^d)
 \longrightarrow G(E)\longrightarrow G(\bar E)\longrightarrow1. \label{eq:core:1.7}
\end{equation}

核位于幺幂根基的中心；在整个群中，它仍按 Tate 角色受共轭作用。不可分解性和框架最小性因而不是同一条件。

\end{theorem}

\begin{theorem}[banana 的两个适用范围]

\label{thm:core:C}

\begin{itemize}
\tightlist
\item
  对 \(m=(1,2,4,8),s=-64,k=\mathbb Q\)，使用第\ref{ch:comparison}节列出的几何及认证输入，\(\kappa=r=1,d=0\)。物理形式周期的最小秩为八，Hodge/Nori 幺幂根基均为十三维；其分级群分别为 \(SO(q)\times\mathbb G_m\)、\(O(q)\times\mathbb G_m\)。
\item
  在四参数族很一般正质量欧氏点，\(\kappa=r=3,d=0\)。最小秩为十，两个幺幂根基均为二十七维，群维数均为四十三。此处并未统一判定 Nori 群的行列式分支。
\item
  在等质量一参数族很一般欧氏点，先剔除三个额外代数方向，得到纯秩三核被顶层扩张的秩四最小对象。其深度为一，Hodge/Nori群均为三维向量群与三维正交相似群的半直积，总维数七，见第\ref{sec:core:8.1}节。
\end{itemize}

``很一般''指排除可数个真解析条件；它不保证某个给定代数点满足结论。全文讨论形式周期对象及其群，不断言复数求值之间的代数独立性。

\end{theorem}


\subsection{正链、有效性与相对对偶}\label{sec:core:2}

对一个排列 p，取 Hepp 扇区坐标

\begin{equation}
\alpha_{p_0}=1,\quad\alpha_{p_1}=t_1,\quad
 \alpha_{p_2}=t_1t_2,\quad\alpha_{p_3}=t_1t_2t_3,
 \qquad 0\leq t_i\leq1.
\end{equation}

记 \(q=t_1t_2t_3\)、\(V=1+t_3+t_2t_3+q\)。直接代入得

\begin{equation}
F_s=t_1^2t_2\,(M_pV-sq),\qquad
 \Omega_{\mathbb P^3}=t_1^2t_2\,dt_1dt_2dt_3,
 \qquad M_pV-sq\geq\mu_{p_0}>0.                        \label{eq:core:2.1}
\end{equation}

因此二十四个扇区上的形式延伸且可积，正链不接触严格变换的图超曲面。由同一估计

\begin{equation}
0<I(s)\leq6\sum_{i=1}^4\mu_i^{-1}.       \label{eq:core:2.2}
\end{equation}

这是几何收敛与正性的证明，不需要先求数值积分。

依次处理坐标顶点、坐标线和额外边节点，并作边界解消，得到光滑射影三维 P，使 X 与 Y 的并为正常交叉除子；X 是图超曲面一侧，Y 是坐标边界一侧，\(Z=X\cap Y\)。解消不改变 \(P\setminus Y\simeq\mathbb T\) 与 \(X\setminus Z\simeq U_s\)。局部形式沿吹起中心的阶数是 \(c- 1- m\)：普通节点给 \(3- 1- 2=0\)，边线给 \(2- 1- 1=0\)；结合\eqref{eq:core:2.1}得到物理链及读出的相对类。

关键不是仅给出 Hodge 空间同构，而是实际构造

\begin{equation}
H^3_{\rm N}(P\setminus Y,X\setminus Z)(3)
       \simeq H^{\rm N}_3(P\setminus X,Y\setminus Z).   \label{eq:core:2.3}
\end{equation}

在正常交叉局部乘积上，交换两种开延拓的自然交换态射在 Betti 实现下为同构。Nori 六运算提供该态射与 Verdier 对偶；忠实正合 Betti 实现使其锥为零，故\eqref{eq:core:2.3}是实际同构。右侧是有效相对同调对象。这说明为什么可以使用有效 Nori 一动机商定理，而不是从左侧的 Tate 扭转外观推断有效性。

这一标准相对对偶的拓扑形式见 Bloch--Kerr--Vanhove §6.1，实际六运算相容性采用 Tubach Theorem 6.5/Corollary 6.8。\cite{bkv,tubach} 这里使用的是上述正常交叉模型中的交换态射和Betti实现的保守性。


\subsection{从边界复形得到权重与实际纯核}\label{sec:core:3}

Hessian 坐标

\begin{equation}
y_i=\mu_i\alpha_i\ (i\leq4),\qquad y_5=-M,\qquad
 \sum_{i=1}^5y_i=0
\end{equation}

把四次曲面写成

\begin{equation}
\sum_{i=1}^5\lambda_i\prod_{j\ne i}y_j=0,
 \qquad (\lambda_1,\ldots,\lambda_5)=(\mu_1,\ldots,\mu_4,s).
                                                               \label{eq:core:3.1}
\end{equation}

其边界由 \(V_i,E_{ij},C_{ij},F_i\) 共二十条有理曲线组成。各自平方为\(- 2\)；非零横截交点为

\begin{equation}
C_{ij}\cdot V_i=C_{ij}\cdot V_j=C_{ij}\cdot E_{ij}=1,
 \qquad E_{ij}\cdot F_k=E_{ij}\cdot F_l=1,
 \quad\{k,l\}=\{1,2,3,4\}\setminus\{i,j\}.             \label{eq:core:3.2}
\end{equation}

对偶图 \(\Gamma \) 连通，有二十顶点、三十边，故 \(b_1(\Gamma)=11\)。四个单位 a/d、b/d、c/d、M/d 的主除子给出四个独立关系。由\eqref{eq:core:3.2}选取十六条非退化边界类，配对行列式\(- 48\)，得到边界张成空间的秩恰十六，而非只知道其 Gram 秩。

对正常交叉边界的 Gysin 复形，曲线的奇上同调为零，遂有

\begin{equation}
0\longrightarrow H_{\rm res}\longrightarrow H^2(U_s)
 \longrightarrow H_1(\Gamma,\mathbb Q)(-2)\longrightarrow0,
 \qquad \dim H^2(U_s)=6+11=17.                          \label{eq:core:3.3}
\end{equation}

\(H^2(\mathbb T)=\mathbb Q(-2)^3\) 的双留数矩阵 Z 在三条指定旗标上有单位子矩阵，因而注入\eqref{eq:core:3.3}的右端。又 \(H^3(U_s)=0\)，因为 \(U_s\) 是光滑仿射曲面。相对长正合列给

\begin{equation}
0\to N_{\rm env}\to E_{\rm env}\to\mathbf1\to0,
 \qquad 0\to A\to N_{\rm env}\to B^8\to0.              \label{eq:core:3.4}
\end{equation}

为把 A 识别为实际动机，取上述十六条边界类 \(D_i\)、Gram 矩阵 G 和有理系数次数一零循环。边界投影子

\begin{equation}
\pi_D=\sum_{i,j}(G^{-1})_{ij}[D_i\times D_j]
\end{equation}

满足 \(\pi_D^2=\pi_D\)，因为复合系数是 \(G^{-1}GG^{-1}\)。剔除次数零、四及边界投影子，得到 \(\pi_{\rm res}\)。它在 Nori 实现中切出 \(H_{\rm res}\)。该因子经限制到 \(U_s\)、再经相对连接映入\eqref{eq:core:3.4}：限制只杀死边界类，而连接的核来自纯权重四的 \(H^2(\mathbb T)\)，与纯权重二因子无交。于是这是一条实际单射，不是事后猜测的 Hodge 同构。

环境权重分级秩为6、8、1，对应权重\(- 4\)、\(- 2\)、0。纯核的 Hodge 型为

\begin{equation}
(-1,-3),\quad(-2,-2)^4,\quad(-3,-1).
\end{equation}

所以活动对象的过滤维数为

\begin{equation}
(\dim F^0,\dim F^{-1},\dim F^{-2},\dim F^{-3})
                  =(1,\kappa+2,\kappa+6,\kappa+7).       \label{eq:core:3.5}
\end{equation}

这是同一几何复形的结论；不能从``相对上同调''四字直接推出预设的三步过滤。完整相对复形还在其他次数有上同调，例如 \(H^2(\mathbb T,U_s)=\mathbb Q(-1)\)。


\subsection{质量 Kummer 商与物理链的活动化}\label{sec:core:4}

由\eqref{eq:core:2.3}及 Ayoub--Barbieri-Viale Theorem 8.12，\(E_{\mathrm{env}}\) 有实际最大一动机商。其 Hodge 实现的核恰为\eqref{eq:core:3.4}的权重\(- 4\)部分。上一节的实际 A 映入此核，二者同秩六，故实际同构。还需确定一动机的点；秩和 Hodge 型本身不提供\eqref{eq:core:1.4}。

取有序符号

\begin{equation}
\xi=\{a/d,b/d,c/d\}\in H^3_{\mathcal M}(U_s,\mathbb Q(3)).
\end{equation}

沿二十条曲线及其交点作两次局部化。顶点---线和边---线旗标的贡献由 Steinberg 关系消去；余下十二条 E---F 旗标的值都是 \([\mu_i]-[\mu_j]\)。例如

\begin{equation}
E_{12}\to F_3:[\mu_1]-[\mu_2],\qquad
 E_{12}\to F_4:-[\mu_1]+[\mu_2].                         \label{eq:core:4.1}
\end{equation}

这些值属于图循环空间 \(\ker\partial\)，模去 Z 的三维像后，质量矩阵秩三。以其前三个独立质量列完成八维留数基 P，坐标是

\begin{equation}
\bigl([\mu_1/\mu_4],[\mu_2/\mu_4],[\mu_3/\mu_4],0^5\bigr).
                                                               \label{eq:core:4.2}
\end{equation}

这里需把符号的局部化与相对对象连接起来。用有序三个单位的杯积给顶层 \(\mathbf1\to R\Gamma(\mathbb T)(3)[3]\)，限制到 \(U_s\) 后取同伦纤维。由于 \(H^3(U_s)=0\)，其第零上同调是由 \(H^2(U_s)(3)\) 扩张顶层的对象；再商去来自 \(H^2(\mathbb T)(3)\) 的三维部分，正是 \(E_{\mathrm{env}}\)。把该构造推出到每个双留数坐标，得到由\eqref{eq:core:4.2}表示的 Kummer 扩张。这条实际态射链同时说明为何不能仅由对数值反推动机比较。

在 P 留数基中使用 \(\nu=e_F-e_B\)，Kummer 类的坐标是正的 \(\log(\mu_i/\mu_4)\)。同调一动机 \([\mathbb Z\to\mathbb G_m,1\mapsto u]\) 的对应约定给 \(-\log u\)，故其余字符基是 \(- P\)；在这个基中得到\eqref{eq:core:1.4}。质量矩阵的完整旗标数据作为附录\ref{app:inputs}所列几何输入保存。

三个单位生成的最小子环面有维数 \(\kappa \)。推出到补商后，一动机扩张分裂，且 \(\operatorname{Hom}(\mathbf1,B)=0\) 使分裂唯一。拉回该活动一动机即得 E，且

\begin{equation}
E_{\rm env}/E\simeq B^{8-\kappa}.    \label{eq:core:4.3}
\end{equation}

所有边界曲线与交点均在 k 上分裂；\eqref{eq:core:4.3}不是根据 Hodge 型排除 Artin 扭曲。复共轭在这个 Tate 商上为\(- 1\)，在正实链上为+1，故物理链在商中的像为零，因而确实属于 E。


\subsection{下层扩张是投影 regulator}\label{sec:core:5}

对图循环 \(c=(c_e)\in H_1(\Gamma,\mathbb Q)\)，在每条有理曲线 \(D_i\) 上取有理函数 \(f_i\)，使

\begin{equation}
\operatorname{div}(f_i)=-\sum_e\partial_{i,e}c_e\,p_e.
                                                               \label{eq:core:5.1}
\end{equation}

图循环条件使右端次数为零；有理曲线上任意此类除子可由函数实现。清分母后，每个交点的两侧系数抵消，故

\begin{equation}
Z_c=\sum_i(D_i,f_i)\in CH^2(X_s,1)_{\mathbb Q}. 
                                                               \label{eq:core:5.2}
\end{equation}

换 \(f_i\) 的常数会增加边界除子乘单位的可分解类，其在 \(\pi_{\rm res}\) 下为零。

\begin{proposition}

\label{thm:core:5.1} 令 \(\eta\in\operatorname{Ext}^1(B^8,A)\) 为\eqref{eq:core:3.4}的下层类。在 P 留数约定下，边界方向 c 的拉回满足

\begin{equation}
(c^*\eta)(-1)
                 =r_{{\rm N},{\rm res}}(Z_c)
       \quad\text{于 }\operatorname{Ext}^1_{\rm N}
                  (\mathbf1,H_{\rm res}(2)).             \label{eq:core:5.3}
\end{equation}

其中先把 c 模去 Z，或者使用与该商相容的代表。

\end{proposition}

\begin{proof}

置 \(\mathcal S=R\Gamma_D(X_s)(2)[3]\)。支持的 Mayer--Vietoris 与纯性给 \(H^0\mathcal S=H_1(\Gamma,\mathbb Q)\)；其相邻次数还含曲线和曲面的 Tate 项，不能删去整个复形。式\eqref{eq:core:5.1}在动机支持复形中给出 \(\alpha_c:\mathbf1\to\mathcal S\)。从支持到整体的态射，再经 Chow 投影，得到

\begin{equation}
g:\mathcal S\longrightarrow H_{\rm res}(2)[1].
\end{equation}

\(H^0\operatorname{Fib}(g)\) 是图循环空间被 \(H_{\rm res}(2)\) 扩张的对象。局部化三角给从 \(H^2(U_s)(2)\) 到它的态射，在两端均为恒等，故是同构。沿 \(\alpha_c\) 拉回时，扩张类正是复合 \(g\alpha_c\)；支持类忘记支持后为\eqref{eq:core:5.2}，所以同一复合又是其投影 regulator。最后商去三项环面二次杯积并作 Tate 扭转，得到\eqref{eq:core:5.3}。

\end{proof}

这个比较使用实际支持态射，不需要从任意 Hodge 扩张中提升动机态射。它还给出三个由全局单位杯积产生的 Tate 直和因子；这属于环境对象，不能与下面活动对象的不可分解性混淆。

因为 \(F^0N=0\)，顶层1有唯一Hodge提升 \(e_F\in F^0E\)。若 \(e_B\) 是有理提升，则

\begin{equation}
\nu_E=[e_F-e_B]\in N_{\mathbb C}/N_{\mathbb Q}
       =\operatorname{Ext}^1_{\rm MHS}(\mathbf1,N).
\end{equation}

它向 \(B^\kappa_{\mathbb C}/B^\kappa_{\mathbb Q}\) 的投影是已算出的Kummer类，但纯核方向仍有自由度。固定N及该边界类、且有提升时，提升类形成 \(\operatorname{Ext}^1_{\rm MHS}(\mathbf1,A)\) 的挠子；这由\(\operatorname{Hom}(1,B)=0\)及Ext长正合列得到。这一Hodge描述不自动提供每个提升的动机来源。实际E选定其中一个提升；取 \(e_B=\sigma\) 时，物理Hodge读出消去 \(e_F\)，故 \(I=-\langle\omega,\nu_E\rangle\)。

相邻扩张 \(\eta\in\operatorname{Ext}^1(B^\kappa,A)\)、 \(\xi\in\operatorname{Ext}^1(\mathbf1,B^\kappa)\) 来自实际 E，故 \(\eta\circ\xi=0\)。理由是 E 已将 \(\xi \) 提升到 N 的扩张。这个零乘积不决定提升的全部比较条目，也不等于宣称所有 Ext\textsuperscript{2} 消失。


\subsection{三步扩张的群：两种秩必须分开}\label{sec:core:6}

本节在有理中性 Tannakian 范畴中证明。假设：A 简单，\(\dim A=n>1\)；B 可逆；三块权重分别为\(- 4\)、\(- 2\)、0；分级生成群 R 约化；\(S=A\otimes B^{-1}\) 与 B 非同构。上层 \(\kappa \) 个扩张类 \(\xi _i\) 有理独立。令

\begin{equation}
D_A=\operatorname{End}(A),\quad
 r=\dim_{D_A}\langle\eta_1,\ldots,\eta_\kappa\rangle,
 \quad V_0=\ker\bigl[L_\eta:\mathbb Q^\kappa\to
             \operatorname{Ext}^1(B,A)\bigr],\quad d=\dim V_0.  \label{eq:core:6.1}
\end{equation}

除环按推出作用于 Ext，使用左系数。以下假定 \(r>0\)。

权重过滤在原范畴中由实际子对象给出，分级群带按这些权重作用的中央权重余特征；本文的Nori/Hodge实现满足这一条件。第\ref{sec:core:7}节另使用相容且严格的Hodge实现。所有最小化均在原有理范畴中进行；de Rham系数域扩大后，不能把任意新增不变子空间当成已经下降的子对象。

选 Levi 因子及相容分裂。下层幺幂像是 \(X\subset S^\kappa\)。对 \(a_i\in D_A\)，有

\begin{equation}
\sum_i a_i\eta_i=0
        \quad\Longleftrightarrow\quad
       \left.\sum_i a_ix_i\right|_X=0.                   \label{eq:core:6.2}
\end{equation}

正向来自分裂扩张；反向中，相应推出扩张的幺幂作用消失，表示降到约化分级群，因而半单并分裂。由于 R 约化、S 简单，X 是若干个 S 的直和；\eqref{eq:core:6.2}的消去关系确定其重数恰为 r。因此 \(X\simeq S^r\)、\(\dim X=nr\)。同理，上层的全部像为 \(Y=B^\kappa\)。

S 与 B 非同构，排除了 X 与 Y 的对角耦合，所以联合像是 \(X\oplus Y\)。严格三角矩阵的交换子为

\begin{equation}
[(x,y,z),(x',y',z')]=(0,0,xy'-x'y).                     \label{eq:core:6.3}
\end{equation}

X 非零、Y 为全部中间层，评价像 \(X\otimes Y\to A\) 非零；A 简单迫使像为整个 A。因而实际幺幂群为

\begin{equation}
U=\left\{\begin{pmatrix}I_A&x&z\\0&I_\kappa&y\\0&0&1\end{pmatrix}
           :x\in X,\ y\in B^\kappa,\ z\in A\right\}.    \label{eq:core:6.4}
\end{equation}

特别地

\begin{equation}
\dim U=nr+\kappa+n,\qquad [U,U]=A,\qquad
 Z(U)=A\oplus(B\otimes V_0).                            \label{eq:core:6.5}
\end{equation}

中心的 x 块必须与所有 y 对易，故 \(x=0\)；其 y 必须被全部 X 消去，由\eqref{eq:core:6.2}的有理关系恰等于 \(B\otimes V_{0}\)。

最大幺幂群的条件是 \(r=\kappa \)。若 \(D_A=\mathbb Q\)，它等价于 \(d=0\)；一般不能互换。例如非零 \(\eta \) 与 \(\alpha\eta\)，\(\alpha\in D_A\setminus\mathbb Q\)，有理独立而在 \(D_A\) 上只张成一维。这个例子说明两种秩的差异，不额外声称任意指定三步几何扩张存在。

一般混合扩张的 Tannakian 框架可比较 Eskandari。\cite{eskandari} 上述相邻像与交换子论证在本文中承担明确假设下的群等号；不以形式上允许这些矩阵作为证明。


\subsection{不可见中心与最小商：进一步的理论结论}\label{sec:core:7}

保留第\ref{sec:core:6}节假设。令 \(\sigma \in E_B\) 在顶层非零，\(\omega\in F^3E_{\rm dR}^\vee\) 且 \(\omega|A\ne0\)。不同实现中的这些条件通过实际子对象比较，不把 Betti 与 de Rham 向量擅自当作同一基下的向量。

本节以同调一侧记最小对象 \(\bar E\)；通常形式周期记号中的最小上同调对象是 \(\bar E^\vee\)，二者秩相同。

\begin{theorem}

\label{thm:core:7.1} 定理\ref{thm:core:B} 的全部结论成立；还满足

\begin{equation}
\operatorname{End}(\bar E)=\mathbb Q,\qquad
 \dim U(\bar E)=nr+(\kappa-d)+n.                         \label{eq:core:7.1}
\end{equation}

最小化只消去幺幂中心中的 \(B\otimes V_{0}\) 方向，保持纯核 A、交换子 \([U,U]=A\) 和下层 \(D_A\) 秩 r。

\end{theorem}

\begin{proof}

首先，\(\eta|_{B\otimes V_0}=0\)，因此 \(B\otimes V_{0}\) 有进入 N 的截面。两个截面之差属于 \(\operatorname{Hom}(B\otimes V_0,A)=0\)，故截面唯一，记像为 \(\widetilde B_0\)。其 Hodge 型为\((- 1,- 1)\)，所以 \(\omega \) 消去它。

其次，\(\sigma \) 生成整个 E。事实上，\eqref{eq:core:6.4}的全部 z 平移作用于顶层非零的 \(\sigma \)，产生 A；再用全部 y 平移，模 A 得到 \(B^\kappa \)，最后保留顶层便得到 E。这一生成性在 \(r>0\) 的群已识别后，不再需要数值正性来证明。

设 \(J\subset E\) 被 \(\omega \) 消去。若 \(J\cap A\ne 0\)，则 A 简单给 \(A\subset J\)，与 \(\omega|A\ne0\) 矛盾。若 J 在顶层非零，同一生成性论证给 \(J=E\)，仍矛盾。所以 \(J\subset N\) 且 \(J\cap A=0\)。其在 \(B^\kappa \) 中的像 B' 给下层扩张沿 B' 的分裂，故 \(B'\subset B\otimes V_{0}\)。唯一提升性质于是给 \(J\subset\widetilde B_0\)。反之该 Tate 提升已被 \(\omega \) 消去，故它正是最大不可见子对象。最小矩阵系数对象因此是 \(\bar E=E/\widetilde B_0\)，得到秩公式。

选 V 的有理补空间 \(V=V_0\oplus V_1\)。X 消去 V\textsubscript{0}，所以在商对象上，x 和 z 完全保留，y 只保留其在 \(B\otimes V_{1}\) 中的像。因 \(r>0\)，\(\dim V_1=\kappa-d>0\)，商对象的分级仍同时包含 A 与 B；分级表示仍忠实于原 R。故 \(G(E)\to G(\bar E)\) 的核没有额外约化部分，其矩阵恰为

\begin{equation}
\begin{pmatrix}I_A&0&0\\0&I_\kappa&y_0\\0&0&1\end{pmatrix},
                 \qquad y_0\in B\otimes V_0.             \label{eq:core:7.2}
\end{equation}

它作为 U 的子群是中心的。由于上层到中间层的共轭角色是 B，它通常不是整个 G 的中心。这证明\eqref{eq:core:1.7}及\eqref{eq:core:7.1}的群维数式。补空间用于显示矩阵；核与商本身由 \(\widetilde B_0\) 规范确定。

最后证明 \(\operatorname{End}(E)=\mathbb Q\)，这一步并不要求 X 最大。与所选 Levi 中的权重环面交换，任意自同态是 \(\operatorname{diag}(a,h,\lambda )\)。与所有 z 交换给 \(a=\lambda I_A\)；与所有 y 交换给 \(h=\lambda I_\kappa \)。故只有标量。商对象仍满足相同 \(r>0\)、全部中心和全部上层条件，同理 \(\operatorname{End}(\bar E)=\mathbb Q\)。

\end{proof}

\begin{corollary}

\label{thm:core:7.2} 若 \(d>0\)，\(\widetilde B_0\subset E\) 是非零不可见子对象，却不可能是直和因子；否则存在非平凡幂等元，与 \(\operatorname{End}(E)=\mathbb Q\) 矛盾。因此``找不到直和分解''不足以证明物理最小性。

\end{corollary}

\begin{corollary}

\label{thm:core:7.3} 这些物理形式周期在最小对象上的幺幂深度仍恰为二，不要求\(r=\kappa \)。确实，三次严格降层的作用必为零。另一方面，所有X的像之和是A：它是非零R子表示，A简单。故对非零 \(\alpha=\omega|A\)，存在\(x\in X\)及\(y\in B^\kappa \)使\(\alpha xy\ne 0\)。对偶表示中先作用x、再作用y，得到非零顶层对偶坐标。\(\sigma \)生成E使相应形式系数不消失，从而深度至少二。商掉 \(\widetilde B_0\) 不改变这个复合，因为全部X都消去该子对象。

因此幺幂深度二、\(\operatorname{End}(E)=\mathbb Q\)、最小对象秩以及U是否最大，是需要分别判断的四个性质。

\end{corollary}

\subsubsection{下层全部分裂的两种结局}\label{sec:core:7.4}

若\(r=0\)于当前实际范畴，\(\operatorname{Hom}(B,A)=0\)使下层\(N=A\oplus B^\kappa \)的分裂唯一。商去唯一B提升，得到Z及规范顶到底类\(\zeta \)；同时\(E=Z\times _{1}K\)。若\(\zeta \ne 0\)，简单性给\(U(E)=A\oplus B^\kappa \)，整个群的z、y作用使\(\sigma \)生成E，最大不可见部分恰\(B^\kappa \)，最小对象Z的秩n+1、深度一。E仍只有标量自同态。

若\(\zeta =0\)，则\(E=A\oplus K\)，物理F\textsuperscript{3}读出消去K。在\(I>0\)的物理框架中，链的A分量非零，最小对象为A、秩n、深度零；\(\operatorname{End}(E)=\operatorname{End}(A)\times Q\)。深度零是纯对象承载，不是代数数判据。\(\kappa =0\)也按\(\zeta \)这两支处理，不能套用\(r>0\)的定理\ref{thm:core:7.1}。

这里还有群核的细节：\(\kappa >0\)时删掉全部B可能遗失一个Tate根。令\(T_{0}=\operatorname{ker}[G(A\oplus B)\to G(A)]\)，则到最小对象群的核是\(B^\kappa \)向量群被T\textsubscript{0}扩张；选Levi后为半直积。极化\(A\otimes A\to B^{4}\)只保证\(T_{0}\subset \mu _{4}\)。六维分级形式\((H,t)=(t^{2}g,t)\)、\(g\in SO\)或O的这个核实际是\(\mu _{4}\)。\(r>0\)时留下至少一个B，故定理\ref{thm:core:7.1}没有这种额外有限核。

\subsubsection{额外除子方向的实际下降}\label{sec:core:7.5}

对Picard跳跃纤维，取有限Galois扩域上完整NS除子基及非退化Gram矩阵\(G_{\mathrm{NS}}\)，构造\(\sum_{i,j}(G_{\mathrm{NS}}^{-1})_{ij}[D_i\times D_j]\)。K3的\(Pic^{0}=0\)使该对应Galois不变；corestriction除以扩域次数给k上投影。限制在有理Chow群中单射，故投影恒等式也下降。

令\(\pi _{\mathrm{alg}}=\pi _{\mathrm{NS}}- \pi _D\)、\(\pi _T=\pi _{2}- \pi _{\mathrm{NS}}\)。由边界空间非退化且包含在NS中，两投影正交，给实际分解\(H_{\mathrm{res}}(3)=T(X)(3)\oplus C_{\mathrm{alg}}\)，其中\(\dim C_{\mathrm{alg}}=\rho - 16\)。\(C_{\mathrm{alg}}\)为Artin表示扭转Q(2)，不必在k上分裂成Tate直和。其对偶全为(2,2)型，故被物理F\textsuperscript{3}形式消去。

实际商\(E_T=E/C_{\mathrm{alg}}\)保存同一物理矩阵系数，纯核变为简单的T(X)(3)，秩\(n=22- \rho \)。将\(\eta \)投向这一纯核后即可应用第\ref{sec:core:6}--\ref{sec:core:7.4}节。这个构造解决理论上的纯核与下降问题，不自动算出任意纤维的\(\rho \)或regulator。


\subsection{纯理论的很一般非消失与秩十}\label{sec:core:8}

令 \(\rho(X_s)=16\) 暂作为结论待得。四个独立 Hessian 参数给四维 K3 周期域。局部 Torelli 的适用在本族可具体核对：十个节点的射影位置固定；若一阶抽象变形平凡，\(H^1(\mathcal O_X)=0\) 使极化线丛唯一延伸，四个截面恢复四次曲面。节点的例外\((- 2)\)曲线刚性迫使诱导射影变换固定节点。四个 \([e_i-e_5]\) 与其余节点强迫该变换为标量；\eqref{eq:core:3.1}中五个单项式独立，参数仅剩共同缩放。因此四个有效参数的 Kodaira--Spencer 映射单射。局部Torelli及本族的模参数验证也见第\ref{ch:deriv}节。

型 IV 域维数也是四，故周期像局部开。每个非零有理向量成为额外(1,1)类的条件是真超平面条件；每个非标量有理算子保持全纯周期线的条件也是真解析条件。后者若含开集，将迫使非退化不可约二次曲面落在有限个真特征子空间之并，矛盾。排除可数并后，\(\rho=16\) 且 \(\operatorname{End}_{\rm Hdg}T=\mathbb Q\)。

在仿射坐标 \(\alpha=(a,b,c,1)\) 中，设

\begin{equation}
f=(\textstyle\sum_i\mu_i\alpha_i)(\sum_i\alpha_i^{-1}),
 \quad\omega=\operatorname{Res}\frac{d\log a\wedge d\log b\wedge d\log c}{f-s},
 \quad\beta_i=-s\frac{\mu_i\alpha_i}{M}\omega.
\end{equation}

在商去来自 \(H^2(\mathbb T)\) 的杯积之后，实际 Cartan 计算给下层对象中的恒等式

\begin{equation}
\sum_i\beta_i=-s\omega,\qquad
             \nabla_s\beta_i=\mu_i\nabla_{\mu_i}\omega. \label{eq:core:8.1}
\end{equation}

\(\beta _i\) 是第 i 质量留数方向的 Hodge 提升。对非零 \(c\in\mathbb Q^3\)，令 \(\beta_c=\sum_{i=1}^3c_i\beta_i\)。若其下层扩张在参数开集有平坦有理 Hodge 截面 q，则因纯层 \(F^2\) 一维，必有 \(q=\beta _c+h\omega \)。沿 s 求导并模 \(F^2\)，\eqref{eq:core:8.1}要求质量方向 \((c_1\mu_1,c_2\mu_2,c_3\mu_3,0,0)\) 与 s 方向、共同缩放在周期切空间中相关。上述局部 Torelli 排除了这一点。

故每个非零有理 c 都不能在非空开集分裂。单连通标记图表中的有理平坦提升只有可数多个；每个成为 Hodge 截面的条件是真解析条件。对 c、提升和图表取可数并，得到很一般点三个下层扩张有理独立。加上 \(\operatorname{End}_{\mathrm{Hdg}} T=Q\)，此时 \(D_A\) 秩也为三。

上层独立性另由 \(\prod_i(\mu_i/\mu_4)^{a_i}=1\) 的可数个真关系超曲面排除得到。复真解析条件不可能含正质量欧氏实开集，否则由全纯唯一性在复开集恒成立。于是同一很一般实参数集合上 \(\kappa =r=3\)、\(d=0\)；定理\ref{thm:core:B} 给 U27 与最小秩十。

Zarhin 的 K3 群定理\cite{zarhin}给权重零纯结构 \(S=T(1)\) 的 Hodge 群 SO(q)。Tate 角色与该半单群独立，故分级群为 \(SO(q)\times \mathbb G_m\)，维数16。Nori 群夹在 Hodge 下界与极化上界之间，两个幺幂根基均已达最大，群维数均为43。不能据此决定每个基域上的行列式分支，也未计算整个秩十绝对比较矩阵。

\subsubsection{等质量子族：最小秩四、深度一与七维群}\label{sec:core:8.1}

归一化\(\mu _i=1\)后，\(t=s/m^{2}<0\)，Kummer活动秩\(\kappa =0\)。BKV §3.2给出的Verrill族一般Picard数为十九，且有十九维除子局部子系统。\cite{bkv} 这项已有几何输入与第\ref{sec:core:7.5}节结合，把原秩七活动对象商成

\begin{equation}
0\to T(X_t)(3)\to Z_t\to\mathbf1\to0,
                    \qquad \operatorname{rank}Z_t=4.
\end{equation}

在局部平坦的除子商上，唯一顶层Hodge提升满足\(\nabla _t e_F=\delta \omega _t\)。非零全纯二形式不被NS投影杀死，故\(e_F\)不在任何非空开集平坦。每个有理平坦提升q与\(e_F\)相等的点集离散；提升及图表可数，所以Hodge分裂只可能落在可数例外集。其外\(\zeta _H\ne 0\)，进而\(\zeta _N\ne 0\)；第\ref{sec:core:7.4}节给最小秩四、深度一。没有把这个很一般结论套用任意指定代数t。

秩三K3超越结构的\(\operatorname{End}_{\mathrm{Hdg}}\)只能是Q：Zarhin 的全实/CM 域分类\cite{zarhin}排除奇次数CM；三次全实域的每个特征空间一维，容不下同一实嵌入下独立的H\textsuperscript{20}和H\textsuperscript{02}。故纯Hodge群含SO\textsubscript{3}，连同权重标量是GO\textsubscript{3}。对三维\(H^{\mathsf t}qH=\lambda q\)，\(c=detH/\lambda \)、\(g=c^{-1}H\)给\(GO(q)=\mathbb G_m\times SO(q)\)，为连通四维群。

实际杯积极化给相同Nori上界，因此纯群亦为GO\textsubscript{3}；\(\zeta _H\)非零使两种幺幂核均为整个三维\(T(X_t)(3)\)。于是选Levi后

\begin{equation}
G_H(Z_t)=G_N(Z_t)=\mathbb V(A_T)\rtimes GO(q),
                    \qquad\dim G=7.
\end{equation}

其Hodge过滤维数为(1,2,3,4)，比较块\((P_A,f\);0,1)的结构式为\(\Delta f=P_A\otimes z+f\otimes 1\)。该结论不计算新的绝对周期矩阵，也不恢复一个不存在的质量对数中间层。外部族几何的输入范围保留上述很一般参数条件。



\endgroup

\begingroup
\providecommand{\Q}{\mathbb Q}
\providecommand{\C}{\mathbb C}
\providecommand{\dR}{\mathrm{dR}}
\providecommand{\id}{\operatorname{id}}
\providecommand{\im}{\operatorname{im}}
\providecommand{\End}{\operatorname{End}}
\providecommand{\Hom}{\operatorname{Hom}}
\providecommand{\Vis}{\operatorname{Vis}}
\providecommand{\soc}{\operatorname{soc}}
\providecommand{\head}{\operatorname{head}}
\providecommand{\rank}{\operatorname{rank}}
\subsection{head 与 socle：保护和最小性的不同问题}
\subsubsection{单源保护只约束非零商}
在有限长度范畴中，$\head R$ 表示最大半单商，
$\soc R$ 表示所有简单子对象之和。对固定有限对象可转到其有限维表示代数，
二者均存在。

\begin{proposition}[范畴版单源 head 判据]\label{ops:prop:head}
设 $R=R_\sigma$，$S$ 为简单对象。
每个保留 $\sigma$ 的非零商都含有 $S$ 作为组成因子，
当且仅当 $\head R$ 的全部简单因子都同构于 $S$。
\end{proposition}
\begin{proof}
一个非零商中的 $\sigma$ 不可能为零，否则商映射的核承载 $\sigma$，
与 $R=R_\sigma$ 矛盾。非零有限长度商有简单商；
它也是 $R$ 的简单商，所以正向保护由 head 的单类型条件得到。
若 head 含 $S'\not\simeq S$，投到 $S'$ 即给出不含 $S$ 的非零商。
\end{proof}

这里量化所有实际商，尚未证明每个商都可由允许物理分子实现。
因此这是充分保护判据及所述\emph{逻辑量化}下的必要条件，
不是任意指定物理读出族的无条件充要分类。

\subsubsection{读出是否完整，由简单子对象检测}
\begin{proposition}[socle 检测]\label{ops:prop:socle}
设 $R=R_\Sigma$。最小化不再损失子对象，即
$\Vis(R;\Sigma,\Omega)=R$，当且仅当不存在非零简单子对象
$S\subset R$ 被全部读出消去。
特别地，若 $\soc R=A$ 是简单对象，则
\begin{equation}\label{ops:eq:socle}
 \Vis(R;\Sigma,\Omega)=R
 \quad\Longleftrightarrow\quad \Omega|A_{\dR}\ne0.
\end{equation}
\end{proposition}
\begin{proof}
任一非零不可见子对象在有限长度范畴中包含简单子对象。
反之，一个被消去的简单子对象已经属于最大不可见核。
若 socle 简单，则每个非零子对象都含 $A$，所以 \eqref{ops:eq:socle} 立即成立。
\end{proof}

非半单对象中，简单纯核可以处于 socle，而顶层 Tate 商处于 head。
``一个输出非零''与``输出在纯核上的限制非零''此时强度完全不同。
这个区别正是下面 banana 应用的关键。


\subsection{指定 banana 对象的全部子对象}
\subsubsection{已建立的几何输入与适用范围}
本节采用已有证明稿给出的实际同调对象 $E$，基域为 $\Q$，
固定 $m=(1,2,4,8)$、$s=-64$。它具有实际过滤
\begin{equation}\label{ops:eq:banana}
 0\subset A\subset N\subset E,
 \qquad \dim A=6,\quad N/A=B=\Q(1),\quad E/N=\mathbf1.
\end{equation}
$A$ 是相应扭转的简单 K3 超越对象。
已有相邻扩张非零证明和实际群论证给出完整幺幂根基
\begin{equation}\label{ops:eq:U}
 U=\left\{
 \begin{pmatrix}I_A&x&z\\0&1&y\\0&0&1\end{pmatrix}
 :x\in\Hom(B,A),\ y\in B,\ z\in A\right\}.
\end{equation}
这是在一个 Levi 分裂中显示的式子；$A,N,E$ 本身不依赖该分裂。
纯分级作用在 $A$ 上不可约。下述证明只需要这些输入，
不再使用新的 Frobenius、周期或路径计算，也不需要重新判别固定纯群的分支。

指定点的扩张非零与固定纯群输入见第\ref{ch:comparison}节；完整幺幂根基由第\ref{thm:core:B}定理给出。
一般三步对象若缺少 \eqref{ops:eq:U} 或纯核简单性，不能套用以下分类。

\begin{theorem}[全部实际子对象]\label{ops:thm:uniserial}
对象 $E$ 的实际子对象恰好为
\begin{equation}\label{ops:eq:chain}
                        0,\ A,\ N,\ E.
\end{equation}
因此 $\soc E=A$，$\head E=\mathbf1$。
\end{theorem}
\begin{proof}
令 $J\subset E$ 非零，在 \eqref{ops:eq:U} 的分裂中检查其 Betti 实现。
若 $J$ 在顶层非零，取 $v=(a,b,c)\in J_B$ 且 $c\ne0$。
仅改变 $z$ 的群元素给出 $(u_z-1)v=(cz,0,0)$，故 $A_B\subset J_B$。
仅改变 $y$ 给 $(u_y-1)v=(0,cy,0)$，故整个中间 $B$ 方向也在 $J_B$ 中。
于是 $N_B\subset J_B$，加上原顶层向量即得 $J=E$。

若顶层为零但 $J$ 在 $N/A$ 中非零，取 $v=(a,b,0)$ 且 $b\ne0$。
由全部 $x$ 得 $(u_x-1)v=(xb,0,0)$，所以 $A_B\subset J_B$，
再加该中间向量得 $J=N$。
最后，若 $J\subset A$，则纯核的简单性给 $J=A$。
上述包含通过忠实正合实现反映回实际范畴，四个候选均已实际存在。
\end{proof}

这个证明没有把不可分解等同于简单。
$E$ 有非平凡子对象 $A,N$，但没有额外直和投影。
由 \eqref{ops:eq:chain} 和扩张的非分裂，非零动机自同态必须同时尊重各层；
已有群作用进一步给 $\End(E)=\Q$。


\subsection{完整双标架分类与分子选择规则}
\subsubsection{链的闭包与读出的核}
对任意 $\sigma\in E_B$，定理\ref{ops:thm:uniserial}立即给
\begin{equation}\label{ops:eq:Rlevels}
 R_\sigma=
 \begin{cases}
 0,&\sigma=0,\\
 A,&0\ne\sigma\in A_B,\\
 N,&\sigma\in N_B\setminus A_B,\\
 E,&\sigma\notin N_B.
 \end{cases}
\end{equation}
对任意 $\lambda\in E_{\dR}^{\vee}$，有
\begin{equation}\label{ops:eq:Ilevels}
 I_\lambda=
 \begin{cases}
 0,&\lambda|A_{\dR}\ne0,\\
 A,&\lambda|A_{\dR}=0,\ \lambda|N_{\dR}\ne0,\\
 N,&\lambda|N_{\dR}=0,\ \lambda\ne0,\\
 E,&\lambda=0.
 \end{cases}
\end{equation}
这里均指向量或限制映射的精确非零，不是十进制求值不为零。

\begin{theorem}[所有双标架的最小对象]\label{ops:thm:allframes}
记 $K=E/A$，它是非分裂的秩二质量 Kummer 对象。
非零输入的完整分类为下表；括号中为秩。
\begin{center}
\small
\begin{tabular}{c|cccc}
\toprule
链的范畴闭包&$\lambda|A\ne0$
 &$\lambda|A=0,\lambda|N\ne0$
 &$\lambda|N=0,\lambda\ne0$&$\lambda=0$\\
\midrule
$A$&$A\ (6)$&$0$&$0$&$0$\\
$N$&$N\ (7)$&$B\ (1)$&$0$&$0$\\
$E$&$E\ (8)$&$K\ (2)$&$\mathbf1\ (1)$&$0$\\
\bottomrule
\end{tabular}
\end{center}
各项都是实际最小对象；在通常形式周期记法中取它的对偶。
\end{theorem}
\begin{proof}
将 \eqref{ops:eq:Rlevels}、\eqref{ops:eq:Ilevels} 代入
$R_\sigma/(R_\sigma\cap I_\lambda)$。
对象识别分别使用 $N/A=B$、$E/A=K$ 和 $E/N=\mathbf1$。
没有剩余未知子对象，故该表既给上界也给最小性。
\end{proof}

\subsubsection{固定物理链的精确秩分层}
原正物理链 $\sigma_{\mathrm{phys}}$ 在 $E/N$ 的坐标已标定为 $1$，
所以其范畴闭包为 $E$。
现在固定它，仅允许改变实际 de Rham 读出。
\begin{corollary}[固定链的四个层]\label{ops:cor:numerators}
设 $V$ 为一个有限维 $\Q$-空间，代表已经映入同一对象的允许分子类；
设线性映射 $L:V\to E_{\dR}^{\vee}$ 给其读出。
定义
\begin{equation}\label{ops:eq:exceptional}
 Z_A=\ker(\operatorname{res}_A L),\qquad
 Z_N=\ker(\operatorname{res}_N L),\qquad Z_E=\ker L.
\end{equation}
则 $Z_E\subset Z_N\subset Z_A\subset V$，且
\begin{equation}\label{ops:eq:rankstrata}
 \rank\Vis(E;\sigma_{\mathrm{phys}},L(v))=
 \begin{cases}
 8,&v\notin Z_A,\\
 2,&v\in Z_A\setminus Z_N,\\
 1,&v\in Z_N\setminus Z_E,\\
 0,&v\in Z_E.
 \end{cases}
\end{equation}
某些层可以为空。该分类不要求先求出任一个完整积分。
\end{corollary}
\begin{proof}
直接应用定理\ref{ops:thm:allframes}的最后一行。
每个例外集合都是限制映射的核，故是有理线性子空间。
若允许读出为空间 $E_{\dR}^{\vee}$ 本身，则
$\dim Z_A=2,\dim Z_N=1,\dim Z_E=0$；
对于物理给定子空间，只取这些空间的逆像，不能默认其全满。
\end{proof}

本推论的前提是分子已给出\emph{同一相对对象}上的类。
若分子增加极点除子、要求新增 sector，或改变边界及链，须先扩大或重构对象。
尤其不能仅按多项式次数将任意分子自动放入这里的八维空间。
多分子线性组合的相消已经由 $L$ 的核精确记录；
积分后一个值恰为零，并不意味着 $v\in Z_E$。

\subsubsection{K3 可见为何迫使混合扩张一起保留}
在保留原链时，$\lambda|A\ne0$ 不仅保护 K3 简单因子，
还使整个 $E$ 最小。这是因为任一非零子对象都包含 $A$。
不能先商去中间层而保留一个同样带原链的纯 K3 商；
\eqref{ops:eq:chain}中没有这种商。

相反，最大割把链变成 $A_B$ 中的非零向量，并使用相应非零纯读出。
定理\ref{ops:thm:allframes}第一行给纯秩六对象。
下层周期取 $N_B\setminus A_B$ 的链，则可能得到秩七。
因此 $6$、$7$、$8$ 的区别来自完整双标架，不能归结为环境中是否存在一个 K3 曲面。

同时，$\head E=\mathbf1$ 表明：仅凭一个未知读出的非零性，
在实际动机范畴中只可强制顶层简单类型，不能单独强制 K3。
顶层商读出 $E\to\mathbf1$ 已给非零但不含 K3 的反例。
原切割单值模的 head 为 K3 型时仍可使用保护定理；
那是另一个对象、另一组链和一个不同的作用范畴，不存在矛盾。


\subsection{幺幂深度与实际割见证}
\subsubsection{三个非零层的深度}
\begin{proposition}\label{ops:prop:depth}
对固定顶层非零链，\eqref{ops:eq:rankstrata}的三个非零层分别具有
形式周期的幺幂深度 $2,1,0$。
\end{proposition}
\begin{proof}
在 $E$ 上，$u-1$ 严格上三角，所以任意三个这样的因子的乘积为零。
当 $\lambda|A\ne0$ 时，选 $x,y$ 使 $\lambda(xy\sigma_{\mathrm{top}})\ne0$；
这是可能的，因为 $\sigma_{\mathrm{top}}\ne0$ 且
$\Hom(B,A)\otimes B\to A$ 的评价满射。
于是一个二重幺幂变化非零，深度恰为二。
此计算在比较扭子的任意平凡化上进行，故是形式余作用的断言。

在商 $K=E/A$ 上，幺幂作用只剩 $y$，其二重变化全为零。
读出在 $B$ 上非零且链顶层非零，故一重变化非零，深度一。
再商为 $\mathbf1$ 后幺幂作用平凡，深度零。
\end{proof}

深度不是数值超越性标签。秩六的纯 K3 对象同样深度零，
而秩一的 Tate 周期也深度零。深度计量扩张中的幺幂方向，
纯核类型必须单独保存。本文也不由深度二推断函数是某种长度二的 MPL。

\subsubsection{用真实割证明读出在纯核上非零}
\begin{proposition}[实际割的充分见证]\label{ops:prop:cutwitness}
设一个实际态射 $f:C\to E$ 的像恰为 $A$，
并有实际割链 $\gamma\in C_B$。
若 $\langle f^*\lambda,\gamma\rangle\ne0$，
则 $\lambda|A\ne0$。
若原链在 $E/N$ 非零，则原标架的最小对象为完整 $E$。
\end{proposition}
\begin{proof}
若 $\lambda|A=0$，则 $f^*\lambda=0$，所有该类割周期都为零，矛盾。
再用定理\ref{ops:thm:allframes}即可。
从一个实际数值\emph{非零}推出形式或向量非零不需要周期猜想；
危险的是反方向，把数值零当成形式零。
\end{proof}

必须证明 $f$ 的像为纯核，而不能仅因某个正的相空间积分非零就假定这一点。
若割对象还含可见边界项，非零也可能完全来自该边界。
此外，函数在一个参数域内不恒零，不保证在任一指定代数点仍非零。
对指定点使用命题\ref{ops:prop:cutwitness}时，应提交该点的见证，
或直接以代数 de Rham 限制证明 $\lambda|A\ne0$。

对于本文已给定的原物理读出，该限制非零来自既有全纯形式与纯核比较，
不需要为本命题再做一次数值周期计算。新分子仍须提交其实际类及限制。


\endgroup

\section{绝对比较、消失球格与固定纯群}\label{ch:comparison}\label{ch:lattice}\label{ch:fixed}\label{ch:coaction}

\begingroup


\subsection{比较矩阵承担的数学任务}
纯 Gauss--Manin 方程给出水平运输；绝对比较矩阵还要求真实 Betti 循环、de Rham 框架及其尺度。没有几何初值时，对基本解矩阵右乘任意常矩阵仍满足同一个系统。因此，恢复微分方程、认证全局单值和识别绝对周期是不同步骤。

固定 $m=(1,2,4,8)$，纯极化系统的有限节点为 $1,9,25,49,81,121,169,225$，另有零点的幺幂退化。规定上半平面路径将消失球送到共同基点。局部反射不能直接相乘求共同基点迹；必须先用同一框架作运输或同时共轭。零点存在共振，不能把未经修正的原留数直接指数化。

现有证书以精确有理中心、Frobenius 残差及逐段包围实现这一过程。论文所需的是其数学界，而不是运行日志。若 $P$ 是近似水平解，$\mathcal R=P'-AP$，则真实误差满足变常数公式；用已证明的传播范数界控制积分项，同时计入初值误差和舍入误差。矩阵求逆另需以残差小于一的界保证可逆性，而非直接使用数值逆。


\subsection{裸留数、Morse 球与 Tate 归一化}
在射影环面图上取
\begin{equation}
 \omega_s^{\mathrm{abs}}=\Res_{f=s}
 \frac{d\log a\wedge d\log b\wedge d\log c}{f-s},\qquad
 f=M\left(\frac1a+\frac1b+\frac1c+1\right).
\end{equation}
节点附近的非退化二次首项提供消失球周期的首项尺度。这里需要同时固定平方根支、球取向和 Poincaré 留数约定。现有 Morse 定标给出相应规范中的
\begin{equation}
 \rho=512\sqrt2\,\pi(1+i),\qquad
 P_D^{\mathsf t}K_DP_D=8\pi^2Q_6.
\end{equation}
因而在本文标准 Tate 约定下，真实代数杯积是公开配对矩阵的 $-2$ 倍。框架中的纯对偶末项应相应乘 $-1/2$；忽略这一变化会使物理读出和比较行列式同时出现尺度错误。

已认证的裸留数六循环矩阵 $P_6$ 采用明确有理 de Rham 框架，在 $s=-64$ 的统一误差小于 $2^{-85}$。原六循环为指数二子格的有理基；后节证明饱和后，右乘相应有理换基矩阵得到完整整 Betti 基的周期矩阵，其误差分别小于 $2^{-84}$ 和标准整基中的 $2^{-81}$。纯整基的结论不自动给出混合对象的完整整基。


\subsection{下层扩张的实分裂障碍}
设 $V_c$ 是未扭转下层七维对象。纯周期之外，还需选定质量方向的 $\beta$ 周期。实际消失球在节点收缩为零给出混合初值；沿已有纯解的积分产生混合条目。这个零初值是几何事实，不能用方便的积分常数替代。

在实 Hodge 分裂的假设下，相关纯与混合周期满足一个实线性相容关系。对三个指定球取 $\operatorname{Re}p,\operatorname{Im}p,\operatorname{Im}b$ 所成的三阶矩阵，其行列式必须为零。现有证书在节点 $1,9,25$ 的对应循环上得到
\begin{equation}
 8.77<\det\begin{pmatrix}
 \Re p_1&\Im p_1&\Im b_1\\
 \Re p_2&\Im p_2&\Im b_2\\
 \Re p_3&\Im p_3&\Im b_3
 \end{pmatrix}<8.78,
\end{equation}
误差小于 $1.2\cdot10^{-27}$。这里 $p_i,b_i$ 依已有几何框架定义；改变归一化必须同步改变判据。该障碍证明指定纤维的实 Hodge 扩张非零，实际 Nori 扩张若分裂则其 Hodge 实现分裂，故 Nori 扩张亦非零。反方向一般不成立。

这一证据与 $\log2\ne0$ 所给的上层 Kummer 非零相互独立。两者共同输入三步扩张定理，得到 $r=\kappa=1,d=0$、物理最小秩八及十三维幺幂根基。不能只计算上层质量对数便宣称完整混合对象最小。


\subsection{第六混合周期与边界 Tate 列}
第六球为 $\delta_6=T_0\delta_1$。其纯周期可由纯 $T_0$ 作用得到，混合周期还必须包括绕零的积分类增量。以既有局部原函数行 $j_*$ 和框架 $W$ 记，完整公式具有
\begin{equation}
 b_6=b_1+j_*W^{-1}(v_6-v_1)
\end{equation}
的形式。它使用零点解析规范和全混合连接，而非仅将纯矩阵作用到一个拼接向量。

边界第七列由双留数为 $+1$ 的小整环面给出：纯列为零，$\beta$ 周期为 $(2\pi i)^2$。这与六个纯循环一起确定完整 $P_7$。现有 $49$ 项比较矩阵误差小于 $2^{-85}$，第六 $\beta$ 条目误差小于 $2^{-88}$。该矩阵是 $V_c$ 的有理 Betti 比较；物理下层为 $N=V_c(3)$。


\subsection{顶层相对膜、物理链与八维比较}
顶层的实际 Cartan 公式为
\begin{equation}
 \nabla_s e=\delta\omega,\qquad
 \mu_i\nabla_{\mu_i}e=\delta\beta_i.
\end{equation}
六个相对消失膜的收缩零初值固定纯核方向的顶层比较。边界环面先作环境像校正，再取与实共轭相容的填充；由此固定 Kummer 项的符号
\begin{equation}
 q_7=-2(2\pi i)^2\log2.
\end{equation}
正物理链标记还需要膜支撑避让及有序三环面交数，不能只根据数值正性猜测链坐标。

令 $z=2\pi i$，采用已经核对的相对对偶整体符号，完整比较为
\begin{equation}
 C_E=z^{-3}P_8^{\mathsf t},\qquad C_M=C_E^{-\mathsf t}=z^3P_8^{-1},
 \qquad M=E^\vee.
\end{equation}
$P_8$ 的 $64$ 项在明确有理基下有误差小于 $2^{-85}$ 的认证。若 $\nu=e_F-\sigma$，几何配对 $C_{\mathrm{geom}}(\omega,w)=-2$ 给出
\begin{equation}
 I=-2\nu_w,\qquad I^{\mathfrak m}=2(C_M^{\mathfrak m})_{8,6},\qquad
 \Delta I^{\mathfrak m}=2\sum_{j=1}^8(C_M^{\mathfrak m})_{8,j}\otimes(g_M)_{j,6}.
\end{equation}
最后一式是选定框架中的实际八项 coaction；它不声称每一项都是一个未加额外标记的图割。指定欧氏积分满足
\begin{equation}
 0.73554702266148<I(-64)<0.73554702266149.
\end{equation}
完整比较不等于所有条目具有经典特殊函数闭式，也不等于已经认证完整混合整 Betti 基。这里的比较以几何膜、实填充和相对对偶符号为输入；其来源和使用位置列于附录\ref{app:inputs}。



\endgroup

\subsection*{消失球格}

\begingroup

本节固定 $m=(1,2,4,8),s=-64$。已认证的六循环 Gram 为下式；前五个循环来自节点 $1,9,25,49,81$，第六个为 $T_0\delta_1$。
\begin{equation}Q_6=\begin{pmatrix}
-2&-2&-2&4&-2&-2\\-2&-2&0&2&0&-4\\-2&0&-2&2&0&-4\\
4&2&2&-2&0&8\\-2&0&0&0&-2&-4\\-2&-4&-4&8&-4&-2
\end{pmatrix}.\label{eq:gram}\end{equation}
设 $R_i=T_i-I$。Picard--Lefschetz 公式给
\begin{equation}\operatorname{tr}(R_iR_j)=q_{ij}^2,\qquad
\operatorname{tr}(R_iR_jR_k)=q_{ij}q_{jk}q_{ki}.\end{equation}
十五个对迹依字典序为
\begin{equation}(4,4,16,4,4,0,4,0,16,4,0,16,0,64,16).\end{equation}
六个三迹按 $(124),(126),(134),(136),(146),(156)$ 的顺序为
\begin{equation}(-16,-16,-16,-16,-64,-16).\end{equation}
拓扑先给整性，严格误差不超过 $2^{-69}$ 的区间再唯一识别这些整数。取 $A_0=T_0-I$，由 $A_0^2=0$ 及正交性得 $A_0^*=-A_0$，从而 $q_{16}=\delta_1^2=-2$；第六球取向并非独立任意选择。

矩阵的秩为六、签名 $(2,4)$、行列式 $192$。完整超越格的既有几何证明给
\begin{equation}T\simeq U\oplus U(2)\oplus A_2(-2),\qquad |\det T|=48.\end{equation}
因此 $L=\sum\Z\delta_i$ 满足 $[T:L]^2=192/48=4$。下文证明实际饱和类唯一为 $w=(\delta_1+\delta_2+\delta_3+\delta_4)/2$。这里完整 $T$ 的输入依赖 Hessian 格和本原性证明，不能仅由秩与行列式推出\cite{dvgeemen}。


\subsection{二阶判别式型的不变量}\label{sec:lattice:2}

对非退化偶整格 K，记

\begin{equation}
A_K=K^\vee/K,\qquad
 q_K(\bar x)=(x,x)\pmod{2\mathbb Z}.
\end{equation}

考虑性质

\begin{equation}
\mathcal P(K):\quad q_K(A_K[2])\subset\mathbb Z/2\mathbb Z.       \label{eq:lattice:2}
\end{equation}

它只要求二阶挠元上的二次型值为整数模2；不要求整个判别群的值都为整数。特别是本例的三阶部分不在断言中。

式 \eqref{eq:lattice:2} 在整格同构下不变。选 Gram 矩阵 G 后，每个二阶挠元唯一写为 \(\varepsilon/2\) 模 K，其中 \(\varepsilon\in\{0,1\}^{\operatorname{rk}K}\)，并满足

\begin{equation}
G\varepsilon\equiv0\pmod2,\qquad
 q_K(\varepsilon/2)=\frac{\varepsilon^{\mathsf t}G\varepsilon}{4}
                      \pmod{2\mathbb Z}.                         \label{eq:lattice:3}
\end{equation}

所以这是至多 \(2^{\operatorname{rk}K}\) 次整系数判定，不需要数值运算或格同构搜索。

\begin{lemma}

\(\mathcal P(T_{\rm std})\) 成立。

\end{lemma}

\begin{proof}

U 是幺模格，没有非零判别元。U(2) 的 Gram 为 \(\left(\begin{smallmatrix}0&2\\2&0\end{smallmatrix}\right)\)，故 \(q((a,b)/2)=ab\)，为整数。\(A2(- 2)\) 的 Gram 为 \(\left(\begin{smallmatrix}-4&2\\2&-4\end{smallmatrix}\right)\)，其所有二阶类由 \((c,d)/2\) 表示，二次型值

\begin{equation}
-c^2+cd-d^2
\end{equation}

也为整数。直和给出 \eqref{eq:lattice:2}。其二阶子群共16个元素。

\end{proof}


\subsection{穷尽所有指数二补全}\label{sec:lattice:3}

由 \([T:L]=2\)，有 \(2T\subset L\)。所以任一可能的 T 都形如

\begin{equation}
K_v=L+\mathbb Z(v/2),\qquad
 0\ne v\in\{0,1\}^6.                                           \label{eq:lattice:4}
\end{equation}

整配对与偶性分别等价于

\begin{equation}
Q_6v\equiv0\pmod2,\qquad v^{\mathsf t}Q_6v\equiv0\pmod8.         \label{eq:lattice:5}
\end{equation}

第一条件因 Q6 所有条目为偶数而自动满足。六十三个非零 v 中，恰十一种满足第二条件。

对每个 v，令 \(p(v)\) 为首个非零坐标，取基矩阵 \(B_v\) 为单位阵把第 p(v) 列替换为 v/2，令 \(G_v=B_v^{\mathsf t}Q_6B_v\)。以下 \(\varepsilon=e_j\) 均是这个特定 \(B_v\) 基中的坐标；并不是原六循环基的第 j 个半球。

{\def\LTcaptype{none} % do not increment counter
\begin{longtable}[]{@{}
  >{\raggedright\arraybackslash}p{(\linewidth - 6\tabcolsep) * \real{0.2500}}
  >{\raggedright\arraybackslash}p{(\linewidth - 6\tabcolsep) * \real{0.2500}}
  >{\raggedright\arraybackslash}p{(\linewidth - 6\tabcolsep) * \real{0.2500}}
  >{\raggedright\arraybackslash}p{(\linewidth - 6\tabcolsep) * \real{0.2500}}@{}}
\toprule\noalign{}
\begin{minipage}[b]{\linewidth}\raggedright
v（按六个坐标连写）
\end{minipage} & \begin{minipage}[b]{\linewidth}\raggedright
判别元见证 \(\varepsilon\)
\end{minipage} & \begin{minipage}[b]{\linewidth}\raggedright
\(G_v\varepsilon\) 模2
\end{minipage} & \begin{minipage}[b]{\linewidth}\raggedright
\(q_{K_v}(\varepsilon/2)\) 模2
\end{minipage} \\
\midrule\noalign{}
\endhead
\bottomrule\noalign{}
\endlastfoot
001100 & e6 & 0 & 3/2 \\
010100 & e6 & 0 & 3/2 \\
011011 & e4 & 0 & 3/2 \\
011101 & e5 & 0 & 3/2 \\
011110 & e6 & 0 & 3/2 \\
100001 & e6 & 0 & 3/2 \\
100010 & e5 & 0 & 3/2 \\
100111 & e6 & 0 & 3/2 \\
101000 & e3 & 0 & 3/2 \\
110000 & e2 & 0 & 3/2 \\
111100 & 无障碍；16个二阶值全部整数 & --- & 0 或 1 \\
\end{longtable}
}

前十行违反 \(\mathcal P(T)\)，不可能是实际 T。实际补全必然存在，故唯一剩余 \(v=(1,1,1,1,0,0)\) 就是其类。这证明了上文所述实际饱和半循环的标记。

计算还显示，被排除的每个候选都有8个非整数二阶值；证明只需表中一个见证。所有十一候选的行列式均为48，说明单用行列式确实看不出这个区别。


\subsection{标记整格和标准基}\label{sec:lattice:4}

以 \(\mathcal B\) 为整基，实际 Gram 为

\begin{equation}
G=
\begin{pmatrix}
0&-1&-1&3&-1&-1\\
-1&-2&0&2&0&-4\\
-1&0&-2&2&0&-4\\
3&2&2&-2&0&8\\
-1&0&0&0&-2&-4\\
-1&-4&-4&8&-4&-2
\end{pmatrix}.                                                   \label{eq:lattice:6}
\end{equation}

取

\begin{equation}
R=\begin{pmatrix}
1&1&4&4&6&4\\
0&-1&1&-1&-1&-1\\
0&0&-2&-1&-2&-2\\
0&0&-1&-1&-2&-1\\
0&0&0&0&-1&1\\
0&0&-2&-1&-2&-1
\end{pmatrix}.
\end{equation}

精确计算给

\begin{equation}
\det R=1,\qquad
 R^{\mathsf t}GR=
 \begin{pmatrix}0&1\\1&0\end{pmatrix}
 \oplus\begin{pmatrix}0&2\\2&0\end{pmatrix}
 \oplus\begin{pmatrix}-4&2\\2&-4\end{pmatrix}.                       \label{eq:lattice:7}
\end{equation}

在确定实际输入整基后，式 \eqref{eq:lattice:7} 的精确等式说明\(\mathcal B R\)是规定几何循环标记下的标准整基。

这一步证明的是整同调类的归属和基，不另外画出 w 的嵌入曲面代表，也不需要这样的代表来验证整格等式。



\endgroup

\subsection*{固定纤维的群}

\begingroup

本节固定 $X=X_{(1,2,4,8),-64}/\Q$，并以第\ref{app:inputs}节所列认证输入中的 $\rho=16$、绝对比较和 regulator 非零为输入。群均明确区分固定纤维与参数族；所有结论保留基域 $\Q$。

\subsection{固定纤维的有限输入}\label{sec:fixed:1}

使用附录\ref{app:inputs}所列良约化输入：\(p=11\) 良约化且几何Picard数为16。其残余六维 Frobenius 多项式为

\begin{equation}
P(Z)=Z^6+10Z^5+55Z^4-484Z^3+6655Z^2+146410Z+1771561.             \label{eq:fixed:3}
\end{equation}

已有有限域点数为 \(288,16568,1794960\)，残余幂迹 \((-10,-10,2102)\)。由这些点数及 Newton 恒等式得到下列纯因子的多项式。

P 模5满足完整六次 Rabin 不可约性判据，故在 \(\mathbb Q\) 上不可约。归一化 \(11^{-6}P(11Z)\) 的系数不全为整数，故 P 没有11倍单位根根。另有

\begin{equation}
P(Z)\equiv Z^5(Z-1)\pmod {11},\qquad P'(1)\equiv1\pmod {11}.     \label{eq:fixed:4}
\end{equation}

因此 P 在 \(\overline{\mathbb Q}_{11}\) 中恰有一个单位根，此处``单位根''指 \textbf{11-adic unit root}，不是有限阶的 root of unity。下文称它为 p-adic 单位特征根，以避免混淆。

实际六循环形成有理基。它们的格虽未饱和，但

\begin{equation}
\det q\equiv\det Q_6=192\equiv3
 \quad\text{于 }\mathbb Q^\times/\mathbb Q^{\times2}.             \label{eq:fixed:5}
\end{equation}

指数二只改变判别式的平方因子，不影响式 \eqref{eq:fixed:5}。


\subsection{排除实乘}\label{sec:fixed:2}

令 \(F=\operatorname{End}_{\rm Hdg}T\)。K3 型的标准定理给出：F 是全实域或 CM 域，其对极化的伴随分别为恒等或复共轭，且 \([F:\mathbb Q]\mid6\)。采用 \href{https://www.imo.universite-paris-saclay.fr/~francois.charles/Picard.pdf}{Charles §2.2 所列 Zarhin 定理}，其

先说明 F 全实时 \(\dim_FT\ge3\)，以免遗漏三次域。

\begin{itemize}
\tightlist
\item
  若 \(\dim_FT=1\)，区别嵌入对应的复特征空间只有一维，却必须同时容纳 \(H^{2,0}\) 与 \(H^{0,2}\)，矛盾。
\item
  若 \(\dim_FT=2\)，写极化为一个 F 上二元对称型的迹，F 上对角化为 \(\langle a,b\rangle\)。每个实嵌入处该二元型均定号：区别嵌入为正定，其余为负定。矩阵
\end{itemize}

\begin{equation}
J=\begin{pmatrix}0&-b\\a&0\end{pmatrix},\qquad J^2=-ab
\end{equation}

在区别嵌入处保持两条复各向同性线，即 \(H^{2,0}\)、\(H^{0,2}\)；其余嵌入全部为 (1,1)。故 J 是有理 Hodge 自同态。又因 \(-ab\) 全负，J 不在全实标量域 F 中，矛盾。

所以非平凡全实 F 唯一可能为二次域，\(\dim_FT=3\) 为奇数。Charles Theorem 1(2) 在 \(p=11\) 给出

\begin{equation}
\rho(X_{\overline{\mathbb F}_{11}})\ge16+2=18,
\end{equation}

与已认证的16矛盾。原定理的特征至少5条件满足。此处采用的是固定良约化，而非很一般参数的特化。

亦可用其 Lemma 16 的群表述理解这个矛盾：某个归一化 Frobenius 幂属于三个 F 维度的特殊正交群，必须有特征值1；这会给 P 的11倍单位根特征值，与第\ref{sec:fixed:1}节矛盾。

因此 F 全实时只能为 \(\mathbb Q\)。


\subsection{\texorpdfstring{CM 候选必含 \(\mathbb Q(\sqrt{-3})\)}{CM 候选必含 \textbackslash mathbb Q(\textbackslash sqrt\{-3\})}}\label{sec:fixed:3}

F 若为 CM 域，其次数只能是2或6。必须谨慎处理次数6的非 Galois 可能性，不能未经证明假定所有 CM 域都含虚二次子域。

\href{https://wwwuser.gwdguser.de/~jjahnel/Arbeiten/Komponentengruppe_v4f.pdf}{Elsenhans--Jahnel Theorem 7.5} 证明：若 \(X\times X\) 的 Hodge 猜想成立，基域 K 与端同态域的合成 KF 在 K 上正规。同文 Theorem 7.7 的证明明确援引 CM 情形下该 Hodge 猜想已知成立；这里 CM 按其 §5.1 指端同态域为 CM 域，不要求 \(\dim_FT=1\)。

仅在当前反证的 CM 情形应用此已知结论。因为 \(X/\mathbb Q\)，F 在 \(\mathbb Q\) 上 Galois。次数6的 Galois CM 域具有中央复共轭，故其六阶 Galois 群不能是中心平凡的 \(S_3\)，只能为 \(C_6\)。其三阶子群固定一个虚二次域 \(K_0\)。次数2时直接取 \(K_0=F\)。

写 \(K_0=\mathbb Q(\sqrt{-d})\)，d 为正无平方因子整数。T 是三维 \(K_0\) 向量空间，伴随在 \(K_0\) 上为复共轭，故极化是一个 Hermitian 型的迹。Hermitian 型可在 \(K_0\) 上对角化，三个对角系数 \(h_j\) 均在 \(\mathbb Q^\times\) 中。相对于每个 \(K_0\) 基向量及其 \(\sqrt{-d}\) 倍，迹型的二阶块为

\begin{equation}
\begin{pmatrix}2h_j&0\\0&2dh_j\end{pmatrix}.
\end{equation}

于是

\begin{equation}
\det q=2^6d^3(h_1h_2h_3)^2\equiv d
          \pmod{\mathbb Q^{\times2}}.
\end{equation}

结合式 \eqref{eq:fixed:5}，\(d=3\)。故每个尚存 CM 候选都包含

\begin{equation}
K_0=\mathbb Q(\sqrt{-3}).              \label{eq:fixed:6}
\end{equation}

这里用的二次域来自极化判别式，不是根据某个数值周期猜测。


\subsection{一个辅助29-adic特征线排除全部 CM 候选}\label{sec:fixed:4}

\subsubsection{引理：P 的任意正整数幂谱都无重复}\label{ux5f15ux7406p-ux7684ux4efbux610fux6b63ux6574ux6570ux5e42ux8c31ux90fdux65e0ux91cdux590d}

设 \(\alpha\ne\beta\) 是 P 的根。若 \(\alpha/\beta\) 为有限阶单位根，因为 P 不可约，可以在其分裂域中取一个 \(\mathbb Q\) 自同构，将 \(\alpha\) 送到第\ref{sec:fixed:1}节唯一的11-adic单位特征根。它将 \(\beta\) 送到另一个与它同11-adic赋值的根，也必须为单位根，因唯一性两根相等，矛盾。

所以任意不同根的比均不是有限阶单位根。特别地，对每个 \(n\ge1\)，所有 \(\alpha^n\) 互不相同。这个论证实际处理了``需要先取 Frobenius 的某个有限幂''的问题；只检查 P 本身无重根并不够。

\subsubsection{29-adic 一维不变空间}\label{sec:fixed:29}

直接模29计算

\begin{equation}
P(Z)\equiv(Z-1)(Z-5)(Z^4-13Z^3+Z^2-7Z-4)\pmod {29},            \label{eq:fixed:7}
\end{equation}

且 \(P'(1)\equiv1\)。Hensel 引理给出一个根 \(\lambda\in\mathbb Q_{29}\)。式 \eqref{eq:fixed:3} 是未扭转 \(T_{29}\) 上 \(p=11\) 几何 Frobenius \(\Phi\) 的特征多项式。因此 \(\Phi\) 有一个一维 \(\mathbb Q_{29}\) 特征空间。前述引理保证，这个空间对每个 \(\Phi^n\) 仍然恰为一维。

K3 的 Hodge 自同态具有相容的29-adic实现，并与 Galois 群的某个有限指数子群交换；采用 Charles §2.2 第6页的绝对 Hodge 实现说明。因此存在 n 使 \(\Phi^n\) 与 \(F\otimes\mathbb Q_{29}\) 交换。若式 \eqref{eq:fixed:6} 成立，则其一维特征线成为

\begin{equation}
K_0\otimes\mathbb Q_{29}=\mathbb Q_{29}(\sqrt{-3})
\end{equation}

的含幺模。但

\begin{equation}
(-3)^{14}\equiv-1\pmod {29},
\end{equation}

故这是一个二次域扩张，不能在一维 \(\mathbb Q_{29}\) 空间上作含幺作用。矛盾。

29只作为系数素数 \(\ell\)；此步骤不要求 X 在29处良约化，也没有新算29处点数。良约化 Frobenius 始终是已认证的 \(p=11\)。

由实乘与CM的排除得到

\begin{equation}
\operatorname{End}_{\rm Hdg}T=\mathbb Q.
                                                                    \label{eq:fixed:8}
\end{equation}


\subsection{固定 Hodge 群}\label{sec:fixed:5}

Zarhin 的 K3 群描述在权重零形式上给出

\begin{equation}
G_H(S)=SO(q).                           \label{eq:fixed:9}
\end{equation}

可采用 Charles 的 K3 Mumford--Tate 定理及 Lemma 16 的连通群表述：\(E=\)\(\mathbb Q\) 时中心化群的恒等分支是 SO。Mumford--Tate 群连通，不能把未标恒等分支的 GO/O 符号机械照抄成不连通的 Hodge 群。

\(A=S\otimes B^{2}\)。S 的群 SO 是连通半单群，B 的群是 \(\mathbb G_m\)。两者没有非平凡共同代数群商，故

\begin{equation}
G_H(S\oplus B)=SO(q)\times\mathbb G_m.
\end{equation}

变回 $A,B$ 的框架即得到 $R_H\simeq SO(q)\times\Gm$。这里证明的是固定点，不使用射线的几何群。


\subsection{比较行列式检测 Nori 群的另一分支}\label{sec:fixed:6}

Hodge 实现以及实际纯核的杯积先给出夹逼

\begin{equation}
SO(q)\times\mathbb G_m
   \subseteq R_N
   \subseteq O(q)\times\mathbb G_m.                               \label{eq:fixed:10}
\end{equation}

两端仅差一个行列式符号，故只须判断一个实际一维对象

\begin{equation}
D=\det S
\end{equation}

是否为单位对象。杯积在 Nori 范畴中给出 \(D^{\otimes2}\simeq\mathbb Q(0)\)。

用已有六循环有理 Betti 基和裸留数 B de Rham 框架记纯周期矩阵为 P。令 \(z=2\pi i\)，则 S 的比较矩阵为 \(C_S=z^{-1}P^{\mathsf t}\)，故 D 的比较系数为

\begin{equation}
c_D=\det(P)/z^6.
\end{equation}

第\ref{ch:comparison}节 的实际定标给出

\begin{equation}
P^{\mathsf t}K_BP=8\pi^2Q_6.
\end{equation}

独立精确代数在 \(s=- 64\) 求得

\begin{equation}
\det K_B=\frac{353067798239177216000}{3},
\end{equation}

从而

\begin{equation}
\boxed{c_D^2=\frac{64\det Q_6}{\det K_B}
        =\frac9{86198192929486625}
        =\frac9{85^2\cdot11930545734185}.}                        \label{eq:fixed:11}
\end{equation}

其中

\begin{equation}
11930545734185=5\cdot13\cdot29\cdot37\cdot73\cdot89\cdot113\cdot233
\end{equation}

为无平方因子整数。因此 \(c_D\notin\mathbb Q\)。

若 D 在基域 \(\mathbb Q\) 的 Nori 范畴中为单位对象，任何有理 Betti/dR 基之间的比较系数都在 \(\mathbb Q^\times\)，与式 \eqref{eq:fixed:11} 矛盾。这只用``实际单位对象必有有理比较''这一方向，没有使用周期映射的单射猜想，也没有从一般数值关系反推出动机同构。

故行列式角色在 \(R_N\) 上非平凡。因其值域为 \(\mu_2\)，式 \eqref{eq:fixed:10} 的右端被全部取到，即

\begin{equation}
R_N=O(q)\times\mathbb G_m.              \label{eq:fixed:12}
\end{equation}

同样得到 \(G_{\rm N,\mathbb Q}(S)=O(q)\)。

基域不能省略：这是 \(\mathbb Q\) 上的 Nori 群；本文没有声称任意扩大基域后仍保留同一个非平凡分支。式 \eqref{eq:fixed:11} 给出一个非平凡二次比较角色，但本文也未仅凭周期数值就宣称 D 与某个指定 Artin 对象同构。


\subsection{完整固定活动群与 coaction 关系}\label{sec:fixed:7}

第\ref{ch:core}节 已证明固定活动 E 的幺幂核为

\begin{equation}
U=\left\{\begin{pmatrix}I_6&u&w\\0&1&v\\0&0&1\end{pmatrix}\right\},
 \qquad \dim U=13.
\end{equation}

它同时证明实际群是其分级群在权重抛物群中的完整逆像。代入已经识别的分级群，在选定有理权重分裂后：

\begin{equation}
\boxed{
 G_{\rm N,\mathbb Q}(E)=
 \left\{\begin{pmatrix}t^2g&u&w\\0&t&v\\0&0&1\end{pmatrix}:
     g\in O(q),\ t\ne0,\ u,w\in\mathbb Q^6,\ v\in\mathbb Q\right\}.
 }                                                               \label{eq:fixed:13}
\end{equation}

式 \eqref{eq:fixed:13} 按任意 \(\mathbb Q\) 代数解释为群概形；Hodge 群由再加 \(\det g=1\) 得到。这些群在适当分裂下分别为 U 与 \(O(q)\times\mathbb G_m\)、\(SO(q)\times\mathbb G_m\) 的半直积。两者维数 \(13+15+1=29\)，Nori 群两个分支，Hodge 群一个。

\textbf{基的提醒。} 原六循环的 q 为 Q6；若使用它们的对偶上同调 Betti 基，则 q 为 \(Q_6^{-1}\)，两者由有理杯积同构联系。de Rham 群须用实际 de Rham 杯积，而不能直接用 Q6 替代。旧 B 框架下为

\begin{equation}
q_{\rm dR}=-2K_B^{-1}.
\end{equation}

物理 \(A=T(3)\) 上的 Tate 值域为 B\textsuperscript{4}，因此 coaction 纯块 H、Tate 块 t 的全部关系可写为

\begin{equation}
H^{\mathsf t}q_{\rm dR}H=t^4q_{\rm dR},
                  \qquad t\ne0.                                \label{eq:fixed:14}
\end{equation}

固定纯关系由正交相似关系给出；不额外强加 \(\det H=t^{12}\)，那会错误地删去 Nori 的另一分支。Hodge 群才增加该行列式等式。

式 \eqref{eq:fixed:14} 连同自由的 u、v、w 和既有矩阵乘法余乘法，给出实际活动对象的完整 Hopf 描述。物理对偶对象仍取 \(C_M=C_E^{-\mathsf t}\)，既有

\begin{equation}
I^{\mathfrak m}=2(C_M^{\mathfrak m})_{8,6},\qquad
 \Delta I^{\mathfrak m}
 =2\sum_{j=1}^8(C_M^{\mathfrak m})_{8,j}
       \otimes(g_{E,\rm dR}^{-\mathsf t})_{j,6}
\end{equation}

不变，且其作用群现在已在指定纤维和基域下识别。这不意味着复周期值已经证明代数独立，也不把数值包围改称闭合特殊函数表达式。



\endgroup



\begingroup

本节以已构造的实际对象为前提讨论 coaction。群的允许块形与实际群等号是不同命题；前节通过扩张像和交换子论证承担后者。

\subsection{形式 coaction、可见子商与幺幂深度}\label{sec:coaction:9}

一般观察资料 \(\Sigma\subset E_B\)、\(\Omega\subset E_{\rm dR}^\vee\) 的规范对象是

\begin{equation}
R_\Sigma=\bigcap_{\Sigma\subset J_B}J,\qquad
 I_\Omega=\sum_{\Omega|J_{\rm dR}=0}J,\qquad
 V_{\Sigma,\Omega}=R_\Sigma/(R_\Sigma\cap I_\Omega).     \label{eq:coaction:9.1}
\end{equation}

交与和取在实际子对象中，有限长度保证存在。固定读出增加链给单射；固定链增加读出给反向满射。两种变化不能拼成未经证明的包含链。单观察的最小对象是标准矩阵系数理论的一部分；这里采用 Brown §2.4的形式周期约定。\cite{brownperiods}

选实际比较矩阵 \(C_E\)，并固定右 coaction 为 \(\Delta C_E=C_Eg_E\)。最大幺幂情形，分裂后

\begin{equation}
g_E=\begin{pmatrix}H&X&z\\0&tI_\kappa&y\\0&0&1\end{pmatrix},
 \qquad H=t^2g,\quad g\in R_{\rm pure}.
\end{equation}

若比较矩阵相应块为 \(P_A,p,f;\ 0,\tau I,q;\ 0,0,1\)，则

\begin{equation}
\Delta p=P_A\otimes X+p\otimes t,\qquad
 \Delta q=\tau\otimes y+q\otimes1,\qquad
 \Delta f=P_A\otimes z+p\otimes y+f\otimes1.              \label{eq:coaction:9.2}
\end{equation}

这些是实际对象的公式；数值求值为零的单个组合不自动给出可在此商去的子动机。实际矩阵关系须在对应 de Rham 框架中写出，不把 Betti 的整数 Gram 矩阵直接挪作 de Rham 配对。

定义标准 coradical 过滤

\begin{equation}
C_{-1}E=0,\qquad C_iE/C_{i-1}E=(E/C_{i-1}E)^U.
\end{equation}

由\eqref{eq:core:6.4}直接计算，在 \(r>0\) 的一般情形

\begin{equation}
C_0E=A\oplus\widetilde B_0,\quad C_1E=N,\quad C_2E=E.
                                                               \label{eq:coaction:9.3}
\end{equation}

对最小商，\(C_0\bar E=A,C_1\bar E=\bar N,C_2\bar E=\bar E\)，恰与三步权重对应。这个对应是实际作用的结论，不是所有混合对象的通则；分裂对象的 C\textsubscript{0}就是整个对象。

对上述最小商取 \(M=\bar E^\vee\)，\(r>0\)时满足 \(C_0M=\mathbf1,C_1M=\operatorname{Ann}(A),C_2M=M\)。\(\sigma \)生成E使系数映射 \(M_{\rm dR}\to\mathcal P^{\mathfrak m}\) 单射：在比较挠子上，恒零系数消去全部群轨道 \(\sigma \)，故消去整个 E。于是 \(\omega|A\ne0\) 给

\begin{equation}
I^{\mathfrak m}\in C_2\setminus C_1. 
                                                               \label{eq:coaction:9.4}
\end{equation}

这也从过滤给出相应的Tate方向下降。以下轨道维数式另假定原对象U最大且\(d=0\)。最大 U 对 \(\omega=(\alpha,\beta,\gamma)\) 的作用为

\begin{equation}
(\alpha,\ \beta-\alpha X,\ \gamma-\beta y+\alpha Xy-\alpha z).
\end{equation}

故幺幂轨道的张成维数为 \(\kappa +2\)，完整群轨道为\(6+\kappa +1\)：指定点分别为3与8，很一般秩十点分别为5与10。纯 K3 表示仍是 coaction 的必要部分，深度二不构成多对数类别判据。

由第\ref{sec:core:7}节的 \(\operatorname{End}(E)=\mathbb Q\)，固定活动对象上的任何有限翻号群都只能作标量作用。四边 banana 的泛能量核有八个非零字符，故这些字符不能同时成为该固定 E 上的八个动机直和投影。可能的参数覆盖、较大对象或跨纤维对应需要另行构造；组合支撑定理本身没有提供它们。



\endgroup

\section{多参数连接与物理导数}\label{ch:connection}\label{ch:deriv}

\begingroup


\subsection{纯连接的框架与配对}
单条质量射线上的微分系统不能代替完整参数空间上的 Gauss--Manin 连接。取四个独立参数 $x_i=\mu_i/s$，纯层的六元有理框架记为 $\mathbf e$，有
\begin{equation}
 \nabla\mathbf e=\sum_{i=1}^4 A_i(x)\mathbf e\,dx_i.
\end{equation}
连接矩阵必须同时满足曲率恒等式与极化相容式
\begin{equation}
 \partial_iA_j-\partial_jA_i+[A_i,A_j]=0,\qquad
 \partial_iK=A_i^{\mathsf t}K+KA_i,
\end{equation}
此处采用列水平解的配对约定；若使用框架列而非解列，须同时转换转置和负号。现有全四参数纯系统从原有理数据恢复，清分母后独立检查了全部曲率与配对条目，并与射线限制一致。全方向平坦性及一个已定标基点把几何杯积为公开配对的 $-2$ 倍延伸到连通光滑族。

这些恒等式的作用是证明同一几何框架的相容性。它们不自动决定其他纤维的绝对周期初值，也不把几何单值群等同固定纤维动机群。由指定射线群包含与全族极化上界，可以得到全参数纯几何闭包的正交群结论；其逻辑仍属于族。


\subsection{三个质量方向与实际秩十对象}
令 $\beta_i$ 为三个质量边界方向的实际 Hodge 提升。通过三坐标质量一动机的拉回，得到未扭转分级秩 $6,3,1$ 的秩十对象。其 Hodge 过滤维数为 $1,5,9,10$。顶层连接以符号 $\nu=e_F-e_B$ 写为
\begin{equation}
 \nabla e=-\Omega\,d\log x_4+
 \sum_{i=1}^3\beta_i\,d\log(x_i/x_4).
\end{equation}
质量提升之间的耦合由实际 Cartan 公式、留数、环面分部积分及纯层配对约束确定。完全对称三张量中的唯一性线性系统行列式非零，消除仅由平坦性猜测耦合的不确定性。所得四个 $10\times10$ 矩阵的 $600$ 个曲率条目清分母后恒为零；这里该有限核对是理论推导的输入检查，不替代实际一动机拉回。

在函数域中，三个质量比单位独立。固定到 $m=(1,2,4,8)$ 的质量射线后，它们全部由 $2$ 生成，活动方向为
\begin{equation}
 \beta_c=3\beta_1+2\beta_2+\beta_3.
\end{equation}
因此族中的秩十环境与该点的秩八最小活动对象并不矛盾。不能把固定点的一个 $\log2$ 方向直接铺成全族秩八连接。


\subsection{纯系统、混合系统与标量方程的不同秩}
一个标量函数的最小微分阶数依赖所选参数方向、生成截面及常数域。纯微分模绝对不可约时，任意非零截面通过协变导数生成整个模，因而纯周期的最小齐次 Picard--Fuchs 阶数等于纯秩。但整个相对对象还包含源项和边界比较，不能从纯六阶方程直接推出完整物理对象只有六维。

同样，一条特殊参数曲线上微分阶数下降，可能来自额外代数类、扩张退化、所选读出不可见或参数方向本身不足以生成全模。必须逐项判别。多参数平坦连接、固定纤维最小形式周期以及单参数标量方程，是三个相关而不同的不变量。

秩十对象的完整绝对比较尚未作为本文结果。若进一步构造，需要三个真实边界 Betti 列、环境像校正、实填充及第六球的混合单值增量；给 $P_8$ 添零不提供任何这样的几何证据。当前更优先的问题是把已知连接用于结构与操作的理论分类。




\endgroup



\begingroup
\providecommand{\Q}{\mathbb Q}
\providecommand{\C}{\mathbb C}
\providecommand{\dR}{\mathrm{dR}}
\providecommand{\Ann}{\operatorname{Ann}}
\providecommand{\Gr}{\operatorname{Gr}}
\providecommand{\Res}{\operatorname{Res}}
\providecommand{\rk}{\operatorname{rk}}
\providecommand{\cE}{\mathcal E}
\providecommand{\cM}{\mathcal M}
\subsection{实际对象、参数与要回答的问题}
\subsubsection{先在全族求导，再限制到特殊活动对象}
设 $\mu_i=m_i^2$。研究二维、四传播子、单位幂次 banana 的正质量欧氏区域 $\mu_i>0$、$s<0$。
在 $T=(\mathbb G_m)^3$ 上取 $\alpha=(a,b,c,1)$，令
\begin{equation}\label{deriv:eq:raw}
 N=\sum_{i=1}^4\alpha_i^{-1},\quad M_\mu=\sum_{i=1}^4\mu_i\alpha_i,
 \quad P=M_\mu N-s,\quad
 \eta=d\log a\wedge d\log b\wedge d\log c.
\end{equation}
项目中已固定无额外动量测度因子的积分
\begin{equation}\label{deriv:eq:I}
 I(s,\mu)=\int_{\sigma_{\rm phys}}\frac{\eta}{P}.
\end{equation}
正实链带既定正向；其边界通过已有紧化及相对对偶模型处理。
质量平方与 $s$ 共同缩放时，$I$ 的齐次次数为 $-1$。

取四个无量纲参数及导数
\begin{equation}\label{deriv:eq:dimensionless}
 x_i=\mu_i/s,\quad D_i=x_i\partial_{x_i},\quad
 \Theta=\sum_iD_i,\quad
 F=\Bigl(\sum_i x_i\alpha_i\Bigr)N-1.
\end{equation}
记 $U_x=\{F=0\}$，以及
\begin{equation}\label{deriv:eq:J}
 J(x)=sI(s,\mu)=\int_{\sigma_{\rm phys}}\eta/F,\qquad
 \Omega=\Res_{U_x}\eta/F.
\end{equation}
此处 $J$ 与原正积分 $I$ 差一个 $s$ 因子；在欧氏区域 $J<0$。
这个归一化有助于使用四参数几何，不能在符号或物理数值比较中省略。

以下工作在带标记的光滑解消开集上。
取已有实际相对对象的质量子对象，并暂以\emph{未扭转}记号写成
\begin{equation}\label{deriv:eq:mass}
 0\subset K\subset W_4\cE\subset\cE,
 \quad \Gr^W_2\cE=K,\quad
 \Gr^W_4\cE=\Q(-2)^3,\quad
 \Gr^W_6\cE=\Q(-3).
\end{equation}
$K$ 为边界正交的秩六 K3 型纯变分，Hodge 数为 $(1,4,1)$。
物理方向的同调对象是 $\cE(3)$，其对偶读出对象为
\begin{equation}\label{deriv:eq:M}
 \cM=\cE(3)^\vee.
\end{equation}
这两个对象不能混用 Hodge 下标。
本文将 $\cM_{\dR}$ 的向量视为 $\cE(3)_{\dR}$ 上的线性读出。
固定 de Rham Tate 生成元后，可将读出的取值写在未扭转的几何框架上。

指定有理点为
\begin{equation}\label{deriv:eq:point}
 m=(1,2,4,8),\quad s=-64,\quad
 x_* =(-1/64,-1/16,-1/4,-1).
\end{equation}
其实际秩八活动对象 $E_*$ 嵌入 $\cE_{x_*}(3)$，并包含原物理链。
只有在求导并取纤维之后才限制到 $E_*$；不假定它对任意质量方向的连接稳定。

\subsubsection{使用的既有输入}
实际对象及以下连接来自项目的四参数混合连接稿及顶层相对比较稿
（文末列出文件名）：
\begin{equation}\label{deriv:eq:beta}
 \beta_i=-x_i\alpha_iN\Omega,\quad
 \sum_i\beta_i=-\Omega,\quad e=(\eta,0),\quad
 \nabla_{D_i}e=\beta_i.
\end{equation}
相对连接注入符号在此省略。令
\begin{equation}\label{deriv:eq:C}
 V_i=\nabla_{D_i}\Omega,\qquad
 C_{ij}=\nabla_{D_j}\beta_i\in F^1K.
\end{equation}
由实际 Cartan 公式、留数平坦性和连接平坦性，
\begin{equation}\label{deriv:eq:sum}
 C_{ij}=C_{ji},\qquad \sum_j C_{ij}=-V_i.
\end{equation}
我们只用这些结构等式，不使用巨大有理连接矩阵的逐项重算。
未扭转 Hodge 过滤为
\begin{align}\label{deriv:eq:filterE}
 F^3\cE&=\langle e\rangle,&
 F^2\cE&=\langle\Omega,\beta_1,\beta_2,\beta_3,e\rangle,\nonumber\\
 F^1\cE&=\langle\Omega,V_1,V_2,V_3,V_4,\beta_1,\beta_2,\beta_3,e\rangle,&
 F^0\cE&=\cE.
\end{align}
因此
\begin{equation}\label{deriv:eq:dualHodge}
 F^p\cM=\Ann(F^{4-p}\cE),\qquad
 (\dim F^3,\dim F^2,\dim F^1,\dim F^0)\cM=(1,5,9,10).
\end{equation}
物理形式给出 $\lambda\in F^3\cM$，其纯限制在既定相对对偶符号下为
\begin{equation}\label{deriv:eq:lambda}
 \lambda|K=-q(\Omega,-),
\end{equation}
其中 $q$ 是实际几何杯积；已有归一化证明给 $q=-2\langle-,-\rangle_{\rm pub}$。
本节保持该符号，不重新定标。


\subsection{导数确实是同一几何环境中的读出}
\begin{lemma}[参数导数及相对类]\label{deriv:lem:actual}
在一个足够小的正质量欧氏参数邻域内，\eqref{deriv:eq:I} 可逐阶对参数求导。
所得类是既有相对几何族的 Gauss--Manin 导数；在固定点限制到 $E_*$ 后，给出该实际对象的有理 de Rham 读出。
一阶原始形式为
\begin{equation}\label{deriv:eq:derforms}
 \partial_s(\eta/P)=\eta/P^2,\qquad
 \partial_{\mu_i}(\eta/P)=-\alpha_iN\eta/P^2.
\end{equation}
\end{lemma}
\begin{proof}
在参数的紧邻域中，所有 $\mu_i$ 和 $-s$ 有共同严格正下界 $c$。
在正实链上 $P=M_\mu N+|s|$，故
\[
 1/P\le 1/c,\qquad \alpha_iN/P\le1/\mu_i\le1/c.
\]
任意有限阶参数导数均为 $1/P$ 乘以这些有界比值的有限线性组合。
邻域内的 $P$ 又与固定参数的 $P_0$ 一致可比。
已有收敛的原积分因而给出可积支配函数，允许在积分号下求导。

几何上，导数仅提高同一极点超曲面的极点阶数，不引入新的支持。
在已建立的光滑相对对族中，以局部拓扑平凡化运输链，代数 de Rham 导数正是 Gauss--Manin 导数。
若某紧化图表需要边界 cochain，使用相对连接完成代表；不能把遗漏边界项的非收敛分部积分当作证明。
正实链在此欧氏邻域没有穿越极点，且其标记边界同步运输，故它是相应局部系统的平坦截面。
因此 $D J=\langle\nabla_D\lambda,\sigma_{\rm phys}\rangle$。

最后，代数连接在有理基点沿有理切向量给有理实现向量。
实际嵌入 $E_*\hookrightarrow\cE_{x_*}(3)$ 的对偶是实现上的满射，故任意这些导数都可限制到同一 $E_*$。
这个步骤不要求在四参数邻域中存在一个水平的秩八子束。
\end{proof}

改变传播子幂次的完整积分族常包含其他 sector。
本引理只处理由明确参数导数得到的这些类及它们的线性组合；它不把所有任意分子或维数平移自动纳入 $E_*$。
同样，连接求导不是固定纤维上的动机自同态，故与上一补篇 $\operatorname{End}(E_*)=\Q$ 的结论不冲突。


\subsection{局部 Torelli 与五维一阶空间}
\subsubsection{为什么这四个参数是真实的模参数}
为使后续秩结论有几何依据，重述项目中的参数单射证明。
Sylvester 坐标满足 $\sum_{i=1}^5y_i=0$，节点四次式写为
\begin{equation}\label{deriv:eq:Sylvester}
 Q_\ell=\sum_{i=1}^5\ell_i\prod_{j\ne i}y_j,
 \qquad [\ell]=[\mu_1:\mu_2:\mu_3:\mu_4:s].
\end{equation}
在仅有标记十节点的开集，节点恒为 $[e_i-e_j]$。
若一阶参数变形在抽象 K3 上平凡，则 $H^1(\mathcal O)=0$ 使超平面线丛唯一平凡化，四个截面给出原模型的一阶射影变换。
每个例外 $(-2)$ 曲线无一阶移动，因为其法丛为 $\mathcal O_{\mathbb P^1}(-2)$。
所以该射影变换必须固定十节点。
固定四个 $[e_i-e_5]$ 迫使线性矩阵对角，固定 $[e_1-e_j]$ 又使四个对角元相同；在 $\operatorname{PGL}_4$ 中变换为零。
五个系数单项式在 $\sum y_i=0$ 上线性独立：在 $y_i=0$ 而其他坐标非零处分别读取即可。
因此唯一平凡方向是共同缩放。

这给参数族 Kodaira--Spencer 映射单射。
K3 局部 Torelli 将其与周期线的微分识别\cite[Chapter 6, Proposition 2.8]{huybrechts}。
于是
\begin{equation}\label{deriv:eq:KS}
 \Omega,V_1,V_2,V_3,V_4\ \text{构成 }F^1K\text{ 的基}.
\end{equation}
这是带标记开集的断言，不以粗模空间没有有限自同构为前提。
指定点处的光滑性和纯秩六输入已由旧几何证明确立，本节不重复计算。

\subsubsection{极化及导数空间}
令
\begin{equation}\label{deriv:eq:Y}
 Y_{ij}=q(V_i,V_j),\qquad Y=(Y_{ij})_{1\le i,j\le4}.
\end{equation}
由 Hodge 型及平坦性，$q(\Omega,\Omega)=q(\Omega,V_i)=0$。
$F^1K=\Omega^\perp$，其配对的根基恰为 $\Q\Omega$；因此 \eqref{deriv:eq:KS} 蕴含 $Y$ 可逆。
这里不需要计算 $Y$ 的十个有理函数值。

\begin{definition}
在局部代数函数环上，令 $\mathcal J^r\subset\cM_{\dR}$ 是所有
$\nabla_{D_{i_1}}\cdots\nabla_{D_{i_j}}\lambda$、$0\le j\le r$ 所生成的子模。
允许系数为在当前点正则的有理函数。
这是\emph{参数导数阶数}过滤，不是选定单参数常微分方程的阶数。
\end{definition}

Griffiths 横截性给
\begin{equation}\label{deriv:eq:upper}
 \mathcal J^r\subset F^{3-r}\cM\qquad(0\le r\le3).
\end{equation}
对偶连接、纯子变分的水平性与 \eqref{deriv:eq:lambda} 给
\begin{equation}\label{deriv:eq:first}
 (\nabla_{D_i}\lambda)|K=-q(V_i,-).
\end{equation}
由 \eqref{deriv:eq:KS}，$\lambda$ 与四个一阶导数在纯限制上已经线性无关。
结合 \eqref{deriv:eq:dualHodge} 的上界即得
\begin{equation}\label{deriv:eq:jet1}
 \mathcal J^0=F^3\cM,\qquad\mathcal J^1=F^2\cM.
\end{equation}


\subsection{二阶导数恰好补齐三个质量方向}
\subsubsection{只需四个独立二次型}
记 $\lambda_i=\nabla_{D_i}\lambda$、$\lambda_{ij}=\nabla_{D_j}\lambda_i$。
由于 $D_i$ 相互交换且连接平坦，$\lambda_{ij}=\lambda_{ji}$。
定义对称矩阵
\begin{equation}\label{deriv:eq:T}
 (T_r)_{ij}=q(V_i,C_{rj}),\qquad r=1,2,3.
\end{equation}
其对称性也可直接由下面的读出等式得到。

\begin{lemma}[二阶读出的四个商坐标]\label{deriv:lem:second}
模去 $F^2\cM$，$F^1\cM$ 可用在 $\Omega,\beta_1,\beta_2,\beta_3$ 上的值作坐标。
在这些坐标下
\begin{equation}\label{deriv:eq:eval2}
 \lambda_{ij}(\Omega)=Y_{ij},\qquad
 \lambda_{ij}(\beta_r)=(T_r)_{ij},\qquad
 \lambda_{ij}(e)=0.
\end{equation}
\end{lemma}
\begin{proof}
纯限制为 $\lambda_{ij}|K=-q(\nabla_{D_j}V_i,-)$。
对 $q(V_i,\Omega)=0$ 求导，第一式随即成立。
因为 $\lambda_i\in F^2\cM$，它消去所有 $\beta_r$，所以
\[
 \lambda_{ij}(\beta_r)=-\lambda_i(C_{rj})=q(V_i,C_{rj}).
\]
最后一式由 $\lambda_{ij}\in F^1\cM=\Ann(e)$ 得到。
商坐标的维数为 $9-5=4$，且 \eqref{deriv:eq:filterE} 说明上述四个测试类恰给全部商坐标。
\end{proof}

\begin{lemma}[无需逐项计算的独立性]\label{deriv:lem:independent}
四个对称二次型 $Y,T_1,T_2,T_3$ 线性无关。
\end{lemma}
\begin{proof}
假定 $aY+\sum_{r=1}^3b_rT_r=0$，并补记 $b_4=0$。
将第二个指标 $j$ 求和，由 \eqref{deriv:eq:sum} 得
\[
 \sum_j(T_r)_{ij}
 =q\Bigl(V_i,\sum_jC_{rj}\Bigr)=-Y_{ir}.
\]
因此每个 $i$ 都满足
\[
 a\sum_jY_{ij}-\sum_rb_rY_{ir}=0,
 \quad\text{即}\quad Y(a\mathbf1-b)=0.
\]
$Y$ 可逆，故 $b=a\mathbf1$；第四个分量 $b_4=0$ 强迫 $a=0$，再得所有 $b_r=0$。
整个证明没有代入连接系数、周期初值或数值近似。
\end{proof}

\begin{proposition}[二阶 Hodge 生成]\label{deriv:prop:jet2}
在上述开集，$\mathcal J^2=F^1\cM$，秩为九。
\end{proposition}
\begin{proof}
引理\ref{deriv:lem:independent}说明对称二阶导数的商坐标 \eqref{deriv:eq:eval2} 张成四维空间。
一阶空间已是 $F^2\cM$，所以二阶空间包含整个 $F^1\cM$；再用 \eqref{deriv:eq:upper} 的上界。
\end{proof}

这也说明二阶纯限制张成全部秩六 $K^\vee$。
但二阶不止恢复纯 K3：另外三个商方向是实际质量 Kummer 方向。
只知道一个纯 Picard--Fuchs 系统不能自动读出这三个方向。


\endgroup
\subsection{高阶读出与指定纤维}
二阶导数补齐质量方向之后，继续求导可提取单位方向。具体算子由有限Gauss--Manin约化和线性代数构造，不要求先求积分。三阶生成、质量对数的符号与二阶算子的构造见\otherref{feynman_periods_supplement}{补充材料S4}。

限制到指定秩八活动对象时，应先在全族求导，再限制读出。沿特殊质量射线求导与先将所有横向方向删去并不等价。补充材料保留指定点的Hodge过滤、一阶保护和最早消去阶数的完整证明。这里各阶导数仍是同一实际对象中的读出；微分生成秩与物理形式周期的最小秩分别由连接和双标架决定。

\part{操作与重整化}

\section{切割、留数过滤与最大割}\label{ch:cut}

\begingroup

本节先构造正常交叉模型中的实际留数态射，再比较割周期与原周期的标架。相关经典留数与变分背景见\cite{muhlbauer}；其一圈 Cutkosky 定理的维数及交数条件不直接用作本文多圈模型的无条件证明。

\subsection{一个光滑除子的实际留数}\label{sec:cut:2}

设X/k光滑，\(D\subset X\)为光滑除子，\(B\subset X\)为光滑闭边界或简单正常交叉边界；D与B的所有交叉层横截，且D不是B的分量。置

\begin{equation}
V=X\setminus D,\quad B_V=B\cap V,\quad B_D=B\cap D.
\end{equation}

边界为空亦允许。记相对复形为 \(R\Gamma(X,B)=\operatorname{Fib}(R\Gamma X\to R\Gamma B)\)。

\begin{proposition}

\label{thm:cut:2.1} 有实际Nori态射及正合段

\begin{equation}
H^n(X,B)\longrightarrow H^n(V,B_V)
 \xrightarrow{\ r_D\ }H^{n-1}(D,B_D)(-1)
 \longrightarrow H^{n+1}(X,B).                         \label{eq:cut:2.1}
\end{equation}

特别地，留数像是原相对对象的商，也是目标割对象的子对象，但留数不必满射或单射。

\end{proposition}

\begin{proof}

对\(D\subset X\)使用局部化三角和余维一纯性

\begin{equation}
R\Gamma(D)(-1)[-2]\to R\Gamma(X)\to R\Gamma(V)\to.
\end{equation}

对\(B_D\subset B\)写相同三角，利用横截基变换连接二者。取相对纤维得到

\begin{equation}
R\Gamma(D,B_D)(-1)[-2]\to R\Gamma(X,B)
 \to R\Gamma(V,B_V)\to R\Gamma(D,B_D)(-1)[-1].         \label{eq:cut:2.2}
\end{equation}

若B有多个正常交叉分量，在其有限交叉的Čech复形上逐项进行，再取总复形。横截性保证每一项纯性相容。取第n上同调得到\eqref{eq:cut:2.1}。所有箭头先在六运算范畴构造，再取Nori上同调，并非由Betti矩阵猜测。

\end{proof}

在局部坐标\(z=0\)定义D时，闭对数形式可写为

\begin{equation}
\omega=\frac{dz}{z}\wedge\alpha+\beta,
 \qquad\operatorname{res}_D\omega=\alpha|_D.           \label{eq:cut:2.3}
\end{equation}

相对形式及其边界分量使用同一映射锥约定。式\eqref{eq:cut:2.3}是未除以\(2\pi i\)的裸代数留数；\eqref{eq:cut:2.1}中的Tate扭转承担比较归一化。取正向小圆为法向取向，拓扑管映射满足

\begin{equation}
\int_{\operatorname{Tub}_D\gamma}\omega
       =2\pi i\int_\gamma\operatorname{res}_D\omega.   \label{eq:cut:2.4}
\end{equation}

相对情形要求管的边界落在\(B_V\)中；横截管邻域可实现这一点。式\eqref{eq:cut:2.4}也由每个小圆上的Cauchy积分和相对Stokes公式直接给出。裸留数公式与含扭转的实际态射必须一同使用，不能省略Tate线后仍称比较矩阵不变。


\subsection{割周期的最小对象：等号需要哪些标记}\label{sec:cut:4}

固定C，令\(F_C=N_C\vee \)，\(\tau _C=r_C\vee :F_C\to E\)。取有理Betti链\(\gamma \in (F_C)_B\)和原读出\(\omega \in E_{\mathrm{dR}}\vee \)。记\(\omega _C=\tau _C^\ast \omega \)，其中Tate线已经包含在\(F_C\)内。则形式周期具有实际态射关系

\begin{equation}
[F_C,\gamma,\omega_C]^{\mathfrak m}
       =[E,\tau_C\gamma,\omega]^{\mathfrak m}.         \label{eq:cut:4.1}
\end{equation}

使用裸积分书写时，\eqref{eq:cut:4.1}就是\otherref{feynman_periods_supplement}{（补充材料S5.3）}，不能把右侧\(2\pi i\)因子漏掉。

\begin{theorem}

\label{thm:cut:4.1} 采用\otherref{feynman_periods_supplement}{补充材料S2}的规范最小对象，有

\begin{equation}
\operatorname{Vis}(F_C;\gamma,\omega_C)
     \simeq\operatorname{Vis}(E;\tau_C\gamma,\omega).  \label{eq:cut:4.2}
\end{equation}

若原物理链\(\sigma \)满足\(R_{\tau _C\gamma }\subset R_\sigma \)，则割的最小同调对象自然单射到原物理最小对象。尤其\(\sigma \)生成E时成立。

\end{theorem}

\begin{proof}

\(\tau _C\)把\(\gamma \)生成的最小子对象满射到\(\tau _C\gamma \)生成的最小子对象。其核被拉回读出消去；任何另一个被消去的子对象，像仍被\(\omega \)消去。反向拉回目标最大不可见子对象，得到两者的最小商同构。若\(R_{\tau _C\gamma }\subset R_\sigma \)，两者使用同一个E内的最大不可见子对象\(I_\omega \)，故

\begin{equation}
R_{\tau _C\gamma }/(R_{\tau _C\gamma }\cap I_\omega)
       \hookrightarrow R_\sigma/(R_\sigma\cap I_\omega).
\end{equation}

\end{proof}

若同时允许所有r重割的全部有理链，生成对象恰\(T_r\)：有限维Betti基张成各\(F_C\)，像之和正是\otherref{feynman_periods_supplement}{（补充材料S5.4）}。因此在固定读出下

\begin{equation}
\mathcal V_r=T_r/(T_r\cap I_\omega),\qquad
 \mathcal V_{r+1}\hookrightarrow\mathcal V_r.          \label{eq:cut:4.3}
\end{equation}

这是经过物理读出的实际割过滤。如果\(\sigma \)不是循环生成元，不得直接把V\textsubscript{0}或全部\(V_r\)当作\(\sigma \)的最小对象；此时只有落在\(R_\sigma \)内的部分受其控制。

式\eqref{eq:cut:4.1}与规范coaction相容，因为矩阵系数关系本来就按实际态射取商。它不推出``所有图割恰好构成coaction第二因子的一个基''：后者还需完备性、线性独立性、对偶配对及收敛/重整化边界的额外证明。


\subsection{单值变分与留数态射必须分开}\label{sec:cut:5}

在一个已构造的解析族中，固定基点处\(T_\ell \)是有理Betti单值。若局部拓扑已经给出

\begin{equation}
(T_\ell-1)\sigma=\sum_C a_C\operatorname{Tub}_C\gamma_C,
 \qquad a_C\in\mathbb Q,                              \label{eq:cut:5.1}
\end{equation}

则积分线性性和\otherref{feynman_periods_supplement}{（补充材料S5.3）}直接给

\begin{equation}
\operatorname{Disc}_\ell I
     =\sum_C a_C(2\pi i)^{|C|}
                      \int_{\gamma_C}\operatorname{res}_C\omega.
                                                               \label{eq:cut:5.2}
\end{equation}

这里\(a_C\)包含实际交数与路径、方向信息，不默认为1。普通简单夹点时，这种表示由局部消失同调提供；但正能量\(\delta _+\)表示还要求识别实际实支、正能量选择以及原链的夹挤。几何复杂性或Landau方程本身不提供这些标记。无穷远贡献、重整化边界及非孤立夹点也不能被\eqref{eq:cut:5.2}自动排除。

引用范围尤其需要准确：\cite{muhlbauer} Theorem41在该文的紧化约定下明确假定奇数D，并保留交数N；不能把该定理直接当作本文二维多圈情形的无条件Cutkosky证明。

\begin{proposition}

\label{thm:cut:5.1} 固定纤维的单值矩阵不一般是该纤维的Hodge或Nori自同态。

\end{proposition}

\begin{proof}

指定banana纤维的纯T满足\(\operatorname{End}_{\mathrm{Hdg}}(T)=Q\)，已由第\ref{ch:fixed}节证明。一个非平凡Picard--Lefschetz节点反射有一维\(- 1\)空间和五维+1空间，因此不是标量，不可能属于\(\operatorname{End}_{\mathrm{Hdg}}(T)\)。若它来自Nori自同态，其Hodge实现也应是自同态，矛盾。

\end{proof}

所以\eqref{eq:cut:5.1}是Betti向量等式；其两侧都可作为形式周期的链标记，但不能将\(T_\ell - 1\)未经证明升级为固定动机中的一个态射或幂等投影。近邻/消失循环的变分态射处在族或退化的范畴中，与固定纤维的断言有不同的源和靶。


\subsection{二维banana的最大割是一个明确代数簇}\label{sec:cut:6}

取h条内线的banana图，圈数\(h- 1\)，二维Minkowski形式为\(q^{2}=(q^{0})^{2}- (q^{1})^{2}\)。在复数域取光锥坐标

\begin{equation}
x_i=q_i^0+q_i^1,\quad y_i=q_i^0-q_i^1,
 \quad x_iy_i=\mu_i\ne0,
 \quad\sum_i x_i=p^+,\quad\sum_i y_i=p^-.
\end{equation}

假设\(s=p^+p^-\ne 0\)。全部传播子在壳的复割簇为

\begin{equation}
Z_h=\{(x_i)\in\mathbb G_m^h:
             \sum_i x_i=p^+,\ \sum_i\mu_i/x_i=p^-\}.   \label{eq:cut:6.1}
\end{equation}

其光滑点的复维数是\(h- 2\)。

\begin{proposition}

\label{thm:cut:6.1} \(Z_h\)与下列射影环面超曲面同构：

\begin{equation}
\Big(\sum_i x_i\Big)\Big(\sum_i\mu_i/x_i\Big)=s,
 \quad[x_1:\cdots:x_h]\in\mathbb P^{h-1},\quad\prod x_i\ne0.
                                                               \label{eq:cut:6.2}
\end{equation}

经\(x_i=\mu _i\alpha _i\)，这正是二维banana的第一、第二Symanzik组合

\begin{equation}
F_h=\Big(\sum_i\mu_i\alpha_i\Big)
               \Big(\sum_i\prod_{j\ne i}\alpha_j\Big)
                      -s\prod_i\alpha_i=0.             \label{eq:cut:6.3}
\end{equation}

\end{proposition}

\begin{proof}

\eqref{eq:cut:6.1}显然给\eqref{eq:cut:6.2}。反向，对\eqref{eq:cut:6.2}的任一射影点，\(s\ne 0\)保证\(\sum x_i\ne 0\)；唯一缩放\(p^+/(\sum x_i)\)把其和归一为\(p^+\)，第二个和随即为\(s/p^+=p^-\)。这给互逆的代数映射。代入\(x_i=\mu _i\alpha _i\)并乘\(\Pi x_i\)，公共因子\(\Pi \mu _i\)消去即得\eqref{eq:cut:6.3}。

\end{proof}

因此，\(h=2\)给一般两个点；\(h=3\)给一般光滑平面三次曲线，亏格由伴随公式为一；\(h=4\)给本文的四次K3曲面模型，解消后为\(X_s\)。\(h=3\)的光滑性可直接检查：边界顶点的一个偏导为非零\(\mu _i\)，边界非顶点处亦至少一个偏导非零；环面内排除临界值。\(s=0\)时三次方程退化为直线乘二次曲线，各不可约正规化分量亏格零。这已经说明，同一边表在不同运动学下并无单一``割曲线亏格''。\(h=4\)甚至根本不是曲线。

在正实物理支，\(x_i>0\)等价于未来向在壳\(q_i\)，且

\begin{equation}
s=\Big(\sum_i x_i\Big)\Big(\sum_i\mu_i/x_i\Big)
          \geq\Big(\sum_i\sqrt{\mu_i}\Big)^2.          \label{eq:cut:6.4}
\end{equation}

这是Cauchy--Schwarz，等号当且仅当\(x_i\propto \)√\(\mu _i\)。因而s小于阈值时没有正能量实最大割，但复割簇仍存在。特别是\(s=- 64\)的复K3周期不能被称为该欧氏点的正相空间积分；两者间还需选择解析延拓路径。


\subsection{四传播子的最大割形式：裸归一化的精确等式}\label{sec:cut:7}

现在\(h=4\)，用前三条内线作独立动量。置\(u_i=q_i^+\)、\(v_i=q_i^-\)（\(i=1,2,3\)），并定义

\begin{equation}
d_i=u_iv_i-\mu_i\quad(i\leq3),\qquad
 d_4=(p^+-u_1-u_2-u_3)(p^--v_1-v_2-v_3)-\mu_4.
\end{equation}

取有序裸六形式

\begin{equation}
\Theta=\frac{du_1\wedge dv_1\wedge du_2\wedge dv_2
                          \wedge du_3\wedge dv_3}
                         {d_1d_2d_3d_4}.              \label{eq:cut:7.1}
\end{equation}

本节始终按d\textsubscript{1},d\textsubscript{2},d\textsubscript{3},d\textsubscript{4}顺序取留数。

\begin{theorem}

\label{thm:cut:7.1} 在光滑最大割上，经命题\ref{thm:cut:6.1}的同构，

\begin{equation}
\boxed{\operatorname{res}_{d_4}\operatorname{res}_{d_3}
 \operatorname{res}_{d_2}\operatorname{res}_{d_1}\Theta
       =\operatorname{Res}_{F_4=0}
                     \frac{\Omega_{\mathbb P^3,\alpha}}{F_4}
       =-\omega_s^{\rm abs}.}
                                                               \label{eq:cut:7.2}
\end{equation}

这里\(\omega _s^{\mathrm{abs}}\)严格采用第\ref{ch:comparison}节的\(d\log (\alpha _{1}/\alpha _{4})\wedge d\log (\alpha _{2}/\alpha _{4})\wedge d\log (\alpha _{3}/\alpha _{4})\)约定。下面定义的标准\(\Omega \)在\(\alpha _{4}=1\)图上等于\(- d\alpha _{1}\wedge d\alpha _{2}\wedge d\alpha _{3}\)，所以本条有明确负号；不能把两种裸形式的取向差默默吸收进旧周期证书。

\end{theorem}

\begin{proof}

\(du_i\wedge dv_i=du_i\wedge dd_i/u_i\)。把(dd\textsubscript{1},dd\textsubscript{2},dd\textsubscript{3})移至所有du之前需要\(1+2+3=6\)次交换，故前三次留数为

\begin{equation}
\frac{du_1\wedge du_2\wedge du_3}{u_1u_2u_3H},\quad
 H=x_4\Big(p^- -\sum_{i=1}^3\mu_i/u_i\Big)-\mu_4,
 \quad x_4=p^+-u_1-u_2-u_3.                           \label{eq:cut:7.3}
\end{equation}

在\(\sum x_i=p^+\)图上，令\(f=(\sum x_i)(\sum \mu _i/x_i)\)。直接代数恒等式给

\begin{equation}
H=-\frac{x_4}{p^+}(f-s),\qquad
 \Omega_x=\sum_{i=1}^4(-1)^{i-1}x_i
             dx_1\wedge\cdots\widehat{dx_i}\cdots\wedge dx_4
            =-p^+du_1\wedge du_2\wedge du_3.           \label{eq:cut:7.4}
\end{equation}

故\eqref{eq:cut:7.3}就是\(\Omega _x/[\Pi x_i(f- s)]\)。最后一次留数可用H或\(f- s\)作局部方程，乘可逆函数不改变所得Poincaré留数。线性替换\(x_i=\mu _i\alpha _i\)使分子\(\Omega _x\)和分母\(\Pi x_i(f- s)\)同时乘\(\Pi \mu _i\)。最后在\(\alpha _{4}=1\)图上，\(\Omega _\alpha /\Pi \alpha _i=- d\log \alpha _{1}\wedge d\log \alpha _{2}\wedge d\log \alpha _{3}\)，得到\eqref{eq:cut:7.2}及其相对于旧绝对框架的负号。

\end{proof}

实际Lorentz测度还有

\begin{equation}
\bigwedge_{i=1}^3(dq_i^0\wedge dq_i^1)
       =-\tfrac18\bigwedge_{i=1}^3(du_i\wedge dv_i).  \label{eq:cut:7.5}
\end{equation}

结合\eqref{eq:cut:7.2}，有序Lorentz测度的四重裸留数为\(+\omega _s^{\mathrm{abs}}/8\)。因此本文没有擅自吸收传播子的i、每圈\(2\pi \)因子或交数。若采用\eqref{eq:cut:7.1}的光锥裸测度，四重管公式的因子就是\((2\pi i)^{4}\)；其他振幅约定必须另外代入\eqref{eq:cut:7.5}及其测度定义。旧Morse球取向、P6/P7/P8以及物理I的符号全部保持不变。

在上同调的Tate账目上，未归一化六维动量对象的留数靶为\(H^{2}(U_s)(- 4)\)。若每圈以一个Tate周期作归一化，整体扭转(3)后靶为\(H^{2}(U_s)(- 1)\)。其纯超越部分\(T(X_s)(- 1)\)由杯积识别为

\begin{equation}
T(X_s)(3)^\vee=A_T^\vee.            \label{eq:cut:7.6}
\end{equation}

这解释纯核的权重为何相容。但本节只构造最大割簇与形式的同构，没有据此宣称已经构造完整动量相对对象与参数相对对象之间的全局同构；后者包括无穷远、相对链及推前的额外工作。


\subsection{最大割能证明什么，不能恢复什么}\label{sec:cut:8}

全纯K3二形式来自\(H^{2}(X_s)\)，又被全部代数除子配对消去。因此\eqref{eq:cut:7.2}的上同调读出通过实际纯超越因子\(T(X_s)\)因子化，不需把整个混合\(H^{2}(U_s)\)当作纯对象。

若所取割循环在T的对偶中非零，T的简单性使该形式周期的最小纯对象就是T，连同\eqref{eq:cut:7.6}所需的Tate扭转。在指定\(\rho 16\)纤维，其秩为六；等质量很一般\(\rho 19\)纤维，其秩为三。二者的幺幂深度均为零。已有实际消失球的非零性可以在指定例中提供所需链输入，无需重做运输。

这和原物理周期的秩八/深度二、等质量很一般秩四/深度一并不矛盾：最大割只看到纯核，不会自动恢复质量Kummer扩张、顶层相对扩张和它们的比较常数。相同纯K3模型可承载不同扩张；留数像的精确几何识别不等于完整混合coaction已由最大割唯一重建。

还有一个不需新微分计算的结果。设单参数纯Gauss--Manin微分模P的秩为r，几何单值绝对不可约，\(\omega \ne 0\)是代数截面。则\(\omega \)的所有协变导数生成P，否则其生成的真微分子模会给单值不变真子系统。相应纯周期满足r阶标量Picard--Fuchs方程；任何非零周期芽在k(s)上的最小齐次方程也恰为r阶：若存在更低阶方程，其所有单值分支仍为该方程的解，但不可约性使这些分支张成r维，与低阶解空间维数矛盾。

应用于已认证几何闭包O(Q\textsubscript{6})的指定质量射线，得到非零纯最大割周期的最小齐次Picard--Fuchs阶数为六。这使用既有单值不可约性，不需要再生成六阶算子；也不把完整混合物理积分的一参数微分阶数与其动机秩直接等同。



\endgroup
\subsection{多个除子的过滤}
对简单正常交叉除子，单次留数在各交叉层上迭代，得到多重留数复形。符号由除子次序固定；过滤应在这一相容复形及其实际像上定义。完整构造与标架条件见\otherref{feynman_periods_supplement}{补充材料S5}。即使某个最大割给出纯核，较低留数及相对边界仍可承载扩张，故最大割一般不足以恢复完整比较矩阵。

\section{子图收缩与森林边界}\label{ch:contraction}

\begingroup

收缩的几何入口是图多项式的法向首项，而非把子图简单从积分中删除。经典图超曲面与边界理论见\cite{bek,browncosmic,bbr}。本节明确区分进入相对对象的面连接与从补空间出去的极点留数。

\subsection{数据和经典结果的位置}\label{sec:contraction:1}

设\(k\subset C\)为特征零域，G为连通图，内边数N、圈数h。每个顶点有外动量\(p_v\)，满足\(\sum p_v=0\)；每条边有质量平方\(\mu _e\)。先在欧氏正质量区域推导，再把有理恒等式延伸到其定义域。固定正整数传播子幂次\(a_e\)，记\(a_H=\sum _{e\in H}a_e\)。自环和平行边允许。外动量在收缩顶点处相加，保留边的质量不变。

使用补树约定

\begin{equation}
U_G=\sum_{T\text{生成树}}\prod_{e\notin T}\alpha_e,\qquad
 F_G=\Phi_G+U_G\sum_e\mu_e\alpha_e,\qquad
 \Phi_G=\sum_{T_1\sqcup T_2}(\sum_{v\in T_1}p_v)^2
                                  \prod_{e\notin T_1\sqcup T_2}\alpha_e.       \label{eq:contraction:1}
\end{equation}

二森林不计两个分支的交换重数。本文F是完整第二Symanzik多项式，对应Brown的Xi；不是该文仅含动量的Phi。

经典来源须明确：U的吹起边界分解见Bloch--Esnault--Kreimer\cite[Proposition 3.5]{bek}；含质量的UV/IR分解及实际边界面关系见Brown\cite[Proposition 2.2, Theorems 2.7, 6.5, 7.8]{browncosmic}。下面给出适合本文幂次、最小标架和RG接口的直接论证，不宣称这些经典分解的原创性。Bergbauer--Brunetti--Kreimer\cite{bbr}提供位置空间嵌套集与重整化的先行框架；不把其位置空间模型直接等同本参数空间模型。

取真连通边子图\(\gamma \)，\(r=|E_\gamma |\ge 2\)，\(h_\gamma \ge 1\)，且\(G/\gamma \)至少有一条边。最初假定

\begin{equation}
F_{G/\gamma}\not\equiv0.                  \label{eq:contraction:2}
\end{equation}

所有保留边质量正时自动满足\eqref{eq:contraction:2}。无质量或退化运动学中必须独立检查。下文的``UV子图对象''会主动忘掉\(\gamma \)的质量及外动量，只保留\(U_\gamma \)；不能把它误认为\(\gamma \)在原物理数据下的完整动机。


\subsection{生成树给出的法向首项}\label{sec:contraction:2}

在\(\gamma \)的边上置\(\alpha _e=t x_e\)，其余边置\(\alpha _e=y_e\)。记\(I_\gamma =(\alpha _e:e\in \gamma )\)。有

\begin{equation}
U_G(tx,y)=t^{h_\gamma}\big(U_\gamma(x)U_{G/\gamma}(y)+tR_U(t,x,y)\big),
\end{equation}

\begin{equation}
F_G(tx,y)=t^{h_\gamma}\big(U_\gamma(x)F_{G/\gamma}(y)+tR_F(t,x,y)\big).       \label{eq:contraction:3}
\end{equation}

\(R_U\)、\(R_F\)是多项式；第二式并不包含\(F_\gamma \)作为UV主因子。

\begin{proof}

对任一生成树T，\(T\cap \gamma \)是覆盖\(\gamma \)顶点的森林。设其分支数为c，则\(\gamma \)边变量在补树单项式中的次数为\(h_\gamma +c- 1\)。因此最低次数为\(h_\gamma \)，且恰在\(T\cap \gamma \)为\(\gamma \)的生成树时达到。这样的T与一对生成树\((T_\gamma ,T_{G/\gamma })\)一一对应：取逆像并并入\(T_\gamma \)即可，连通性及边数保证无圈。这证明第一式及余项界。

二森林中相同计数成立。达到\(h_\gamma \)时，\(\gamma \)完整位于同一分支；收缩不改变该分支总外动量，因此对应\(G/\gamma \)的二森林及其同一权重。故Phi的最低项为\(U_\gamma  Phi_{G/\gamma }\)。质量和分成\(\sum _{e\notin \gamma }\mu _e y_e+t\sum _{e\in \gamma }\mu _e x_e\)，后一项额外含t，得到第二式。\(U_\gamma  U_{G/\gamma }\)非零；\eqref{eq:contraction:2}保证F的指定首项也非零。

\end{proof}

这也是\(I_\gamma -adic\)分级环中的等式。令\(L_\gamma ={\alpha _e=0:e\in \gamma }\subset P^{N- 1}\)，吹起\(\pi :X=Bl_{L_\gamma }P^{N- 1}\to P^{N- 1}\)。因\(N_{L/P}=O(1)^{\oplus r}\)，例外除子

\begin{equation}
E_\gamma\simeq P^{r-1}\times P^{N-r-1}.          \label{eq:contraction:4}
\end{equation}

设Y是\(V(U_G F_G)\)的严格变换之并。由\eqref{eq:contraction:3}在E上去掉总变换的t幂后，得到

\begin{equation}
E_\gamma\setminus Y
  \simeq\big(P^{r-1}\setminus V(U_\gamma)\big)
       \times\big(P^{N-r-1}\setminus V(U_{G/\gamma}F_{G/\gamma})\big).
                                                                    \label{eq:contraction:5}
\end{equation}

这里用的是约化补空间，乘积中\(U_\gamma \)在总方程中出现几次不改变\eqref{eq:contraction:5}。几何上，严格变换与E的交是相应超曲面沿L的射影法向锥；不是把\(t=0\)直接代入原U、F，后者会同时得到零。

若\eqref{eq:contraction:2}失败，只能从\eqref{eq:contraction:3}得\(\operatorname{ord}_{I\gamma }F_G\ge h_\gamma +1\)，不能仍用\eqref{eq:contraction:5}或下节的F阶数。一个明确例子是无质量一圈二点bubble，\(U=\alpha _1+\alpha _2\)、\(F=q^{2}\alpha _1\alpha _2\)，缩放\(\gamma ={1}\)。虽然\(h_\gamma =0\)，但\(G/\gamma \)为无尺度自环，\(F_{G/\gamma }=0\)；F实际一阶消失。这个例子说明退化分支必须由物理数据识别。它只用于首项失效，不满足本文例外除子\(r\ge 2\)的假设。


\subsection{收缩面连接与UV留数是两个不同箭头}\label{sec:contraction:3}

令\(\mathcal U=X\setminus Y\)，B为坐标超平面严格变换在\(\mathcal U\)中的并，不把E重复算入B。X光滑，E与坐标边界正常交叉；删去Y后仍成立。置\(E^\circ=E\cap \)\(\mathcal U\)，\(B_E=B\cap E^\circ\)，\(n=N- 1\)。单次吹起时，\eqref{eq:contraction:5}与边界一起是乘积对：

\begin{equation}
(E^\circ,B_E)=(V_\gamma\times V_c,
                    B_\gamma\times V_c\ \cup\ V_\gamma\times B_c),            \label{eq:contraction:6}
\end{equation}

其中\(V_\gamma =P^{r- 1}\setminus V(U_\gamma )\)，\(V_c=P^{N- r- 1}\setminus V(U_c F_c)\)，\(B_\gamma \)、\(B_c\)为各自坐标边界。定义

\begin{equation}
M_\gamma^{UV}=H^{r-1}(V_\gamma,B_\gamma),\quad
 M_c=H^{N-r-1}(V_c,B_c),\quad
 M_G^{\mathrm{face}}=H^n(\mathcal U,B\cup E^\circ).                \label{eq:contraction:7}
\end{equation}

若需要进一步解消，只能使用已经构造、且在E上与两个因子的解消相容的模型；不能假定任意新增吹起不改变所有标架。Brown的motic模型提供这样的系统性选择\cite{browncosmic}。本节单次吹起的实际对已足够构造下列箭头。

\begin{theorem}[相对面连接]

\label{thm:contraction:3.1} 外积与三重对的连接给实际Nori态射

\begin{equation}
M_\gamma^{UV}\otimes M_c
 \xrightarrow{\ \times\ }H^{n-1}(E^\circ,B_E)
 \xrightarrow{\ \partial_E\ }M_G^{\mathrm{face}}.                \label{eq:contraction:8}
\end{equation}

其中没有Tate扭转；不保证复合单射。本文把连接号固定为Betti对偶提取E面的正乘积链：在正法向\(t\ge 0\)的领口中，通常向外边界取向与此相反，须同时调整连接的整体负号。这个选择不改变正合性。对提升后避开Y的正单纯形链\(\sigma _G\)，有

\begin{equation}
\langle\partial_E(\alpha\times\beta),\sigma_G\rangle
    =\langle\alpha,\sigma_\gamma\rangle
       \langle\beta,\sigma_c\rangle.                             \label{eq:contraction:9}
\end{equation}

如进一步边界仍碰到Y，\eqref{eq:contraction:9}只在完成所需解消并得到实际相对链后使用。

\end{theorem}

\begin{proof}

三重对\(B\subset B\cup E^\circ\subset \)\(\mathcal U\)的纤维三角给

\begin{equation}
H^{n-1}(\mathcal U,B)\longrightarrow H^{n-1}(E^\circ,B_E)
 \xrightarrow{\partial_E}H^n(\mathcal U,B\cup E^\circ).            \label{eq:contraction:10}
\end{equation}

中间项由闭并的相对Mayer--Vietoris/切除得到。复合相对外积即\eqref{eq:contraction:8}，不需要把全部\(H^{n- 1}\)强行等同单个Kunneth分量。连接的Betti对偶按上述号约定取\(\sigma _G\)的正E面；正实吹起在该面记录\(\gamma \)参数之比及剩余参数，正好是\(\sigma _\gamma \times \sigma _c\)。相对Stokes给\eqref{eq:contraction:9}。这些步骤先在实际几何对与六运算中完成，再取Nori上同调；\cite{tubach}保证与Hodge实现相容。

\end{proof}

式\eqref{eq:contraction:10}也给出精确可见性限制：外积类若从\(H^{n- 1}(\)\(\mathcal U\),B)限制而来，则面连接消去它。另一方面，\eqref{eq:contraction:9}右边非零时，连接类必非零。即使非零，它仍可能对应与原物理形式不同的de Rham类；``整个图对象含这个周期''不等于``给定振幅读出必须含它''。

\begin{theorem}[法向极点留数]

\label{thm:contraction:3.2} 另有正合段

\begin{equation}
H^n(\mathcal U,B)\longrightarrow
 H^n(\mathcal U\setminus E^\circ,B\setminus E^\circ)
 \xrightarrow{\operatorname{Res}_E}
 H^{n-1}(E^\circ,B_E)(-1)
 \xrightarrow{\operatorname{Gys}_E}H^{n+1}(\mathcal U,B).          \label{eq:contraction:11}
\end{equation}

它来自余维一纯性及相对局部化，证明已在第\ref{ch:cut}节展开。因此留数像精确为Gysin核。未扭转的裸形式留数满足

\begin{equation}
\int_{\operatorname{Tub}_E\tau}\omega
                     =2\pi i\int_\tau\operatorname{res}_E\omega.              \label{eq:contraction:12}
\end{equation}

\eqref{eq:contraction:8}从面对象进入相对对象；\eqref{eq:contraction:11}从删去面的补空间出去并带(-1)。两者来源、方向、次数、Tate因子及物理链均不同。divergent正单纯形不能仅凭\eqref{eq:contraction:8}当作删面补空间中的普通有限相对周期；需先给正规化和减除。

\end{theorem}


\subsection{幂次计数与对数边界形式}\label{sec:contraction:4}

取偶整数\(D_0\)。去掉整体Gamma与测度常数后的标量射影形式为

\begin{equation}
\omega_G(D)=\Big(\prod_e\alpha_e^{a_e-1}\Big)
       \frac{U_G^{A_G(D)}}{F_G^{B_G(D)}}\Omega_G,
 \quad A_G=a_G-(h_G+1)D/2,\quad B_G=a_G-h_GD/2.                  \label{eq:contraction:13}
\end{equation}

取标准\(\Omega _H=\sum _i(- 1)^{i- 1}\alpha _i d\alpha _1\wedge \)\ldots{}\(\wedge \)省略\(d\alpha _i\wedge \)\ldots，子图和收缩图继承边顺序；留数按法向在前。\(D=D_0\)时是有理形式；D的非整数变化仅作解析正规化，并不称其为普通Nori对象。

置

\begin{equation}
d_\gamma=a_\gamma-D_0h_\gamma/2.        \label{eq:contraction:14}
\end{equation}

在收缩图某一非零坐标设为1的图中，再用\(\gamma \)的一条非零方向坐标定义t，Jacobian是\(t^{r- 1}\)；幂次分子额外给\(t^{a\gamma - r}\)，\eqref{eq:contraction:3}的分母净贡献为\(t^{- D_0h\gamma /2}\)。故

\begin{equation}
\operatorname{ord}_{E_\gamma}\pi^*\omega_G(D_0)
                                  =d_\gamma-1.                 \label{eq:contraction:15}
\end{equation}

等号指E的泛点，且\eqref{eq:contraction:2}成立。\(a_e\)正整数保证单项式分子在方向泛点不消失；它不排除交叉层上的额外阶数。

\begin{theorem}[对数收缩留数]

\label{thm:contraction:4.1} 若\(d_\gamma =0\)，则

\begin{equation}
\operatorname{res}_{E_\gamma}\pi^*\omega_G(D_0)
       =s_\gamma\eta_\gamma^{UV}\boxtimes\omega_{G/\gamma}(D_0),\qquad
 \eta_\gamma^{UV}=\frac{\prod_{e\in\gamma}x_e^{a_e-1}\Omega_\gamma}
                                      {U_\gamma^{D_0/2}}.        \label{eq:contraction:16}
\end{equation}

其中\(s_\gamma \)为体积形式的显式取向号。将\(\gamma \)的r条边移到最前、保持两块内部顺序，设该置换符号为\(sgn(\pi _\gamma )\)，则\(s_\gamma =(- 1)^r sgn(\pi _\gamma )\)。采用相容正密度约定时可把这个号同时吸收入链取向，不能只改形式。

\end{theorem}

\begin{proof}

法向幂由\eqref{eq:contraction:15}为dt/t。\(U_\gamma \)净指数为\(A_G- B_G=- D_0/2\)；且\(a_\gamma =D_0h_\gamma /2\)、\(h_G=h_\gamma +h_{G/\gamma }\)，给\(A_G=A_{G/\gamma }\)、\(B_G=B_{G/\gamma }\)。在\(\gamma \)已排最前、\(\alpha _r=t\)、\(x_r=1\)、\(\alpha _N=1\)的图上，直接楔积得到\(\pi \)*\(\Omega _G=(- 1)^r t^{r- 1}dt\wedge \Omega _\gamma \wedge \Omega _c\)；恢复原边序再乘\(sgn(\pi _\gamma )\)。这给\eqref{eq:contraction:16}。\(\eta _\gamma \)的总齐次次数为零，确实是\(\gamma \)方向射影形式。

\end{proof}

若\(\gamma \)的UV方向积分收敛，定义其形式周期\(p_\gamma \)。它的几何对、形式与正链都不依赖原\(\gamma \)质量及外运动学。因此\eqref{eq:contraction:16}给一个常值的算术对象乘上收缩图对象。这具体实现了局部性的一种来源：外运动学被忘掉，而\(p_\gamma \)的算术不必是Tate。在Tannakian层面，常值变分不等于Tate对象；不能把外动量导数为零当作混合Tate判据。

分子会改变判据。对齐次有理形式\(\Omega _G P/(U_G^u F_G^v)\)，若P的\(\gamma \)法向次数为\(\nu _\gamma (P)\)，則泛点阶数是

\begin{equation}
r-1+\nu_\gamma(P)-h_\gamma(u+v).                 \label{eq:contraction:17}
\end{equation}

其首项\(P_0\)一般是方向变量与收缩变量的和式，并不分解为一个乘积。因而它只给张量通道之和，且可使\eqref{eq:contraction:16}的候选留数消失；单位标量的乘积结论不原样覆盖任意分子或Ward投影。


\subsection{极点由哪一阶法向数据决定}\label{sec:contraction:5}

令\(D=D_0- 2\epsilon \)，先固定避开其余奇异面的紧致方向区域，或已经对其余面作出有效减除。则某个法向积分可写成

\begin{equation}
J(\epsilon)=\int_0^1 t^{d_\gamma-1+h_\gamma\epsilon}
                                A(t,\epsilon)\,dt,             \label{eq:contraction:18}
\end{equation}

A在\(t=0\)充分光滑，在\(\epsilon =0\)解析，所需余项一致可积。这里\(\Gamma (B_G)\)、额外正规化因子和其他尺度的极点尚未乘入。

若\(d_\gamma =0\)，减加\(A(0,\epsilon )\)立即得到

\begin{equation}
J(\epsilon)=\frac{A(0,\epsilon)}{h_\gamma\epsilon}
      +\int_0^1 t^{-1+h_\gamma\epsilon}(A(t,\epsilon)-A(0,\epsilon))dt,
 \qquad\operatorname{Res}_{\epsilon=0}J=A(0,0)/h_\gamma.          \label{eq:contraction:19}
\end{equation}

对数情况下，法向锥主项正好够用。但\textbf{这只是指定法向端点的裸积分极点}；反项采用\(C=- K\)时另有负号，整图的Gamma极点、森林和物理求和需按实际方案另行组合。

若\(d_\gamma =- k<0\)，展开\(A=\sum _{j=0}^k A_j(\epsilon )t^j+t^{k+1}\widetilde A\)，逐项积分给

\begin{equation}
J(\epsilon)=\sum_{j=0}^k\frac{A_j(\epsilon)}{h_\gamma\epsilon+j-k}
                  +H(\epsilon),\qquad
 \operatorname{Res}_{\epsilon=0}J=A_k(0)/h_\gamma.                \label{eq:contraction:20}
\end{equation}

H在\(\epsilon =0\)解析。首项\(A_0\)控制截断意义下最严重的幂次发散，但\(\epsilon =0\)的留数由第k阶法向喷射给出。比如\(A=1+c t\)、\(d_\gamma =- 1\)时，\(J=(h_\gamma \epsilon - 1)^{-1}+c/(h_\gamma \epsilon )\)：同一首项可有不同的MS极点。这是直接反例，说明``只看法向锥即可给所有反项''不成立。

\eqref{eq:contraction:20}是Schwinger尺度坐标中的公式，不直接等于外动量Taylor算符。两者在指定场论中的比较还需分子、局部算符插入和IR控制；不能把这里的k机械读成外动量微分阶数。一般正规化\(a_e(\epsilon )=a_e+\epsilon r_e\)时，分母\(h_\gamma \)改为\(h_\gamma +\sum _{e\in \gamma }r_e\)；若它为零，这个\(\epsilon \)方向根本不调节该端点。


\subsection{相容森林的多尺度因子化}\label{sec:contraction:6}

采用明确而足够的相容条件：有限集合\(\mathscr F\)由真连通子图组成，任意两个或边集包含，或顶点集不交。对每个\(H\in \)\(\mathscr F\)\(\cup {G}\)，令child(H)为其最大真子成员，定义

\begin{equation}
H^\circ=H/\bigcup_{K\in child(H)}K.      \label{eq:contraction:21}
\end{equation}

要求每个H°至少一条边，根G°满足\(F_{G^\circ}\ne 0\)。顶点不交条件保证子成员的圈数可加；可以放宽，但必须另证相应秩可加性。仅说``边不交''并不足以支持下面的证明。

在每个H°选一条标记边\(e_H\)。置

\begin{equation}
\alpha_e=u_e\prod_{\substack{\gamma\in\mathscr F\\e\in\gamma}}t_\gamma,
                  \qquad u_{e_H}=1.                             \label{eq:contraction:22}
\end{equation}

根标记固定射影尺度。每层标记先定义该层相对父层的尺度，故\eqref{eq:contraction:22}在非零标记图上可逆；\(t_\gamma =0\)是相应迭代吹起图的边界。体积形式的Jacobian因子为\(\prod _\gamma  t_\gamma ^{|E_\gamma |- 1}\)，可按最外层到最内层逐次计算；嵌套子成员的所有边都计入外层指数。

\begin{theorem}

\label{thm:contraction:6.1} 去掉法向单项式后在\(t_\gamma =0\)对所有\(\gamma \)取值，有

\begin{equation}
U_G=\Big(\prod_\gamma t_\gamma^{h_\gamma}\Big)
       \left(\prod_{H\in\mathscr F\cup\{G\}}U_{H^\circ}+R\right),
\end{equation}

\begin{equation}
F_G=\Big(\prod_\gamma t_\gamma^{h_\gamma}\Big)
       \left(F_{G^\circ}\prod_{\gamma\in\mathscr F}U_{\gamma^\circ}
                                                       +R'\right),          \label{eq:contraction:23}
\end{equation}

R、R'属于由全部\(t_\gamma \)生成的理想。若所有\(d_\gamma =0\)，则

\begin{equation}
\operatorname{res}_{\mathscr F}\pi^*\omega_G(D_0)
        =s_{\mathscr F}\mathop{\boxtimes}_{\gamma\in\mathscr F}\eta_{\gamma^\circ}^{UV}
                                                \boxtimes\omega_{G^\circ}.
                                                                    \label{eq:contraction:24}
\end{equation}

这里\(s_F\)由选定消去次序中\eqref{eq:contraction:16}的各取向号及把法向移到统一顺序的Koszul号相乘得到。法向顺序及因子顺序同时固定；交换法向次序的反号不被省略。

\end{theorem}

\begin{proof}

最内层生成树可先在顶点不交的子成员中独立选取。其并是森林，收缩后选父商图生成树，再提升，逐层得到同时适应全部成员的树。反过来，每棵适应树唯一恢复这些商图树。对每个\(\gamma \)，非适应树的补树单项式中\(\gamma \)次数至少\(h_\gamma +1\)；所有树的该次数至少\(h_\gamma \)。故可同时提出\(\prod t_\gamma ^{h_\gamma }\)，常数项精确为各商图U之积。根用二森林并保留外运动量及剩余质量，给F公式；子成员内部质量每次进入余项。这个论证也证明首项不依赖选取哪一个包含关系的线性延伸次序。

圈数和a在收缩下可加；全部\(d_\gamma =0\)蕴含每个非根商图\(\gamma ^\circ\)仍对数。将\eqref{eq:contraction:23}及Jacobian代入\eqref{eq:contraction:13}，每个法向给\(dt_\gamma /t_\gamma \)，各非根U指数为\(- D_0/2\)，根指数与其自身\eqref{eq:contraction:13}相同，证明\eqref{eq:contraction:24}。

\end{proof}

在多尺度\(A(t,\epsilon )\)解析且余下积分一致可积的单一乘积图上，\eqref{eq:contraction:24}还给

\begin{equation}
[\epsilon^{-|\mathscr F|}]\int\prod_{\gamma\in\mathscr F}
                t_\gamma^{h_\gamma\epsilon-1}A(t,\epsilon)\,dt
                       =\frac{A(0,0)}{\prod_{\gamma\in\mathscr F}h_\gamma}.
                                                                    \label{eq:contraction:25}
\end{equation}

若还有方向变量，应对右侧继续积分。证明是对各变量反复使用\eqref{eq:contraction:19}；含至少一个\(t_\gamma \)的余项最多有少一个法向极点。式\eqref{eq:contraction:25}不含整体Gamma、其他图、其他扇区及其重叠计数，因此不是未经拼接的整图最高极点公式。

\textbf{最高留数不足以重建RG简单极点。} 两个相容尺度就有反例：

\begin{equation}
\int_0^1\!\int_0^1
 t_1^{h_1\epsilon-1}t_2^{h_2\epsilon-1}(1+c t_1)\,dt_1dt_2
 =\frac1{h_1h_2\epsilon^2}
                  +\frac{c}{h_2\epsilon(1+h_1\epsilon)}.         \label{eq:contraction:25a}
\end{equation}

所有c的最深角点首项都等于1，但简单极点相差\(c/h_2\)。即使所有端点均对数，也必须保留较浅面上的数据及其余变量积分；一般还须保留其\(\epsilon \)展开。这个有理积分反例与\otherref{feynman_periods_supplement}{补充材料S10}中``有限项先插入后可进入极点''的代数反例相容。它不声称任意c对应实际Feynman图；它严格排除了仅凭最深森林乘积首项推出全部局部RG系数的形式论证。


\subsection{重叠子图为什么不能同时作为独立残数}\label{sec:contraction:7}

令A、C为两个不可比边集，\(L_A\)、\(L_C\)为坐标中心。若\(A\cup C\ne E_G\)，则交为\(L_{A\cup C}\)。先吹起该交后，\(L_A\)、\(L_C\)的严格变换不交：它们在交中心法向空间中的方向子空间交为零，射影化后没有公共点。后续吹起二者产生的例外除子也无共同层。若\(A\cup C=E_G\)，则原射影中心已经不交。

这说明在包含该交中心的模型里，不能给这两个原标签捏造共同的二重留数面。对于相交边集，这是处理重叠的一条具体规则；对边不交的横截中心，也可以选择再吹起其交，此时模型把同时面细分为旗标面。因此\textbf{边界层由选定building set决定，不能无条件说``所有Zimmermann森林就是某一个模型的全部层''}。不同细分必须通过实际拉回、链分解和取向比较。

例：四维三线sunset的三个二边bubble均对数发散，但对应P\textsuperscript{2}中的三个不同坐标点。吹起后所得三个例外曲线两两不交。因此一个局部图最多含其中一个proper-bubble留数；三项须在不同面/区域求和，不能把三bubble相乘当作共同三重极点。整体图另有幂次发散及整体尺度，故这里也没有断言整图只有简单\(\epsilon \)极点。

在已经得到的正常交叉对数图上，局部Taylor零阶算符\(T_i A=A|_{t_i=0}\)彼此交换，

\begin{equation}
\prod_i(1-T_i)A=\sum_{S}(-1)^{|S|}\prod_{i\in S}T_i A          \label{eq:contraction:26}
\end{equation}

消去所有对应端点的对数不可积性。更高幂次用第\ref{sec:contraction:5}节所需有限阶Taylor算符。这个局部几何包含--排除式解释森林项的来源；要等同某个全局物理BPHZ/MS方案，仍需给全局拼接、重叠计数、外动量Taylor/局部插入比较及IR处理。\cite{bbr}在其位置空间、至多对数的主范围内完成了相应方案理论，不能未经比较就移植为本文所有参数积分的定理。


\subsection{本文banana的检验性推论}\label{sec:contraction:8}

二维h线banana，单位传播子、正质量，任取k条边组成的真连通子图\(\gamma \)，有\(h_\gamma =k- 1\)、\(a_\gamma =k\)，因此

\begin{equation}
d_\gamma=k-(k-1)=1.                \label{eq:contraction:27}
\end{equation}

每个这样的UV面上参数形式泛点阶数为零，\textbf{没有对数UV留数}。整体也有\(a_G- D_0h_G/2=1\)，无整体UV Gamma极点。四线情形\eqref{eq:contraction:13}恰为\(\Omega _G/F_G\)；其K3及已证混合周期来自完整几何与边界扩张，而不是UV极点。

所以``UV筛选为空''绝不推出原物理最小对象是Tate或零。本文指定四线对象的纯K3核与秩八最小性是既有独立证据；\eqref{eq:contraction:27}只用圈数和法向计数，未重算这些证据。

另一方面，若一个对数子图的\(p_\gamma \)已经在实际形式周期中证明含某算术方向，\eqref{eq:contraction:16}只证明该方向可沿对应收缩通道出现。完整森林与物理求和仍可能消去它。只有它在规定MS型的完整局部RG数据中非零后，\otherref{feynman_periods_supplement}{补充材料S10}的反演才保证它不会在所有临界系数中同时消失。这样连接三层时，每一次``可见''都有自己的实际映射与非零条件。



\endgroup

\section{本原减除与共同森林提升}\label{ch:primitive}\label{ch:forest}

\begingroup

以下先处理没有子发散的范围，并完整固定物理减除顺序。含子发散的扩展在第\ref{ch:zero}、\ref{ch:finitemom}节讨论。参数减除的背景见\cite{bk2008}，Kummer 周期关系见\cite{brownperiods}。

\subsection{精确输入与物理归一化}\label{sec:primitive:1}

沿用第\ref{ch:contraction}节的U、完整含质量F以及标准射影体积形式\(\Omega \)。G连通，内边数N、圈数\(h\ge 1\)，\(n=N- 1\)，\(D_0\)为正偶整数，传播子幂次\(a_e\)为正整数。给定

\begin{equation}
a_G=D_0h/2,\qquad
 d_S:=a_S-D_0h_S/2>0\quad(\varnothing\ne S\subsetneq E_G).       \label{eq:primitive:1}
\end{equation}

\(h_S\)按边子图的所有连通分支计算；不只检查连通子图。\eqref{eq:primitive:1}是本文采用的严格UV原始性条件。\(N>1\)时它排除桥和自环：删桥后的d为负；一个自环与其余边的d相加为零，不能同时正。

两组数据\(\kappa \)、\(\kappa _0\)均有每边质量平方\(\mu _e>0\)和欧氏外动量，动量守恒；其系数属于固定\(k\subset R\)。在数域/代数参数纤维上得到k上的Nori对象；在参数族上先限制到上同调局部常值的光滑开集。\(\kappa _0\)可以与\(\kappa \)相同。

取\(D=D_0- 2\epsilon \)，欧氏动量测度每圈\(\pi ^{- D/2}\)归一化，暂不含耦合、图对称因子及Minkowski的i因子。设

\begin{equation}
c_a=\frac1{\prod_e\Gamma(a_e)}\in\mathbb Q_{>0},\qquad
 \eta=\frac{\prod_e\alpha_e^{a_e-1}\Omega_G}{U_G^{D_0/2}},\qquad
 p_G=c_a\int_\sigma\eta.                                        \label{eq:primitive:2}
\end{equation}

\(\sigma \)取使Schwinger密度为正的取向；若用\(\sum \alpha =1\)且\(\alpha _1,\)\ldots{}\(,\alpha _{N- 1}\)递增的坐标，标准\(\Omega \)拉回有\((- 1)^{N- 1}\)号，须同时计入链取向。本文不改动旧banana框架。

定义裸积分及减除

\begin{equation}
I_\epsilon(\kappa)=c_a\,\mu_R^{2h\epsilon}\Gamma(h\epsilon)
 \int_\sigma\eta\,U_G^{(h+1)\epsilon}F_\kappa^{-h\epsilon},\qquad
 R_G(\kappa,\kappa_0)=\lim_{\epsilon\to0}
                   (I_\epsilon(\kappa)-I_\epsilon(\kappa_0)).      \label{eq:primitive:3}
\end{equation}

减除顺序明确为当前点减参考点。某些文献振幅含相反整体符号；不能把其log比值直接替换\eqref{eq:primitive:3}。


\subsection{收敛与取消极限有统一证明}\label{sec:primitive:2}

取任意Hepp扇区，边序\(e_1,\)\ldots{}\(,e_N\)满足\(\alpha _{e_1}\le \)\ldots{}\(\le \alpha _{e_N}\)，设\(\alpha _{e_N}=1\)及

\begin{equation}
\alpha_{e_i}=t_i\cdots t_{N-1},\quad 0\le t_i\le1,\quad
 S_j=\{e_1,\ldots,e_j\}.                                        \label{eq:primitive:4}
\end{equation}

Kruskal逐边扩展森林给一棵同时在每个\(S_j\)上取最大森林的生成树。故每个补树单项式的\(t_j\)次数至少\(h_{S_j}\)，且有一个单项式对所有j同时达到该下界。因此

\begin{equation}
U_G=\prod_{j=1}^{N-1}t_j^{h_{S_j}}\widehat U,
             1\le\widehat U\le\tau(G)\quad\hbox{在实单位立方体上}.           \label{eq:primitive:5}
\end{equation}

Jacobian及幂次分子给

\begin{equation}
\eta=\pm\prod_jt_j^{d_{S_j}-1}\widehat U^{-D_0/2}\,dt_1\cdots dt_{N-1}.
                                                                    \label{eq:primitive:6}
\end{equation}

由\eqref{eq:primitive:1}，每个指数非负。于是\(p_G\)有限且正。

把任意二森林延伸为生成树，其补单项式是该树的补单项式乘一条边参数。故在\(\alpha \)最大值为1的图上，Phi/U有有限上界。质量项则给\(F_\kappa \ge \mu _{e_N}U\)，因而

\begin{equation}
F_\kappa=\prod_jt_j^{h_{S_j}}\widehat F_\kappa,\qquad
       0<c\le\widehat F_\kappa,\widehat F_{\kappa_0}\le C<\infty.             \label{eq:primitive:7}
\end{equation}

常数可在任意给定的紧致正质量欧氏参数区域统一选取。

这里还有一个不依赖扇区的比较：正质量保证F的单项式支撑等于U·\(\sum \alpha \)的支撑；动量二森林项不增加新的单项式类型。因此\(F_\kappa \)和\(F_{\kappa 0}\)系数逐项可比较，\(F_{\kappa 0}/F_\kappa \)在整个正单纯形上有正的上下界。

\begin{theorem}

\label{thm:primitive:2.1} \eqref{eq:primitive:3}存在，且

\begin{equation}
R_G(\kappa,\kappa_0)=c_a\int_\sigma\eta
                            \log\frac{F_{\kappa_0}}{F_\kappa},
 \qquad \operatorname{Res}_{\epsilon=0}I_\epsilon=p_G/h.          \label{eq:primitive:8}
\end{equation}

\end{theorem}

\begin{proof}

\eqref{eq:primitive:5}--\eqref{eq:primitive:7}把正规化后各端点幂改为\(d_{S_j}- 1+h_{S_j}\epsilon \)。对足够小的\(|Re \epsilon |\)，所有积分及有限次\(\epsilon \)导数由可积幂乘有限对数次幂支配。记\(F_0=F_{\kappa 0}\)，则

\begin{equation}
\frac{F_\kappa^{-h\epsilon}-F_0^{-h\epsilon}}{h\epsilon}
                    \longrightarrow\log(F_0/F_\kappa).
\end{equation}

比值的对数有界，剩余公共因子也满足上述支配。又\(h\epsilon \Gamma (h\epsilon )\to 1\)，得到第一式。未作减除时\(\Gamma (h\epsilon )\)的简单极点给第二式。

\end{proof}

这比在积分号内形式取消\(1/\epsilon \)更强：给出了极限交换所需的端点控制。任意公共解析正规化因子\(A(\epsilon )\)、\(A(0)=1\)，都不改变\(R_G\)；它可以改变单个有限部分。


\subsection{一个全局正链避让的模型}\label{sec:primitive:3}

令X为\(P^{N- 1}\)沿所有非空坐标中心作的标准扇区细分模型，坐标边界\(D_X\)正常交叉。等价地使用完整坐标旗标的光滑toric模型，\eqref{eq:primitive:4}给其正扇区图。\(N=1\)时X为点。记\(Y_U\)、\(Y_F\)、\(Y_{F0}\)为相应非边界严格变换，置

\begin{equation}
V=X\setminus(Y_U\cup Y_F\cup Y_{F0}),\quad D=D_X\cap V,
                         r=F_0/F_\kappa.                         \label{eq:primitive:9}
\end{equation}

\eqref{eq:primitive:5}--\eqref{eq:primitive:7}证明提升后的正单纯形\(\widetilde\sigma\)是V内的紧致流形带角，其边界在D。F和\(F_0\)在每个坐标除子上具有相同阶数，故其法向单项式取消；r在V上是代数单位。\(\eta \)由\eqref{eq:primitive:6}延伸为V上的正则n形式，且在D每个分量上的拉回为零。这些是全局构造所需的正链与形式，不只是扇区上的渐近记号。

V还删去F和\(F_0\)，可能承载超出UV周期的几何。因此不能仅根据\(p_G\)的最小对象判断\(H^n(V,D)\)的全部结构。UV周期本身来自不依赖\(\kappa \)、\(\kappa _0\)的较大开空间\(X\setminus Y_U\)；限制到V是实际拉回，不改变给定\(\eta \)、\(\sigma \)的形式周期关系。


\subsection{分开端点的实际相对对象}\label{sec:primitive:4}

若直接使用(Gm,\{1,r\})，\(r=1\)时两个端点合并，普通相对H\textsuperscript{1}由秩二降为秩一。这不是保持标记的Kummer家庭，不能用于无条件的对角特殊化。

引入额外坐标w，只用来保留端点标记：

\begin{equation}
Y=V\times\mathbb G_{m,z}\times\mathbb A^1_w,\qquad
 s_0(x)=(x,1,0),\quad s_1(x)=(x,r(x),1),
\end{equation}

\begin{equation}
B_R=(D\times\mathbb G_m\times\mathbb A^1)
                          \cup s_0(V)\cup s_1(V),\qquad
                M_R=H^{n+1}(Y,B_R).                              \label{eq:primitive:10}
\end{equation}

两个截面总不相交，包括\(r=1\)。\(B_R\)不必纯余维一；一般相对Nori上同调容许这种闭边界，不能为套用除子留数而假装它是光滑除子。

定义

\begin{equation}
\Omega_R=c_a\,\operatorname{pr}_V^*\eta\wedge\frac{dz}{z},
\quad\Gamma_R(x,t)=(x,1+t(r(x)-1),t),\quad
 (x,t)\in\widetilde\sigma\times[0,1].                            \label{eq:primitive:11}
\end{equation}

链取``基底在前，t在后''的取向。\(r>0\)保证\(z>0\)，所以\(\Gamma _R\)紧致且避开\(z=0\)；其端点落在\(s_0\)、\(s_1\)，其侧边落在D的逆像。

\begin{theorem}[实际减除周期]

\label{thm:primitive:4.1} \(\Omega _R\)定义\(M_R\)的de Rham类，\(\Gamma _R\)定义其有理Betti对偶标架，且

\begin{equation}
R_G^{\mathfrak m}(\kappa,\kappa_0)
           =[M_R,[\Gamma_R],[\Omega_R]]^{\mathfrak m},\qquad
          \operatorname{per}R_G^{\mathfrak m}=R_G(\kappa,\kappa_0).           \label{eq:primitive:12}
\end{equation}

\end{theorem}

\begin{proof}

\(\Omega _R\)闭；在\(s_0\)上拉回为零，在\(s_1\)上为\(\eta \wedge d\log  r=0\)，因为V维数n；在D逆像上\(\eta \)已为零。在边界闭并的相对de Rham复形中取其余分量为零，得到实际相对类。链条件已由\eqref{eq:primitive:11}验证。拉回dz/z时，基底微分项与\(\eta \)楔积为零，剩下

\begin{equation}
\Gamma_R^*\Omega_R=c_a\eta\wedge
                  \frac{r-1}{1+t(r-1)}dt.
\end{equation}

对t积分恰为log r，按\eqref{eq:primitive:8}得到\eqref{eq:primitive:12}。整个构造都在代数几何对上完成；Nori/Hodge比较使用实际六运算\cite{tubach}，没有由数值周期相等倒推动机态射。

\end{proof}

式\eqref{eq:primitive:11}也把减除振幅写成收敛的有理形式积分。额外w坐标没有新增物理自由度，其作用是使端点为分离的代数截面。


\subsection{Kummer系数及推前的精确核商}\label{sec:primitive:5}

沿\(f:Y\to V\)取相对端点上同调。纤维\((\mathbb G_m\times A^{1},{(1,0),(r,1)})\)的相对上同调只在一次非零。得到秩二Kummer系数对象\(\mathcal K_r\)，有

\begin{equation}
0\longrightarrow\mathbb Q(0)\longrightarrow\mathcal K_r
                      \longrightarrow\mathbb Q(-1)\longrightarrow0.       \label{eq:primitive:13}
\end{equation}

这是实际相对家庭的序列；并非凭扩张类猜测一个矩阵。设\(e_0=[dw]\)、\(e_1=[dz/z]\)为纤维相对de Rham基，Betti链取\(\gamma _0\)为\eqref{eq:primitive:11}的纤维段，\(\gamma _1\)为正向z圆，则在正实支

\begin{equation}
\left(\int_{\gamma_j}e_i\right)_{i,j=0,1}
       =\begin{pmatrix}1&0\\\log r&2\pi i\end{pmatrix},\qquad
 \nabla e_0=0,\quad \nabla e_1=d\log r\,e_0.                    \label{eq:primitive:14}
\end{equation}

\(e_0\)虽然在绝对复形中是恰当形式，在端点相对复形中不是，因为w在两端取值0、1。Hodge层F\textsuperscript{1}由\(e_1\)生成，\(W_0\)由\(e_0\)生成。r为常数时\(d\log  r=0\)仍可能有非平凡算术Kummer扩张；连接式不能取代比较条目log r。

记\(H^j=H^j(V,D)\)。相对推前与边界映射锥相容，给

\begin{equation}
M_R\simeq H^n(V,D;\mathcal K_r).               \label{eq:primitive:15}
\end{equation}

可直接用Y的乘积投影及两个截面的Čech复形验证：相对纤维复形是\(K_r[- 1]\)，限制到D给同一构造，再取相对纤维。此处不依赖未验证的一般非proper基变换。

令\(\delta _j:H^j(-1)\to H^{j+1}\)为\eqref{eq:primitive:13}产生的连接，则

\begin{equation}
0\longrightarrow\operatorname{coker}\delta_{n-1}
       \longrightarrow M_R\longrightarrow\ker\delta_n\longrightarrow0.
                                                                    \label{eq:primitive:16}
\end{equation}

式\eqref{eq:primitive:16}也可完全绕开系数对象的记号而从实际几何对导出。令\(A=D\times \mathbb G_m\times A^{1}\)、\(C=s_0(V)\)\ensuremath{\coprod}\(s_1(V)\)。三重对和闭并的Mayer--Vietoris给

\begin{equation}
\cdots\to H^j(Y,A)\to H^j(C,C\cap A)
             \to H^{j+1}(Y,A\cup C)\to\cdots .
\end{equation}

以\(z=1\)分裂Gm的常值因子，Kunneth给\(H^j(Y,A)=H^j\oplus H^{j-1}(-1)\)，而\(H^j(C,C\cap A)=H^j\oplus H^j\)。两个截面的拉回把(a,b)送至\((a,a+\delta _{j-1}b)\)，其中\(\delta \)由\(r:V\to \mathbb G_m\)的类与b杯积得到，按``基底形式在前''固定符号。在\(j=n\)、n+1处分别取余核和核，便得到\eqref{eq:primitive:16}。所用外积、截面拉回及正合列均为实际Nori态射。

这是需要保留的真实扩张结构，而不是无条件的\(M_R=H^n\otimes \)某个固定秩二对象。其Betti连接由r的Kummer拓扑类杯积给出，de Rham中对应\(d\log  r\)；连接之外还保留算术扩张数据。

\begin{corollary}

\label{thm:primitive:5.1} 设\(b_n=\operatorname{dim} H^n_B\)，则

\begin{equation}
\dim(M_R)_B=2b_n-\operatorname{rank}\delta_{n-1}
                         -\operatorname{rank}\delta_n\le2b_n.               \label{eq:primitive:17}
\end{equation}

\(M_R\)的简单组成来自\(H^n\)及\(H^n(-1)\)的子商。因此\(H^n\)若已证明mixed Tate，则\(M_R\)及减除周期的最小对象也mixed Tate；反向不成立。这里只给载体的精确核商和充分判据，没有声称由边表即可求得\(b_n\)或\(\delta _j\)。

在\(r=1\)时，\(K_1\)实际分裂为\(Q(0)\oplus Q(-1)\)：投影到(Gm,\{1\})分裂Q(-1)因子，投影到(A\textsuperscript{1},\{0,1\})给Q(0)因子；两截面仍不同。更严格地，\(\Gamma _R\)由代数子对

\begin{equation}
(V\times\{1\}\times\mathbb A^1,
 D\times\{1\}\times\mathbb A^1\ \cup\ V\times\{1\}\times\{0,1\})
                       \longrightarrow(Y,B_R)
\end{equation}

中的相对链推前而来，\(\Omega _R\)沿此态射的de Rham拉回恒为零。因此形式周期关系本身给\(R_G^m(\kappa _0,\kappa _0)=0\)，而非仅由数值积分为零倒推形式周期为零。对象没有因端点合并而降秩。

\end{corollary}


\subsection{减除点变换是形式周期恒等式}\label{sec:primitive:6}

对第三组正质量欧氏数据\(\kappa _1\)，有

\begin{equation}
R_G^{\mathfrak m}(\kappa_2,\kappa_0)
 =R_G^{\mathfrak m}(\kappa_2,\kappa_1)
                         +R_G^{\mathfrak m}(\kappa_1,\kappa_0).              \label{eq:primitive:18}
\end{equation}

\begin{proof}

在同时删去三个F零集的共同基底上，\(r_{20}=r_{10}r_{21}\)。用标记为\(w=0\)、1、2的三个截面，分别取\(z=1\)、\(r_{10}\)、\(r_{20}\)。第一段是\(1\to r_{10}\)；第二段由\(1\to r_{21}\)经代数映射\(z\mapsto r_{10}(x)z\)、\(w\mapsto w+1\)得到。该映射将\(\eta \wedge d\log  z\)拉回为自身，因为\(\eta \wedge d\log  r_{10}=0\)。两段的和与\(1\to r_{20}\)的正路径在\(z>0\)的纤维条带中相对同调，侧边在D。三个标记点的几何对及实际变换给\eqref{eq:primitive:18}的形式周期关系，不需要假设period映射单射。

\end{proof}

文献\cite{brownperiods}给普通Kummer对数的同类乘法关系；本文把端点标记、随基底变化的r以及\(\eta \)的高次形式一同保留。对固定参考点\(\kappa _0\)、\(\kappa _1\)，\eqref{eq:primitive:18}的差与可变外数据\(\kappa _2\)无关，正是本原对数图的有限常数减除变化。


\subsection{运动学导数及微分系统上界}\label{sec:primitive:7}

在共同光滑参数开集上，固定\(\kappa _0\)，\(\eta \)来自不变的U几何。\eqref{eq:primitive:14}的相对de Rham模型给

\begin{equation}
\nabla_v[\Omega_R]=\iota\!\left[c_a\eta\,v\log r\right],\qquad
 vR_G=-c_a\int_\sigma\eta\,\frac{vF_\kappa}{F_\kappa},          \label{eq:primitive:19}
\end{equation}

其中\(\iota:H^n(V,D)\to M_R\)来自\eqref{eq:primitive:13}的Q(0)项。第一式是类的等式：对\(\eta e_1\)应用系数连接，\(\eta \)在参数方向水平；同一相对复形中得到\(\eta (vlog r)e_0\)。第二式再取实际链周期；第\ref{sec:primitive:2}节的正下界给微分与积分交换。

注意\eqref{eq:primitive:19}的右端一般依赖运动学，不能当成常值UV局部投影。它消去一层对数，但留下F几何。

\begin{theorem}[单参数阶数界]

\label{thm:primitive:7.1} 限制到一个光滑代数曲线参数s，设\(b_n\)为该族\(H^n(V,D)\)的秩。减除周期满足一个k(s)系数齐次微分方程，阶数至多\(b_n+1\)。

\end{theorem}

\begin{proof}

\(\operatorname{im}\iota\)为水平子对象，所有\(\nabla _s^j[\Omega _R]\)、\(j\ge 1\)均在其中。设这些导数生成的微分子模秩为\(d\le b_n\)，则连同\([\Omega _R]\)的循环子模秩至多d+1；d+2个连续导数必线性相关。对水平Betti链取周期即可。也可先找阶数\(\le d\)的P使\(P\partial _{\mathrm{sR}}=0\)，再写\(P\partial _s\)。无须先算R的完整积分。

\end{proof}

这是有明确定义的上界，不是声称最小阶数等于\(b_n+1\)，也没有给任意图的\(b_n\)求取算法。它与\eqref{eq:primitive:17}的完整相对对象秩上界不同；一个标架生成的微分子模可以远小于整个动机载体。


\endgroup



\begingroup

\subsection{同一个辅助参数的留数判据}\label{sec:forest:1}

设X为k上的光滑n维代数簇，\(k\subset R\)；\(D=\cup D_i\)正常交叉。\(K\subset X(R)\)为紧致有理半代数n链，边界在D。设\(\eta _j\)为有限个有理n形式，沿D至多有对数极点，其余极点均避开K；\(r_j\)为K邻域中的代数单位，在K上严格正。可从X删去这些非边界极点及\(r_j\)的零极点而不改变K。

采用共同变量u，置

\begin{equation}
\ell_u(r)=\frac{r-1}{1+u(r-1)},\qquad
 \Xi=\sum_j\eta_j\,\ell_u(r_j)\wedge du.                         \label{eq:forest:1}
\end{equation}

从\(X\times A^{1}_u\)删去\(1+u(r_j- 1)=0\)的超曲面，得到Y。它们避开\(K\times [0,1]\)，因为分母为\(1- u+ur_j>0\)。设B为\((D\times A^{1})\)及\(u=0,1\)两端与Y的交。

对每个不可约\(D_i\)，在函数域\(k(D_i)\)中记

\begin{equation}
\theta_{ij}=\operatorname{res}_{D_i}\eta_j,\qquad
 \bar r_{ij}=r_j|_{D_i}.
\end{equation}

将j按相同的\(\bar r_{ij}\)分组；相等指函数域中的精确代数相等，不是某个数值点相等。

\begin{theorem}[共同线性提升的充要条件]

\label{thm:forest:1.1} 下列条件等价：

\(1. \Xi \)沿\(D\times A^{1}\)的所有分量没有极点。 2. 对每个\(D_i\)和每个出现的\(q\ne 1\)，有

\begin{equation}
\sum_{j:\bar r_{ij}=q}\theta_{ij}=0.          \label{eq:forest:2}
\end{equation}

\(q=1\)这一组不受条件\eqref{eq:forest:2}约束。满足条件时，

\begin{equation}
M_{\rm ren}=H^{n+1}(Y,B),\qquad
 [M_{\rm ren},[K\times[0,1]],[\Xi]]^{\mathfrak m}                \label{eq:forest:3}
\end{equation}

为实际相对周期，其数值为合并后的绝对收敛积分

\begin{equation}
\int_K\sum_j\eta_j\log r_j.                                    \label{eq:forest:4}
\end{equation}

\end{theorem}

\begin{proof}

\(\Xi \)的\(D_i\)留数是

\begin{equation}
\left(\sum_j\theta_{ij}
       \frac{\bar r_{ij}-1}{1+u(\bar r_{ij}-1)}\right)\wedge du.   \label{eq:forest:5}
\end{equation}

\(q=1\)项为零。其余项是\(1/(u+1/(q- 1))\)，不同q给不同简单极点；比较\(k(D_i)(u)\)中的部分分式留数，\eqref{eq:forest:5}恒零当且仅当\eqref{eq:forest:2}成立。这是有理恒等式判据，不能用近似数值替代。

\(\Xi \)至多有对数边界极点。局部将它写成f·\(dt_1\)\ldots{}\(dt_n du/(t_1\)\ldots{}\(t_s)\)，其余分母为单位。每个边界留数恒零使f分别被\(t_i\)整除；正常交叉坐标于是给\(\Pi t_i\)整除f。故\(\Xi \)在交叉角点也正则，不只在各面的泛点正则。

\(\Xi \)是Y上的最高次正则形式，因而闭，且拉回B的每个分量为零；B为正常交叉闭并，可在相对复形中取其余分量为零。\(K\times [0,1]\)为实际相对链。紧致性、正分母和正则性给绝对收敛；先对u积分等于\(log r_j\)，得到\eqref{eq:forest:4}。相对上同调及比较由实际几何对定义。

\end{proof}

这个充要条件针对给定的共同线性提升\eqref{eq:forest:1}，不是所有可能周期表示的普适判据。例如

\begin{equation}
\frac{dt}{t}(2\log2-\log4)=0,
\end{equation}

但直接提升得到

\begin{equation}
\frac{dt}{t}
 \left(\frac2{1+u}-\frac3{1+3u}\right)du,                        \label{eq:forest:6}
\end{equation}

\(t=0\)的留数不为零，双积分不绝对收敛。先对u积分为零不足以证明原有理提升正则。应先用实际对数函数方程改变表示，或采用能保持边界关系的其他辅助链。


\endgroup
\subsection{共同提升的物理检验}
上述辅助参数判据要求森林各项在同一个几何模型中比较。对两圈dunce's cap，子图与整图的减除条件分别固定；共同Schwinger参数给出森林式，全部坐标面及交角的留数相消后，正链才定义所需相对周期。改变子图减除点时，差为子图的有限常数乘以收缩图的锚定周期。这解释了局部反项为什么可以产生非局部的有限响应。

该例的图多项式、共同核、逐面检查及乘积律证明收于\otherref{feynman_periods_supplement}{补充材料S7}。下面的零动量森林将相同构造推广到所列至多对数发散图类；相容非零参考点另在第\ref{ch:finitemom}节处理。

\section{零动量森林与实际端点下降}\label{ch:zero}\label{ch:descent}

\begingroup

准备映射的代数机制采用 Brown--Kreimer\cite{bk2011}。本节在正质量、欧氏数据和所有子图至多对数的范围内，说明共同质量核、实际相对周期和未正规化零点 BPHZ 之间的比较。

\subsection{输入、方案和结论}\label{sec:zero:1}

固定连通标量图G，边集E带全序，单位传播子幂次，\(D=4\)，\(N=|E|=2h_G>0\)。允许重边；对每个边子集J，定义

\begin{equation}
h_G(J)=b_1(G|J),\qquad d_G(J)=|J|-2h_G(J).
\end{equation}

假设

\begin{equation}
d_G(J)\geq0\quad(J\subseteq E),\qquad d_G(E)=0.       \label{eq:zero:1}
\end{equation}

所以没有幂次UV子发散，所有需减除的子图都是对数型。孤立顶点不计作有边图，空边图为单位。

每边质量平方\(\mu _e\)严格正。外动量欧氏、守恒；以二森林动量平方为系数的无质量第二Symanzik多项式记\(\Phi _G\)，所以系数非负。写

\begin{equation}
M_E=\sum_{e\in E}\mu_e\alpha_e,\qquad F_G=U_GM_E+\Phi_G.
\end{equation}

为给固定纤维的数域动机，先取\(\mu \)及外Gram数据在一个带指定实嵌入的数域k中；变化参数可另取相应代数族。对任意实数据，下面收敛积分的解析结论仍成立，但不据此称其为Q上的数值周期。

采用以下完整减除约定：

\begin{itemize}
\tightlist
\item
  所有子图反项的Taylor算子均为在其全部外动量等于零处取值，质量保持原边质量；
\item
  整图也在零外动量减除；
\item
  每圈采用\(\pi ^{-D/2}\)的欧氏测度，耦合与图对称因子另计；
\item
  收缩携带原外动量之和，质量只随保留下来的原边继承。
\end{itemize}

\begin{theorem}

\label{thm:zero:1} 对满足\eqref{eq:zero:1}的图，以上零外动量BPHZ振幅有一个显式、绝对收敛的有理相对周期表示

\begin{equation}
R_G(\kappa,0)=\int_{\widetilde\sigma_G\times[0,1]}\Xi_G,
 \qquad M_G^R=H^N(Y_G,B_G),                                  \label{eq:zero:2}
\end{equation}

其中\(Y_G\)、\(B_G\)及\(\Xi _G\)由有限森林准备映射和图多项式构造。\(\Xi _G\)在完整坐标边界及全部角点正则，正链避开其余极点。结论允许图含嵌套或重叠的对数子图；不要求先求出完整积分。

这里证明的是指定方案的实际标架存在及物理比较。它不决定\(M_G^R\)的全部纯因子、扩张或最小秩，也不完成任意非零参考点、幂次发散、IR或规范理论。

\end{theorem}


\subsection{一个所有森林共用的质量核}\label{sec:zero:4}

对\eqref{eq:zero:6}中的项定义\(U_A\)为A各块U的乘积，\(U_1=1\)；B的\(\Phi _B\)按继承外动量计算。由于每个左块在零外动量，

\begin{equation}
\sum_{\text{左块 }a}F_a^0/U_a+F_B^0/U_B=M_E.                 \label{eq:zero:8}
\end{equation}

这是真正的同一个线性多项式，与森林项无关。

将外运动学平方不变量乘\(\tau \)，\(0\le \tau \le 1\)，质量不变。于是

\begin{equation}
q_B=\Phi_B/U_B,\qquad
 D_{A|B}(\tau)=U_BM_E+\tau\Phi_B,
\end{equation}

\begin{equation}
\Psi_\tau(A\otimes B)
  =\frac{\Omega_E}{U_A^2U_B^2}\,
              \frac{\Phi_B}{U_BM_E+\tau\Phi_B}
  =\frac{\Omega_E}{U_A^2U_B^2}\,
              \frac{q_B}{M_E+\tau q_B},                       \label{eq:zero:9}
\end{equation}

并规定\(\Psi \tau (A\otimes 1)=0\)。各项的总齐次次数为零，故\eqref{eq:zero:9}为射影\((N- 1)\)形式。最终取

\begin{equation}
\Xi_G=-\Psi_\tau(\mathcal R(G))\wedge d\tau.                  \label{eq:zero:10}
\end{equation}

在内部形式上积分\(\tau \)给

\begin{equation}
\sum_{A,B}c_{A,B}\frac{\Omega_E}{U_A^2U_B^2}
              \log\frac{M_E}{M_E+\Phi_B/U_B}.                 \label{eq:zero:11}
\end{equation}

\eqref{eq:zero:10}是先合并森林再积分的定义；不能把\eqref{eq:zero:11}每个尚不收敛的项单独当成周期。

注意\(\Psi \tau (1\otimes G)\)一般不是\(\Omega _E/U_G^{2}\)。原文抽象设置的这一附加等式在此不适用；后面的相消只使用右单位消失、UV单位为零及实际留数协变性，不假装满足原文的全部输入条件。


\endgroup
\subsection{共同核的边界控制}
零动量森林定理使用三项事实：图类对对数子图和收缩封闭；森林准备映射消去每个对数面的法向首项；Hepp扇区估计同时控制相交UV面与整体尺度。共同核的正性保证正链避开极点。由这些事实得到实际相对对象，并与同一零外动量BPHZ规则比较。完整逐面证明及径向积分见附录\ref{app:zero-proof}。

以下端点下降只调用这一实际对象和共同核，不重新扩大图类。关键问题是：对外部参数求导后，能否在同一个正常交叉模型上构造有理原函数，并以实际连接态射取其两端点？



\begingroup
\providecommand{\Q}{\mathbb Q}
\providecommand{\dR}{\mathrm{dR}}
\providecommand{\Res}{\operatorname{Res}}
\providecommand{\im}{\operatorname{im}}
\providecommand{\rk}{\operatorname{rk}}
\providecommand{\cR}{\mathcal R}
\providecommand{\cH}{\mathcal H}
\subsection{输入、方案和已建立的共同核}
\subsubsection{图与物理数据}
设 $G$ 连通，边集 $E$ 有序，$N=|E|=2h_G>0$，并置 $n=N-1$。
对任意边子集 $J$ 要求
\begin{equation}\label{descent:eq:power}
 d_G(J)=|J|-2b_1(G|J)\geq0.
\end{equation}
这排除了幂次 UV 子发散，允许嵌套和重叠的对数子图。
所有边质量平方 $\mu_e>0$；外数据为守恒的欧氏动量。
每圈采用 $\pi^{-D/2}$ 的欧氏测度，$D=4$，图权和耦合另计。
所有子图以及整图均在其全部外动量为零处减除，保留原边质量。
固定代数参数属于带指定实嵌入的数域 $k$。

设 $U$ 为第一 Symanzik 多项式，$\Phi$ 为不含质量项的第二多项式，
\begin{equation}\label{descent:eq:mass}
 M=\sum_{e\in E}\mu_e\alpha_e,\qquad F_G=U_GM+\Phi_G.
\end{equation}
在所选欧氏域 $\Phi$ 非负。
允许先把其动量不变量系数作为仿射参数，再限制到实际运动学子空间；这不声称任意系数都可由四维外动量实现。

\subsubsection{森林准备映射}
在保留边标签的对数图 Hopf 代数中，记反极为 $S_H$，
$\pi_+=\mathrm{id}-u\varepsilon$。定义
\begin{equation}\label{descent:eq:prep}
 \cR=(S_H\otimes\pi_+)\Delta_H,
 \qquad \cR(G)=\sum_{A,B}c_{A,B}A\otimes B,
 \qquad B\ne1.
\end{equation}
$A$ 是反项块的乘积，$B$ 是带外运动学的根商图。
各项的边集仍分割原 $E$。
准备恒等式为
\begin{equation}\label{descent:eq:prepidentity}
 \mu_{13}\bigl(\Delta_H\otimes((\mathrm{id}\otimes\pi_+)\Delta_H)\bigr)\cR
 =1\otimes\cR.
\end{equation}
这是经典森林相消机制\cite{bk2011}；其在本方案中的具体核和收敛证明已在项目零点森林稿给出。

写
\begin{equation}\label{descent:eq:wq}
 w_{A,B}=\frac{\Omega_E}{U_A^2U_B^2},\qquad
 q_B=\frac{\Phi_B}{U_B},\qquad H_B(\tau)=M+\tau q_B.
\end{equation}
因所有左块在零外动量，森林各项共享同一个 $M$。
既有实际相对形式为
\begin{equation}\label{descent:eq:Xi}
 \Xi_G=-\sum_{A,B}c_{A,B}w_{A,B}
                     \frac{q_B}{H_B(\tau)}\wedge d\tau,
 \qquad 0\leq\tau\leq1,
\end{equation}
其正链周期为零点 BPHZ 振幅 $R_G$。
在正内部积分 $\tau$ 可写成
\begin{equation}\label{descent:eq:log}
 R_G=\int_{\widetilde\sigma_G}
       \sum_{A,B}c_{A,B}w_{A,B}\log\frac{M}{M+q_B}.
\end{equation}
式\eqref{descent:eq:log}始终指\emph{先合并森林再积分}；单项可以发散。
本文不另行对每个发散对数赋予形式周期。

\subsubsection{既有边界结论和本节新增责任}
令 $X_G$ 为射影空间沿全部坐标旗标的光滑 toric 模型，$D_X$ 为坐标边界。
森林各项沿 $D_X$ 至多有对数极点。
在真边子集 $I$ 的法向面上，其留数分解为 $\omega_{\rm UV}(I)$ 与相应收缩核的外积；右块全部收缩时，$q_B$ 多出一阶法向参数。
式\eqref{descent:eq:prepidentity}使合并后的每个留数分别为零，正常交叉整除使全部角点正则。
正质量及生成树正性保证正链避开非边界极点。

这些既有结论证明\eqref{descent:eq:Xi}的存在，却尚未说明其参数导数来自哪个实际对象。
下面给出这个端点对象及实际态射；尤其证明所用有理原函数也满足森林相消。


\subsection{在两端点之间构造有理原函数}
令 $v$ 为质量和仿射运动学参数的一个常系数方向。
图多项式 $U$ 与这些参数无关，故 $v(w_{A,B})=0$。
定义 $n$ 次形式
\begin{align}\label{descent:eq:Hprimitive}
 \mathcal A_v(\tau)
 &=\sum_{A,B}c_{A,B}w_{A,B}
       \left(\frac{vM}{M}-\frac{vH_B(\tau)}{H_B(\tau)}\right)\nonumber\\
 &=\sum_{A,B}c_{A,B}w_{A,B}
       \frac{\tau(q_BvM-Mvq_B)}{M H_B(\tau)}.
\end{align}
其中未合并的括号在 $B=1$ 时规定 $q_1=0$。
由此该括号在右单位上恒零，并且
\begin{equation}\label{descent:eq:endpointzero}
 \mathcal A_v(0)=0,\qquad
 v\Xi_G=(\partial_\tau\mathcal A_v)\wedge d\tau
             =(-1)^n d_Y\mathcal A_v.
\end{equation}
最后一个等号是相对参数纤维上的外微分：$\mathcal A_v$ 已是 $X_G$ 方向最高次形式。

\begin{remark}
若只取 $-\sum c_{A,B}w_{A,B}vH_B/H_B$ 作为原函数，它在右单位上一般不为零。
此时不能直接调用森林准备恒等式消除反项新增面的极点。
式\eqref{descent:eq:Hprimitive}减去初端点值，是几何证明所必需的规范化，不是略去一个可任意丢弃的常数。
\end{remark}

\begin{lemma}[原函数的全部边界相消]\label{descent:lem:cancel}
在删去 $M$、各 $U$ 块及各 $H_B(\tau)$ 的非边界零集后，$\mathcal A_v(\tau)$ 沿全部坐标面和正常交叉角点正则。
结论包括 $\tau=0,1$，并且对任意有限阶参数导数仍成立。
\end{lemma}
\begin{proof}
单项的 $w_{A,B}$ 至多有简单对数极点。
在一个坐标面提取共同法向单项式后，$M$ 与 $H_B$ 的余因子在正邻域非零，故式\eqref{descent:eq:Hprimitive}的附加因子不产生新的法向负阶。

若 $I_B\ne B$ 且该面有对数留数，则树及二森林因子化给
\[
 q_B\big|_I=q_{B/I_B},\qquad M\big|_I=M_{E\setminus I}.
\]
参数方向 $v$ 在保留下来的质量及运动学上作相应限制。
属于 $I$ 的质量项及其参数导数均含法向因子，在该面消失。
于是\eqref{descent:eq:Hprimitive}的括号限制恰好是收缩后同一个括号。
$\omega_{\rm UV}$ 不依赖所求导的参数，因此原函数具有与原核相同的留数协变性。

若 $I_B=B$，则 $q_B$ 和 $vq_B$ 均具有额外法向因子；其括号限制为零。
这也由 $q_1=0$ 的右单位约定表达。
因此将原函数核代入\eqref{descent:eq:prepidentity}，全部带边标签留数分别为零。

在旗标坐标 $t_1,\ldots,t_n$ 中，总和可写为
$f(t,\tau)\,dt_1\cdots dt_n/(t_1\cdots t_n)$，其中 $f$ 正则。
每个留数为零意味着 $t_j\mid f$。
光滑正常交叉坐标环中各 $t_j$ 是不同素因子，所以 $\prod_jt_j\mid f$。
因此角点也正则。

上述恒等式在参数代数环上成立，并未先固定某个数值。
继续对参数求导仍保持初端点零值、右单位零值和留数协变性；至多提高非边界分母的幂。
同一整除证明逐阶适用。
\end{proof}

定义端点形式
\begin{equation}\label{descent:eq:Theta}
 \Theta_v=\mathcal A_v(1)
   =\sum_{A,B}c_{A,B}w_{A,B}
       \left(\frac{vM}{M}-\frac{v(M+q_B)}{M+q_B}\right).
\end{equation}
它在正链上的每个角点正则；其定义不含辅助参数 $\tau$。


\subsection{实际相对对及连接态射}
\subsubsection{共同端点几何}
从 $X_G$ 删去所有森林项的 $U_A,U_B,M,M+q_B$ 的非边界零集，得到 $V_G$；有理函数的零集按清分母后的严格变换理解。
记 $D_V=D_X\cap V_G$。
同一正质量下界保证 $\widetilde\sigma_G$ 及其边界仍在这对中。
置
\begin{align}\label{descent:eq:pairs}
 Y&=\{(x,\tau)\in V_G\times\mathbb A^1:
                         H_B(\tau)\ne0\text{ 对所有森林项}\},\nonumber\\
 D_Y&=Y\cap(D_V\times\mathbb A^1),\qquad
 E_i=V_G\times\{i\}\quad(i=0,1),\nonumber\\
 B&=D_Y\cup E_0\cup E_1,\qquad
 H_G=H^n(V_G,D_V),\qquad M_G=H^{n+1}(Y,B).
\end{align}
边界处使用提取法向单项式后的非边界方程。
由于已删去两个端点的零集，$E_0,E_1$ 确为整个 $V_G$ 的两份。
这是原零点森林模型的共同开细化；实际限制态射与同一正链保持其形式周期，不声称较大的载体与原载体同构。

三重对 $D_Y\subset B\subset Y$ 的连接态射及端点分解给
\begin{equation}\label{descent:eq:connecting}
 H^n(B,D_Y)\simeq H_G\oplus H_G
       \longrightarrow M_G.
\end{equation}
取来自 $E_1$ 的分量，并按乘积链取向归一化，记为 $\iota:H_G\to M_G$。
它是实际相对上同调态射，也是在共同平滑参数开集上的水平态射。
这里没有留数管的 Tate 扭转或额外 $2\pi i$ 因子。

\subsubsection{符号和相对 Stokes}
为固定\eqref{descent:eq:connecting}的符号，采用相对复形
\[
 C^j(Y,B)=\Omega_Y^j\oplus\Omega_B^{j-1},\qquad
 d(\alpha,\beta)=(d\alpha,i^*\alpha-d\beta),
\]
对链的配对为 $\int_\Gamma\alpha-\int_{\partial\Gamma}\beta$。
$B$ 的并分量可用相容的相对或 simplicial 模型解释。
对 $\Gamma=\widetilde\sigma_G\times[0,1]$，终端面在 $\partial\Gamma$ 中的符号为 $(-1)^n$。
于是 $\iota[\theta]$ 在局部代表中为
\begin{equation}\label{descent:eq:cone}
 \bigl[(0,(-1)^{n+1}\theta\text{ 位于 }E_1)\bigr],
 \qquad
 \langle\iota[\theta],\Gamma\rangle
     =\langle[\theta],\widetilde\sigma_G\rangle.
\end{equation}
端点形式在 $D_V$ 上的拉回为零；必要的相对 cochain 按同一模型补齐。

取 $\eta_v=(-1)^n\mathcal A_v$。
式\eqref{descent:eq:endpointzero}给 $d\eta_v=v\Xi_G$。
其在 $E_0$ 上为零，在 $E_1$ 上为 $(-1)^n\Theta_v$，在 $D_Y$ 上亦为零，因为它没有 $d\tau$ 且是 $X_G$ 方向的 $n$ 次形式。
减去相对 coboundary $d(\eta_v,0)$，得到
\begin{equation}\label{descent:eq:GM}
 \nabla_v[\Xi_G]=\iota[\Theta_v].
\end{equation}
这一步是实际相对复形中的等式，不是从数值导数相等倒推动机态射。


\subsection{任意正阶导数的降维定理}
在质量严格正、欧氏参数的一个紧邻域中，生成树余因子有统一正下界，$M$ 以及 $H_B(\tau)$ 的余因子也有统一正下界。
引理\ref{descent:lem:cancel}中的整除在参数族中成立。
故合并后的形式及其任意给定有限阶参数导数在有限个闭旗标立方体上都有统一有界密度。
这允许在合并森林积分内逐阶求导。
它不允许分别积分尚未相消的发散项。

\begin{theorem}[实际端点下降]\label{descent:thm:descent}
在\eqref{descent:eq:power}的图类及本节规定的零点方案中，任意正阶参数导数均为端点对象 $H_G$ 的实际周期。
在共同光滑且上同调常秩的参数开集上，式\eqref{descent:eq:connecting}给出实际水平态射 $\iota$，满足
\begin{equation}\label{descent:eq:descent}
 \nabla_v[\Xi_G]=\iota[\Theta_v],\qquad
 vR_G=\int_{\widetilde\sigma_G}\Theta_v.
\end{equation}
所有更高正阶协变导数均在 $\im\iota$ 中。
在代数参数点，以该实际微分标架定义的形式周期满足
\begin{equation}\label{descent:eq:motivic}
 (vR_G)^{\mathfrak m}
     =[H_G,[\Theta_v],[\widetilde\sigma_G]]^{\mathfrak m}.
\end{equation}
因而其最小形式周期对象是 $H_G$ 的子商。
\end{theorem}
\begin{proof}
原函数的边界正则性由引理\ref{descent:lem:cancel}保证；这使\eqref{descent:eq:GM}成为实际相对类等式。
上述统一有界性给积分导数，式\eqref{descent:eq:cone}给正链及符号。
两者的形式周期比较来自同一个实际态射，故不需要一般周期映射的单射性。
继续求导时，$\iota$ 的水平性保证所得类仍在其像中。
实际像是子对象，最小矩阵系数对象的构造再给最后结论。
\end{proof}

这里 ``光滑参数开集'' 包括对代数对的必要分层：若删去的零集碰并导致上同调秩跳跃，不能把邻近纤维的秩界直接搬到该点。
显式有理形式及其收敛积分在适当欧氏边界点仍可有意义；连接与秩的陈述只用于其已定义的常秩开集。

\subsubsection{可直接写出的全部导数核}
若 $v_1,\ldots,v_r$ 是仿射参数空间上的常系数方向，$r\geq1$，则 $M$ 和 $H_B(1)$ 均为仿射函数，故
\begin{align}\label{descent:eq:allorders}
 \Theta_{v_1\cdots v_r}
 ={}&(-1)^{r-1}(r-1)!\sum_{A,B}c_{A,B}w_{A,B}\nonumber\\[-2pt]
 &\quad\cdot\left[
 \frac{\prod_{j=1}^r v_jM}{M^r}
 -\frac{\prod_{j=1}^r v_j(M+q_B)}{(M+q_B)^r}
 \right].
\end{align}
重复方向也适用。对非线性物理参数曲线按链式法则取这些形式的线性组合；不能遗漏向量场系数的导数。

例如质量多重指标 $a=(a_e)$、$|a|=r\geq1$ 给
\begin{equation}\label{descent:eq:massderivative}
 \partial_\mu^a R_G
 =(-1)^{r-1}(r-1)!\int_{\widetilde\sigma_G}
   \sum_{A,B}c_{A,B}w_{A,B}\alpha^a
       \left(M^{-r}-(M+q_B)^{-r}\right).
\end{equation}
一个纯外运动学方向 $v$ 则给
\begin{equation}\label{descent:eq:kinematic}
 vR_G=-\int_{\widetilde\sigma_G}
             \sum_{A,B}c_{A,B}w_{A,B}\frac{vq_B}{M+q_B}.
\end{equation}
所有表达式仍先作有限森林总和。


\subsection{无需先积分的微分系统上界}
\begin{corollary}[曲线上的阶数界]\label{descent:cor:ode}
将参数限制到一条定义在 $k$ 上的光滑代数曲线，局部函数域记 $k(s)$。
设端点相对局部系统的秩为 $b=\dim(H_G)_B$。
振幅所有正阶 $s$ 导数所生成的微分模秩 $d$ 满足
\begin{equation}\label{descent:eq:rankbound}
 d\leq\rk(\im\iota)\leq b.
\end{equation}
原振幅具有一个 $k(s)$ 系数齐次线性微分方程，阶数至多 $d+1\leq b+1$；其一阶导数具有阶数至多 $d$ 的齐次方程。
\end{corollary}
\begin{proof}
令 $\omega=[\Xi_G]$，$\mathcal D_+$ 为 $\nabla_s^j\omega$、$j\geq1$ 生成的 $k(s)$ 向量空间。
它对连接稳定，由定理\ref{descent:thm:descent}包含于 $\im\iota$，故\eqref{descent:eq:rankbound}成立。
加入 $\omega$ 最多增加一维，所得循环微分模秩至多 $d+1$。
其连续导数必在该维数内线性相关；与水平 Betti 链配对便得到方程。
若 $d=0$，导数恒零，原振幅的常数方程按同一约定处理。
\end{proof}

这是端点几何控制的上界，不是一个已经给出所有图 $b$ 的算法，更不是最小阶数恰为 $b+1$ 的断言。
多参数导数空间也包含于同一端点像；其实际秩仍需几何求取。

\begin{proposition}[算术范围]\label{descent:prop:arithmetic}
设 $\mathcal C$ 是所用实际动机范畴中对子商封闭的满张量子范畴。
若 $H_G\in\mathcal C$，则每个代数点的正阶参数导数形式周期的最小对象属于 $\mathcal C$。
特别地，端点对象 mixed Tate 是这些导数具有 mixed Tate 载体的充分条件。
\end{proposition}
\begin{proof}
直接用式\eqref{descent:eq:motivic}及子商封闭性。该证明不依赖导数的具体数值。
\end{proof}

反向不成立：特定读出或链可能看不见端点对象的大部分结构。
本命题也没有从导数载体的类型自动推出原振幅的完整动机类型。
微分模中新增的一个水平商方向可以承载非平凡的\emph{常数算术}；一维微分模并不意味着一维固定算术对象。
本原尺度振幅 $-p_G\log\lambda$ 的参数方程阶数只有二，而 $p_G$ 的算术 coaction 必须完整保留，正是此前证明的例子。


\subsection{质量插入怎样穿过森林减除}
\subsubsection{交换的是继承的减除算子}
记 $\mathbf R_G$ 为原图规定的零点森林算子，包含整图零点减除，并按嵌套顺序施加各 $t_\gamma$。
$t_\gamma$ 将子图的全部外动量取零，保留所有边质量。
因此在被积函数层面
\begin{equation}\label{descent:eq:commute}
 [\partial_{\mu_e},t_\gamma]=0,\qquad
 \partial_\mu^a\mathbf R_G f_G
                =\mathbf R_G\partial_\mu^a f_G.
\end{equation}
这个代数恒等式在相容的收敛正规化区域同样成立；最终有限值的交换由前述合并森林的统一估计保证。

单位传播子给
\begin{equation}\label{descent:eq:insert}
 \partial_\mu^a\prod_e\frac1{k_e^2+\mu_e}
 =(-1)^{|a|}\left(\prod_e a_e!\right)
              \prod_e(k_e^2+\mu_e)^{-1-a_e}.
\end{equation}
故质量导数确实实现提高传播子幂次的插入，但右端必须带\emph{从原图继承}的森林算子。
若插入改善发散度后重新仅对仍发散子图作减除，一般会得到另一个有限振幅。

\begin{proposition}[不能自动改用插入图的最少减除]\label{descent:prop:oversub}
即使原图本原、质量插入图已经 UV 收敛，原重整化振幅的质量导数仍可能含一个非零的零点减除值。
因此式\eqref{descent:eq:commute}不能无条件写成 ``质量导数等于插入图的最少重整化值''。
\end{proposition}
\begin{proof}
取等质量一圈 bubble，$z=Q^2\geq0$，$\mu=m^2>0$。
同一归一化下
\[
 R(z,\mu)=\int_0^1\log\frac{\mu}{\mu+zx(1-x)}\,dx,
\]
所以
\begin{equation}\label{descent:eq:bubble}
 \partial_\mu R
   =\frac1\mu-\int_0^1\frac{dx}{\mu+zx(1-x)}.
\end{equation}
未作额外零点减除的共同质量插入之和为
$-\int_0^1dx/(\mu+zx(1-x))$，已经 UV 收敛。
原振幅的导数另外减去该插入在 $z=0$ 的值 $-1/\mu$。
两者相差明确非零的 $1/\mu$，无须求完整积分。
\end{proof}

对子图反项作质量求导同样不能遗漏。
在一个具有多块乘积的森林项中，共同质量或一般方向导数按 Leibniz 法则分配到每块；高阶导数带相应多项式系数。
若质量由原边独立标记，每个 $\partial_{\mu_e}$ 只作用于含该边的块。
反项块虽对外动量局部，却通常依赖质量。

\subsubsection{外动量导数不同于质量导数}
对一变量零点取值 $t_0f=f(0)$，有
\begin{equation}\label{descent:eq:kincommutator}
 [\partial_z,t_0]f=-f'(0).
\end{equation}
所以外参数导数不能逐项穿过减除而忽略局部项。
式\eqref{descent:eq:kinematic}已经将所有原反项及总体减除处理在合并核内。
它仅让根商图的 $q_B$ 发生变化；没有假定对子图内部外动量作任意变化都可与原减除交换。


\endgroup

\section{相容非零MOM与二次减除sunset}\label{ch:finitemom}\label{ch:sunsetjet}

\begingroup

本节补齐零点森林之后的相容非零参考比较，并为第\ref{ch:unified}节的新组成方案保持提供实际输入。全部参考与锚定条件保留。

\providecommand{\Q}{\mathbb Q}
\providecommand{\cP}{\mathcal P}
\providecommand{\cC}{\mathcal C}
\providecommand{\ev}{\operatorname{ev}}
\providecommand{\per}{\operatorname{per}}
\providecommand{\eps}{\varepsilon_H}

\subsection{输入与相容参考条件}
\subsubsection{图类和零点输入}
沿用项目已有零点森林定理的范围：$D=4$，单位标量传播子，
\begin{equation}\label{finitemom:eq:class}
 |E(G)|=2h_G>0,\qquad |J|-2b_1(G|J)\geq0\quad(J\subseteq E(G)).
\end{equation}
所有质量平方严格正，外数据为守恒欧氏动量。
图的边、外腿及收缩后的外结构均保留标记。
抽取非空对数子图并收缩后，仍属于这一图类。
记相应连通分次 Hopf 代数为 $H$，余积写为
\begin{equation}\label{finitemom:eq:coproduct}
 \Delta G=G\otimes1+1\otimes G+
       \sum_{\varnothing\ne\gamma\subsetneq G}\gamma\otimes G/\gamma.
\end{equation}
求和含嵌入重数，$\gamma$ 可为许可的不交连通子图乘积。
次数取圈数；约化余积两侧的每个非单位因子次数均严格降低。

已有零点规则记为 $\Phi_0(G;\kappa)$，对应所有子图和整体均在零外动量减除。
在代数的正质量欧氏数据上，已有实际形式周期
\begin{equation}\label{finitemom:eq:zero}
 \Phi_0^{\mathfrak m}(G;\kappa)\in\cP_k^{\mathfrak m},
 \qquad \Phi_0^{\mathfrak m}(G;0)=0,
\end{equation}
并且其数值周期与零点 BPHZ 振幅相同。
不连通图取实际外积，单位图的值为 $1$。
这里 $k$ 是包含本节全部质量和参考 Gram 数据、带指定实嵌入的数域。
任意实数据可先得到解析结论，但不据此称其为 $k$ 上的周期。

\subsubsection{参考点分配的精确定义}
给每个需要出现的连通带标记图 $X$ 指定一个欧氏参考点 $\rho_X$，质量保持该图原边质量。
同一带标记图在不同余积分解中出现时使用同一分配；不连通图按分量乘法延拓。
外数据在收缩下的继承映射记为
\begin{equation}\label{finitemom:eq:projection}
 \pi_\gamma:\mathcal K_G\longrightarrow\mathcal K_{G/\gamma}.
\end{equation}
它通过合并顶点上的外动量之和定义，连续收缩必须满足相应复合关系。
继承的欧氏数据仍然欧氏。

\emph{不要求} $\rho_{G/\gamma}=\pi_\gamma(\rho_G)$。
左边是商图自身选定的减除点，右边是将整图参考点限制到该商图的评价点，两者在递归中有不同作用。
把它们混同会错误排除独立子图参考条件。

本文的方案是逐图的零次局部 MOM 条件。
若要来自某个固定拉格朗日量的一组耦合重整化条件，还需参考分配尊重外腿对称性、顶点类型和图求和；这些条件并不由任意逐图分配自动满足。
这里也不包含幂次发散所需的多阶 Taylor 反项。

\subsubsection{局部 character 的作用}
局部 character $f$ 对不交并乘法，$f(1)=1$，其系数可依赖质量和参考数据，却不依赖当前外动量 $\kappa$。
它对振幅规则的作用定义为
\begin{equation}\label{finitemom:eq:action}
 (f*\Phi)(G;\kappa)=\Phi(G;\kappa)+f(G)
   +\sum_{\varnothing\ne\gamma\subsetneq G}
        f(\gamma)\Phi(G/\gamma;\pi_\gamma\kappa).
\end{equation}
所有子图系数的质量变量按原边标签拉回。
不同图的运动学空间不必相同；式\eqref{finitemom:eq:action}明确这些拉回，避免把它们当成未经说明的同一变量空间。
局部 character 之间的卷积由\eqref{finitemom:eq:coproduct}定义，余结合性及\eqref{finitemom:eq:projection}的复合性给
\begin{equation}\label{finitemom:eq:associative}
 (f*g)*\Phi=f*(g*\Phi),\qquad f^{-1}=f\circ S_H.
\end{equation}
有限圈数上这些公式均为有限多项式，不涉及无限级数收敛。


\subsection{有限局部反项的递归构造}
\begin{theorem}[非零参考点的实际周期提升]\label{finitemom:thm:construction}
对上述参考分配 $\rho$，存在唯一局部 character $f_\rho^{\mathfrak m}$，使
\begin{equation}\label{finitemom:eq:lift}
 \Phi_\rho^{\mathfrak m}=f_\rho^{\mathfrak m}*\Phi_0^{\mathfrak m},
 \qquad \Phi_\rho^{\mathfrak m}(G;\rho_G)=0
\end{equation}
对每个非单位连通图成立。其递归为
\begin{align}\label{finitemom:eq:recursion}
 f_\rho^{\mathfrak m}(G)={}&-\Phi_0^{\mathfrak m}(G;\rho_G)\nonumber\\
 &-\sum_{\varnothing\ne\gamma\subsetneq G}
 f_\rho^{\mathfrak m}(\gamma)
 \Phi_0^{\mathfrak m}(G/\gamma;\pi_\gamma\rho_G).
\end{align}
每个 $f_\rho^{\mathfrak m}(G)$ 和 $\Phi_\rho^{\mathfrak m}(G;\kappa)$ 都有由有限直和、张量积及零点相对对象构成的实际 Nori 载体。
该构造允许任意有限层嵌套与重叠，不需先求出周期数值。
\end{theorem}
\begin{proof}
给单位赋值 $1$，再按圈数递增使用\eqref{finitemom:eq:recursion}，对不交并乘法延拓。
右侧的 $f_\rho(\gamma)$ 只涉及更低圈数，所以递归终止且唯一。
代入\eqref{finitemom:eq:action}，在 $\kappa=\rho_G$ 处三个部分恰好相消；这是形式周期环中的精确等式。
反过来，条件\eqref{finitemom:eq:lift}逐阶强制\eqref{finitemom:eq:recursion}，所以唯一性也成立。

每个零点值由既有实际几何对给出。
乘积由张量积载体、和由直和载体、负号由有理标架给出。
逐次应用这些实际操作即可构造本定理的载体。
全部参考及继承评价点在有限的子图--收缩闭包内出现，固定最大圈数后只需要有限个几何输入。
此处的载体未必最小，最小对象还需按实际标架取子商。
\end{proof}

\begin{remark}
这里从既定零点实际周期构造了一个相容的非零点形式周期提升。
不能把另一个独立提出的非零点几何模型仅凭数值相等就宣布与它具有相同形式周期。
不同几何模型之间仍需实际态射、变量变换或相对 Stokes 比较。
本定理证明实际载体的存在及下节的物理读出，不声称已将所有方案合成一个最小的共同有理核。
\end{remark}


\subsection{与规定的物理 BPHZ 规则比较}
\subsubsection{一个可乘且相容的截断}
采用逐边 Schwinger 截断 $\alpha_e\geq a>0$。对连通图定义
\begin{equation}\label{finitemom:eq:cutoff}
 \phi_a(G;\kappa)=\int_{[a,\infty)^{E(G)}}
    U_G^{-2}\exp\left(-M_G-\frac{\Phi_G(\kappa)}{U_G}\right)
                                  \prod_e d\alpha_e.
\end{equation}
这采用前述欧氏测度，不含耦合与对称因子。
固定 $a>0$ 时，生成树正系数及 $\alpha_e\geq a$ 给 $U_G\geq a^{h_G}$；质量指数控制无穷远。
故每个裸积分绝对收敛。
边标签使子图与商图乘积使用互不重叠的同一组截断变量，因此规则对不交并乘法。
这是理论上的辅助截断，不需要执行任何截断积分。

在截断层定义局部反项
\begin{equation}\label{finitemom:eq:counterterm}
 C_{\rho,a}(G)=-\phi_a(G;\rho_G)
  -\sum_{\varnothing\ne\gamma\subsetneq G}
       C_{\rho,a}(\gamma)\phi_a(G/\gamma;\pi_\gamma\rho_G).
\end{equation}
这正是先减子发散、再对准备后图作 $\rho_G$ 点零次取值的递归。
记
\begin{equation}\label{finitemom:eq:regulated}
 \Phi_{\rho,a}=C_{\rho,a}*\phi_a,
 \qquad f_{\rho,a}=C_{\rho,a}*C_{0,a}^{-1}.
\end{equation}
由式\eqref{finitemom:eq:associative}，
\begin{equation}\label{finitemom:eq:regulatedaction}
 \Phi_{\rho,a}=f_{\rho,a}*\Phi_{0,a}.
\end{equation}
各 $C_{\rho,a}$ 可随 $a\downarrow0$ 发散；我们并不对它们分别赋予有限周期。

\subsubsection{只需零点的绝对收敛作为解析输入}
零点反项为零点评价规则与反极的复合。
展开其准备映射 $\mathcal R(G)=\sum c_{A,B}A\otimes B$，将各块放回原边变量，得到
\begin{equation}\label{finitemom:eq:zerodensity}
 \Phi_{0,a}(G;\kappa)=
 \int_{[a,\infty)^{E(G)}}
 \sum_{A,B}c_{A,B}
 \frac{e^{-M_E-q_B(\kappa)}-e^{-M_E}}{U_A^2U_B^2}
                                  \prod_e d\alpha_e.
\end{equation}
本文零点森林证明已给出括号中\emph{合并密度}在全正正交体上的绝对可积性，含所有 UV 角点及整体尺度。
因此以该绝对可积密度为支配函数，直接有
\begin{equation}\label{finitemom:eq:zerolimit}
 \lim_{a\downarrow0}\Phi_{0,a}(G;\kappa)
              =\Phi_0(G;\kappa)=\per\Phi_0^{\mathfrak m}(G;\kappa).
\end{equation}
这里没有未经证明地换用维数正规化的极限，也没有逐项取发散森林积分的极限。

\begin{theorem}[物理比较及有限性]\label{finitemom:thm:physical}
对每个图及规定的代数欧氏参考分配，$f_{\rho,a}(G)$ 有有限极限，且
\begin{equation}\label{finitemom:eq:physical}
 \lim_{a\downarrow0} f_{\rho,a}(G)=\per f_\rho^{\mathfrak m}(G),
 \qquad
 \lim_{a\downarrow0}\Phi_{\rho,a}(G;\kappa)
                         =\per\Phi_\rho^{\mathfrak m}(G;\kappa).
\end{equation}
所以第\ref{finitemom:thm:construction}定理构造的实际周期给出了该非零点对数 BPHZ 方案的振幅。
\end{theorem}
\begin{proof}
对任意 $a>0$，$\Phi_{\rho,a}(G;\rho_G)=0$。
将\eqref{finitemom:eq:regulatedaction}在该点展开，得到与\eqref{finitemom:eq:recursion}完全相同的递归，只把 $\Phi_0$ 换成 $\Phi_{0,a}$。
由\eqref{finitemom:eq:zerolimit}及对圈数的归纳，右侧每个低阶系数和零点评价都有有限极限。
有限和与乘积可逐项取这些\emph{已经有限的}极限，得到第一个等式。
随后在任意当前点展开\eqref{finitemom:eq:regulatedaction}，便得第二个等式。
局部性在每一步都明确：$f_\rho(G)$ 无当前外动量依赖，反项只有零次局部插入。
\end{proof}

这给出规定截断和递归下的全图类物理比较。
若要与另一种正规化的有限部分或 MS 约定比较，仍须那一方案自身的归一化关系；本定理不把各正规化的裸有限部分自动视为同一个周期。


\subsection{实际方案之间的复合律}
对方案 $a,b,c$，以下下标表示参考分配，与上一节已取极限的截断变量不再混用。
定义
\begin{equation}\label{finitemom:eq:transition}
 g_{ba}^{\mathfrak m}=f_b^{\mathfrak m}*(f_a^{\mathfrak m})^{-1}.
\end{equation}
连通分次性使逆 character 的每个分量都是有限多项式。
从而
\begin{equation}\label{finitemom:eq:cocycle}
 \Phi_b^{\mathfrak m}=g_{ba}^{\mathfrak m}*\Phi_a^{\mathfrak m},
 \quad g_{cb}^{\mathfrak m}*g_{ba}^{\mathfrak m}=g_{ca}^{\mathfrak m},
 \quad g_{aa}^{\mathfrak m}=\eps.
\end{equation}
这是图 character 的方案复合，不是对固定费曼动机随意规定一个自同态。

转换系数也可直接从方案 $a$ 递归恢复：
\begin{align}\label{finitemom:eq:directtransition}
 g_{ba}^{\mathfrak m}(G)={}&-\Phi_a^{\mathfrak m}(G;\rho_G^b)\nonumber\\
 &-\sum_{\varnothing\ne\gamma\subsetneq G}
 g_{ba}^{\mathfrak m}(\gamma)
 \Phi_a^{\mathfrak m}(G/\gamma;\pi_\gamma\rho_G^b).
\end{align}
因为 $\Phi_b(G;\rho_G^b)=0$，此式由\eqref{finitemom:eq:action}直接得到。
注意 $\Phi_a(G;\rho_G^b)$ 一般不为零。

实际动机 coaction 是形式周期环的代数同态，故上述每个有限多项式恒等式均与其相容。
例如对每项乘积要使用 $\Delta^{\mathfrak m}(pq)=\Delta^{\mathfrak m}(p)\Delta^{\mathfrak m}(q)$。
这不将图余积 $\Delta_H$ 与动机 coaction 识别为同一个映射。


\subsection{固定锚点之后，哪些算术结构不随方案改变}
\subsubsection{整体常数为什么需要单独处理}
式\eqref{finitemom:eq:action}含 $g_{ba}(G)$，其中\eqref{finitemom:eq:directtransition}可能出现整图本身在另一个参考点的周期。
这个项与当前外动量无关，但其圈数仍为 $h_G$。
局部常数不意味着它由低圈图的周期构成。

对一个本原图，转换已经简化为
\begin{equation}\label{finitemom:eq:primitive}
 \Phi_b(G;\kappa)-\Phi_a(G;\kappa)
   =\Phi_0(G;\rho_G^a)-\Phi_0(G;\rho_G^b).
\end{equation}
右侧没有真子图；它是明确的同圈数参考周期差。

\subsubsection{锚定差及严格低圈三角关系}
给整图选共同的代数欧氏锚点 $\kappa_*$，定义
\begin{equation}\label{finitemom:eq:anchor}
 \widehat\Phi_a(G;\kappa,\kappa_*)
       =\Phi_a(G;\kappa)-\Phi_a(G;\kappa_*).
\end{equation}
商图锚点始终是 $\pi_\gamma\kappa_*$，不强行改成商图自身的参考点。

\begin{theorem}[锚定后的低圈保持]\label{finitemom:thm:anchor}
在形式周期环中有
\begin{align}\label{finitemom:eq:triangular}
 \widehat\Phi_b^{\mathfrak m}(G;\kappa,\kappa_*)
 ={}&\widehat\Phi_a^{\mathfrak m}(G;\kappa,\kappa_*)\nonumber\\
 &+\sum_{\varnothing\ne\gamma\subsetneq G}
 g_{ba}^{\mathfrak m}(\gamma)
 \widehat\Phi_a^{\mathfrak m}
       (G/\gamma;\pi_\gamma\kappa,\pi_\gamma\kappa_*).
\end{align}
右侧增量只使用圈数严格小于 $h_G$ 的图及其产品。
若两方案的整图参考点本来相同，取该共同点为 $\kappa_*$，则原振幅本身满足同一个低圈关系。
\end{theorem}
\begin{proof}
将\eqref{finitemom:eq:cocycle}分别在 $\kappa,\kappa_*$ 展开相减。
局部常数 $g_{ba}(G)$ 消失，单位项给原振幅差，约化余积项给\eqref{finitemom:eq:triangular}。
真非空对数子图及非空商图均严格低圈；$g_{ba}(\gamma)$ 的递归只需 $\gamma$ 及更低圈图，所以全部增量具有所述范围。
若 $\rho_G^a=\rho_G^b=\kappa_*$，两项在锚点均为形式周期零，锚定不改变当前振幅。
\end{proof}

令 $\mathcal A_{<G}$ 为包含此次转换所需全部低圈形式周期的代数。
具体只需有限子图--收缩闭包、两方案参考点及所用当前/锚点评价点；必要时加入反向转换的同类数据。
式\eqref{finitemom:eq:triangular}及其反向关系给
\begin{equation}\label{finitemom:eq:algebra}
 \mathcal A_{<G}[\widehat\Phi_a^{\mathfrak m}(G)]
       =\mathcal A_{<G}[\widehat\Phi_b^{\mathfrak m}(G)].
\end{equation}
这是有基底代数的等式，不是两个数值生成域未经条件的等式。

\begin{corollary}[稳定的实际张量范畴]\label{finitemom:cor:category}
令 $\cC_{<G}$ 为上述低圈实际载体生成的满张量子范畴，$M(p)$ 为形式周期 $p$ 的最小对象。
则
\begin{equation}\label{finitemom:eq:category}
 \langle\cC_{<G},M(\widehat\Phi_a^{\mathfrak m}(G))\rangle_\otimes
 =\langle\cC_{<G},M(\widehat\Phi_b^{\mathfrak m}(G))\rangle_\otimes.
\end{equation}
\end{corollary}
\begin{proof}
式\eqref{finitemom:eq:triangular}将右方案的形式周期写成左方案周期和低圈周期乘积的有限和。
这些和、乘积由实际直和与张量积承载，其最小矩阵系数对象为该载体的子商\cite{brownperiods}。
故右侧生成范畴包含于左侧；反向方案公式给另一包含。
没有使用数值周期映射单射性。
\end{proof}

若所关心的结构子范畴对子商和张量操作封闭，并且包含 $\cC_{<G}$，则锚定振幅的最小对象在两方案中属于该子范畴的条件等价。
例如全部相关低圈载体均 mixed Tate 时，锚定振幅的 mixed Tate 归属性具有这个方案保持性。
这不证明未锚定振幅的孤立对象同构，也不排除某个具体数值的偶然抵消。

\subsubsection{一圈反例说明未锚定对象可改变}
等质量 bubble 的零点值为
\begin{equation}\label{finitemom:eq:bubble}
 \Phi_0(z)=\int_0^1\log\frac{\mu}{\mu+z x(1-x)}\,dx,
 \quad \mu>0,\quad z\geq0.
\end{equation}
在当前点 $z=0$，零点方案的形式周期为零。
另取参考点 $z_0>0$，其振幅为 $\Phi_b(0)=-\Phi_0(z_0)>0$。
严格正性由被积函数的符号直接给出，不求完整积分。
因此后者形式周期非零，最小对象非零，而前者最小对象为零。
这已经反驳了 ``任意改变整图减除点仍保持孤立最小对象''；无需提出任何尚未证明的超越性结论。


\endgroup



\begingroup
\providecommand{\Q}{\mathbb Q}
\providecommand{\A}{\mathbb A}
\providecommand{\PP}{\mathbb P}
\providecommand{\per}{\operatorname{per}}
\providecommand{\im}{\operatorname{im}}
\providecommand{\dR}{\mathrm{dR}}
\providecommand{\SR}{S_{\mathrm R}}
\subsection{图、减除条件与归一化}
考虑两顶点、三条内边的 sunset 二点图。每边传播子为
$(k_e^2+\mu_e)^{-1}$，其中 $\mu_e>0$ 是质量平方。
外欧氏动量记为 $p$，$z=p^2\geq0$。
本节单个圈采用测度 $d^4k/\pi^2$；耦合和图对称因子到第\ref{sunsetjet:sec:field}节才加入。
令
\begin{equation}\label{sunsetjet:eq:poly}
 U=ab+ac+bc,\quad M=\mu_1a+\mu_2b+\mu_3c,\quad
 r=\frac{abc}{U},\quad F_z=UM+zabc,\quad H_z=\frac{F_z}{U}=M+zr.
\end{equation}
这里 $a,b,c$ 为 Schwinger 参数或其齐次坐标，不用它们表示减除点。
在正射影单形 $\Delta\subset\PP^2(\mathbb R)$ 上，取正向形式
\[
 \Omega=a\,db\wedge dc-b\,da\wedge dc+c\,da\wedge db.
\]
所有本节几何对象在包含 $\mu_i,z,z_0$ 的实代数数域 $K$ 上定义；
参数变化时则先取相应的代数族。解析收敛结论允许任意正实质量。

固定 $z_0\geq0$，整体减除算子为
\begin{equation}\label{sunsetjet:eq:taylor}
 t_{z_0}^1 f(z)=f(z_0)+(z-z_0)f'(z_0).
\end{equation}
这是旋转不变二点函数的 MOM 值与斜率条件。
当 $z_0=0$，它等于 $p$ 变量中至二次的 Taylor 减除；
当 $z_0>0$，不把它称为在任意选定非零向量 $p_0$ 处的完整笛卡尔 Taylor 算子。
适用条件明确是 \eqref{sunsetjet:eq:taylor}。

\begin{proposition}[三个重叠子反项的消去]\label{sunsetjet:prop:forest}
先给所有内边施加同一类截断 $a,b,c\geq\varepsilon>0$。
对三个两边 bubble 以任意固定标量参考条件作局部常数减除，
再对整图施加 $1-t_{z_0}^1$。
在每个固定截断处，所得完整森林规则精确等于
\begin{equation}\label{sunsetjet:eq:regulated}
 S_\varepsilon(z)-S_\varepsilon(z_0)
          -(z-z_0)S'_\varepsilon(z_0).
\end{equation}
该结论不依赖三个子图的有限常数选择。
\end{proposition}
\begin{proof}
图有 $N=3$、$h=2$，表面发散度 $4h-2N=2$。
每个两边子图有一个圈，表面发散度零，故其局部反项与子图外动量无关。
三个子图两两重叠，任何许可森林至多包含其中一个。
收缩任意一个子图后，两条外腿位于同一顶点，剩余边是 tadpole，
因而商图积分不依赖 $z$。
所以准备式中每项 $C_{\gamma,\varepsilon}S_{G/\gamma,\varepsilon}$
均被 $1-t_{z_0}^1$ 消去。
逐边 Schwinger 截断使这些积分、参数导数和乘积在固定 $\varepsilon$ 时均有限。
等式先在此处成立，再取共同截断极限；没有将三个发散裸值分别赋予有限部分。
\end{proof}


\subsection{从二阶余项直接得到有理周期}
\begin{theorem}[实际二次减除表示]\label{sunsetjet:thm:period}
设 $z,z_0\geq0$，$\mu_i>0$。命题\ref{sunsetjet:prop:forest}的极限存在且绝对收敛。
记 $d=z-z_0$、$z_t=z_0+td$，则
\begin{align}
 \SR(z;z_0)
 &=\int_\Delta\frac{\Omega}{U^2}
       \left[H_z\log\frac{H_z}{H_{z_0}}-d r\right],\label{sunsetjet:eq:log}\\
 &=d^2\int_{\Delta\times[0,1]}
       \frac{(1-t)(abc)^2\Omega\wedge dt}{U^3F_{z_t}}.\label{sunsetjet:eq:rational}
\end{align}
它满足 $\SR(z_0;z_0)=\partial_z\SR(z_0;z_0)=0$。
在代数参数处，\eqref{sunsetjet:eq:rational} 是实际代数几何对的相对周期。
\end{theorem}
\begin{proof}
先作两个圈的 Gaussian 积分。采用本节测度后，裸 Schwinger 被积式是
\[
 \frac{e^{-M-zr}}{U^2}\,da\,db\,dc.
\]
令 $(a,b,c)=\lambda(\alpha,\beta,\gamma)$，
$\alpha+\beta+\gamma=1$。
径向因子为 $\lambda^{-2}d\lambda$。
式\eqref{sunsetjet:eq:regulated}合并后的径向核是
\begin{equation}\label{sunsetjet:eq:radial}
 \lambda^{-2}e^{-\lambda H_{z_0}}
       \left(e^{-\lambda d r}-1+\lambda d r\right)
 =d^2r^2\int_0^1(1-t)e^{-\lambda H_{z_t}}\,dt.
\end{equation}
这是指数函数的带积分余项 Taylor 公式，对 $d$ 正负均成立。
由于 $z_t\geq0$，在闭单形上 $H_{z_t}\geq\min_i\mu_i>0$。
积分 $\lambda\in(0,\infty)$ 立即给出 \eqref{sunsetjet:eq:rational}。
再对 $t$ 作初等积分得到 \eqref{sunsetjet:eq:log}；
实际周期的定义使用有理式，不以求出对数表达式为前提。

下面证明包括全部角点的可积性，而不只检验各个开边。
将正射影单形分成六个坐标排序扇区。
在其中一个扇区令 $c=1,b=y,a=xy$，$0\leq x,y\leq1$。
置
\[
 A=1+x+xy,\quad U=yA,\quad
 M=\mu_3+\mu_2y+\mu_1xy,\quad r=\frac{xy}{A}.
\]
相应正向密度精确化为
\begin{equation}\label{sunsetjet:eq:sector}
 \frac{\Omega}{U^2}\frac{r^2}{H_{z_t}}
   =\frac{x^2y}{A^4H_{z_t}}\,dx\,dy.
\end{equation}
分母在整个闭立方上有正下界，因而该式在 $x=0$、$y=0$ 及其交点同时正则。
其余五扇区只需置换质量与坐标。
这也证明相应吹起后坐标边界上的代数正则性。

逐边截断在尺度变量中仅加入指示函数
$\mathbf1_{\lambda\alpha_e\geq\varepsilon\ (e=1,2,3)}$。
式\eqref{sunsetjet:eq:radial}非负，并由刚刚证明可积的完整余项控制。
因此共同截断趋零时可取极限，所得正是 \eqref{sunsetjet:eq:rational}。
端点两个条件由因子 $d^2$ 和统一可微性得到。
实际相对对象的构造将在下一节明确给出。
\end{proof}

\begin{corollary}[凸性和所有高阶参数导数]\label{sunsetjet:cor:derivatives}
在正质量欧氏区域，
\begin{align}
 \partial_z\SR(z;z_0)
 &=d\int_{\Delta\times[0,1]}
             \frac{(abc)^2\Omega\wedge dt}{U^3F_{z_t}}
   =\int_\Delta\frac{\Omega}{U^2}r\log\frac{H_z}{H_{z_0}},\label{sunsetjet:eq:first}\\
 K(z):=\partial_z^2\SR(z;z_0)
 &=\int_\Delta\frac{(abc)^2\Omega}{U^3F_z}>0,\label{sunsetjet:eq:second}\\
 \partial_z^k\SR(z;z_0)
 &=(-1)^{k-2}(k-2)!
       \int_\Delta\frac{\Omega}{U^2}\frac{r^k}{H_z^{k-1}}\qquad(k\geq2).
                                                        \label{sunsetjet:eq:higher}
\end{align}
特别地 $K$ 与参考点无关；$\SR(z;z_0)>0$ 当且仅当 $z\ne z_0$。
\end{corollary}
\begin{proof}
在任意紧的正质量欧氏参数集，\eqref{sunsetjet:eq:sector}中的正下界统一成立。
有限阶求导仅提高 $H$ 的分母次数并加入有界的 $r$，
所以可在积分号内求导。
对 \eqref{sunsetjet:eq:log}求导得第一式，再求导得其余各式。
严格正性来自正单形内部的严格正核。
\end{proof}


\subsection{实际相对对象与二阶端点下降}\label{sunsetjet:sec:motive}
取 $X\to\PP^2$ 为坐标顶点的吹起，并按需要作共同 toric 细化。
令 $D$ 为坐标边界的全变换。
六扇区式\eqref{sunsetjet:eq:sector}说明下面使用的形式在正链边界附近均代数正则。
删去 $U$ 及 $F$ 的严格变换零集，而不因齐次表达式中的公共单项式因子删去整个例外面。
具体设
\[
 V_z=X\setminus\bigl(\widetilde{\{U=0\}}\cup
        \widetilde{\{F_{z_0}=0\}}\cup\widetilde{\{F_z=0\}}\bigr),
 \qquad D_z=D\cap V_z.
\]
取 $Y_z\subset V_z\times\A^1_t$，再删去 $F_{z_t}$ 的严格变换零集；
记 $D_Y=Y_z\cap(D_z\times\A^1)$，
$B_z=D_Y\cup(V_z\times\{0,1\})$。
这给出实际对象
\begin{equation}\label{sunsetjet:eq:objects}
 \mathcal E_z=H^3(Y_z,B_z),\qquad
 \mathcal V_z=H^2(V_z,D_z).
\end{equation}
如需进一步解消，仅在不穿过正链的代数极点处作共同细化。
正链 $\widetilde\Delta\times[0,1]$ 的全部边界都在 $B_z$；
正性保证它不遇被删去的严格变换。
因此定理\ref{sunsetjet:thm:period}确实给出实际相对 Nori 对象的矩阵系数，
其 Hodge 和 de Rham 实现由同一对产生。
这不声称 \eqref{sunsetjet:eq:objects} 已经最小。

以下在固定 $z_0,\mu_i$ 的代数参数直线上工作，
限制到上述对给出常秩局部系统的光滑开集。
记 $w=\Omega/U^2$；$w$ 单独不是齐次次数零形式，
但下列每个完整组合均是射影形式。
定义
\begin{align}
 \Xi_z&=d^2(1-t)w\frac{r^2}{H_{z_t}}\wedge dt,&
 \Psi_z&=d\,w\frac{r^2}{H_{z_t}}\wedge dt,\label{sunsetjet:eq:forms}\\
 A_1&=d\,t(1-t)w\frac{r^2}{H_{z_t}},&
 A_2&=t\,w\frac{r^2}{H_{z_t}},&
 \theta_z&=w\frac{r^2}{H_z}.\label{sunsetjet:eq:primitives}
\end{align}

\begin{theorem}[二阶实际端点下降]\label{sunsetjet:thm:descent}
在上述族中存在由三重对给出的实际水平态射
$\iota:\mathcal V\to\mathcal E$，使
\begin{equation}\label{sunsetjet:eq:descent}
 \nabla_z[\Xi_z]=[\Psi_z],\qquad
 \nabla_z^2[\Xi_z]=\iota[\theta_z].
\end{equation}
所有二阶及更高阶导数的最小形式周期对象都来自端点对象的子商。
若 $b=\operatorname{rank}(\im\iota)$，则原周期在该参数直线上的齐次微分方程阶数至多为 $b+2$；
一阶导数至多 $b+1$，二阶导数至多 $b$。
\end{theorem}
\begin{proof}
在 $Y_z$ 上直接微分得到
\begin{equation}\label{sunsetjet:eq:exact}
 \partial_z\Xi_z=\Psi_z+d_YA_1,\qquad
 \partial_z\Psi_z=d_YA_2.
\end{equation}
例如第一式的纯 $t$ 系数之差为
\[
 d\left(\frac{1-2t}{H_{z_t}}
          -\frac{d\,t(1-t)r}{H_{z_t}^2}\right)
   =\partial_t\left(\frac{d\,t(1-t)}{H_{z_t}}\right),
\]
第二式使用
$\partial_t(t/H_{z_t})=H_{z_0}/H_{z_t}^2$。
底空间形式次数为二，故这里 $d_Y(A)=\partial_t A\wedge dt$ 没有额外负号。
两个原函数都含 \eqref{sunsetjet:eq:sector} 的正则核，
在坐标边界的拉回为零；$A_1$ 还在 $t=0,1$ 同时为零。
所以第一式在相对 de Rham 复形中给出第一条等式。

明确取 cone 微分
$D(\omega,\eta)=(d\omega,i^*\omega-d\eta)$，
相对配对为 $\int_\sigma\omega-\int_{\partial\sigma}\eta$。
于是 $[d_YA_2,0]=[0,-i^*A_2]$；
其初端点为零，终端点为 $\theta_z$。
定义 $\iota[\theta]=[0,-(0,\theta)]$。
终端面在 $\widetilde\Delta\times[0,1]$ 上取正向，
该类的周期正是 $\int_{\widetilde\Delta}\theta$，
从而得到第二式及 \eqref{sunsetjet:eq:second}，不产生 Tate 或 $2\pi i$ 因子。
三重对的长正合序列给出该实际态射；它与 Gauss--Manin 连接相容。

其像是水平子对象，故所有继续导数仍在该像中。
将 $[\Xi]$ 生成的微分子模模去该像后，
至多由 $[\Xi]$ 与 $\nabla[\Xi]$ 两个类张成；
因此秩至多 $b+2$，其余两项同理。
这只是微分模阶数界，不是固定算术纤维的秩等号，
也没有假定该端点态射单射。
\end{proof}


\subsection{参考喷射的变换与椭圆几何}
\begin{proposition}[参考点变换的实际形式周期律]\label{sunsetjet:prop:reference}
在同一正质量欧氏区域，任意参考 $\rho,\sigma\geq0$ 满足
\begin{equation}\label{sunsetjet:eq:change}
 \SR(z;\sigma)=\SR(z;\rho)-\SR(\sigma;\rho)
                -(z-\sigma)\partial_z\SR(\sigma;\rho).
\end{equation}
对于代数参数，该等式可提升为实际形式周期等式；
两方案差是至一次的局部多项式，系数由同一图的值与斜率周期给出。
\end{proposition}
\begin{proof}
在共同截断处直接使用两个 Taylor 投影便得到数值恒等式，
前述绝对收敛允许取极限。
形式提升另有直接几何依据：\eqref{sunsetjet:eq:rational}与\eqref{sunsetjet:eq:first}
经代数变量变换 $s=\rho+t(z-\rho)$ 分别成为
\begin{equation}\label{sunsetjet:eq:weighted}
 \SR(z;\rho)=\int_\rho^z(z-s)K(s)\,ds,\qquad
 \partial_z\SR(z;\rho)=\int_\rho^zK(s)\,ds.
\end{equation}
这里的 $K(s)$ 代表实际有理二维被积式与正链，未先作数值求值。
沿有向区间 $[\rho,z]=[\rho,\sigma]+[\sigma,z]$ 分解，
再使用 $z-s=(\sigma-s)+(z-\sigma)$，就得到 \eqref{sunsetjet:eq:change}。
若三点次序不同，用反向链抵消；所有区间都位于无物理极点的非负轴。
端点重合时，带标签端点的相对链恒等式取零链，结论仍成立。
这些操作是实际链加法、代数拉回与相对配对双线性，
所以不是依靠周期映射的单射性推断形式等式。
\end{proof}

\begin{remark}
局部多项式并不意味着低圈周期：\eqref{sunsetjet:eq:change}的系数是两圈 sunset 自身的参考值。
更换任意两个点之后再复合，\eqref{sunsetjet:eq:change}自动给出同一参考变换；
这种喷射复合律尚不是完整物理 RG 流。
\end{remark}

\begin{proposition}[端点曲线的几何类型]\label{sunsetjet:prop:elliptic}
固定正质量后，$E_z=\{F_z=0\}\subset\PP^2$ 在一般 $z$ 处是光滑三次曲线，
因此亏格为一。它与二维 sunset 的图三次曲线相同，
欧氏变量与文献 Minkowski 变量的关系为 $z=-K^2$。
\end{proposition}
\begin{proof}
三次铅笔由 $UM$ 和 $abc$ 生成，没有公共分量。
其基点为三个坐标顶点及各坐标边上 $M=0$ 的点。
在顶点 $[1:0:0]$，$F_z$ 对另两个坐标的线性系数均为 $\mu_1\ne0$；
另外两顶点相同。
在 $c=0$、$ab\ne0$、$M=0$ 处，沿此坐标边的限制为
$ab(\mu_1a+\mu_2b)$，该根单，因此该点也对每个 $z$ 光滑。
特征零 Bertini 定理保证一般成员在基点之外光滑。
故一般成员处处光滑；平面三次的亏格公式给出一。
图多项式的识别直接比较 \eqref{sunsetjet:eq:poly} 与文献\cite[式(2.5)--(2.6)]{bv}。
\end{proof}

对 $E_z$ 的实际定位序列提供其开曲线 $H^1$ 的 Tate 扭转作为几何候选来源；
坐标边界与 conic $U=0$ 另外贡献边界信息。
但定位序列仍有映射和核像，不能单靠出现了三次曲线，
就宣布某一被积形式与指定链一定读出了完整椭圆核。
尤其 $z=0$ 时 $F_0=UM$ 已退化，
正欧氏积分的正则性并不意味着每个参数处的三次曲线都光滑。
本文没有给出固定数值不属于 MZV 的超越性结论。
文献\cite{bv}主要处理二维等质量闭合表达式；
这里的四维减除和实际形式提升来自前述独立推导。


\subsection{两圈场、质量与耦合的有限匹配}\label{sunsetjet:sec:field}
现取单实标量欧氏作用量，相互作用为 $\lambda\phi^4/4!$，
并定义 $u=\lambda/(16\pi^2)$。
这与第\ref{ch:physics}节的物理基准的 $\mathcal M_{\rm tree}=-g$ 相容于 $g=\lambda$，
但本节使用欧氏离壳 MOM 条件，不把它与极点质量、在壳单位留数方案混同。
固定非负二点参考 $\rho$，令
\begin{equation}\label{sunsetjet:eq:2conditions}
 \Gamma^{(2)}_\rho(\rho)=\rho+\mu_\rho,
 \qquad \partial_z\Gamma^{(2)}_\rho(\rho)=1.
\end{equation}
一圈 tadpole 独立于外动量，其整体质量减除使一圈二点校正为零。

\begin{lemma}[二点函数的非局部两圈项]\label{sunsetjet:lem:weight}
在这些条件下，
\begin{equation}\label{sunsetjet:eq:gamma2}
 \Gamma^{(2)}_\rho(z)=z+\mu_\rho
       -\frac{u_\rho^2}{6}\SR(z;\rho;\mu_\rho,\mu_\rho,\mu_\rho)
       +O(u_\rho^3).
\end{equation}
\end{lemma}
\begin{proof}
具有两顶点的四价二点图有三条内部线。
每顶点各接一条外腿时，唯一动量相关的 1PI 拓扑是三条线全连接两顶点的 sunset。
若两条外腿接同一顶点，1PI 图的另一顶点带 tadpole，整个图不携带穿行外动量。
单连接边加两侧 tadpole 的图是一粒子可约，属于传播子迭代而非 $\Gamma^{(2)}$。
一圈耦合反项及质量反项在这一阶的二点插入也只给动量无关项。
所以值与斜率条件消去其余局部项。

sunset 的 Wick 计数可手算：展开系数是
$1/[2!(4!)^2]$；
两条带标签外腿在两顶点的分配有 $2$ 种，
每顶点外腿端口有 $4$ 种，三个剩余端口的配对有 $3!$ 种。
故图系数为 $2\cdot4^2\cdot3!/[2!(4!)^2]=1/6$。
欧氏 $e^{-S_{\rm int}}$ 的二阶项给连通传播子正的 sunset 修正；
取逆传播子后，1PI 非局部项带负号。
每圈从 $d^4k/\pi^2$ 换为 $d^4k/(2\pi)^4$ 带来 $1/(16\pi^2)$，
两圈与 $\lambda^2$ 合成 $u^2$，得到 \eqref{sunsetjet:eq:gamma2}。
\end{proof}

取另一个非负二点参考 $\sigma$，在旧方案参数下记
\begin{equation}\label{sunsetjet:eq:AB}
 A=\partial_z\SR(\sigma;\rho;\mu_\rho,\mu_\rho,\mu_\rho),\qquad
 B=\SR(\sigma;\rho;\mu_\rho,\mu_\rho,\mu_\rho).
\end{equation}
这里 $A$ 无量纲，$B$ 的量纲为质量平方。
定义
\[
 \zeta=\partial_z\Gamma^{(2)}_\rho(\sigma),\qquad
 \phi_\sigma=\zeta^{1/2}\phi_\rho.
\]
平方根按常数项为一的形式幂级数理解，不引入额外代数分支选择。
于是 $\Gamma^{(n)}_\sigma=\zeta^{-n/2}\Gamma^{(n)}_\rho$
是同一有效作用量在新场坐标下的顶点。

\begin{theorem}[有限二点及四点匹配]\label{sunsetjet:thm:matching}
固定参考坐标 $\rho,\sigma$，则
\begin{align}
 \zeta&=1-\frac{u_\rho^2}{6}A+O(u_\rho^3),\label{sunsetjet:eq:zeta}\\
 \mu_\sigma&=\mu_\rho+
       \frac{u_\rho^2}{6}\bigl[(\sigma+\mu_\rho)A-B\bigr]
       +O(u_\rho^3).\label{sunsetjet:eq:mass}
\end{align}
再令旧方案的欧氏四点顶点为
\begin{equation}\label{sunsetjet:eq:vertex}
 \frac{\Gamma^{(4)}_\rho(\kappa)}{16\pi^2}
   =u_\rho+u_\rho^2V_{1,\rho}(\kappa)
            +u_\rho^3V_{2,\rho}(\kappa)+O(u_\rho^4),
\end{equation}
其耦合参考为固定代数欧氏四点配置 $\kappa_\rho$，
$V_{j,\rho}(\kappa_\rho)=0$。
在新的固定参考 $\kappa_\sigma$ 定义耦合后，
\begin{equation}\label{sunsetjet:eq:coupling}
 u_\sigma=u_\rho+u_\rho^2V_{1,\rho}(\kappa_\sigma)
    +u_\rho^3\left[V_{2,\rho}(\kappa_\sigma)+\frac A3\right]
    +O(u_\rho^4).
\end{equation}
$A,B$ 以及在本节统一条件下的 $V_1,V_2$ 都有实际形式周期提升，
所以这些有限匹配系数由实际周期的有限和与乘积给出。
\end{theorem}
\begin{proof}
第一式是 \eqref{sunsetjet:eq:gamma2} 的斜率。
新逆传播子为 $\zeta^{-1}\Gamma^{(2)}_\rho$，
因此 $\mu_\sigma=\zeta^{-1}\Gamma^{(2)}_\rho(\sigma)-\sigma$；
展开到 $u^2$ 得到 \eqref{sunsetjet:eq:mass}。
新四点顶点乘上 $\zeta^{-2}=1+u_\rho^2 A/3+O(u_\rho^3)$，
与 \eqref{sunsetjet:eq:vertex} 相乘就给出 \eqref{sunsetjet:eq:coupling}。
尤其不能遗漏 $A/3$：它来自场的有限归一化，虽始于两圈，却在三阶耦合匹配中出现。

$A,B$ 的实际载体和参考律已在前面构造。
一圈四点 bubble 与两圈非局部双 bubble、wineglass 都在已有对数森林实际周期定理的图类中。
两圈四点中的一圈 tadpole 插入与其质量反项相消，
因为本节的一圈二点减除是全动量无关 tadpole 的完整减除。
剩余局部四点项由同一耦合参考条件固定，
因此 $V_1,V_2$ 可由已证明的非零 MOM 递归及实际直和、张量积实现。
场变换中出现的逆数或平方根，在任意固定微扰阶的系数都是已有系数的有理多项式。
最后，将内部质量由 $\mu_\rho$ 换写为 $\mu_\sigma$，
只会从四点一圈项的 $u^2\cdot O(u^2)=O(u^4)$ 起产生新项，
不会改变本定理到 $u^3$ 的公式。
\end{proof}

同一理论在第三个参考点的匹配，由有效作用量的场坐标变换和
\eqref{sunsetjet:eq:change}复合，故满足逐阶复合律。
这避免将彼此不相容的逐图减除常数拼作一个拉格朗日量方案。
本节特意固定参考坐标；若 $\rho,\sigma,\kappa_\rho,\kappa_\sigma$
依赖正在求导的质量或尺度，必须加入对应的链式法则项。


\subsection{实际场反常维数与有理紫外出口}
上述构造补上了此前纯对数图范围不能提供的实际二点输入。
固定旧方案参数，改变新二点参考 $\sigma$，有
\begin{equation}\label{sunsetjet:eq:scale-response}
 \partial_\sigma\log\zeta=-\frac{u_\rho^2}{6}K(\sigma)+O(u_\rho^3).
\end{equation}
现在明确 RG 条件，令 $\sigma=Q^2>0$，以 $\log Q$ 为 RG 时间。
旧参考及其全部重整化参数固定，即固定同一裸理论；
新质量和耦合由前节的匹配同时定义，不人为冻结它们。
令 $\phi_{\rm bare}=Z_\rho^{1/2}\phi_\rho=Z_Q^{1/2}\phi_Q$，
于是 $Z_Q=Z_\rho/\zeta(Q^2)$。定义
\begin{equation}\label{sunsetjet:eq:gamma-def}
 \gamma_\phi=\frac12\left.Q\frac{d}{dQ}\log Z_Q\right|_{\rm bare}
 =-\frac12\left.Q\frac{d}{dQ}\log\zeta(Q^2)\right|_{\rho,\mu_\rho,u_\rho}.
\end{equation}
这些比值可先在同一正规化下定义，再取已证明存在的有限差；
不需要把单独的发散 $Z_\rho$ 称为有限周期。

\begin{theorem}[两圈场响应与紫外有理值]\label{sunsetjet:thm:gamma}
按 \eqref{sunsetjet:eq:gamma-def} 的约定，
\begin{equation}\label{sunsetjet:eq:gamma}
 \gamma_\phi(Q)=\frac{u_Q^2}{6}\,\mathcal T(Q^2;\mu_Q)+O(u_Q^3),
 \qquad
 \mathcal T(z;\mu)=zK(z;\mu,\mu,\mu).
\end{equation}
对任意固定正质量，$\mathcal T(z;\mu)$ 在 $z>0$ 严格递增，且
\begin{equation}\label{sunsetjet:eq:threshold}
 0<\mathcal T(z;\mu)<\frac12,\qquad
 \lim_{z\to0}\mathcal T(z;\mu)=0,\qquad
 \lim_{z\to\infty}\mathcal T(z;\mu)=\frac12.
\end{equation}
因此在两圈系数层面，其紫外值为
\begin{equation}\label{sunsetjet:eq:uvgamma}
 \gamma_{\phi,\rm UV}(u)=\frac{u^2}{12}+O(u^3).
\end{equation}
有限质量的 $\mathcal T$ 是实际端点周期；
紫外系数所对应的极限有理积分，其形式周期恰等于有理数 $1/2$。
\end{theorem}
\begin{proof}
将 \eqref{sunsetjet:eq:scale-response} 代入定义，
得到 $u_\rho^2Q^2K(Q^2;\mu_\rho)/6$。
由 $u_Q=u_\rho+O(u_\rho^2)$、$\mu_Q=\mu_\rho+O(u_\rho^2)$，
在两圈阶可换写为 \eqref{sunsetjet:eq:gamma}。
这是逐阶恒等式，允许高阶余项的系数含尺度对数。

为证明极限，以下甚至允许三个不同的正质量。
用式 \eqref{sunsetjet:eq:second} 可写成
\begin{equation}\label{sunsetjet:eq:monotone}
 zK(z)=\int_\Delta\frac{\Omega\,abc}{U^3}
                    \frac{zr}{M+zr}.
\end{equation}
内部的最后因子随 $z$ 严格递增，由零趋向一，且严格介于零和一。
其控制积分可以完全用有理原函数计算。
在 $c=1$ 图表中，
\begin{align}
 C&=\int_0^\infty\int_0^\infty
          \frac{ab}{(a+b+ab)^3}\,da\,db \notag\\
  &=\frac12\int_0^\infty\frac{db}{(1+b)^2}
   =\frac12. \label{sunsetjet:eq:C}
\end{align}
第一步只需关于 $a$ 的原函数
\[
 \frac{b}{(1+b)^2}
 \left[-\frac{1}{(1+b)a+b}
       +\frac{b}{2((1+b)a+b)^2}\right],
\]
第二步原函数为 $-1/[2(1+b)]$。
也可先作代数换元 $v=(1+b)a/b$，将二重积分分成
$\int_0^\infty v(1+v)^{-3}dv$ 与
$\int_0^\infty(1+b)^{-2}db$ 的乘积；
其两条紧化区间的端点均代数，角点经同一吹起处理。
所以不需要椭圆特殊函数或数值求积。
正控制核可积，支配收敛给出两端极限与严格不等式；
差分的严格正性给单调性。
在代数紧化后的两个区间端点应用同一有理 Stokes 公式，
还得到 $C^{\mathfrak m}=\tfrac12\,1^{\mathfrak m}$；
不是从其数值恰好有理反推形式周期有理。
将 $C=1/2$ 代入 \eqref{sunsetjet:eq:gamma}，得到 \eqref{sunsetjet:eq:uvgamma}。
\end{proof}

这里的紫外极限针对微扰系数，等价于该阶阈值函数的质量/尺度比趋零；
并未把固定非零耦合下的大对数微扰级数当成一致收敛函数，
也未声称实际跑动质量沿任意轨道自动满足这一极限。
式 \eqref{sunsetjet:eq:uvgamma} 是经典结果，和文献
\cite[式(1)、(2)、(5)、(9)，$n=1$]{kp} 的归一化吻合；
本文给出的贡献是把它从实际二点周期和明确场坐标中直接导出。
有限尺度的椭圆环境并不迫使这个紫外系数含椭圆常数。
不过上述结论只识别此极限的有理周期，
不建立从每个固定椭圆纤维到单位对象的无条件动机投影。


\endgroup

\clearpage\part{物理出口}

\section{两圈MOM的RG出口}\label{ch:momrg}

\begingroup

\otherref{feynman_periods_supplement}{补充材料S10}说明MS型形式临界反演；本节给实际有限质量MOM出口并证明两者间不能省略的质量坐标条件。下一节的在壳散射基准使用另行声明的物理规范。

\providecommand{\Q}{\mathbb Q}
\providecommand{\SR}{S_{\mathrm R}}
\providecommand{\gphi}{\gamma_\phi}
\providecommand{\dd}{\mathcal D}
\renewcommand{\Z}{\mathcal Z}

\subsection{同一方案中的二点、四点及尺度}
考虑四维单实标量 $\lambda\phi^4/4!$，令
$u=\lambda/(16\pi^2)$，内部质量平方记为 $\mu>0$。
单图的 Schwinger 核沿用每圈 $d^4k/\pi^2$ 的归一化；
$u$ 已吸收通常的圈因子。
本节取欧氏非例外对称四点配置 $\kappa_z$：
\begin{equation}\label{momrg:eq:symmetric}
 p_i^2=\frac{3z}{4},\quad p_i\cdot p_j=-\frac z4\ (i\ne j),
 \quad \sum_i p_i=0,\qquad z=Q^2>0.
\end{equation}
于是任意两腿通道的平方均为 $z$。
例如 $p_i$ 可取 $\sqrt z/2$ 乘以四个正四面体的符号向量。
这是离壳欧氏配置，不是旧散射基准的在壳交叉对称点。

二点条件和耦合条件同时定义为
\begin{equation}\label{momrg:eq:conditions}
 \Gamma_Q^{(2)}(z)=z+\mu_Q,\qquad
 (\Gamma_Q^{(2)})'(z)=1,\qquad
 \Gamma_Q^{(4)}(\kappa_z)=16\pi^2u_Q.
\end{equation}
所有外腿对称性和局部四点反项均按这些共同条件处理。
RG 时间是 $\log Q$；取导数时固定同一裸理论，
因此新 $u_Q,\mu_Q$ 随参考点同时改变。
在微扰语言中，可以固定某个旧参考及其全部重整化参数来实现这一操作。
这不意味求导时把新 $\mu_Q,u_Q$ 人为冻结。

第\ref{ch:sunsetjet}节的二点构造给出实际相对周期 $\SR(z;\rho;\mu)$，
其中 $\rho$ 是固定旧二点参考，且
\begin{equation}\label{momrg:eq:sunset-input}
 \Gamma_\rho^{(2)}(z)
 =z+\mu-\frac{u_\rho^2}{6}\SR(z;\rho;\mu)+O(u_\rho^3),
 \qquad
 K(z;\mu)=\partial_z^2\SR(z;\rho;\mu)>0.
\end{equation}
三个内部质量在此相同。
明确地，对三边参数 $\alpha,\beta,\gamma$，
\begin{equation}\label{momrg:eq:K}
 K(z;\mu)=\int_\Delta
 \frac{(\alpha\beta\gamma)^2\Omega}
 {U_s^3\{\mu(\alpha+\beta+\gamma)U_s+z\alpha\beta\gamma\}},
 \quad U_s=\alpha\beta+\alpha\gamma+\beta\gamma.
\end{equation}
它是已有二维端点对象的实际周期，而不只是候选数值表示。
记质量/尺度比 $r=\mu/z$，定义
\begin{equation}\label{momrg:eq:threshold-sunset}
 T(r)=zK(z;\mu).
\end{equation}
齐次性使右侧仅依赖 $r$。已知 $0<T(r)<1/2$，
$T(0^+)=1/2$，$T(+\infty)=0$；
该紫外端点的形式周期为 $\frac12\,1^{\mathfrak m}$。


\subsection{质量不是自动乘性跑动的坐标}\label{momrg:sec:mass}
\begin{theorem}[二点值/斜率的精确质量流]\label{momrg:thm:mass}
固定裸理论，令旧场坐标中的二点逆传播子为 $F(z)$，
并假设 $F'(z)$ 可逆；形式微扰中其常数项为一。
定义新场因子 $\zeta=F'(z)$，则
\begin{equation}\label{momrg:eq:mass-coordinate}
 \mu_Q=\frac{F(z)}{F'(z)}-z,\qquad
 \gphi=-z\frac{F''(z)}{F'(z)}.
\end{equation}
因此有精确恒等式
\begin{equation}\label{momrg:eq:mass-flow}
 Q\frac{d\mu_Q}{dQ}=2(z+\mu_Q)\gphi,\qquad
 Q\frac{dr_Q}{dQ}=-2r_Q+2(1+r_Q)\gphi,
 \quad r_Q=\mu_Q/Q^2.
\end{equation}
在第\ref{ch:sunsetjet}节的实际两圈范围内，它给出
\begin{equation}\label{momrg:eq:mass-two}
 \gphi=\frac{u_Q^2}{6}T(r_Q)+O(u_Q^3),\qquad
 \dot r_Q=-2r_Q+\frac{u_Q^2}{3}(1+r_Q)T(r_Q)+O(u_Q^3).
\end{equation}
\end{theorem}
\begin{proof}
新二点函数为 $F/F'(z)$，在参考点的斜率为一；
值条件立即给出 $\mu_Q$。
由 $\phi_{\rm bare}=Z_Q^{1/2}\phi_Q$ 和
$Z_Q=Z_\rho/\zeta$，定义
$\gphi=\frac12Qd\log Z_Q/dQ|_{\rm bare}$，即得第二式。
直接求导，
\[
 \frac{d}{dz}\left(\frac F{F'}-z\right)
 =-\frac{FF''}{(F')^2}
 =-(z+\mu_Q)\frac{F''}{F'}.
\]
乘上 $2z$ 并除去无量纲质量的经典尺度项，得到
\eqref{momrg:eq:mass-flow}。
代入 \eqref{momrg:eq:sunset-input}，再在本阶换写为新参数，
得到 \eqref{momrg:eq:mass-two}。
精确恒等式是条件性运动学结论，不主张已经构造任意高圈的实际周期。
\end{proof}

式 \eqref{momrg:eq:mass-flow} 含 $2Q^2\gphi$，故不是
$\dot\mu=\gamma_m\mu$ 的乘性形式。
仅对 $\mu_Q\ne0$ 强行定义对数导数，会产生
$2(1+1/r_Q)\gphi$，它不能自动认作 MS 方案中的质量反常维数。
本节未改变旧基准的极点质量；这里的 $\mu_Q$ 是逆传播子在参考点切线的截距。


\subsection{四点森林的完整非局部部分}\label{momrg:sec:forests}
\subsubsection{无需程序的拓扑与系数}
对四条带标签外腿，含 $V$ 个四价顶点的图有
$I=2V-2$、$L=V-1$。
一圈无 tadpole 的 1PI 图是三个通道 bubble，各有平行边对称因子 $1/2$。
两圈、三个顶点的外腿分配，除去 tadpole 和一粒子可约图，
只有 $(2,2,0)$ 与 $(2,1,1)$。
前者由度数方程唯一给出双 bubble，三个通道各权重 $1/(2!2!)=1/4$；
后者唯一给出 wineglass，双外腿顶点有 $\binom42=6$ 种选择，
双边的对称因子给每图 $1/2$。
一圈 tadpole 自能与其质量反项在本方案完整相消；
外腿可约项不属于 1PI 顶点。
因此以下表达穷尽所需非局部四点部分，局部常数由共同耦合条件固定。

\subsubsection{显式正质量核}
固定旧参考 $\rho>0$。置
\begin{equation}\label{momrg:eq:one-loop}
 X=x(1-x),\quad Y=y(1-y),\quad
 L_\rho(z)=\int_0^1\log\frac{\mu+zX}{\mu+\rho X}\,dx.
\end{equation}
重整化标量 bubble 本身是 $-L_\rho$；欧氏四点一圈顶点带相反号。
对 wineglass 取边 $a=AB,b=AC,c=BC,d=BC$，A 接两条外腿。
生成树及二森林直接给出
\begin{equation}\label{momrg:eq:wine-polynomial}
 U=cd+(a+b)(c+d),\quad
 F_z=\mu(a+b+c+d)U
       +z\left[ab(c+d)+\frac34(a+b)cd\right].
\end{equation}
其经典图多项式可与文献\cite[第5节]{bk2012}逐项比较；
本节的 $3/4$ 来自 \eqref{momrg:eq:symmetric}，不能照搬在壳值。

采用 $a=\ell x,b=\ell(1-x),c=\ell t y,d=\ell t(1-y)$，定义
\begin{align}
 A_t&=1+tY,& D_z&=\mu+zX,& C_\rho&=\mu+\rho Y, \notag\\
 Q_z&=\mu(1+t)A_t+z\left(X+\frac34tY\right).
                                                        \label{momrg:eq:Q}
\end{align}
按同一个 Schwinger 截断先作子 bubble 减除、再作总体减除，
积分总尺度后得到
\begin{equation}\label{momrg:eq:W}
 W_\rho(z)=\int_0^1dx\int_0^1dy\int_0^\infty\frac{dt}{t}
 \left[
  \log\frac{D_z+tC_\rho}{D_\rho+tC_\rho}
  -\frac1{A_t^2}\log\frac{Q_z}{Q_\rho}
 \right].
\end{equation}
符号由“当前裸积分减参考裸积分”及子反项决定，
与本节欧氏顶点约定一起使用；不直接沿用任意文献的整体图符号。

式 \eqref{momrg:eq:W} 是前述正质量单子图森林实际周期在对称族上的拉回。
总尺度积分前，裸项指数为 $-\ell Q_z/A_t$，
子反项乘积指数为 $-\ell(D_z+tC_\rho)$，
相对密度分别为 $1/(tA_t^2)$ 与 $1/t$；
这一检查直接固定 \eqref{momrg:eq:W} 的符号和系数。
正质量的合并收敛及实际四维相对载体沿用已证单子图模型，
不分别把两个发散项解释成周期。

\begin{proposition}[有限顶点与场匹配]\label{momrg:prop:vertex}
在旧二点、四点参考均为 $\rho$ 的方案中，
\begin{equation}\label{momrg:eq:vertex}
 \frac{\Gamma_\rho^{(4)}(\kappa_z)}{16\pi^2}
 =u_\rho+\frac32u_\rho^2L_\rho(z)
    +u_\rho^3\left[\frac34L_\rho(z)^2+3W_\rho(z)\right]
    +O(u_\rho^4).
\end{equation}
令 $A_s(z)=\partial_z\SR(z;\rho;\mu)$。
变到新场坐标并在 $\kappa_z$ 定义耦合，得到
\begin{equation}\label{momrg:eq:matching}
 u_Q=u_\rho+u_\rho^2 f_1(z)+u_\rho^3 f_2(z)+O(u_\rho^4),
 \quad
 \begin{cases}
 f_1=\frac32L_\rho,\\
 f_2=\frac34L_\rho^2+3W_\rho+\frac13A_s.
 \end{cases}
\end{equation}
\end{proposition}
\begin{proof}
双 bubble 的准备式减去总体参考后成为 $L_\rho^2$；
三个图给 $3/4$。六个 wineglass 权重和为三。
二圈欧氏 1PI 四点顶点的符号为正，一圈为负，
与标量 bubble 的 $-L_\rho$ 结合给出 \eqref{momrg:eq:vertex}。
这些系数也可由有效作用量的微扰展开确定。
一圈子反项插入已包含在各准备式中，不另加一次。
最后，新四点顶点乘以
$\zeta^{-2}=1+u_\rho^2 A_s/3+O(u_\rho^3)$，
即得 \eqref{momrg:eq:matching}。
这保留了此前已证明的场归一化项。
\end{proof}


\subsection{先求导再拉回参考点：两圈实际端点}\label{momrg:sec:beta}
令 $\dd=Q\partial_Q=2z\partial_z$，求导保持旧 $\rho,\mu,u_\rho$ 固定。
形式反演 \eqref{momrg:eq:matching} 给
\begin{equation}\label{momrg:eq:beta-old}
 \beta_Q=u_Q^2\dd f_1+
 u_Q^3\{\dd f_2-2f_1\dd f_1\}+O(u_Q^4).
\end{equation}
于是三阶系数为
\begin{equation}\label{momrg:eq:beta-combination}
 6z\{W_\rho'(z)-L_\rho(z)L_\rho'(z)\}
                   +\frac{2z}{3}K(z;\mu).
\end{equation}

\begin{lemma}[旧参考的精确消去]\label{momrg:lem:reference}
组合 $W_\rho'-L_\rho L_\rho'$ 不依赖旧参考 $\rho$。
它等于先对 $W_\rho(z)$ 的当前变量求导、再置 $\rho=z$ 的结果，
而不是恒零函数 $W_z(z)$ 的全导数。
\end{lemma}
\begin{proof}
改变旧参考为 $\sigma$ 时，
$L_\sigma(z)=L_\rho(z)-L_\rho(\sigma)$，而
\[
 W_\sigma'(z)-W_\rho'(z)
 =\int dx\,dy\int_0^\infty\frac{dt}{t}\,
 X\left[\frac1{D_z+tC_\sigma}-\frac1{D_z+tC_\rho}\right]
 =-L_\rho(\sigma)L_\rho'(z).
\]
最后一步是收敛的有理 Frullani 恒等式；
也可以在相对 Kummer 模型中用有理变量替换和链加法证明。
两项因而恰好相消。
先求导后的 $L_\rho(z)$ 在 $\rho=z$ 为零，给出所述拉回。
若一开始沿对角求全导数，会错误地连物理响应也消去。
\end{proof}

令 $r=\mu/z$，保持 $X,Y,A_t$ 的定义，置
\[
 q_r=r(1+t)A_t+X+\frac34tY,\qquad
 j_r(x,y,t)=\frac1t
 \left[
  \frac{X}{r+X+t(r+Y)}
       -\frac{X+\frac34tY}{A_t^2q_r}
 \right].
\]
定义实际端点积分
\begin{equation}\label{momrg:eq:J}
 J(r)=\int_0^1dx\int_0^1dy\int_0^\infty j_r(x,y,t)\,dt.
\end{equation}

\begin{theorem}[有限正质量的两圈RG核]\label{momrg:thm:beta}
在 \eqref{momrg:eq:conditions} 的同一方案中，
\begin{align}
 \beta(u,r)&=b_1(r)u^2+b_2(r)u^3+O(u^4),\label{momrg:eq:beta}\\
 b_1(r)&=3\int_0^1\frac{X}{r+X}\,dx,\qquad
 b_2(r)=6J(r)+\frac23T(r),\label{momrg:eq:b12}\\
 \gphi(u,r)&=\frac{u^2}{6}T(r)+O(u^3),\qquad
 \dot r=-2r+\frac{u^2}{3}(1+r)T(r)+O(u^3).
                                                        \label{momrg:eq:vector}
\end{align}
这里 $r$ 最终取新质量比 $r_Q$。
在任意正代数 $r$，所显示的系数均为已构造的实际形式周期；
不需预先求出整个四点函数或 sunset 的特殊函数表达式。
\end{theorem}
\begin{proof}
将引理\ref{momrg:lem:reference}代入 \eqref{momrg:eq:beta-combination}，
按 \eqref{momrg:eq:Q} 对当前 $z$ 求导，便得到
$zW_\rho'|_{\rho=z}=J(r)$。
一圈来自 $\dd f_1=3zL_\rho'$，其余两式是定理\ref{momrg:thm:mass}。
旧质量换成新质量的差从 $O(u^2)$ 开始，
故其代入一圈四点项只影响 $O(u^4)$；
耦合换参则已由 $-2f_1\dd f_1$ 计入。

实际性需在几何层面解释：$W_\rho$ 的共同有理对数提升，
在当前参数求导后给端点有理三形式；
子图面相消在求导及 $\rho=z$ 的代数拉回下保持。
这是已有单子图相对对的态射与限制，不是从数值等式逆推。
积分 $J(r)$ 的合并核正是该拉回形式。
$T(r)$ 来自二点端点，$b_1$ 来自 bubble 端点。
有限和与有理系数在直和载体上实现。
下一节还直接证明这些合并积分的绝对收敛及紫外边界控制。
\end{proof}


\subsection{同时角点的控制及紫外精确求值}\label{momrg:sec:bounds}
把质量直接置零之前，必须排除 $x,y,t$ 端点与 $r$ 同时趋边界产生的额外项。
以下分解完成这一点，亦证明正质量 $J$ 的合并绝对收敛。
记 $d=r+X$、$c=3/4$，则
\begin{align}
 j_r={}&\frac{X}{d}\frac1t
 \left[\frac1{1+t(r+Y)/d}-\frac1{A_t^2}\right]
 +R_r,\label{momrg:eq:decomposition}\\
 R_r={}&\frac{r\{X(1+Y+tY)-cY\}}{d A_t^2q_r}.
                                                        \label{momrg:eq:R}
\end{align}
这只是通分后的有理恒等式，保留了小 $t$ 的相消。

\begin{lemma}[统一可积性与质量极限]\label{momrg:lem:domination}
当 $r\downarrow0$ 时，$J(r)$ 收敛到将 $r=0$ 代入合并核的积分；
没有额外边界常数。并且当 $r\to\infty$，$J(r)=O(1/r)$。
\end{lemma}
\begin{proof}
先将 \eqref{momrg:eq:decomposition} 第一项对 $t$ 积分。结果是
\begin{equation}\label{momrg:eq:main-integrated}
 \frac{X}{r+X}
       \left[1+\log\frac{(r+X)Y}{r+Y}\right].
\end{equation}
其绝对积分可由 $1-\log X-\log Y$ 控制：
事实上 $XY\leq(r+X)Y/(r+Y)<1$。
将两个有理分式以 $1/(1+tY)$ 为中间项拆开，
第一部分定号，第二部分非负，
还给出原 $t$ 积分绝对值的同一控制。
所以对 $x,y$ 的支配收敛成立，极限是 $1+\log X$。

对余项，分别控制其 $X$ 和 $cY$ 两部分。
因 $q_r\geq r(1+t)A_t$、$X/d\leq1$，
正部分的绝对值不超过
\begin{equation}\label{momrg:eq:Rplus-bound}
 \frac{2}{(1+t)(1+tY)^2}.
\end{equation}
它的 $t$ 积分至多 $3-2\log Y$：
分成 $0<t<1$、$1<t<1/Y$ 及 $t>1/Y$ 即可。
另一方面，$q_r\geq X+ctY$、$r/d\leq1$ 给
\begin{equation}\label{momrg:eq:Rminus-bound}
 \frac{cY}{(1+tY)^2(X+ctY)}.
\end{equation}
令 $v=tY$，其积分由
$\log(1+c/X)+1/2$ 控制；这在 $x$ 上可积。
每个内部点的 $R_r$ 均趋零，故上述统一控制给 $\int R_r\to0$。
这同时控制了耦合角点，不是仅逐个开面取极限。

最后当 $r\geq1$，第一项的 $X/(r+X)\leq1/(4r)$，
而对数仍由 $-\log X-\log Y$ 控制。
余项改用 $q_r\geq r(1+t)A_t$ 及 $d\geq r$，
得到 $1/r$ 乘同类可积函数。
于是 $J(r)=O(1/r)$。
\end{proof}

\begin{theorem}[有理紫外系数与低能衰减]\label{momrg:thm:UV}
实际系数满足
\begin{equation}\label{momrg:eq:UV}
 b_1(0^+)=3,\qquad J(0^+)=-1,\qquad
 b_2(0^+)=-\frac{17}{3}.
\end{equation}
紫外端点的实际形式周期分别为这些有理数乘单位对象。
在另一端，$b_1(r)\sim1/(2r)$，
$b_2(r)=O(1/r)$，$T(r)=O(1/r)$。
因此耦合与场的这些量子响应在固定正质量、低减除尺度的极限衰减；
无量纲质量流还含 $1+r$ 因子，不能据此说其量子项也趋零。
\end{theorem}
\begin{proof}
一圈用支配收敛；$X$ 的区间积分为 $1/6$，也给出大 $r$ 首项。
两圈由引理\ref{momrg:lem:domination}，可以使用
\begin{equation}\label{momrg:eq:massless}
 J(0)=\int dx\,dy\int_0^\infty\frac{dt}{t}
       \left[\frac{X}{X+tY}-\frac1{(1+tY)^2}\right].
\end{equation}
作代数换元 $v=tY$ 后，内层积分恰为 $1+\log X$。
故
\[
 J(0)=\int_0^1[1+\log x+\log(1-x)]\,dx=1-1-1=-1.
\]
形式周期等式亦有实际证明：$v$ 积分分成一个有理端点周期 $1$
和 Kummer 周期 $\log^{\mathfrak m}X$；
Kummer 乘法律把它拆成 $\log^{\mathfrak m}x+\log^{\mathfrak m}(1-x)$。
再以实际有向三角链
$\log x=-\int_x^1dt/t$ 得
$\int_0^1\log^{\mathfrak m}x\,dx=-1^{\mathfrak m}$，
交换两维链顺序后内层长度为 $t$，其相消是代数推前或 Stokes 关系。
另一个对数同理。
这无需假定数值周期映射单射。
二点已证 $T(0)=1/2$，故 $6(-1)+(2/3)(1/2)=-17/3$。

大 $r$ 的 $J$ 已在引理中处理；
由 \eqref{momrg:eq:K}，固定单形质量和后
$T(r)\leq C_s/r$，
其中 $C_s$ 是正质量端点正则性已保证有限的有理核积分。
其余结论随之成立。
\end{proof}

因此在本规范的逐阶紫外边界，
\begin{equation}\label{momrg:eq:classical}
 \beta_{\rm UV}=3u^2-\frac{17}{3}u^3+O(u^4),
 \qquad \gamma_{\phi,\rm UV}=\frac{u^2}{12}+O(u^3).
\end{equation}
这些是经典系数，和文献\cite[第II节及式(19)]{kp}的 $n=1$ 约定一致；
本节不主张发现了新系数。
新接合在于从实际相对周期、先微分后拉回和受控退化得到它们。
特别地，若丢弃场项，就会得到错误的两圈值 $-6$；
正确的 $1/3$ 补项具有明确二点来源。
这展示了图级几何与最终局部读出之间的真实筛选，
不等于每个固定椭圆纤维都有到单位对象的普遍投影。


\subsection{临界坐标的障碍及与MS反演的边界}\label{momrg:sec:critical}
\subsubsection{零截距面并不自动不变}
由 \eqref{momrg:eq:mass-two} 的受控 $r\to0$ 系数极限，
\begin{equation}\label{momrg:eq:r-zero}
 \left.\dot r\right|_{r=0}
 =\frac{u^2}{6}+O(u^3).
\end{equation}
所以不能把 $r=0$ 直接当作此方案的不变无质量流。
对本节显示的两圈截断向量场，若存在 $C^2$ 不变图 $r=h(u)\geq0$，
$h(0)=h'(0)=0$，则 $h'=O(u)$。
其不变性条件
$h'(u)\beta=\dot r|_{r=h(u)}$ 强制
\begin{equation}\label{momrg:eq:critical-graph}
 h(u)=\frac{u^2}{12}+o(u^2).
\end{equation}
因为左端为 $O(u^3)$，而右端为 $-2h+u^2/6+o(u^2)$。
这只是截断场的必要首项，不是完整临界流形存在性定理；
推广到全场还需控制高圈余项沿该质量边界的行为。

也可直接检验逆易感率的扰动条件。
在参考 $z=Q^2$，第\ref{ch:sunsetjet}节的二点公式给
\[
 \Gamma_Q^{(2)}(0)=\mu_Q-\frac{u_Q^2}{6}\SR(0;z;\mu_Q)+O(u_Q^3),
 \qquad
 \SR(0;z;\mu)=\int_0^z sK(s;\mu)\,ds.
\]
由于 $0<sK(s;\mu)<1/2$，当 $\mu/z\to0$ 时
$\SR(0;z;\mu)/z\to1/2$。
因此截断到已证明阶数的 $\Gamma_Q^{(2)}(0)=0$ 条件也给
$\mu_Q/Q^2=u_Q^2/12$ 的首项。
更高阶在临界边界的统一控制、真实临界流形及稳定指数仍需独立证明。

\subsubsection{质量边界并非普通二变量形式幂级数}
问题还不止是坐标平移。一圈已出现
\begin{equation}\label{momrg:eq:b1-nonanalytic}
 b_1(r)=3\left[1-r\int_0^1\frac{dx}{r+x(1-x)}\right]
       =3-6r\log(1/r)+O(r),\qquad r\downarrow0.
\end{equation}
这可把积分在 $x=1/2$ 对称分开，再作一次有理积分证明；
精确表达为
\[
 \int_0^1\frac{dx}{r+x(1-x)}
 =\frac4{\sqrt{1+4r}}
       \operatorname{artanh}\frac1{\sqrt{1+4r}}.
\]
故向量场虽在每个 $r>0$ 解析，不能直接当作
$\Q[[u,r]]$ 中的原点正则形式向量场。
有限质量系数存在连续紫外极限，不等于所有质量导数都能与该极限交换。

\otherref{feynman_periods_supplement}{补充材料S10}的RG--临界展开反演有明确前提：
一个归一化单耦合截面、
$\beta=-\epsilon u+uB(u)$、量子系数无额外 $\epsilon$ 依赖，
以及已经识别的乘性热方向。
本节的 $\mu_Q$ 是切线截距，满足 \eqref{momrg:eq:mass-flow}；
不能把其对数导数塞入该反演式就得到热指数 $\nu$。
真正的比较还需正则临界坐标或适当的复合算符插入，
并证明它与固定方案的实际周期及流相容。
相同的紫外 $\beta$、$\gamma_\phi$ 首项也不足以替代这项证明。


\endgroup

\section{独立物理基准与开放问题}\label{ch:physics}

\begingroup


\subsection{理论、方案与有限验收范围}
取一个实标量场，未破缺对称性为 $\phi\mapsto-\phi$，拉格朗日量为
\begin{equation}
 \mathcal L=\frac12(\partial\phi)^2-\frac12m^2\phi^2
 -\frac{g}{4!}\phi^4+\mathcal L_{\mathrm{ct}}.
\end{equation}
采用度规 $(+---)$、$S=1+iT$ 及树振幅 $\mathcal M_{\mathrm{tree}}=-g$。二点函数在 $p^2=m^2$ 减去值与一阶导数，使质量为极点质量、留数为一。四点耦合在解析对称点 $s=t=u=4m^2/3$ 定义；四个外腿在壳，故 $s+t+u=4m^2$。

外部态归一化为 $2E(2\pi)^3\delta^3$，本节 $d\Phi_n$ 已包含 $(2\pi)^4\delta^4$。同种末态的阶乘另计。与标准运动学公式\cite{pdg}比较时必须避免重复乘入四维动量守恒因子。这里不存在规范 Ward 恒等式；分子只有标量常数及规定的二点反项。


\subsection{全图、反项与 LSZ}
对四条带标签外腿，$V$ 个四价顶点满足 $I=2V-2$、$L=V-1$。现有穷尽及独立 Wick 配对给树级一张、一圈七张、两圈四十二张连通图。两圈非局部四点 1PI 部分为三张双泡（每张权重 $1/4$）与六张 wineglass（每张权重 $1/2$）。三张 tadpole 插入与质量反项相消，三十张外腿可约图由单位留数的 LSZ 处理。

外腿图不能再次作为四点 1PI 贡献加入。泡子发散的有限减除、一圈耦合反项插入及两圈总体局部反项，都由同一耦合定义固定。这个步骤使有限表达对应一个明确物理方案，而非裸图有限部分的任意相加。


\subsection{有限振幅的解析表示}
置 $x_0=4/3$，定义
\begin{equation}
 B(x)=-\int_0^1\log[1-xu(1-u)-i0]\,du,\qquad b(x)=B(x)-B(x_0).
\end{equation}
在 $x>4$ 令 $\beta=\sqrt{1-4/x}$，则
\begin{equation}
 B(x)=2-\beta\log\frac{1+\beta}{1-\beta}+i\pi\beta,\qquad
 B_0=2-2\sqrt2\arctan(1/\sqrt2).
\end{equation}
交叉通道角平均和 wineglass 减除函数为
\begin{equation}
 A(x)=\frac1{x-4}\int_0^{x-4}B(-v)\,dv,
\end{equation}
\begin{equation}
 J(x)=(x-x_0)\int_4^\infty
 \frac{\beta(r)[A(r)-B_0]}{(r-x_0)(r-x-i0)}\,dr.
\end{equation}
于是
\begin{equation}
\begin{aligned}
 \mathcal M&=-g+g^2\mathcal A_1+g^3\mathcal A_2+O(g^4),\\
 \mathcal A_1&=\frac1{32\pi^2}\sum_{\xi=s,t,u}b(\xi/m^2),\\
 \mathcal A_2&=\frac1{1024\pi^4}\sum_{\xi=s,t,u}
 \bigl[-b(\xi/m^2)^2-4J(\xi/m^2)\bigr].
\end{aligned}
\end{equation}
三个通道同时在 $x_0$ 减除，所以交叉对称与耦合条件直接成立。

色散式以外，$J$ 具有实际森林减除得到的 Schwinger 表示。记 $Y=y(1-y)$、$D_x=1-xu(1-u)$、$C_0=1-x_0Y$、$Q_x=D_x+t+t^2Y$，则
\begin{equation}
 J(x)=\int_0^1du\int_0^1dy\int_0^\infty\frac{dt}{t}
 \left[\log\frac{D_x+tC_0}{D_{x_0}+tC_0}
 -\frac1{(1+tY)^2}\log\frac{Q_x}{Q_{x_0}}\right].
\end{equation}
$t=0$ 两项常数相消；$t\to\infty$ 分区控制 $tY$ 后可积。低于阈值时正分母允许交换积分。切口上切开外层两条线，剩余泡处于非正动量转移，得到 $\Im J=\pi\beta(A-B_0)$。一次减除的 Cauchy 公式及已固定的局部常数给出前述色散表示。两种表示的相容不只依赖数值吻合。


\subsection{匹配阶数的光学检验与方案变化}
两粒子中间态带因子 $1/2!$。由树与一圈的完整干涉得
\begin{equation}
 \Im\mathcal A_1=\frac{\beta}{32\pi},\qquad
 \Im\mathcal A_2=-\frac{\beta}{512\pi^3}
 [\Re B(s/m^2)+2A(s/m^2)-3B_0].
\end{equation}
双泡与 wineglass 的吸收部分别给前一方括号中的对应项，有限反项 $B_0$ 也必须保留。这是两圈振幅 $g^3$ 阶的检验；四体态从 $g^4$ 阶光学右端出现，所以不能把它称为三圈前向 $g^4$ 全检验。

全角域、同种末态已除二的弹性密度为
\begin{equation}
 \frac{d\sigma_{\mathrm{el}}}{d\Omega}
 =\frac{g^2-2g^3\Re\mathcal A_1+
 g^4(|\mathcal A_1|^2-2\Re\mathcal A_2)}{128\pi^2s}+O(g^5).
\end{equation}
阈值耦合 $g_{\mathrm{th}}=-\mathcal M(4m^2,0,0)$ 与对称点耦合满足 $g_{\mathrm{th}}=g-ag^2-bg^3+O(g^4)$。反演并代回给
\begin{equation}
 \mathcal A_1^{\mathrm{th}}=\mathcal A_1-a,\qquad
 \mathcal A_2^{\mathrm{th}}=\mathcal A_2+2a\mathcal A_1-2a^2-b.
\end{equation}
截面逐阶相同。这里比较的是同时换参后的可观测量，不是要求各方案的单图有限部分相同。


\subsection{四体产生、正相空间与阈值律}
当 $s>16m^2$，四体通道开通。令六个外部动量全取向外，$\ell=(-p_1,-p_2,k_1,k_2,k_3,k_4)$。两四价顶点的树图由十个无序 $3|3$ 分割给出：
\begin{equation}
 \mathcal M_{24}^{(0)}=-g^2\sum_{\substack{|A|=3\\0\in A}}
 \frac1{Q_A^2-m^2+i0},\qquad Q_A=\sum_{i\in A}\ell_i.
\end{equation}
每张树的 Wick 权重为一；相补三元组不可重复计数。四项类时分母至少为 $8m^2$，余六项可写 $q=K-p_1=p_2-R$，其中 $K,R$ 为两未来向末态之和。若 $q$ 因果，质量不等式将迫使某入射粒子的质量至少为 $2m$，矛盾。因此这六项严格类空，整个物理积分域内没有内传播子极点。

记分母和为 $S_6$。首阶四体截面是
\begin{equation}
 \sigma_{24}=\frac{g^4}{2s\beta\,4!}\int d\Phi_4\,|S_6|^2+O(g^5).
\end{equation}
必须先相加再平方，保留全部干涉。将相空间拆成两对，或拆成逐次 $3+1$，均给独立的有限积分表示。例如 $a=(k_1+k_2)^2,b=(k_3+k_4)^2$ 时
\begin{equation}
 d\Phi_4=\frac{da\,db}{(2\pi)^2}
 d\Phi_2(P;Q,R)d\Phi_2(Q;k_1,k_2)d\Phi_2(R;k_3,k_4).
\end{equation}
它同时固定了相空间归一化和积分区域 $\sqrt a+\sqrt b\le\sqrt s$。

\begin{proposition}[阈值首项]
设 $\epsilon=\sqrt s/m-4\to0^+$，则
\begin{equation}
 \sigma_{24}\sim\frac{g^4}{m^2}
 \frac{\sqrt{2/3}}{5160960\pi^4}\epsilon^{7/2}.
\end{equation}
\end{proposition}
\begin{proof}
阈值时四个类时分母为 $8m^2$，六个类空分母为 $-4m^2$，故 $S_6\to-1/m^2$。两对质量变量写成 $\sqrt a=2m+m\epsilon u$、$\sqrt b=2m+m\epsilon(1-u)v$；精确 Jacobian 提出 $\epsilon^{7/2}$。余下正权积分为 $\int_0^1(1-u)^2\sqrt u\,du=16/105$ 和 $\int_0^1\sqrt{v(1-v)}\,dv=\pi/8$，其光滑因子一致趋于阈值常数。代入通量及 $4!$ 即得。正紧参数域和无极点性保证取极限合法。
\end{proof}


\subsection{结果的精度与未完成项}
现有高精度积分在 $s/m^2=8$ 给
\begin{equation}
 \sigma=\frac{g^2+0.00356759925535171g^3
 -0.0000461288386699841g^4}{256\pi m^2}+O(g^5).
\end{equation}
在 $s/m^2=20$ 合并弹性与四体后，得到近似包容结果
\begin{equation}
 \sigma_{\mathrm{incl}}\approx
 \frac{g^2+0.0180991989932455g^3+0.00014250842g^4}{640\pi m^2}.
\end{equation}
四体系数约为 $m^2\sigma_{24}/g^4\approx5.75125\cdot10^{-11}$。这些显示小数来自普通求积与独立相空间参数化，不是全部位数的严格认证。另有真正严格、但明显较粗的区间
\begin{equation}
 8.8\cdot10^{-13}<m^2\sigma_{24}/g^4<6.2\cdot10^{-10}
 \qquad(s/m^2=20)
\end{equation}
用于证明正性与有限范围。它不能冒充上述近似小数的窄误差界。

本基准的后续责任是选定点的窄积分误差认证，以及若要求同阶总截面光学闭合时独立计算三圈前向全切割。给出光学右端的正确组合不等于验证左端。banana 的指定积分、两圈弹性基准和四体产生承担不同任务，三者不可互相替代。



\endgroup
\subsection{开放问题}\label{ch:roadmap}
结构判据把尚待解决的问题定位在具体的比较上。对一般图，首先需要构造完整相对推前，并保留分支局部系统和边界。逐纤维为Tate的结果不足以控制推前；补充材料中的二次覆盖反例说明了这一点。有限骨架也不保证可见秩有统一界，因为质量列表和分子读出仍可引入任意多的椭圆组成。

第二个问题是一般操作下的混合粘合。相邻扩张的Yoneda乘积为零，只说明存在相容提升；完整对象及其比较仍取决于所选提升。本文的banana构造给出一个实际实例，推广到任意图仍需几何输入。全局物理割也需要管态射、物理链和解析延拓之间的比较，不能只把单值矩阵解释成固定纤维的动机自同态。

重整化方面，正质量对数森林与相容MOM已有本文所列构造。一般幂次森林需要更高法向喷射及其重叠相容性，二次减除sunset只解决一个具体实例。有限质量MOM的临界极限还须确定适当质量坐标；本文的两圈核不能直接替代完整临界流。

物理基准的近期剩余问题是$s/m^2=20$的窄误差认证。必须先固定被认证量及其无量纲归一化，再给出整个积分和余项的包围。现有四体粗界虽然在绝对单位下很小，却不能认证所列近似小数，也没有同时认证包容截面的全部系数。若进一步要求$g^4$阶光学闭合，还须独立计算三圈前向振幅的全部切割。全阶重求和与非微扰连续极限不作为本文结论。

\appendix
\section{基础输入与定理依赖}\label{app:inputs}\label{app:evidence}
本附录将容易被“前述结果”掩盖的五类输入集中写出。它们分别属于基础文献、本文证明和指定纤维的有限认证。有限输入的数值精度不替代其几何标记；后续证明只在相应输入成立时调用其结论。

\subsection{锚定方案差的低圈性}
设$\Phi_a,\Phi_b$是第\ref{ch:finitemom}节规定的相容MOM角色，$g_{ba}$为实际有限局部角色，且$\Phi_b=g_{ba}\star\Phi_a$。对连通正圈数图$G$，余积展开为
\[
 \Phi_b(G)=\Phi_a(G)+g_{ba}(G)
 +\sum_{\varnothing\ne\gamma\subsetneq G}g_{ba}(\gamma)\Phi_a(G/\gamma).
\]
在共同锚点$\kappa_*$相减，整图常数$g_{ba}(G)$消失。剩余每项是严格低圈子图常数和严格低圈收缩图周期的乘积；基域、参考点评价和所需Tate因子一并纳入背景。因而差由该背景中的有限直和与张量积承载。这是第\ref{unified:thm:ren}定理调用的全部低圈输入。

这个论证的前提是实际角色关系已经构造。零动量部分的收敛和几何模型见第\ref{ch:zero}节及附录\ref{app:zero-proof}；相容非零MOM的递归及BPHZ比较在第\ref{ch:finitemom}节给出。锚定保护没有被用于证明这些前提，故依赖不成环。

\subsection{实际变分类别与物理管}
从代数相对对的上同调出发，缩小光滑参数基底并删去判别式，使其Betti实现为局部系统，de Rham实现为带正则连接的向量丛。在本文的变分陈述中，取几何产生的可容许、分级可极化有理混合Hodge变分，以及它们在这一范畴中的子商。对象有有限秩，实际子对象严格保留权重与Hodge过滤，因而最小链子对象和最大读出核存在。

若实际管$f:\mathcal C\to\mathcal E$满足$(T_\ell-1)\sigma=c f_B\gamma$，$c\ne0$，则链生成子变分被$T_\ell$保持，故含$f_B\gamma$。由第\ref{unified:prop:lower}命题即得割最小对象的单射。管的存在、取向及交数属于具体几何输入；不可约性只在它们建立之后将非零方向扩大为整个纯子变分。固定纤维的Nori最小性另用实际扩张与群，不由此处的单值论证自动得到。

\subsection{六运算、纯性与相对留数}
采用Tubach已发表版本的Theorem 6.5与Corollary 6.8\cite{tubach}：特征零有限型概形上的Nori实现及Hodge实现与六运算相容。本文只在这些范围及有理系数下使用其局部化、基变换、纯性和对偶。对光滑除子$Y\subset X$，余维一纯性给
\[
 R\Gamma_Y(X)\simeq R\Gamma(Y)(-1)[-2].
\]
若$Y$横截正常交叉边界$B$的各层，则各层上的局部化与限制相容。取有限交叉复形，再取相对映射纤维，得到第\ref{unified:thm:pencil-separation}定理的相对三角。第$n$上同调给出留数正合段；这一步无需假设一般Nori变分类别已经构造。

banana中还使用正常交叉对的相对对偶，将扭转的相对上同调识别为有效相对同调。随后才应用有效Nori一动机商定理\cite{abv}。权重型的相同本身不能提供这些实际识别。

\subsection{二点对象与两圈MOM的共同归一化}
两圈RG的二点输入是第\ref{ch:sunsetjet}节的四维二次减除sunset。先固定值和一阶导数的减除条件，再用二阶Taylor余项写出收敛参数核。该核的有理提升给实际相对对象，二阶端点下降给场归一化的导数读出。参考喷射改变时，应同时变换质量和场因子。

四点部分采用同一MOM定义，依次组合双泡、wineglass、子反项与整图局部条件。第\ref{ch:momrg}节先在独立运动学变量中求尺度导数，再拉回参考点；倒置这个顺序会丢掉链式项。两个核$b_1(r),b_2(r)$及其角点控制在该节完整给出，紫外极限为$3,-17/3$。因此这里的二点输入是实际sunset与场匹配，不是二维banana的K3秩八结论。独立散射基准又有自己的在壳条件，第\ref{ch:physics}节逐项声明。

\subsection{固定banana的几何与有限认证输入}
固定$m=(1,2,4,8)$、$s=-64$和$k=\mathbb Q$。所用输入分为四组。

首先，Hessian正常交叉模型给二十条边界曲线、三十个横截交点和秩十六边界子格。第\ref{ch:core}节的投影子从这些类切出秩六纯核，质量双留数确定最大一动机商。完整边界矩阵和旗标计算是几何证书的一部分；仅有分级维数不足以恢复质量点。

其次，共同基点的六个几何循环、全局闭路径及其严格运输给第\ref{eq:gram}式。十五个对迹与六个三迹的包围误差至多$2^{-69}$；整数性来自Picard--Lefschetz拓扑，区间负责唯一选出整数。Hessian格的本原性另给完整超越格$U\oplus U(2)\oplus A_2(-2)$。两者结合后，第\ref{ch:lattice}节的判别式型证明才识别唯一指数二补全。

第三，$p=11$的良约化及点数$288,16568,1794960$给第\ref{eq:fixed:3}式的六次残余Frobenius多项式。其不可约性和非单位根条件排除额外除子，结合已有十六个除子给Picard数十六。第\ref{ch:fixed}节由此排除实乘与CM，得到固定纯Hodge群。这里的素数、基点及残余多项式必须属于同一纤维。

最后，裸留数的Morse定标、边界环面及相对膜固定比较矩阵的尺度和符号。下层实分裂障碍的行列式落在$(8.77,8.78)$，与上层$\log2\ne0$分别证明两个扩张非零。纯与八维比较的误差界、Tate扭转及物理链标记见第\ref{ch:comparison}节。由此调用第\ref{thm:core:B}定理得到最小秩八。完整特殊函数闭式以及混合整Betti基不包含在这些输入中。

以上有限认证继承自原研究档案，本次长论文重构未重新运行它们。可复核交付应同时保存源算子、基和路径、误差不等式及证书数据；本文列出的结果是其数学接口，不能作为替代原证书的运行记录。依赖关系如图\ref{fig:dependencies}，关键结构定理的证明不依赖这些指定点的数值输入。

\section{零动量森林的边界与收敛证明}\label{app:zero-proof}

\begingroup

\subsection{对数子图与收缩的封闭性}\label{sec:zero:2}

若\(\gamma \)为紧子集，即\(d_G(\gamma )=0\)，则对\(T\subset E\\gamma \)，

\begin{equation}
h_{G/\gamma}(T)=h_G(T\cup\gamma)-h_G(\gamma),\qquad
 d_{G/\gamma}(T)=d_G(T\cup\gamma)-d_G(\gamma)\geq0.     \label{eq:zero:3}
\end{equation}

第一式是图循环空间在收缩下的秩公式；把边数减去两倍圈数即得第二式。原图的紧子图及其紧收缩商仍满足\eqref{eq:zero:1}。对子图自身取更小边集，圈数由该边集决定，故同样满足非负条件。

因此可使用有边标记的、以对数UV子图为抽取对象的图Hopf代数H：

\begin{equation}
\Delta_H G=\sum_{\gamma\subseteq E,\ d_G(\gamma)=0}
                         \gamma\otimes G/\gamma.             \label{eq:zero:4}
\end{equation}

空和全子图均包括；不连通\(\gamma \)解释为其有边连通分支的乘积，每个分支都是对数型，因为各分支的d非负且总和为零。不同嵌入保留各自原边标签，不省掉多重系数。

\eqref{eq:zero:3}给抽取与收缩的双射，从而给\eqref{eq:zero:4}的余结合性。乘法为不交并，分次为边数，零次仅单位，故H连通且有反极\(S_H\)。把H限制到G反复抽取/收缩所需的有限类型即可；全体图的形式框架不等于一个宣称高效的枚举算法。

公式只在图和边标记层面使用Hopf结构。物理外动量保留在最后的根商图；左侧反项块全部置零外动量。这样不需要为子图原来的内部半边猜测额外运动学。


\subsection{准备映射及其恒等式}\label{sec:zero:3}

令\(\pi _+=id- u\epsilon _H\)，取

\begin{equation}
\mathcal R=(S_H\otimes\pi_+)\Delta_H,\qquad\mathcal R(1)=0.     \label{eq:zero:5}
\end{equation}

对正次图展开反极，得到有限的交替旗标和。写成

\begin{equation}
\mathcal R(G)=\sum_{A,B}c_{A,B}\,A\otimes B,\qquad B\ne1.      \label{eq:zero:6}
\end{equation}

A是被抽取/收缩后的反项块的乘积，B是非空根商。每项的边标签仍分割原E，没有复制积分变量。首项是\(1\otimes G\)。

准备映射的基本恒等式为

\begin{equation}
\mu_{13}\bigl(\Delta_H\otimes((id\otimes\pi_+)\Delta_H)\bigr)
           \mathcal R
       =1\otimes\mathcal R,                                  \label{eq:zero:7}
\end{equation}

其中\(\mu_{13}(a\otimes b\otimes c\otimes d)=ac\otimes b\otimes d\)。

为避免单位约定含糊，给出直接验证。将四重余积写成\(\sum x_1\otimes x_2\otimes x_3\otimes x_4\)，利用\(\Delta S=(S\otimes S)\Delta ^{\mathrm{op}}\)，左侧为

\begin{equation}
\sum S_H(x_2)x_3\otimes S_H(x_1)\otimes\pi_+(x_4)
   =1\otimes\sum S_H(x_1)\otimes\pi_+(x_2).
\end{equation}

等号只用反极公理和余结合性。\(\pi _+\)杀掉右单位，故该证明也涵盖\(x=1\)。正次图上的\eqref{eq:zero:7}对应Brown--Kreimer Theorem66；\eqref{eq:zero:5}明确了本文的单位约定。

重叠子图不需要同时成为一个森林的成员。\eqref{eq:zero:6}取所有允许抽取旗标的总和，\eqref{eq:zero:7}控制其全部项；这是后面处理重叠的方式。


\subsection{全部坐标面上的极点和留数}\label{sec:zero:5}

令\(X_G\)为\(P^{N- 1}\)沿全部真坐标面的完整旗标toric模型；坐标边界D为SNC。每个真非空\(I\subset E\)给一个面，局部将\(\alpha _e(e\in I)\)同乘t。

对每项写\(I_A=I\cap E(A)\)、\(I_B=I\cap E(B)\)，并逐块计算圈数。令

\begin{equation}
d_{A,B}(I)=d_A(I_A)+d_B(I_B)\geq0.
\end{equation}

适应生成树给

\begin{equation}
\frac{\Omega_E}{U_A^2U_B^2}
       =t^{d_{A,B}(I)-1}\,dt\wedge(\text{正则形式}).            \label{eq:zero:12}
\end{equation}

\(\Phi _B\)沿\(I_B\)的阶至少为\(h_B(I_B)\)。若\(I_B=B\)，阶为\(h_B+1\)。I为真子集时\(M_E\)在面上保留至少一条边的质量项；因此在相应正开集，\eqref{eq:zero:9}附加因子没有负阶。各项最多有简单的对数边界极点。

\subsubsection{非对数面和右块全收缩}\label{sec:zero:5.1}

\(d_{A,B}(I)>0\)时留数为零。

若\(d_{A,B}(I)=0\)而\(I_B=B\)，则\(q_B\)多出一阶t；I真且B非空意味着A仍有未缩边，其质量使分母的面值非零。因此留数同样为零。这正是\(\Psi (A'\otimes 1)=0\)的几何原因，必须纳入审计，不能只检查原G的发散子图。

\subsubsection{剩余对数面}\label{sec:zero:5.2}

设\(d_{A,B}(I)=0\)且\(I_B\ne B\)。由各d非负，\(I_A\)、\(I_B\)在各自块中都是紧子集。树/二森林因子化给

\begin{equation}
U_B=t^{h_{I_B}}\bigl(U_{I_B}U_{B/I_B}+O(t)\bigr),\qquad
 \Phi_B=t^{h_{I_B}}\bigl(U_{I_B}\Phi_{B/I_B}+O(t)\bigr),         \label{eq:zero:13}
\end{equation}

\begin{equation}
M_E=M_{E\setminus I}+O(t).
\end{equation}

\(\Phi _{B/I_B}=0\)时首项自动为零，不需假设非退化外运动学。

固定原边全序，按法向dt/t在前的约定取裸留数。令\(s_I\)为由这一边序决定的置换号；它对同一I的所有森林项相同。去掉共同\(s_I\)后，\eqref{eq:zero:12}--\eqref{eq:zero:13}给

\begin{equation}
\operatorname{res}_I\Psi_\tau(A\otimes B)
 =\omega_{\rm UV}(I_A I_B)\wedge
             \Psi_\tau(A/I_A\otimes B/I_B),\qquad
 \omega_{\rm UV}(H)=\Omega_{E(H)}/U_H^2,\quad\omega_{\rm UV}(1)=0.       \label{eq:zero:14}
\end{equation}

在面的两个射影因子上理解右式，后续旗标吹起取其拉回。即使H不连通，\(\omega UV(H)\)仍是其全部边变量上的一个射影形式；此处没有把它当成已经收敛的UV数值周期或简单周期乘积。

每个留数以实际子集I为标签。将\eqref{eq:zero:14}汇总，是

\begin{equation}
\operatorname{Res}\Psi_\tau
   =(\omega_{\rm UV}\otimes\Psi_\tau)
                         \mu_{13}(\Delta_H\otimes\Delta_H).   \label{eq:zero:15}
\end{equation}

当用\eqref{eq:zero:15}计算真面的留数时，空左因子由\(\omega UV(1)=0\)消去，全右收缩由\(\Psi (...\otimes 1)=0\)消去。共同取向号不影响零结论；这里没有引入或遗漏管积分的\(2\pi i\)因子。


\subsection{森林同时消去全部面及角点}\label{sec:zero:6}

把\eqref{eq:zero:15}作用于\eqref{eq:zero:5}，右单位项先消失，再用\eqref{eq:zero:7}，得到

\begin{equation}
\operatorname{Res}\Psi_\tau\mathcal R
 =(\omega_{\rm UV}\otimes\Psi_\tau)(1\otimes\mathcal R)=0.       \label{eq:zero:16}
\end{equation}

每个边标签对应的留数分别为零，不只是所有面的数值求和为零。

还需检查真正的角点。选旗标坐标\(t_1,...,t_{N- 1}\)。在删掉非边界极点的开集内，各项至多对数，因此总形式可写成

\begin{equation}
\frac{f(t,\tau)}{t_1\cdots t_{N-1}}\,
                 dt_1\cdots dt_{N-1}.
\end{equation}

\eqref{eq:zero:16}意味着f在每个\(t_j=0\)上为零；在光滑正常交叉坐标环中，它被每个\(t_j\)整除，因而被其乘积整除。总形式在所有角点正则。\(\tau \)是代数参数，恒等式包括\(\tau =0\)、1，并非只对内部某个数值成立。

这证明了相容和重叠情形的共同结论。它不声称两个重叠的原坐标中心在同一局部旗标内相交；完整模型的坐标与准备映射处理的是两个不同层面的问题。


\subsection{正链避让和实际相对对象}\label{sec:zero:7}

对每个Hepp扇区，以最大Schwinger参数为1，取其余参数为旗标单项式，\(t_j\in [0,1]\)。适应树给每个U块提取一个共同最小单项式，剩下的多项式\(\widehat U\)在闭立方体上至少为1；这是树多项式的正系数及适应树单项式给出的下界。

\(\Phi _B\)可提取同样的\(U_B\)边界单项式，剩下\(\widehat \Phi _B\)非负。于是

\begin{equation}
\widehat D_{A|B}
    =\widehat U_B M_E+\tau\widehat\Phi_B
       \geq\mu_{e_{\max}}\widehat U_B>0
           \quad (t,\tau\in[0,1]^{N}).                        \label{eq:zero:17}
\end{equation}

所有质量严格正使\(\tau =0\)仍有统一下界。有限森林和有限旗标取共同最小常数即可；不用先求任何周期数值。

从\(X_G\times A^{1}_\tau \)删去全部森林项的U块和\(D_{A|B}\)非边界严格变换的零集，得到\(Y_G\)；取

\begin{equation}
B_G=Y_G\cap\bigl((D\times\mathbb A^1)
                  \cup(X_G\times\{0,1\})\bigr).              \label{eq:zero:18}
\end{equation}

边界单项式不作为被删掉的极点；它们已由\eqref{eq:zero:16}消去。由于\eqref{eq:zero:17}，提升正链\(K_G=\widetilde\sigma _G\times [0,1]\)紧致且避开删除集。\(\Xi _G\)是\(Y_G\)上的正则N形式，在\(B_G\)的拉回为零，故给相对de Rham类；\(K_G\)给相对Betti类。因此

\begin{equation}
R_G^{\mathfrak m}
  =[H^N(Y_G,B_G),[\Xi_G]_{dR},[K_G]_B]^{\mathfrak m}           \label{eq:zero:19}
\end{equation}

是实际Nori对象的标架周期，规范coaction随之存在。其最小对象仍须按标架另取子商，不能将整个\(H^N\)的环境类型等同物理可见类型。

\(\kappa =0\)时所有根\(\Phi _B\)恒零，\(\Xi _G\)恒零；故形式周期本身为零，不借助周期映射单射性。没有据此要求环境\(H^N\)在该点消失。


\subsection{与物理零外动量BPHZ的比较}\label{sec:zero:8}

有实际收敛几何并不自动识别物理振幅。本节单独补齐这个箭头。

\subsubsection{反项递归与共同Schwinger变量}\label{sec:zero:8.1}

对数BPHZ的子反项是外动量无关常数。把裸Feynman规则记为\(\phi \kappa \)，零外动量取值为\(\phi 0\)；可先在一个可乘的UV正规化中写代数递归。零点反项C0满足

\begin{equation}
C_0(G)=-\phi_0(G)-\sum_{\emptyset\ne\gamma\subsetneq G}
                                C_0(\gamma)\phi_0(G/\gamma),
\end{equation}

所以\(C0=\phi 0\)\ensuremath{\circ}\(S_H\)。由于C0*\(\phi 0=\epsilon _H\)，对正次图

\begin{equation}
(C_0*\phi_\kappa)(G)
   =\sum_{A,B}c_{A,B}\,
                  \phi_0(A)\bigl(\phi_\kappa(B)-\phi_0(B)\bigr).       \label{eq:zero:20}
\end{equation}

这正是先减所有子发散再作整图零次Taylor减除的森林规则；反极的嵌入重数和符号不是另行拟合。

把各块的高斯动量积分转为Schwinger密度。由于各项边标签分割E，全部项可写到同一个正正交体上。\(D=4\)的合并被积密度为

\begin{equation}
\sum_{A,B}c_{A,B}\,
 \frac{e^{-M_E-q_B}-e^{-M_E}}{U_A^2U_B^2}\,d^N\alpha.          \label{eq:zero:21}
\end{equation}

单项不作不合法的无正规化积分。下面直接证明\eqref{eq:zero:21}的总和绝对可积，从而它定义未正规化的BPHZ森林积分；无需把某个数值正规化有限部分假定为普通周期。

\subsubsection{同时控制UV角点和整体尺度}\label{sec:zero:8.2}

在每个Hepp扇区写\(\alpha =r\beta \)，最大的\(\beta \)等于1。所有块总圈数为\(h_G\)，\(N=2h_G\)，故径向密度是dr/r。恒等式

\begin{equation}
e^{-r(M+q_B)}-e^{-rM}
      =-r\int_0^1q_Be^{-r(M+\tau q_B)}\,d\tau                 \label{eq:zero:22}
\end{equation}

先在内部逐点成立。

对固定\(r\ge 0\)和\(\tau \in [0,1]\)，考虑扇区形式

\begin{equation}
\Theta_{r,\tau}(A\otimes B)
   =\frac{\Omega_E}{U_A^2U_B^2}\,
                         q_Be^{-r(M_E+\tau q_B)}.             \label{eq:zero:23}
\end{equation}

这里最大\(\beta \)等于1的截面已固定，\eqref{eq:zero:23}不冒充单独齐次的射影形式。其UV首项有与\eqref{eq:zero:14}相同的因子化：\(q_B\)限制为\(q_{B/I_B}\)，质量限制为\(M_{E\setminus I}\)；\(I_B=B\)时\(q_B\)多一阶。右单位仍为零。因此同一\eqref{eq:zero:7}给总和\(\Theta _{r,\tau }(R(G))\)的每个坐标留数为零。

正立方体上，提取边界单项式后的\(q_B=\widehat\Phi _B/\widehat U_B\)为光滑非负函数，其各有限阶导数有界；\(M_E\ge \mu _{e_{\mathrm{max}}}\)。把全部项写成

\begin{equation}
\frac{F(t,r,\tau)}{t_1\cdots t_{N-1}}\,
                      dt_1\cdots dt_{N-1}.
\end{equation}

F及其t混合导数的绝对值至多为\(C(1+r)^L e^{-cr}\)，\(c>0\)，L有限。理由是U余因子有正下界，所有其余代数因子在闭立方体光滑；对指数求有限次t导数只增加r的多项式因子。

全部面留数为零即\(F|_{t_j=0}=0\)。重复使用一维积分余项，

\begin{equation}
\frac{F(t,r,\tau)}{\prod_jt_j}
 =\int_{[0,1]^{N-1}}
   (\partial_1\cdots\partial_{N-1}F)
              (v_1t_1,\ldots,v_{N-1}t_{N-1},r,\tau)\,dv.     \label{eq:zero:24}
\end{equation}

故消极点后的总密度也被同一个可积型函数\(C(1+r)^L e^{-cr}\)控制，且在\(r=0\)有界。该估计同时控制所有角点、\(r\to 0\)、\(r\to \infty \)及\(\tau \)端点，没有逐项绝对可积的假设。

所以\eqref{eq:zero:21}的合并密度绝对可积，可用Fubini交换r、\(\tau \)和角变量。此处给的是未正规化森林总和的收敛证明；没有声称已对所有复\(\epsilon \)及所有MS约定完成统一极限比较。

\subsubsection{\texorpdfstring{径向积分给\eqref{eq:zero:10}}{径向积分给}}\label{sec:zero:8.3}

利用

\begin{equation}
\int_0^\infty q_Be^{-r(M_E+\tau q_B)}\,dr
                           =\frac{q_B}{M_E+\tau q_B},         \label{eq:zero:25}
\end{equation}

\eqref{eq:zero:22}带来的整体负号恰给\eqref{eq:zero:10}。因此\eqref{eq:zero:2}既是实际Nori周期，又是\eqref{eq:zero:20}规定的物理零外动量BPHZ值。没有通过计算样例或高精度相等代替这一比较。

若要把另一个正规化程序的极限与\eqref{eq:zero:21}比较，仍需该程序自己的森林定理或相应一致估计；本定理不把这一责任抹去。


\endgroup

\begin{thebibliography}{99}
\addcontentsline{toc}{section}{参考文献}
\bibitem{bek} S. Bloch, H. Esnault and D. Kreimer, \emph{On Motives Associated to Graph Polynomials}, \href{https://arxiv.org/abs/math/0510011}{arXiv:math/0510011}.
\bibitem{bkv} S. Bloch, M. Kerr and P. Vanhove, \emph{A Feynman integral via higher normal functions}, Compositio Mathematica \textbf{151} (2015), 2329--2375, \href{https://arxiv.org/abs/1406.2664}{arXiv:1406.2664}.
\bibitem{bv} S. Bloch and P. Vanhove, \emph{The elliptic dilogarithm for the sunset graph}, \href{https://arxiv.org/abs/1309.5865}{arXiv:1309.5865}.
\bibitem{brownperiods} F. Brown, \emph{Notes on Motivic Periods}, \href{https://arxiv.org/abs/1512.06410}{arXiv:1512.06410}.
\bibitem{browncosmic} F. Brown, \emph{Feynman amplitudes, coaction principle, and cosmic Galois group}, \href{https://arxiv.org/abs/1512.06409}{arXiv:1512.06409}.
\bibitem{abv} J. Ayoub and L. Barbieri-Viale, \emph{Nori 1-motives}, \href{https://arxiv.org/abs/1206.5923}{arXiv:1206.5923}, Theorem 8.12.
\bibitem{tubach} S. Tubach, \emph{On the Nori and Hodge realisations of Voevodsky motives},
Compositio Math. \textbf{161} (2025), 2155--2201.
\href{https://doi.org/10.1112/S0010437X25102431}{doi:10.1112/S0010437X25102431}.
Published Theorem 6.5 and Corollary 6.8 (Theorem 4.5 and Corollary 4.8 in arXiv:2309.11999v4).
\bibitem{eskandari} P. Eskandari, \emph{Tannakian fundamental groups of blended extensions}, \href{https://arxiv.org/abs/2407.01379}{arXiv:2407.01379}.
\bibitem{charles} F. Charles, \emph{On the Picard number of K3 surfaces over number fields}, \href{https://www.imo.universite-paris-saclay.fr/~francois.charles/Picard.pdf}{author's manuscript}, especially Section 2.2.
\bibitem{zarhin} Yu. G. Zarhin, \emph{Hodge groups of K3 surfaces}, Journal für die reine und angewandte Mathematik \textbf{341} (1983), 193--220. \href{https://doi.org/10.1515/crll.1983.341.193}{doi:10.1515/crll.1983.341.193}.
\bibitem{ej} A.-S. Elsenhans and J. Jahnel, \emph{On the component group of the algebraic monodromy group of a K3 surface}, \href{https://wwwuser.gwdguser.de/~jjahnel/Arbeiten/Komponentengruppe_v4f.pdf}{author's manuscript, 10 January 2024}, Theorems 7.5 and 7.7. 本文仅在所述 CM 条件下使用其正规性结论。
\bibitem{dvgeemen} E. Dardanelli and B. van Geemen, \emph{Hessians and the moduli space of cubic surfaces}, \href{https://arxiv.org/abs/math/0409322}{arXiv:math/0409322}, Section 1.6. 完整 Hessian 整格公式及其在本族的适用性见项目的格本原性证明。
\bibitem{dualcritical} Z. Király and S. Kisfaludi-Bak, \emph{Notes on dual-critical graphs}, \href{https://arxiv.org/abs/1410.1653}{arXiv:1410.1653} (2014).
\bibitem{kiefer} S. Kiefer, \emph{Notes on equivalence and minimization of weighted automata}, \href{https://arxiv.org/abs/2009.01217}{arXiv:2009.01217}.
\bibitem{muhlbauer} M. Mühlbauer, \emph{Cutkosky's Theorem for Massive One-Loop Feynman Integrals---Part 1}, \href{https://arxiv.org/abs/2206.08402}{arXiv:2206.08402}. 本文使用留数及变分背景，不跨越其一圈物理定理的假设。
\bibitem{bbr} C. Bergbauer, R. Brunetti and D. Kreimer, \emph{Renormalization and resolution of singularities}, \href{https://arxiv.org/abs/0908.0633}{arXiv:0908.0633}.
\bibitem{bk2008} S. Bloch and D. Kreimer, \emph{Mixed Hodge Structures and Renormalization in Physics}, \href{https://arxiv.org/abs/0804.4399}{arXiv:0804.4399}.
\bibitem{bk2011} F. Brown and D. Kreimer, \emph{Angles, scales and parametric renormalization}, \href{https://arxiv.org/abs/1112.1180}{arXiv:1112.1180}, especially Proposition 34 and Theorems 40, 66.
\bibitem{bk2012} F. Brown and D. Kreimer, \emph{Decomposing Feynman rules}, \href{https://arxiv.org/abs/1212.3923}{arXiv:1212.3923}.
\bibitem{guozhang} L. Guo and B. Zhang, \emph{Differential Birkhoff decomposition and the renormalization of multiple zeta values}, \href{https://arxiv.org/abs/0710.0432}{arXiv:0710.0432}, Theorem 3.2.
\bibitem{antipin} O. Antipin, A. Pinoy, F. Sannino and S. Vatani, \emph{Consistency Relation for Fixed Point Dynamics}, \href{https://arxiv.org/abs/2504.05988}{arXiv:2504.05988}.
\bibitem{kp} M. V. Kompaniets and E. Panzer, \emph{Minimally subtracted six loop renormalization of O(n)-symmetric phi4 theory and critical exponents}, \href{https://arxiv.org/abs/1705.06483}{arXiv:1705.06483}, equations (16)--(19).
\bibitem{brownmzv} F. Brown, \emph{On the decomposition of motivic multiple zeta values}, \href{https://arxiv.org/abs/1102.1310}{arXiv:1102.1310}.
\bibitem{pdg} Particle Data Group, \emph{Kinematics}, Review of Particle Physics, \href{https://pdg.lbl.gov/2025/reviews/rpp2025-rev-kinematics.pdf}{2025 review}. 本文单独固定同种末态阶乘和 $d\Phi_n$ 约定。

\bibitem{stacks}The Stacks Project,
\href{https://stacks.math.columbia.edu/tag/02MS}{Tag 02MS},
Lemma 12.10.6. Serre商及精确函子。

\bibitem{KR}E. Kani and M. Rosen,
\emph{Idempotent relations and factors of Jacobians},
Math. Ann. 284 (1989), 307--327.
\href{https://eudml.org/doc/164555}{出版记录}。


\bibitem{HHMZ}J. D. Hauenstein, R. Huang, D. Mehta and Y. Zhang,
\emph{Global Structure of Curves from Generalized Unitarity Cut of Three-loop Diagrams},
\href{https://arxiv.org/abs/1408.3355}{arXiv:1408.3355}.
链式White-house亏格为先行结果，相同亏格不证明星式图与其同构。

\bibitem{bkone}S. Bloch and D. Kreimer,
\emph{Feynman Amplitudes and Landau Singularities for 1-loop graphs},
\href{https://arxiv.org/abs/1007.0338}{arXiv:1007.0338v2},
\S7--10, 特别是Proposition10.1。一圈Hodge结构及双对数的经典来源。

\bibitem{svone}M. Spradlin and A. Volovich,
\emph{Symbols of One-Loop Integrals From Mixed Tate Motives},
\href{https://arxiv.org/abs/1105.2024}{arXiv:1105.2024}。
偶数维多边形的二次型/符号递归背景；本节未宣称实现该算法。

\bibitem{abwsun}L. Adams, C. Bogner and S. Weinzierl,
\emph{The two-loop sunrise graph with arbitrary masses},
\href{https://arxiv.org/abs/1302.7004}{arXiv:1302.7004},
\S4--5，式(24)--(32)。本文采用不乘任意质量尺度的归一化。

\bibitem{bkvsunset}S. Bloch, M. Kerr and P. Vanhove,
\emph{Local mirror symmetry and the sunset Feynman integral},
\href{https://arxiv.org/abs/1601.08181}{arXiv:1601.08181v3},
\S3--4。该文参数$s$是外部动量平方的倒数，本文不沿用它的变量约定。

\bibitem{dhv2023}C. F. Doran, A. Harder and P. Vanhove,
\emph{Motivic Geometry of Two-Loop Feynman Integrals},
\href{https://arxiv.org/abs/2302.14840}{arXiv:2302.14840v1},
\S7.3, 式(120)--(121), Proposition7.6。
绝对曲面/局部系统有此先行研究；本文的相对面与标架论证单独给出，
不宣称它们的文献优先权已检索完毕。

\bibitem{internalproofs} 费曼周期研究项目，内部几何与理论证明稿（2026）。本文与补充材料给出实际对象、操作图表、有限MOM、参数导数及RG端点证明；来源映射见随附整合清单。有限认证仍分别依赖正文列出的原证书。本条不表示外部同行评议或新颖性检索已完成。
\bibitem{delignemilne}P. Deligne and J. S. Milne,
\emph{Tannakian Categories}, in \emph{Hodge Cycles, Motives, and Shimura Varieties},
Lecture Notes in Mathematics 900, Springer (1982), pp. 101--228.
使用作者维护版本 Theorem 2.11：
\url{https://www.jmilne.org/math/xnotes/tc.pdf}.

\bibitem{huybrechts}D. Huybrechts,
\emph{Lectures on K3 Surfaces}, author manuscript,
Chapter 6, \S2, especially the period-map differential and Proposition 2.8.
\url{https://www.math.uni-bonn.de/people/huybrech/K3Global.pdf}.
周期映射微分与局部 Torelli 的一般理论归于此文献。

\bibitem{conneskreimer}A. Connes and D. Kreimer,
\emph{Renormalization in quantum field theory and the Riemann--Hilbert problem I:
the Hopf algebra structure of graphs and the main theorem},
Commun. Math. Phys. \textbf{210} (2000), 249--273.
\url{https://arxiv.org/abs/hep-th/9912092}.

\end{thebibliography}
\end{document}
